% !TEX program = xelatex
% NOTE: Для корректного Unicode используйте XeLaTeX (или LuaLaTeX).
\documentclass[11pt]{article}
\usepackage{fontspec}
\setmainfont{CMU Serif}
\usepackage{polyglossia}
\setdefaultlanguage{russian}
\setotherlanguage{english}
\newfontfamily\cyrillicfonttt{CMU Typewriter Text}
\usepackage[margin=1in]{geometry}
\usepackage{amsmath}
\usepackage{amssymb}
\usepackage{amsthm}
\usepackage{hyperref}
\usepackage{microtype}
\usepackage{longtable}
\usepackage{etoolbox}
\usepackage{caption}
\usepackage{float}
\usepackage{tabularx}
\usepackage{array}
\usepackage{tikz}
\usepackage{pdflscape}
\usetikzlibrary{arrows.meta,positioning}

% Унификация оформления таблиц по всему документу
\newcolumntype{Y}{>{\raggedright\arraybackslash}X}
\newcolumntype{P}[1]{>{\raggedright\arraybackslash}p{#1}}
\newcolumntype{C}[1]{>{\centering\arraybackslash}p{#1}}
\newcolumntype{T}[1]{>{\raggedright\arraybackslash\ttfamily\small}p{#1}}
\setlength{\tabcolsep}{4pt}
\renewcommand{\arraystretch}{1.22}
\captionsetup[table]{font=small,skip=6pt}
\setlength{\extrarowheight}{1pt}
\AtBeginEnvironment{table}{\centering}
\newcommand{\smarts}[1]{\nolinkurl{#1}}

\emergencystretch=2em

\newcommand{\warnnote}[1]{\begin{quote}\textbf{Внимание.} #1\end{quote}}
\newcommand{\keyidea}[1]{\begin{quote}\textbf{Ключевая идея.} #1\end{quote}}
\newcommand{\validated}[1]{\begin{quote}\textbf{Подтверждено.} #1\end{quote}}
\newcommand{\openquestion}[1]{\begin{quote}\textbf{Гипотеза для проверки.} #1\end{quote}}

\frenchspacing
\setlength{\parindent}{0pt}
\setlength{\parskip}{0.6\baselineskip}

\begin{document}

\title{Отчёт: Физически-информированная графовая нейронная сеть\\
для предсказания растворимости твёрдых веществ в жидкостях}
\author{Поломошнов Никита Львович}
\date{Апрель 2026}
\maketitle

\tableofcontents
\newpage

\section*{Часть I. Постановка задачи и термодинамический фундамент}
\addcontentsline{toc}{section}{Часть I. Постановка задачи и термодинамический фундамент}

\section{Введение}

\subsection{Практическая важность задачи}

Предсказание равновесной растворимости твёрдого растворяемого вещества (компонент 2) в жидком растворителе (компонент 1) при заданной температуре $T$ является центральной задачей для проектирования процессов в химии и фармацевтике.

Растворимость существенно зависит от структурных особенностей молекулы: размера и формы каркаса, распределения полярности, способности образовывать водородные связи, наличия ароматических фрагментов и гибкости конформации. Поэтому близкие по формуле соединения могут существенно отличаться по растворимости в одном и том же растворителе. Для исследователя это означает, что модель должна учитывать как химическую структуру, так и физику межмолекулярных взаимодействий. Целевая величина в работе --- логарифм мольной доли растворённого вещества в насыщенном растворе:

\begin{equation*}
y = \ln x_2.
\end{equation*}

Диапазон значений в практических задачах широк: типично $\ln x_2 \in [-25, 0]$ (то есть $x_2 \in [10^{-11}, 1]$). Такая шкала означает, что модель должна одинаково хорошо работать как для практически нерастворимых, так и для хорошо растворимых систем.

Для наглядности: $\ln x_2=-10$ соответствует мольной доле $x_2\approx 4.5\times10^{-5}$, что типично для плохо растворимых фармацевтических соединений класса BCS II. При $\ln x_2=-2$ получаем $x_2\approx 0.14$ --- умеренная растворимость, характерная для многих органических соединений в полярных органических растворителях.

Минимальные рабочие определения (полный словарь --- в §\ref{sec:glossary}): \textbf{SLE} --- равновесие «твёрдое--жидкость», \textbf{GNN} --- графовая нейронная сеть для молекулярных графов, \textbf{NRTL} --- модель активности на основе гипотезы локального состава.

Практическая значимость задачи проявляется в нескольких отраслях.

\begin{enumerate}
  \item \textbf{Фармацевтика.} Растворимость активных фармацевтических ингредиентов (API) определяет биодоступность, выбор технологического маршрута и параметры кристаллизации. На этапе разработки лекарственной формы требуется быстрый отбор растворителей и антирастворителей, а также оценка температурных окон кристаллизации.
  \item \textbf{Химическая инженерия.} В процессах экстракции, перекристаллизации и очистки растворимость задаёт достижимые концентрации, материальные потоки и энергозатраты. Ошибка даже на уровне десятых долей в $\ln x_2$ может приводить к заметной ошибке в расчёте выходов и режимов.
\end{enumerate}

Отдельной прикладной проблемой остаётся дефицит экспериментальных данных: для подавляющего большинства пар «растворяемое вещество--растворитель» отсутствуют прямые измерения при нужных температурах. Эксперимент остаётся дорогим и медленным, поэтому возникает потребность в предиктивных моделях, способных работать в условиях разреженных данных.

\subsection{Существующие подходы и их ограничения}

Подходы к прогнозу растворимости можно условно разделить на две большие группы: физико-химические модели на основе групповых вкладов и модели машинного обучения, основанные на данных.

\textbf{Групповые методы} (например, семейства UNIFAC-подобных подходов и COSMO-ориентиро-ванных реализаций, включая NIST-COSMO) обеспечивают интерпретируемую параметризацию межмолекулярных взаимодействий. Однако их точность и переносимость ограничены полнотой набора функциональных групп, качеством параметров взаимодействий и способностью схемы учитывать тонкие структурные различия конкретных молекул.

\textbf{Модели машинного обучения без физического блока} (RF на дескрипторах, GNN с прямым предсказанием $\ln x_2$) нередко показывают высокую подгонку на обучающей выборке, но имеют принципиальные ограничения:

\begin{enumerate}
  \item не предоставляют интерпретируемых промежуточных термодинамических величин;
  \item не гарантируют термодинамическую консистентность;
  \item склонны к ухудшению качества при экстраполяции по температуре вне диапазона обучения;
  \item дают ограниченные возможности переноса обучения, поскольку не выделяют физически осмысленные латентные параметры процесса.
\end{enumerate}

Таким образом, между интерпретируемостью и предсказательной мощностью возникает классический компромисс: строго физические модели часто недостаточно гибки, а сугубо статистические --- недостаточно физически надёжны.

\subsection{Цель работы}

Цель данной работы --- построить физически-информированную GNN-модель для SLE, которая объединяет преимущества физического и ML-подходов. В отличие от прямой регрессии $\ln x_2$, модель предсказывает физические параметры системы (в частности, $T_{\mathrm{m}}$ в $\mathrm{K}$, $\Delta H_{\mathrm{fus}}$ в $\mathrm{J/mol}$ и параметры NRTL), после чего эти параметры подаются в аналитически заданный термодинамический решатель.

Сформулируем целевые требования к модели:

\begin{enumerate}
  \item \textbf{Физическая корректность:} использование точного уравнения SLE как вычислительного ядра;
  \item \textbf{Интерпретируемость:} сохранение физического смысла всех промежуточных предсказаний;
  \item \textbf{Точность:} качество не хуже сильных ориентиров из машинного обучения (и, желательно, выше при сопоставимом бюджете обучения);
  \item \textbf{Обобщаемость по температуре:} устойчивое поведение при интерполяции и умеренной экстраполяции по $T$.
\end{enumerate}

\textbf{Основные вклады данной работы:}
\begin{enumerate}
  \item Архитектура TGNN-Solv, объединяющая графовый энкодер с дифференцируемым термодинамическим решателем SLE.
  \item Корректное неявное дифференцирование через итеративный физический слой без накопления вычислительного графа по шагам решателя.
  \item Систематическая диагностика источников ошибки в физически-информированных моделях, включая линейную пробу и декомпозицию ошибок.
  \item Анализ идентифицируемости параметров и условий, при которых физическая декомпозиция становится неустойчивой.
  \item Практические рекомендации по балансировке многокомпонентных функций потерь в физически-информированном обучении.
\end{enumerate}

\subsection{Структура отчёта}

Дальнейшее изложение организовано от физического основания к архитектуре модели и затем к эмпирической валидации.

\begin{enumerate}
  \item В \textbf{Части I} вводится термодинамический фундамент задачи: уравнение SLE, декомпозиция идеального и неидеального вкладов, а также роль параметров плавления и активности.
  \item В \textbf{Части II} рассматриваются вычислительные аспекты: постановка в виде неподвижной точки, численный решатель и неявное дифференцирование.
  \item В \textbf{Части III} описывается архитектура TGNN-Solv, включая построение молекулярных признаков, блоки взаимодействия и физический слой.
  \item В \textbf{Части IV} формализуются протокол обучения, функции потерь, разбиения данных и настройка оптимизации.
  \item В \textbf{Части V} анализируются эмпирические результаты, гипотезы отказа, фальсификационные эксперименты и дорожная карта модификаций.
\end{enumerate}

Такое построение позволяет последовательно связать математическую постановку, инженерную реализацию и экспериментальные выводы в единую воспроизводимую исследовательскую логику.

\section{Глоссарий и обозначения}
\label{sec:glossary}

Раздел предназначен для читателей из смежных областей (химия, фармацевтика, химическая технология), чтобы все сокращения и технические термины были определены до обсуждения моделей и экспериментов.

Поскольку в тексте местами используются устоявшиеся англоязычные термины из литературы по машинному обучению, хемоинформатике и термодинамике, в глоссарий включены не только русские определения, но и рабочие английские формы терминов, встречающиеся далее без дополнительного перевода.

\begin{longtable}{p{0.28\textwidth}p{0.67\textwidth}}
\textbf{Термин} & \textbf{Пояснение} \\
\hline
SLE (равновесие «твёрдое--жидкость») & Равновесие между кристаллической фазой вещества и насыщенным раствором. Именно в этой постановке рассчитывается растворимость. \\
GNN (графовая нейронная сеть) & Модель, которая обрабатывает молекулу как граф: вершины --- атомы, рёбра --- химические связи. \\
TGNN-Solv & Физически-информированная модель из отчёта: графовый энкодер + термодинамический решатель. \\
TIMP (термодинамически-информированная передача сообщений) & Архитектурная модификация графового энкодера, в которой обмен сообщениями разделяется на физически мотивированные каналы, чтобы лучше сохранять дисперсионную и полярно-ассоциативную информацию. \\
IDAC (коэффициент активности при бесконечном разбавлении) & Экспериментально измеренный коэффициент активности $\gamma_2^{\infty}$ для пары \textit{solute--solvent}, используемый как вспомогательный обучающий сигнал для ограничения параметров NRTL. Не следует путать с величиной $\gamma_2$, вычисляемой внутри решателя при конечном составе; стартовый внешний корпус опубликован в Zenodo~\cite{idac_zenodo}. \\
Растворяемое вещество / растворитель & Растворяемое твёрдое вещество и растворитель соответственно; базовая роль компонентов в бинарной SLE-постановке. \\
Энкодер & Часть модели, которая преобразует структуру молекулы в вектор признаков. \\
Эмбеддинг & Векторное представление молекулы, атома или пары молекул, в котором модель хранит релевантную информацию для дальнейшего предсказания. \\
Узкое место (bottleneck) & Ограничивающее звено модели, которое определяет верхнюю границу качества всей системы. \\
Латентные параметры & Скрытые переменные модели, которые напрямую не наблюдаются в данных, но используются для объяснения или предсказания физических величин. \\
Выходной блок (readout) & Финальный блок агрегации и проекции, переводящий графовые эмбеддинги в целевые физические параметры или итоговый прогноз. \\
Парное представление & Совместное векторное представление пары «растворяемое вещество--растворитель» после слияния признаков, перекрёстного внимания или конкатенации. \\
Тёплый старт & Инициализация обучения из предобученного чекпойнта вместо случайного старта. \\
Чекпойнт & Сохранённое состояние весов модели и оптимизатора в конкретный момент обучения, пригодное для восстановления или дальнейшего дообучения. \\
Начальное зерно (seed) & Начальное значение генератора случайных чисел; фиксируется для воспроизводимости разбиения, инициализации и порядка батчей. \\
Одиночный запуск / серия запусков & Результат одного запуска модели по сравнению со статистикой по нескольким независимым запускам с разными начальными зёрнами. \\
Предобучение / этап 0 & Предварительное обучение на вспомогательных задачах до основной регрессии растворимости; используется для улучшения качества инициализации энкодера. \\
Дообучение / обучение на прикладной задаче & Дообучение модели на целевой задаче после предобучения; именно на этом этапе оптимизируется основная метрика по $\ln x_2$. \\
Прикладная задача & Конечная прикладная задача, ради которой строится предобучение или вспомогательные головы; в отчёте это прогноз растворимости. \\
Резервный решатель & Резервный численный режим (например, бисекция), который включается при нестабильности основной итерационной схемы. \\
Связующий штраф (bridge-loss) & Дополнительный член функции потерь, связывающий Hansen-оценку и предсказанные термодинамические параметры. \\
Регуляризация & Любое дополнительное ограничение или штраф, стабилизирующий обучение и уменьшающий переобучение либо нефизичные решения. \\
Поэтапная учебная схема & Учебная схема, в которой сложность задачи возрастает по фазам; в отчёте реализована как трёхфазное обучение с разными весами потерь. \\
Термометрический энкодер (ThermometerEncoder) & Дискретизированное температурное встраивание, то есть разбиение температуры на интервалы с one-hot-подобным кодированием, подаваемое в головы, чувствительные к $T$~\cite{buckman2018}. \\
Голова предсказания (выходной блок) & Выходной блок для конкретной физической величины: например, $T_{\mathrm{m}}$, $\Delta H_{\mathrm{fus}}$ или параметров модели активности. \\
Дескриптор & Фиксированный числовой признак молекулы (например, MolLogP, TPSA, MolWt), вычисляемый внешней хемоинформатической библиотекой. \\
Отпечаток молекулы (fingerprint) & Обычно бинарное или счётное векторное представление локальных подструктур молекулы (например, Morgan/ECFP), используемое как табличный признак. \\
Жёстко заданный решатель SLE & Вычислительный блок без обучаемых параметров; использует аналитические термодинамические формулы. \\
GC-априори (априори групповых вкладов) & Априорные оценки по методу групповых вкладов (например, Joback), используемые как физически обоснованная регуляризация. \\
Оракульная проверка с истинными параметрами & Диагностический режим, когда часть параметров берётся из эксперимента, чтобы разделить ошибку энкодера и ошибку физического слоя. \\
Принудительная подстановка истинных промежуточных значений (teacher forcing) & Режим обучения, в котором на часть шагов подаются истинные промежуточные величины вместо предсказанных; близок по духу к оракульной подстановке. \\
Линейная проба (linear probe) & Простая, обычно линейная диагностическая модель поверх замороженных эмбеддингов, используемая для оценки того, какая информация уже содержится в представлении. \\
График паритета (parity plot) & Диагностический график «предсказано по сравнению с истинным», позволяющий быстро увидеть смещение, разброс и систематические ошибки модели. \\
Режим применения модели (inference) & Режим использования уже обученной модели для получения предсказаний без обновления весов. \\
Оценка неопределённости & Оценка неопределённости прогноза; в отчёте обсуждаются MC-dropout, deep ensembles и калибровка интервалов. \\
Калибровка / калиброванный режим & Постобработка или дополнительная настройка, улучшающая согласование между предсказанной неопределённостью или значением и фактической ошибкой на валидации. \\
Применение без дообучения (zero-shot) & Режим применения внешней модели без дообучения на целевой выборке. \\
Полное дообучение на целевых данных & Полное переобучение или серьёзное дообучение внешней модели на целевых данных для оценки её верхней достижимой точности. \\
Базовая модель (baseline) & Базовая модель-сравнение, которая служит ориентиром для оценки пользы новой архитектуры или модификации. \\
Бенчмарк & Стандартизованный блок сравнений моделей при фиксированном протоколе, метриках и разбиениях. \\
Полный вычислительный бюджет & Полноценный сценарий обучения: больше эпох, больше повторных запусков и более тщательный подбор гиперпараметров. \\
Прокси-бюджет (укороченный) & Быстрый сценарий для инженерной проверки гипотез перед полным обучением. \\
Разбиение по каркасам (scaffold split) & Строгое разбиение, при котором молекулярные каркасы разделяются между обучением и тестом. \\
Случайное разбиение & Обычное случайное разбиение выборки без жёсткого контроля химических каркасов; обычно даёт более оптимистичную оценку качества, чем разбиение по каркасам. \\
Сдвиг домена & Ситуация, в которой тестовые пары молекул, температуры или растворители статистически отличаются от обучающих данных. \\
Вне области применимости (OOD) & Примеры, лежащие вне или на краю области применимости модели; именно на них особенно важны неопределённость и физическая устойчивость. \\
Абляционное исследование (ablation study) & Эксперимент с удалением или отключением конкретного компонента модели, чтобы количественно оценить его вклад в итоговое качество. \\
MAE & Средняя абсолютная ошибка прогноза. \\
$R^2$ & Коэффициент детерминации: доля вариации целевой величины, объясняемая моделью. \\
Ретросинтез & Планирование последовательности химических превращений для получения целевой молекулы из доступных прекурсоров. \\
NRTL (модель локального состава) & Полуэмпирическая модель коэффициентов активности, в которой неидеальность смеси описывается через параметры локального состава и контактных энергий между компонентами. \\
Wilson & Более простая модель коэффициентов активности с меньшим числом параметров; часто используется как контрольная альтернатива NRTL при анализе идентифицируемости. \\
UNIQUAC & Модель коэффициентов активности, разделяющая вклад на комбинаторную часть (размер и форма) и остаточную энергетическую часть. \\
UNIFAC & Групповой метод прогноза коэффициентов активности, в котором взаимодействия оцениваются через функциональные группы, а не через индивидуальные молекулы целиком. \\
COSMO-RS & Квантово-статистическая модель растворов, использующая поверхностные сегменты молекулы и их экранированные заряды для оценки межмолекулярных взаимодействий. \\
Параметры Хансена & Триада параметров $(\delta_d,\delta_p,\delta_h)$, разделяющая когезионную энергию на дисперсионный, полярный и водородно-связывающий вклады. \\
MPNN (сеть передачи сообщений) & Базовый класс графовых нейронных сетей, где атомы многократно обмениваются сообщениями вдоль химических связей, а затем эти состояния агрегируются в представление молекулы. \\
GPS (гибридный графовый трансформер) & Архитектура, сочетающая локальную передачу сообщений на графе с глобальным механизмом внимания для дальнодействующего обмена информацией. \\
MLP (многослойная перцептронная сеть) & Небольшая полносвязная нейронная сеть, используемая как универсальный аппроксиматор в сообщениях, выходных блоках и проекциях признаков. \\
Нормализация по слоям (LayerNorm) & Нормализация скрытого вектора по его компонентам внутри одного объекта; стабилизирует обучение глубоких архитектур. \\
Остаточная связь & Соединение вида ``вход + преобразование входа'', которое облегчает обучение глубоких сетей и уменьшает деградацию сигнала по слоям. \\
Перекрёстное внимание & Механизм, в котором элементы одного множества (например, атомы растворяемого вещества) вычисляют веса важности по отношению к элементам другого множества (например, атомам растворителя). \\
Самовнимание & Механизм внимания внутри одного и того же множества элементов; позволяет каждому атому или токену учитывать все остальные элементы того же объекта. \\
Агрегирование / pooling & Операция свёртки множества атомных состояний в один молекулярный вектор, например средним, вниманием или более сложной схемой. \\
Set2Set & Рекуррентная схема агрегирования графовых признаков с вниманием, возвращающая информативное молекулярное представление фиксированной размерности. \\
Неподвижная точка & Формулировка уравнения, в которой решение представляется как объект, не меняющийся после применения определённого оператора; используется при записи численного решателя. \\
Неявное дифференцирование & Вычисление градиентов через решение уравнения без разворачивания всех итераций решателя в вычислительном графе. \\
RDKit & Открытая хемоинформатическая библиотека, используемая для разбора молекул, вычисления дескрипторов, отпечатков и некоторых эвристических физических признаков. \\
BigSolDB & Рабочий корпус данных по растворимости, агрегирующий экспериментальные наблюдения для пар ``растворяемое вещество--растворитель'' при различных температурах. \\
RF (случайный лес) & Ансамбль деревьев решений, используемый в отчёте как сильный табличный ориентир на фиксированных дескрипторах. \\
QSPR & Подход ``количественная связь структура--свойство'', в котором свойства вещества предсказываются по инженерным описаниям химической структуры. \\
Optuna & Библиотека подбора гиперпараметров, автоматизирующая поиск конфигураций обучения по заданной целевой метрике. \\
TPE-сэмплер & Метод байесовского подбора гиперпараметров, оценивающий, какие конфигурации с большей вероятностью приведут к хорошему значению целевой функции. \\
Прунинг запусков & Раннее прекращение заведомо слабых запусков обучения на основании промежуточной метрики, чтобы перераспределить вычислительный бюджет. \\
Разогрев скорости обучения (warmup) & Краткий начальный этап, на котором шаг оптимизатора постепенно увеличивается от малого значения к рабочему. \\
Ограничение нормы градиента & Приём стабилизации обучения, при котором слишком большие градиенты принудительно сокращаются до заданной нормы. \\
SASA (доступная растворителю площадь поверхности) & Площадь поверхности молекулы или атома, которая доступна молекулам растворителя; служит прокси межмолекулярной контактной ``видимости''. \\
MolLogP & Дескриптор гидрофобности, оценивающий распределение молекулы между неполярной и полярной средой; часто используется как грубая мера липофильности. \\
TPSA & Топологическая полярная площадь поверхности; дескриптор, тесно связанный с полярностью, водородным связыванием и транспортными свойствами молекулы. \\
MolWt & Молекулярная масса вещества как дескриптор; часто коррелирует с размером молекулы и некоторыми транспортными свойствами. \\
MolMR & Молярная рефракция; приближённый дескриптор поляризуемости и объёма электронной оболочки молекулы. \\
Morgan/ECFP-отпечаток & Кольцевой отпечаток молекулы, кодирующий локальные окрестности атомов и широко используемый в табличных ML-моделях для химии. \\
One-hot-кодирование & Представление дискретного признака в виде вектора, где ровно одна позиция равна единице, а остальные --- нулю. \\
DirectGNN & Контрольная архитектура без физического решателя, которая предсказывает растворимость напрямую из графового и парного представления. \\
Ограниченная коррекция & Небольшой корректирующий блок поверх физического прогноза, который компенсирует систематическую погрешность модели, не разрушая физическую структуру основного решателя. \\
FusionHead & Выходной блок для предсказания кристаллических свойств растворяемого вещества, прежде всего $T_{\mathrm{m}}$ и $\Delta H_{\mathrm{fus}}$. \\
NRTLHead & Выходной блок, предсказывающий параметры модели активности NRTL из парного представления молекул. \\
Полный агрегированный корпус (merged frame) & Все записи после объединения источников, очистки и служебных соединений таблиц, независимо от наличия конечной целевой метки. \\
Подмножество с целевой меткой (supervised subset) & Строки корпуса, в которых присутствует целевая растворимость и которые действительно участвуют в основной задаче обучения. \\
GAT & Архитектура графовой нейронной сети, использующая механизм внимания для взвешивания вкладов соседних вершин. \\
FFN-блок & Позиционно-независимый полносвязный блок преобразования признаков, обычно применяемый после внимания или графового подслоя. \\
Позиционное кодирование & Способ ввести в модель информацию о положении, структуре графа или спектральной геометрии, чтобы различать узлы с одинаковыми локальными признаками. \\
RWPE & Вид позиционного кодирования на графе, основанный на статистиках случайных блужданий. \\
Адаптер & Небольшой дополнительный блок малой ёмкости, который донастраивает общее представление под конкретную роль, не меняя радикально весь энкодер. \\
Заряды Гастейгера--Марсили & Быстрая приближённая схема оценки парциальных зарядов на атомах без квантово-химического расчёта. \\
Сигмоида & Гладкая функция активации, ограничивающая значение в интервале $(0,1)$; удобна для вероятностей и мягких ограничений диапазона. \\
Softplus & Гладкая неотрицательная функция активации $\ln(1+e^z)$, часто используемая там, где требуется положительный выход без жёсткого отсечения. \\
Логит & Значение до применения нормировки softmax или сигмоиды; именно логиты затем преобразуются в вероятности или веса внимания. \\
Многоголовое внимание & Механизм внимания, в котором несколько независимых ``голов'' параллельно вычисляют разные шаблоны зависимости между элементами. \\
Агрегирование вниманием & Способ свёртки множества признаков, в котором вклад каждого элемента взвешивается обучаемым коэффициентом важности. \\
Двудольный граф & Граф, в котором вершины разбиты на две непересекающиеся доли, а рёбра допускаются только между долями. \\
Блок взаимодействия & Часть архитектуры, где признаки растворяемого вещества и растворителя начинают обмениваться информацией напрямую. \\
\end{longtable}

\subsection{Ключевые обозначения}

\begin{itemize}
  \item $x_2$ --- мольная доля растворённого вещества в насыщенном растворе;
  \item $y=\ln x_2$ --- целевая величина регрессии;
  \item $T$ --- температура в К;
  \item $T_{\mathrm{m}}$ --- температура плавления;
  \item $\Delta H_{\mathrm{fus}}$ --- энтальпия плавления;
  \item $\Delta H_{\mathrm{sol}}$ --- энтальпия растворения (в первом приближении $\Delta H_{\mathrm{sol}}\approx\Delta H_{\mathrm{fus}}+\Delta H_{\mathrm{mix}}$; в van 't Hoff-блоке используется её эффективное значение);
  \item $\gamma_2$, $\ln\gamma_2$ --- коэффициент активности растворённого вещества и его логарифм при \textbf{конечном составе}, вычисляемые внутри SLE-решателя;
  \item $\gamma_2^{\infty}$, $\ln\gamma_2^{\infty}$ --- коэффициент активности и его логарифм в пределе бесконечного разбавления; это \textbf{внешняя вспомогательная IDAC-метка}, которая связана с величиной внутри решателя только через предел $x_2\to 0$, но не тождественна ей при произвольном составе;
  \item $\Delta C_p^{\mathrm{fus}}$ --- разность теплоёмкостей между жидкой и кристаллической фазой при плавлении;
  \item $\Phi(T)$ --- идеальный термодинамический вклад в уравнении SLE, зависящий от $T$, $T_{\mathrm{m}}$, $\Delta H_{\mathrm{fus}}$ и, при необходимости, $\Delta C_p^{\mathrm{fus}}$;
  \item $G^E$ --- избыточная энергия Гиббса смеси относительно идеального раствора;
  \item $\tau_{12}, \tau_{21}$ --- безразмерные параметры взаимодействия в модели NRTL, отвечающие за несимметричность контактных энергий между компонентами;
  \item $\tau_{12}^{\mathrm{ref}}, \tau_{21}^{\mathrm{ref}}$ --- значения параметров $\tau$ при опорной температуре $T_{\mathrm{ref}}$;
  \item $\beta_{12}, \beta_{21}$ --- коэффициенты температурной зависимости параметров $\tau$ в параметризации по $1/T$;
  \item $\alpha_{12}$ --- параметр неслучайности в NRTL, определяющий степень локальной структурированности жидкости;
  \item $(\delta_d,\delta_p,\delta_h)$ --- параметры Хансена: дисперсионный, полярный и водородно-связывающий вклады в когезионные свойства вещества;
  \item $V_m$ --- молярный объём вещества;
  \item $\mathbf{h}_i^{(\ell)}$ --- скрытое состояние атома $i$ после $\ell$-го слоя графового энкодера;
  \item $\mathbf{g}_{\mathrm{mol}}$ --- молекулярное представление после агрегирования атомных признаков;
  \item $\mathbf{g}_{\mathrm{pair}}$ --- совместное представление пары ``растворяемое вещество--растворитель'', подаваемое в блок активности или прямой прогностический путь.
\end{itemize}

\subsection{Состав данных}

\begin{table}[H]
\centering
\small
\begin{tabularx}{\textwidth}{|P{4.1cm}|Y|}
\hline
Показатель & Значение \\
\hline
Полный merged frame & 120\,197 записей, 19\,878 unique solute, 212 unique solvent, 29\,826 уникальных пар; доля строк с меткой растворимости --- 90.09\% \\
\hline
Основной корпус & \textasciitilde108\,000 supervised записей (текущая water-inclusive конфигурация; точное значение 108\,287). Исторически до включения воды использовалась water-excluded конфигурация на 101\,763 строках, 211 solvent и 11\,392 парах. \\
\hline
Температурное покрытие supervised subset & $T \in [243.15,\, 425.77]$ K, медиана около 303 K; медианное число температурных точек на пару --- 9 \\
\hline
Покрытие по $\ln x_2$ & $\ln x_2 \in [-23.6237,\, -0.0466]$; медиана корпуса до включения воды составляла $-4.5282$, а water-inclusive конфигурация смещается в сторону меньшей растворимости из-за 6\,524 водных строк \\
\hline
Вспомогательные метки & $T_{\mathrm{m}}$: 51.04\% строк исторической water-excluded конфигурации; $\Delta H_{\mathrm{fus}}$: 1.00\%; Hansen: 0.85\%; $\gamma^\infty$: 0.00\% внутри базового supervised subset \\
\hline
Каноническое разбиение & \texttt{solute\_scaffold}; текущая water-inclusive конфигурация сохраняет режим 80/10/10. Исторически до включения воды это соответствовало 90\,808 / 5\,385 / 5\,570 строкам \\
\hline
\end{tabularx}
\caption{Краткая сводка по корпусу.}
\label{tab:data-summary-brief}
\end{table}

Подробная предобработка, статистика и смещения корпуса разобраны в §\ref{subsec:data-preprocessing}.

\subsection{Данные: предобработка, статистика и смещения корпуса}
\label{subsec:data-preprocessing}

Этот подраздел фиксирует, из каких источников собран корпус, какие шаги предобработки были применены и какие структурные ограничения данных необходимо учитывать при интерпретации результатов модели. Ниже принципиально различаются два среза: \textbf{полный merged frame} (все агрегированные записи после объединения источников) и \textbf{solubility-supervised subset} (строки, в которых присутствует целевая метка $\ln x_2$ и которые реально участвуют в основной задаче регрессии).

\subsubsection{Источник данных и BigSolDB}

Основной источник наблюдений --- \textbf{BigSolDB v2.1}. Корпус агрегирует измерения растворимости из разнородных экспериментальных источников и после унификации хранится в строковом формате
\begin{equation*}
(\mathrm{solute\_smiles},\; \mathrm{solvent\_smiles},\; T,\; \ln x_2),
\end{equation*}
где при необходимости присутствуют и вспомогательные поля (идентификаторы источника, признаки наличия дополнительных меток, служебные поля join-операций и контроля качества).

BigSolDB в настоящем отчёте рассматривается именно как рабочий публичный корпус для воспроизводимых сравнений. Доступна как публичный ресурс; формальная цитата уточняется.

\validated{Корпус данных в отчёте не является внутренним проприетарным активом: все количественные выводы ниже относятся к публичному ресурсу BigSolDB v2.1 и его явно описанной предобработке.}

\subsubsection{Pipeline предобработки}

Предобработка строится как детерминированный pipeline, одинаковый для всех последующих разбиений и задач:

\begin{enumerate}
  \item \textbf{Каноникализация SMILES.} Для каждого SMILES выполняется парсинг и последующая каноникализация через RDKit: \texttt{Chem.MolToSmiles(Chem.MolFromSmiles(smi))}. Это уменьшает число искусственных различий между эквивалентными строковыми представлениями одной и той же структуры.
  \item \textbf{Удаление невалидных структур.} Если \texttt{Chem.MolFromSmiles(smi)} возвращает \texttt{None}, запись исключается до любых операций объединения и агрегации.
  \item \textbf{Агрегация дубликатов.} Для записей с одинаковой тройкой $(\mathrm{solute},\, \mathrm{solvent},\, T)$ итоговое значение формируется как медиана по $\ln x_2$. Такое правило снижает чувствительность к единичным выбросам и согласуется с инженерной целью получить один целевой таргет на одну пару-температуру.
  \item \textbf{Фильтрация физических диапазонов.} Рабочее окно предобработки задаётся как $\ln x_2 \in [-30,\, 0]$ и $T \in [200,\, 600]$ K. Фактически после фильтрации и объединения наблюдаемые диапазоны оказываются уже: см. таблицу~\ref{tab:data-global-stats}.
  \item \textbf{Присоединение вспомогательных меток.} К solute-объектам дополнительно присоединяются источники по $T_{\mathrm{m}}$, $\Delta H_{\mathrm{fus}}$, Hansen-параметрам $(\delta_d, \delta_p, \delta_h)$ и, при наличии, $\gamma^\infty$. В текущей конфигурации использованы Bradley melting points для $T_{\mathrm{m}}$, NIST-источники для $\Delta H_{\mathrm{fus}}$ и справочник Hansen для $(\delta_d, \delta_p, \delta_h)$.
\end{enumerate}

Отдельно важно текущее поведение фильтра \texttt{filter\_for\_sle}. В исторической конфигурации вода исключалась из supervised subset неявно через условие \texttt{min\_atoms = 2}, поскольку молекула \smarts{O} содержит лишь один тяжёлый атом. В текущей water-inclusive конфигурации этот шаг дополнен whitelist для критически важных малых растворителей, поэтому вода целенаправленно включается в supervised subset и больше не рассматривается как артефакт фильтрации.

\validated{После агрегации по $(\mathrm{solute},\, \mathrm{solvent},\, T)$ в solubility-supervised subset не осталось дублированных pair-temperature entries: 0.}

\subsubsection{Общая статистика корпуса}

Таблица~\ref{tab:data-global-stats} разделяет полный merged frame и основной supervised-срез. Это различие критично: именно смешение этих двух уровней обычно и порождает некорректные формулировки вида «в корпусе порядка 10k уникальных solute». Для задач прямой регрессии растворимости релевантны не все 19\,878 unique solute из merged frame, а лишь существенно меньший water-inclusive supervised subset; в исторической water-excluded конфигурации это соответствовало 1\,444 unique solute.

\begin{table}[H]
\centering
\small
\setlength{\tabcolsep}{2.5pt}
\begin{tabularx}{\textwidth}{|P{2.4cm}|C{1.45cm}|C{1.25cm}|C{1.25cm}|C{1.35cm}|Y|Y|}
\hline
Срез & Строки & Solute & Solvent & Пары & $T$ (min / med / max), K & $\ln x_2$ (min / med / max) \\
\hline
Merged frame & 120\,197 & 19\,878 & 212 & 29\,826 & 243.15 / 298.15 / 425.77 & $-23.6237$ / $-3.9575$ / $0.0$ \\
\hline
Supervised subset (water-inclusive) & 108\,287 & --- & 212 & 12\,129 & $243.15 / \approx 303 / 425.77$ & $-23.6237$ / $\lesssim -4.53$ / $-0.0466$ \\
\hline
\end{tabularx}
\caption{Общая статистика полного merged frame и основного solubility-supervised subset.}
\label{tab:data-global-stats}
\end{table}

\validated{После включения воды доля строк с прямой supervision по $\ln x_2$ составляет 90.09\% от merged frame (108\,287 из 120\,197 строк). Историческая water-excluded конфигурация соответствовала 84.66\% (101\,763 из 120\,197 строк).}

\subsubsection{Температурная плотность: число точек на пару}

Для идентифицируемости температурных параметров важен не только суммарный размер корпуса, но и то, сколько температурных наблюдений приходится на одну пару \textit{solute--solvent}. В этом смысле текущий корпус заметно лучше распространённого в литературе упрощения «1--2 точки на пару».

\begin{table}[H]
\centering
\small
\begin{tabularx}{\textwidth}{|P{2.2cm}|C{2.3cm}|C{2.3cm}|C{2.3cm}|Y|}
\hline
Бин по числу точек & Число пар & Доля пар, \% & Число записей & Доля записей, \% \\
\hline
1 & 380 & 3.34 & 380 & 0.37 \\
\hline
2 & 14 & 0.12 & 28 & 0.03 \\
\hline
3 & 6 & 0.05 & 18 & 0.02 \\
\hline
4--5 & 1\,233 & 10.82 & 6\,074 & 5.97 \\
\hline
6--10 & 7\,830 & 68.73 & 69\,265 & 68.07 \\
\hline
11--20 & 1\,870 & 16.42 & 24\,554 & 24.13 \\
\hline
21+ & 59 & 0.52 & 1\,444 & 1.42 \\
\hline
\end{tabularx}
\caption{Распределение числа температурных точек на одну пару в solubility-supervised subset.}
\label{tab:data-points-per-pair}
\end{table}

Среднее число точек на пару равно 8.93, медиана --- 9. Следовательно, постановка в основном является \textbf{мультитемпературной}, а не однотемпературной.

\keyidea{Корпус преимущественно мультитемпературный: 68.73\% пар имеют 6--10 температурных точек, а медиана числа точек на пару равна 9. Это существенно повышает шансы на идентификацию температурных коэффициентов по сравнению со сценарием «1--2 точки на пару», хотя устойчивость всё равно ограничивается узостью температурного окна и коллинеарностью вкладов; см. §\ref{sec:identifiability}.}

\subsubsection{Покрытие вспомогательных меток}

Ключевое ограничение текущего корпуса --- резкая неравномерность покрытия по вспомогательным физическим меткам. Для инженерной интерпретации важно различать покрытие на уровне строк supervised subset и покрытие на уровне уникальных solute.

\begin{table}[H]
\centering
\small
\begin{tabularx}{\textwidth}{|P{2.0cm}|C{1.9cm}|C{1.9cm}|C{1.8cm}|C{1.8cm}|Y|}
\hline
Метка & Строк с меткой & Row-level, \% & Solute с меткой & Solute-level, \% & Практический статус \\
\hline
$T_{\mathrm{m}}$ & 51\,944 & 51.04 & 609 & 42.17 & Достаточно плотная вспомогательная супервизия; полезна для oracle teacher forcing при обучении, см. §\ref{subsec:teacher-forcing} \\
\hline
$\Delta H_{\mathrm{fus}}$ & 1\,016 & 1.00 & 8 & 0.55 & Крайне разрежённый канал; практически применим только как train-only supervision \\
\hline
Hansen & 863 & 0.85 & 5 & 0.35 & Ещё более разрежённый канал; информативен скорее как слабый регуляризатор, чем как полноценная задача \\
\hline
$\gamma^\infty$ & 0 & 0.00 & 0 & 0.00 & В базовом solubility-supervised subset отсутствует; фактическая supervision подключается через внешний IDAC-корпус (404 строки, 138 пар) \\
\hline
\end{tabularx}
\caption{Покрытие вспомогательных меток в solubility-supervised subset на уровне строк и уникальных solute.}
\label{tab:data-aux-coverage}
\end{table}

Для $T_{\mathrm{m}}$ покрытие достаточно велико, чтобы этот сигнал реально влиял на обучение и использовался в режиме принудительной подстановки известных кристаллических свойств. Напротив, $\Delta H_{\mathrm{fus}}$ и Hansen-параметры настолько разрежены, что их корректнее трактовать как дополнительные \textit{train-only} источники слабой супервизии.

\warnnote{В базовом supervised subset наблюдения IDAC/$\gamma^\infty$ действительно отсутствуют полностью, поэтому таблица~\ref{tab:data-aux-coverage} показывает 0\% покрытия именно для \emph{основного solubility-supervised корпуса}. Однако в текущей IDAC-augmented конфигурации канал $\mathcal{L}_{\gamma^\infty}$ уже активен за счёт отдельного внешнего IDAC-корпуса (404 строки, 138 пар), который подключается как \texttt{aux\_only\_gamma} и не обязан совпадать по строкам с наблюдениями $\ln x_2$.}

Дополнительно важно, что покрытие вспомогательных меток концентрируется преимущественно в train-части канонического split: для $T_{\mathrm{m}}$ оно составляет 54.41\% в train, 25.48\% в validation и 20.86\% в test; для $\Delta H_{\mathrm{fus}}$ --- 1.10\%, 0.33\% и 0.00\%; для Hansen --- 0.95\%, 0.00\% и 0.00\%. Это усиливает различие между «обучаемостью» физического прокси и его реальной проверяемостью на отложенных данных.

\subsubsection{Потенциальные смещения корпуса}

Наблюдаемая статистика корпуса позволяет зафиксировать несколько смещений, которые важно учитывать при интерпретации результатов.

\validated{Хотя полный merged frame содержит 19\,878 unique solute, в основной задаче регрессии растворимости используется лишь существенно меньший supervised subset. Следовательно, значительная часть молекулярного разнообразия доступна модели только через этап 0 и/или вспомогательные источники меток; см. §\ref{sec:stage0-pretraining}.}

\warnnote{Медианная температура supervised subset равна 303.15 K. Это означает, что модель в основном обучается около комнатных условий; прогнозы при $T < 260$ K и $T > 380$ K следует интерпретировать как выход за основную область обучающего распределения.}

\warnnote{Диапазон $\ln x_2 \in [-23.6237,\, -0.0466]$ покрывает более десяти порядков по мольной доле, однако левый хвост с $\ln x_2 < -15$ разрежён. Для практически нерастворимых систем там вероятны повышенные шумы и нестабильность оценки ошибки.}

\subsubsection{Распределение растворителей}

Распределение по растворителям в supervised subset выраженно скошено, но не монопольно: индекс Херфиндаля--Хиршмана равен $\mathrm{HHI}=0.0457$, коэффициент Джини --- $0.8696$, а доля top-5 / top-10 / top-20 растворителей составляет 37.05\% / 62.10\% / 79.07\% соответственно. Это означает заметную концентрацию наблюдений в ограниченном числе популярных сред, но не режим, в котором один-единственный растворитель фактически определяет весь корпус.

Вода входит в supervised subset на 6-м месте по числу записей (6\,524 строки, 6.02\%, 737 уникальных solute). Медианная водная растворимость ($\ln x_2 = -7.82$) существенно ниже медианы корпуса ($-4.53$), что отражает типичную плохую водную растворимость фармацевтически релевантных органических соединений. Дополнительно 18\,397 водных строк присутствуют как \texttt{aux\_only} для property pretraining.

\keyidea{Включение водных данных в supervised subset делает модель принципиально пригодной для предсказания водной растворимости, что критично для BCS-классификации (§VII.7.4) и PK-релевантного профиля (§VII.7.5). При этом водные системы остаются одним из наиболее сложных режимов из-за аномальных свойств воды как растворителя и выраженной ассоциации.}

Длинный хвост по растворителям сохраняется: 82 растворителя имеют менее 10 записей, а 13 представлены ровно одной записью. Следовательно, корпус одновременно содержит и хорошо представленные технологически важные среды, и широкий редкий хвост, что делает обобщение на малоизученные solvent самостоятельной задачей.

С химической точки зрения top-10 задаётся типичным для фармацевтической и тонкой органической химии набором: спирты (5 из top-10), вода, сложные эфиры, кетоны, нитрилы и амиды. Такое покрытие хорошо соответствует реальным лабораторным и process-development сценариям, но означает, что модель в первую очередь калибруется на наиболее частых фармацевтических solvent-классах.

Полный ранжированный список растворителей следует выносить в приложение: top-50 --- в основной текст приложения, а оставшиеся 163 растворителя --- при необходимости в дополнительную таблицу или внешний артефакт корпуса.

\section{Обзор литературы: существующие методы предсказания растворимости и позиционирование TGNN-Solv}

Этот раздел систематизирует классы методов, применяемых для прогноза растворимости, и фиксирует, где именно находится TGNN-Solv относительно каждого класса. В обзоре используется единая целевая переменная $y=\ln x_2$ (где $x_2$ --- мольная доля растворённого вещества в насыщенном растворе), а для работ на шкале $\log S$ явно отмечается ограниченная сопоставимость.

\subsection{Термодинамические модели (UNIFAC, COSMO, GC-корреляции)}
\label{subsec:thermo-models-review}

\textbf{UNIFAC/mod-UNIFAC.} В групповых моделях коэффициент активности строится из комбинаторного и остаточного вкладов:
\begin{equation*}
\ln \gamma_i = \ln \gamma_i^{C} + \ln \gamma_i^{R},
\end{equation*}
где $\ln \gamma_i^{C}$ описывает энтропийный вклад размера/формы (Flory--Huggins-подобная часть), а $\ln \gamma_i^{R}$ --- энергетические взаимодействия функциональных групп. В связке с уравнением SLE это даёт
\begin{equation*}
\ln x_2 = -\Phi(T) - \ln \gamma_2^{\text{UNIFAC}},
\end{equation*}
где $\Phi(T)$ включает вклад плавления (через $T_{\mathrm{m}}$, $\Delta H_{\mathrm{fus}}$, при необходимости $\Delta C_p$). Репрезентативные источники: \cite{fredenslund1975,weidlich1987,ddb2024}.

Для систем внутри матрицы параметров обычно сообщают MAE порядка $0.8$--$1.5$ по $\ln x_2$; при появлении новых или плохо параметризованных групп возможны резкие деградации точности. Сильные стороны --- интерпретируемость, нулевой этап обучения, физическая согласованность. Ограничения --- конечность словаря групп, низкая чувствительность к изомерии и внутримолекулярным эффектам, отсутствие механизма дообучения на новых данных.

\textbf{COSMO-RS/COSMO-SAC.} Эти подходы используют квантово-химический этап (обычно DFT) для оценки $\sigma$-profile/сегментных поверхностных характеристик, после чего активность вычисляется в рамках статистико-термодинамической модели сегментов. Часто используемая запись:
\begin{equation*}
\ln\gamma_i = \ln\gamma_i^{\text{comb}} + \sum_{\nu} n_{\nu}^{(i)}\!\left[\ln\Gamma_{\nu}-\ln\Gamma_{\nu}^{(i)}\right].
\end{equation*}
Репрезентативные работы: \cite{klamt1995,lin2002,klamt2017}. В задачах растворимости в ряде публикаций сообщается MAE около $0.5$--$1.0$ по $\ln x_2$ (обычно с ухудшением на выраженных системах с выраженными водородными связями и сложных водных смесях). Преимущества --- широкая применимость и меньшая зависимость от экспериментальной калибровки. Ограничения --- высокая вычислительная стоимость (минуты на молекулу при типичном DFT-конвейере), зависимость от выбранных приближений и параметризации сегментных взаимодействий.

\textbf{Групповые вклады для кристаллических свойств.} Методы Joback/van Krevelen дают априорные оценки $T_{\mathrm{m}}$, $\Delta H_{\mathrm{fus}}$, $\Delta C_p$ как сумму вкладов групп. Например,
\begin{equation*}
T_{\mathrm{m}}^{\text{GC}} = 122.5 + \sum_k n_k\,\Delta T_{\mathrm{m},k}.
\end{equation*}
Репрезентативные источники: \cite{joback1987,marrero2001}. Обычно ошибка составляет около $30$--$50$ K по $T_{\mathrm{m}}$ и $20$--$30\%$ по $\Delta H_{\mathrm{fus}}$.

\textbf{Позиционирование TGNN-Solv относительно класса.} TGNN-Solv не заменяет термодинамику непрозрачной регрессией: расчёт растворимости остаётся в явной SLE-форме с коэффициентами активности из параметризованной модели взаимодействий. Отличие от UNIFAC/COSMO --- обучаемая параметризация на данных, а не только табличные параметры или чистый \textit{ab initio} конвейер. Одновременно GC-априори для кристаллических голов напрямую опираются на инженерные подходы типа Joback и служат структурным априори для стабилизации обучения.

\subsection{Методы машинного обучения для растворимости: дескрипторные модели}

Классический QSPR-подход использует фиксированные признаки молекул и обучает регрессию на уровне пары «растворяемое вещество--растворитель» (обычно с температурой как отдельным входом):
\begin{equation*}
\hat{y} = f_\theta\!\left([\mathbf{d}_{\text{solute}}\,\|\,\mathbf{d}_{\text{solvent}}\,\|\,T]\right), \quad y=\ln x_2\;\text{или}\;\log S.
\end{equation*}
Здесь $\mathbf{d}$ --- RDKit-дескрипторы (например, MolLogP, TPSA, MolWt, NumHDonors), отпечатки молекулы (Morgan/ECFP) или их комбинация.

\textbf{RF/градиентный бустинг.} Например, работы \cite{lusci2013,boobier2017,vermeire2021} показывают, что деревья решений остаются сильным ориентиром: типично MAE $\approx 0.8$--$1.3$ по $\ln x_2$ (или сопоставимый уровень на шкале $\log S$ в зависимости от протокола разбиения), при приемлемых $R^2$ на тестах в пределах области обучения. На практике XGBoost/LightGBM нередко дают ещё $0.05$--$0.1$ выигрыша в MAE относительно RF за счёт бустинга.

Отдельно отметим сравнение протоколов: в работе Vermeire et al.~\cite{vermeire2021} на случайном разбиении сообщается уровень около MAE $\approx 0.71$ (в сопоставимой шкале), тогда как при более строгом разбиении по каркасам качество обычно заметно снижается. В текущем отчёте именно разбиение по каркасам рассматривается как основной реалистичный критерий переносимости на новые химические каркасы.

\textbf{Линейные QSPR.} Исторические модели (классического типа, Abraham-подобные формы) интерпретируемы, но по точности обычно уступают нелинейным методам: типично MAE около $1.5$--$2.0$.

\textbf{Плюсы и ограничения класса.} Плюсы: быстрое обучение, низкий вычислительный бюджет, сильное качество на табличных данных, отсутствие необходимости в GPU. Ограничения: слабая физическая декомпозиция предсказания, ограниченная переносимость по температуре вне обучающего диапазона и уязвимость к сдвигу домена признаков.

\textbf{Позиционирование TGNN-Solv относительно класса.} В наших экспериментах RF на RDKit-дескрипторах даёт MAE $\approx 1.20$ (при разбиении по каркасам) и выступает ключевым практическим ориентиром. TGNN-Solv без дескрипторного дополнения пока хуже (MAE $\approx 1.93$), а DirectGNN на той же базовой архитектуре даёт MAE $\approx 1.88$. Линейная диагностическая модель (зонд) показывает медианный $R^2=0.505$ при восстановлении дескрипторов из эмбеддингов, что указывает на информационное бутылочное горлышко энкодера; по декомпозиции ошибок около $93\%$ разрыва с RF объясняется именно энкодером и лишь $7\%$ --- физическим слоем (подробная методика декомпозиции приведена в \hyperref[subsec:error-decomposition]{Части~V, §\ref*{subsec:error-decomposition}}). Следовательно, расширение признаков дескрипторами выглядит не второстепенной опцией, а наиболее обоснованным направлением сокращения разрыва с сильным QSPR-ориентиром при сохранении физической структуры модели.

\subsection{Внешние базовые модели для сравнения: SolProp и FastSolv}

Помимо классических QSPR-базовых ориентиров в отчёте используется отдельный внешний бенчмарк-блок из двух прикладных моделей: \textbf{SolProp} и \textbf{FastSolv}. Их роль --- не заменить внутренние сравнения TGNN/DirectGNN/RF, а дать дополнительную внешнюю точку отсчёта относительно современных подходов на основе предобучения и ансамблей.

\textbf{SolProp.} В рамках этого отчёта SolProp рассматривается как внешняя базовая модель для прогнозирования растворимости и связанных физико-химических характеристик~\cite{solprop2024}. Практически используются три режима:
\begin{enumerate}
  \item \textbf{режим без дообучения:} запуск без адаптации на целевой выборке;
  \item \textbf{калиброванный режим:} посткалибровка предсказаний и неопределённости на валидационном подмножестве;
  \item \textbf{полное переобучение:} полноценное переобучение на целевых данных текущего отчёта.
\end{enumerate}

Содержательно эти режимы отвечают на разные вопросы: режим без дообучения проверяет переносимость «как есть», калиброванный режим --- насколько быстро модель подстраивается без полной переоценки весов, а полное переобучение показывает её потолок качества в тех же данных и при сопоставимом протоколе.

\textbf{FastSolv.} FastSolv относится к семейству быстрых предикторов растворимости на основе предобучения и ансамблирования дескрипторных и нейросетевых компонентов~\cite{osanloo2024}. Сильная сторона подхода --- высокая практическая точность при умеренной стоимости применения модели; ограничения в нашем контуре связаны с тем, что интеграция чувствительна к версиям модельных артефактов, формату признаков и согласованию протокола данных.

\textbf{Почему именно эти две модели взяты во внешний бенчмарк.}
\begin{enumerate}
  \item Они представляют два практически востребованных внешних направления: SolProp как воспроизводимый прикладной конвейер с несколькими режимами адаптации, FastSolv как сильный ориентир на предобучение и ансамблирование.
  \item Они имеют иные индуктивные предпосылки по сравнению с TGNN-Solv (меньше жёстких термодинамических ограничений), что полезно для сопоставления подходов.
  \item Они позволяют отделить вопрос «насколько хороша физически-информированная декомпозиция» от вопроса «насколько силён внешний универсальный предиктор на тех же данных».
\end{enumerate}

Именно \textbf{SolProp и FastSolv} используются как внешние модели сравнения в бенчмарк-контуре этого отчёта и далее явно фигурируют в разделах про протокол сравнения, риски и прикладную интерпретацию результатов.

\subsection{Методы машинного обучения для растворимости: графовые модели}

\textbf{MPNN-парадигма.} Базовые GNN для молекул строятся по схеме обмена сообщениями:
\begin{align*}
\mathbf{m}_{ij}^{(t)} &= \phi\!\left(\mathbf{h}_i^{(t)},\mathbf{h}_j^{(t)},\mathbf{e}_{ij}\right),\\
\mathbf{h}_i^{(t+1)} &= \psi\!\left(\mathbf{h}_i^{(t)},\sum_{j\in\mathcal{N}(i)}\mathbf{m}_{ij}^{(t)}\right),\\
\mathbf{g} &= R\!\left(\{\mathbf{h}_i^{(T)}\}\right).
\end{align*}
Репрезентативный базис --- \cite{gilmer2017}; для D-MPNN см. \cite{yang2019}. Известное ограничение выразительности: стандартные MPNN ограничены мощностью 1-WL-теста.

\textbf{3D-ориентированные семейства.} SchNet использует свёртку с непрерывно параметризуемыми фильтрами \cite{schutt2018}; DimeNet/SphereNet вводят направленные сообщения с углами и торсионами \cite{gasteiger2020,liu2022}. Эти модели обычно точнее на геометрически чувствительных задачах, но тяжелее вычислительно и требовательнее к качественным 3D-конформациям.

\textbf{GNN в задачах растворимости.} Большинство работ предсказывают один скаляр напрямую (logS или близкие прокси), например \cite{francoeur2021,ye2021}. Часто заявляемые уровни MAE на водной растворимости находятся в диапазоне $0.8$--$1.2$ по $\log S$; однако это не эквивалентно температурно-зависимой SLE-постановке по $\ln x_2$. Для FastSolv-подходов (предобучение и ансамблирование на дескрипторах, \cite{osanloo2024}) сообщаются значения около MAE $\approx 0.73$ в сопоставимой постановке задачи, но обычно на иной схеме данных и/или фиксированном наборе растворителей.

\textbf{Позиционирование TGNN-Solv относительно класса.} Наш энкодер относится к стандартному семейству MPNN (6 слоёв), то есть основной вклад работы состоит не в архитектуре графового блока. Ключевое отличие находится после энкодера: блок попарного внимания (cross-attention), выходные блоки физически интерпретируемых параметров и последующее вычисление растворимости через термодинамический решатель. С практической точки зрения результаты DirectGNN (MAE $\approx 1.88$) показывают, что «прямой скалярный» путь для нашего датасета недостаточен; одновременно данные анализа зондом подтверждают необходимость усиления представления (расширение признаков дескрипторами и/или более информативных априорных ограничений), а не отказа от физического слоя.

\subsection{Физически-информированное машинное обучение: PINN, DEQ, Neural ODE и дифференцируемые решатели}

\textbf{PINNs.} В PINNs сеть обучается одновременно по данным и по невязке физического уравнения:
\begin{equation*}
\mathcal{L}=\mathcal{L}_{\text{data}}+\lambda\,\mathcal{L}_{\text{PDE}}.
\end{equation*}
Классические источники --- \cite{raissi2019}. Для задач SLE-предсказания растворимости этот шаблон напрямую применяется редко, но общий принцип «физика как ограничение» здесь напрямую релевантен.

\textbf{Гамильтоновы (structure-preserving) нейронные сети.} В работах типа \cite{greydanus2019} сеть параметризует энергию (гамильтониан), а динамика задаётся жёстко уравнениями. Это близко по философии к схемам, где нейросеть предсказывает параметры, а физический закон выполняет вычисление наблюдаемой величины.

\textbf{Neural ODE и дифференцируемые симуляторы.} Идея «модель + стандартный решатель + дифференцирование через решатель» (например, через сопряжённые или неявные производные) систематизирована в \cite{chen2018} и последующих работах.

\textbf{DEQ-постановки.} В DEQ выход определяется как неподвижная точка
\begin{equation*}
\mathbf{z}^*=f_\theta(\mathbf{z}^*,\mathbf{x}),
\end{equation*}
а градиенты считаются по теореме о неявной функции \cite{bai2019}. Эта математика напрямую переносится на итеративные физические блоки, где состояние определяется самосогласованным решением.

В контексте TGNN-Solv DEQ-аналогия особенно полезна: роль неподвижной точки играет равновесное значение $x_2^*$, а отображение задаётся SLE-оператором. Принципиальное отличие от классического DEQ --- оператор фиксированной точки в TGNN-Solv не содержит обучаемых параметров; обучаемость сосредоточена в предсказании физически интерпретируемых параметров, поступающих в аналитически заданный решатель. Это гарантирует термодинамическую консистентность по построению и снижает риск «скрытой» нефизичности промежуточных переменных.

\textbf{Физически-информированное машинное обучение в термодинамике смесей.} Существуют гибридные направления (например, \cite{lim2022,winter2022,sanchezmedina2023}), где GNN/NN сочетаются с ограничениями термодинамики. Однако в ряде случаев решатель либо однопроходный, либо неполный относительно полной SLE-итерации.

\textbf{Позиционирование TGNN-Solv относительно класса.} TGNN-Solv совмещает пять линий: (i) обучаемый молекулярный энкодер, (ii) жёстко заданное термодинамическое ядро SLE (0 обучаемых параметров в решателе), (iii) имплицитное дифференцирование через постановку в виде неподвижной точки, (iv) априори групповых вкладов для кристаллических параметров, (v) расширение признаков дескрипторами для компенсации ограничений представления. По отношению к ближайшим аналогам важны именно самосогласованная итеративность, ограниченная поправка с повторным решением SLE и фокус на идентифицируемости параметров, а не только на итоговом скалярном MAE.

\subsection{Сводная таблица сравнения}

\begin{table}[H]
\centering
\tiny
\setlength{\tabcolsep}{0.5pt}
\begin{tabularx}{0.93\textwidth}{P{3.0cm}Y C{1.05cm}P{1.8cm}P{2.1cm}}
\hline
Метод (год) & Обучающие данные и учёт температуры & MAE & \shortstack[l]{Интерпретируе-\\мость} & \shortstack[l]{Термодин.\\согласованность} \\
\hline
\shortstack[l]{UNIFAC/mod-\\UNIFAC\\(1975/1987)} & Табличные параметры групп; температура учитывается через явные уравнения модели & 0.8--1.5 & Высокая & Да \\
\shortstack[l]{COSMO-RS/\\COSMO-SAC\\(1995/2002/2017)} & Квантово-химические расчёты и сегментные параметры; температура входит в термодинамический расчёт & 0.5--1.0 & Высокая & Да \\
\shortstack[l]{RF на\\дескрипторах\\(2013--2021)} & Табличные молекулярные дескрипторы; температура подаётся как отдельный входной признак & 0.8--1.3 & Низкая & Нет \\
\shortstack[l]{Direct GNN\\(2021+)} & Графы молекул; температура подаётся как вход сети & 0.8--1.2$^*$ & Низкая & Нет \\
\shortstack[l]{SolvGNN\\(2022)} & Графы и признаки компонентов смеси; температурный учёт зависит от постановки данных & 0.6--1.0$^*$ & Средняя & Частично \\
\textbf{\shortstack[l]{TGNN-Solv target\\(2026)}} & \textbf{Графы + дескрипторы + физические вспомогательные метки; температура учитывается через SLE/NRTL решатель} & \textbf{1.1--1.4}$^{**}$ & \textbf{Высокая} & \textbf{Да} \\
\hline
\end{tabularx}
\caption{$^*$ Для водной растворимости ($\log S$), не для полной SLE-постановки по $\ln x_2$. $^{**}$ Проектный целевой диапазон после расширения признаков и стабилизации обучения; это ориентир, а не уже достигнутый результат текущей версии.}
\end{table}

\subsection{Позиционирование TGNN-Solv в контексте обзора литературы}

Систематический обзор показывает, что существующие семейства методов хорошо закрывают отдельные требования, но редко объединяют их одновременно. Групповые и \textit{ab initio} подходы обеспечивают физическую осмысленность и устойчивую температурную интерполяцию, но ограничены параметрическим покрытием или вычислительной стоимостью. Дескрипторные ML-модели дают сильное качество и низкую цену обучения, но не предоставляют термодинамически интерпретируемую декомпозицию. Direct-GNN расширяют гибкость представлений, однако в режимах с малым объёмом данных часто сталкиваются с проблемой ограниченной информативности латентных признаков и ухудшением out-of-domain поведения.

На этом фоне вклад TGNN-Solv состоит не в утверждении «физика всегда лучше» и не в попытке заменить сильные базовые модели одной новой нейросетью. Центральная идея работы --- инженерно корректная композиция физики и ML с проверяемыми диагностическими критериями. Эмпирически это поддерживается следующими наблюдениями: (i) RF на дескрипторах уже формирует высокий порог качества (MAE $\approx 1.20$), (ii) графовый вариант без физического блока без дескрипторов (MAE $\approx 1.93$) и DirectGNN (MAE $\approx 1.88$) фиксируют разрыв в качестве представлений, (iii) probe-анализ ($R^2=0.505$) и декомпозиция ошибок ($93\%$ энкодер, $7\%$ физический слой; методика --- в \hyperref[subsec:error-decomposition]{Части~V, §\ref*{subsec:error-decomposition}}) указывают, где именно находится bottleneck. Следовательно, наиболее рациональная траектория улучшения --- не отказ от физически-информированной постановки, а усиление входной информативности и структурных априорных ограничений.

С методологической точки зрения работа вносит три исследовательских акцента. Во-первых, проблема идентифицируемости параметров рассматривается как первичная, а не вторичная: many-to-one соответствие «параметры $\rightarrow$ растворимость» требует специальных ограничений и диагностик при обучении. Во-вторых, конструирование функции потерь рассматривается как часть физической постановки (контроль доли растворимости, стабилизация van't Hoff-компоненты, баланс задач), а не как исключительно численная эвристика. В-третьих, расширение признаков дескрипторами и GC-априори трактуются как согласованный мост между классической инженерной химией и современными представлениями GNN.

Итоговое позиционирование относительно конкурирующих классов можно сформулировать так: TGNN-Solv стремится достичь как минимум уровня сильных дескрипторных ориентиров по MAE, одновременно сохраняя то, чего эти ориентиры обычно не дают --- температурную экстраполяцию, явные промежуточные физические величины и термодинамическую консистентность вычислительного ядра. Если конфигурация с расширением признакового пространства дескрипторами достигает MAE не хуже RF при этих дополнительных свойствах, это является практически значимым результатом для проектирования растворителей и для переноса моделей на новые температурные режимы.
\section{Термодинамический фундамент}

В основном тексте оставлены только формулы, необходимые для чтения результатов и архитектуры. Развёрнутые промежуточные выводы и технические доказательства вынесены в приложение C.

\subsection{Химический потенциал}

Химический потенциал компонента $i$ в смеси:

\begin{equation*}
\mu_i \equiv \left(\frac{\partial G}{\partial n_i}\right)_{T,\, P,\, n_{j \neq i}}
\end{equation*}

Для компонента $i$ в жидкой смеси:

\begin{equation*}
\mu_i^{\mathrm{liq}}(T, P, \mathbf{x}) = \mu_i^{\mathrm{liq},\circ}(T, P) + RT \ln(\gamma_i x_i)
\end{equation*}

где $\mu_i^{\mathrm{liq},\circ}$ --- химический потенциал чистой жидкости $i$, $x_i$ --- мольная доля, $\gamma_i$ --- коэффициент активности, $R$ --- универсальная газовая постоянная.

Для чистого кристалла:

\begin{equation*}
\mu_2^{\mathrm{sol}}(T, P) = \mu_2^{\mathrm{sol},\circ}(T, P)
\end{equation*}

\subsection{Условие фазового равновесия твёрдое--жидкое}

При равновесии кристалл $\leftrightarrow$ насыщенный раствор:

\begin{equation*}
\mu_2^{\mathrm{sol},\circ}(T) = \mu_2^{\mathrm{liq},\circ}(T) + RT \ln(\gamma_2 x_2)
\end{equation*}

Определяя свободную энергию Гиббса плавления:

\begin{equation*}
\Delta G_{\mathrm{fus}}(T) \equiv \mu_2^{\mathrm{liq},\circ}(T) - \mu_2^{\mathrm{sol},\circ}(T)
\end{equation*}

получаем мастер-уравнение \eqref{eq:master_sle}:

\begin{equation}
\boxed{\ln x_2 = -\frac{\Delta G_{\mathrm{fus}}(T)}{RT} - \ln \gamma_2}
\label{eq:master_sle}
\end{equation}

Это уравнение точное (единственные допущения: чистая твёрдая фаза, несжимаемость при $P \approx 1$ атм).

\subsection{Структурная декомпозиция}

\begin{equation*}
\ln x_2 = \underbrace{-\frac{\Delta G_{\mathrm{fus}}(T)}{RT}}_{\equiv\; -\Phi(T)\;\text{(идеальный вклад)}} + \underbrace{(-\ln \gamma_2)}_{\text{вклад неидеальности}}
\end{equation*}

Первый член $(-\Phi)$ зависит только от свойств чистого solute ($T_{\mathrm{m}}$, $\Delta H_{\mathrm{fus}}$, $\Delta C_p^{\mathrm{fus}}$) и температуры $T$. Второй член $(-\ln \gamma_2)$ зависит от взаимодействия solute--solvent и от состава раствора $x_2$.

\section{Свободная энергия плавления}

\subsection{Интегрирование Кирхгофа}
Классический вывод через термодинамические тождества и интегрирование по температуре (форма Кирхгофа) приведён в стандартных курсах по молекулярной термодинамике~\cite{prausnitz1998}.

Энтальпия плавления при произвольной температуре:

\begin{equation*}
\Delta H_{\mathrm{fus}}(T) = \Delta H_{\mathrm{fus}}(T_{\mathrm{m}}) + \int_{T_{\mathrm{m}}}^{T} \Delta C_p^{\mathrm{fus}}(T') \, dT'
\end{equation*}

Энтропия плавления:

\begin{equation*}
\Delta S_{\mathrm{fus}}(T) = \frac{\Delta H_{\mathrm{fus}}(T_{\mathrm{m}})}{T_{\mathrm{m}}} + \int_{T_{\mathrm{m}}}^{T} \frac{\Delta C_p^{\mathrm{fus}}(T')}{T'} \, dT'
\end{equation*}

где использовано условие равновесия при $T_{\mathrm{m}}$: $\Delta G_{\mathrm{fus}}(T_{\mathrm{m}}) = 0 \Rightarrow \Delta S_{\mathrm{fus}}(T_{\mathrm{m}}) = \Delta H_{\mathrm{fus}}(T_{\mathrm{m}}) / T_{\mathrm{m}}$.

\subsection{Приближение постоянной dCp плавления}

При $\Delta C_p^{\mathrm{fus}} = \mathrm{const}$:

\begin{equation*}
\Delta G_{\mathrm{fus}}(T) = \Delta H_{\mathrm{fus}}(T_{\mathrm{m}})\!\left(1 - \frac{T}{T_{\mathrm{m}}}\right) + \Delta C_p^{\mathrm{fus}}\!\left[(T - T_{\mathrm{m}}) - T \ln\frac{T}{T_{\mathrm{m}}}\right]
\end{equation*}

Деление на $RT$:

\begin{equation*}
\Phi(T) = \frac{\Delta G_{\mathrm{fus}}(T)}{RT} = \frac{\Delta H_{\mathrm{fus}}(T_{\mathrm{m}})}{R}\!\left(\frac{1}{T} - \frac{1}{T_{\mathrm{m}}}\right) - \frac{\Delta C_p^{\mathrm{fus}}}{R}\!\left[\left(\frac{T_{\mathrm{m}}}{T} - 1\right) - \ln\frac{T_{\mathrm{m}}}{T}\right]
\end{equation*}

\subsection{Анализ поправочного члена}

Определим $\psi(u) \equiv (u - 1) - \ln u$ для $u = T_{\mathrm{m}} / T$.

\textbf{Утверждение.} $\psi(u) \geq 0$ для всех $u > 0$, причём $\psi(u) = 0$ тогда и только тогда, когда $u = 1$.

\textbf{Доказательство.}

\begin{equation*}
\psi'(u) = 1 - \frac{1}{u} = \frac{u - 1}{u}
\end{equation*}

\begin{equation*}
\psi'(u) = 0 \iff u = 1, \quad \psi''(u) = \frac{1}{u^2} > 0
\end{equation*}

Следовательно, $u = 1$ --- единственный минимум, $\psi(1) = 0$, и $\psi$ строго выпукла. \qed

\textbf{Следствие.} При $\Delta C_p^{\mathrm{fus}} > 0$ поправочный член отрицателен:

\begin{equation*}
-\frac{\Delta C_p^{\mathrm{fus}}}{R} \cdot \psi\!\left(\frac{T_{\mathrm{m}}}{T}\right) \leq 0
\end{equation*}

и увеличивает идеальную растворимость. Физический смысл: при бо́льшей теплоёмкости жидкой фазы рост $\Delta G_{\mathrm{fus}}$ при удалении от $T_{\mathrm{m}}$ замедляется.

\subsection{Упрощение Хильдебранда (dCp = 0)}

Это приближение соответствует классической regular-solution постановке~\cite{hildebrand1970}.

\begin{equation*}
\Phi(T) = \frac{\Delta H_{\mathrm{fus}}}{R}\!\left(\frac{1}{T} - \frac{1}{T_{\mathrm{m}}}\right)
\end{equation*}

Двухпараметрическая модель. Текущая имплементация использует это приближение по умолчанию (\texttt{fixed\_dCp\_fus = 0}).

\subsection{Количественная оценка ошибки при dCp = 0}

Для парацетамола ($T_{\mathrm{m}} = 442$ К, $\Delta H_{\mathrm{fus}} = 26400$ Дж/моль, $\Delta C_p^{\mathrm{fus}} \approx 80$ Дж/(моль$\cdot$К)) при $T = 298.15$ К (типичные литературные параметры для расчётов идеальной растворимости~\cite{yalkowsky1979,acree2010}):

\begin{itemize}
  \item С $\Delta C_p$: $x_2^{\mathrm{ideal}} = 0.0735$
  \item Без $\Delta C_p$: $x_2^{\mathrm{ideal}} = 0.0313$
  \item Ошибка в $\ln x_2^{\mathrm{ideal}}$: $\approx 0.85$
\end{itemize}

Эта систематическая ошибка сопоставима с целевой точностью модели (MAE $\sim 0.5$--$1.0$).

\section{Идеальная растворимость}

\subsection{Определение и свойства}

\begin{equation*}
\ln x_2^{\mathrm{ideal}} = -\Phi(T)
\end{equation*}

\textbf{Свойства:}

\begin{enumerate}
  \item При $T = T_{\mathrm{m}}$: $\Phi(T_{\mathrm{m}}) = 0 \Rightarrow x_2^{\mathrm{ideal}} = 1$.
  \item При $T \to 0$: $\Phi \to +\infty \Rightarrow x_2^{\mathrm{ideal}} \to 0$.
  \item Монотонность: $\dfrac{d\ln x_2^{\mathrm{ideal}}}{dT} = \dfrac{\Delta H_{\mathrm{fus}}(T)}{RT^2} > 0$ при $\Delta H_{\mathrm{fus}} > 0$.
\end{enumerate}

Физическая интерпретация предельных режимов: при $T\to T_{\mathrm{m}}$ идеальная растворимость стремится к единице ($x_2^{\mathrm{ideal}}\to 1$), что соответствует практически полной совместимости расплава solute с жидкой фазой вблизи точки плавления. При $T\ll T_{\mathrm{m}}$ идеальная растворимость экспоненциально убывает, отражая термодинамический штраф за плавление кристалла при низкой температуре.

\subsection{Чувствительность}

\begin{equation*}
\frac{\partial \Phi}{\partial T_{\mathrm{m}}} = \frac{\Delta H_{\mathrm{fus}}}{RT_{\mathrm{m}}^2} > 0
\end{equation*}

\begin{equation*}
\frac{\partial \Phi}{\partial \Delta H_{\mathrm{fus}}} = \frac{1}{R}\!\left(\frac{1}{T} - \frac{1}{T_{\mathrm{m}}}\right) > 0 \quad (T < T_{\mathrm{m}})
\end{equation*}

При типичных значениях ($T_{\mathrm{m}} = 400$ К, $\Delta H_{\mathrm{fus}} = 25000$ Дж/моль, $T = 298$ К):

\begin{table}[H]
\centering
\begin{tabular}{|c|c|c|}
\hline
Параметр & Возмущение & $\Delta(\ln x_2)$ \\
\hline
$T_{\mathrm{m}}$ & $\pm 10$ К & $\pm 0.19$ \\
\hline
$\Delta H_{\mathrm{fus}}$ & $\pm 2000$ Дж/моль & $\pm 0.25$ \\
\hline
\end{tabular}
\caption{Чувствительность идеальной растворимости к погрешностям кристаллических параметров при типичных условиях.}
\end{table}

\section{Коэффициент активности и модель NRTL}

\subsection{Определение коэффициента активности}

\begin{equation*}
RT \ln \gamma_i = \left(\frac{\partial (n_{\mathrm{tot}} G^E)}{\partial n_i}\right)_{T, P, n_{j \neq i}}
\end{equation*}

где $G^E$ --- мольная избыточная энергия Гиббса.

Физический смысл: $\gamma_i > 1$ --- молекула $i$ «менее комфортна» в растворе, чем среди себе подобных; $\gamma_i < 1$ --- наоборот.

\subsection{Гипотеза локального состава (NRTL)}

Renon и Prausnitz (1968) постулировали распределение Больцмана для локальных составов~\cite{renon1968}:

\begin{equation*}
\frac{x_{12}}{x_{22}} = \frac{x_1}{x_2} \exp(-\alpha_{12}\tau_{12})
\end{equation*}

где

\begin{equation*}
\tau_{12} = \frac{g_{12} - g_{22}}{RT}, \quad \tau_{21} = \frac{g_{21} - g_{11}}{RT}
\end{equation*}

$g_{ij}$ --- энергия парного взаимодействия, $\alpha_{12} \in [0.1, 0.6]$ --- параметр неслучайности.

\subsection{Избыточная энергия Гиббса NRTL}
\label{subsec:nrtl-ge}

Ниже используется функциональная форма из исходной публикации NRTL и последующих сводок~\cite{renon1968,prausnitz1998}; далее используем выражения \eqref{eq:nrtl_ge} и \eqref{eq:nrtl_lngamma2}.

Обозначим $G_{12} = \exp(-\alpha_{12}\tau_{12})$, $G_{21} = \exp(-\alpha_{12}\tau_{21})$.

\begin{equation}
\frac{G^E}{RT} = x_1 x_2 \left(\frac{\tau_{21} G_{21}}{x_1 + x_2 G_{21}} + \frac{\tau_{12} G_{12}}{x_2 + x_1 G_{12}}\right)
\label{eq:nrtl_ge}
\end{equation}

\subsection{Коэффициент активности растворённого вещества}
\label{subsec:nrtl-gamma2}

Дифференцированием $G^E$ по $n_2$:

\begin{equation}
\ln \gamma_2 = x_1^2 \left[\tau_{12}\!\left(\frac{G_{12}}{x_2 + x_1 G_{12}}\right)^{2} + \frac{\tau_{21} G_{21}}{(x_1 + x_2 G_{21})^2}\right]
\label{eq:nrtl_lngamma2}
\end{equation}

\textbf{Верификация пределов.}

При $x_2 \to 1$ ($x_1 \to 0$):

\begin{equation*}
\ln \gamma_2 \to 0 \cdot [\ldots] = 0 \quad \Rightarrow \quad \gamma_2 \to 1
\end{equation*}

Нормировка Рауля выполняется автоматически.

При $x_2 \to 0$ ($x_1 \to 1$):

\begin{equation*}
\ln \gamma_2^\infty = \tau_{12} + \tau_{21} \exp(-\alpha_{12}\tau_{21})
\end{equation*}

\subsection{Термодинамическая консистентность}

\textbf{Утверждение.} NRTL удовлетворяет уравнению Гиббса--Дюгема по построению~\cite{prausnitz1998}:

\begin{equation*}
x_1 \frac{\partial \ln \gamma_1}{\partial x_1} + x_2 \frac{\partial \ln \gamma_2}{\partial x_1} = 0 \quad \forall \, x_1 \in [0, 1]
\end{equation*}

\textbf{Доказательство.} $\ln \gamma_1$ и $\ln \gamma_2$ получены дифференцированием одной и той же функции $G^E(x_1, x_2)$. По теореме Эйлера для однородных функций первого порядка:

\begin{equation*}
n_{\mathrm{tot}} G^E = n_1 \bar{G}_1^E + n_2 \bar{G}_2^E
\end{equation*}

где $\bar{G}_i^E = RT \ln \gamma_i$ --- парциальные мольные избыточные энергии. Дифференцируя тождество по $x_1$ при фиксированных $T$, $P$:

\begin{equation*}
\frac{dG^E}{dx_1} = \bar{G}_1^E - \bar{G}_2^E + x_1 \frac{d\bar{G}_1^E}{dx_1} + x_2 \frac{d\bar{G}_2^E}{dx_1}
\end{equation*}

С другой стороны, из определения парциальных мольных величин и тождества Гиббса--Дюгема:

\begin{equation*}
x_1 \, d\bar{G}_1^E + x_2 \, d\bar{G}_2^E = 0
\end{equation*}

Деление на $RT$ даёт требуемое. \qed

\subsection{Число параметров NRTL и их диапазоны}

\begin{table}[H]
\centering
\begin{tabularx}{0.99\textwidth}{|P{1.9cm}|P{3.1cm}|Y|}
\hline
Параметр & Диапазон & Физический смысл \\
\hline
$\Delta g_{12}$ & $[-20000, +30000]$ Дж/моль & Энергетическая разница контакта 1--2 относительно 2--2 \\
\hline
$\Delta g_{21}$ & $[-20000, +30000]$ Дж/моль & Энергетическая разница контакта 2--1 относительно 1--1 \\
\hline
$\alpha_{12}$ & $[0.1, 0.6]$ & Неслучайность расположения \\
\hline
\end{tabularx}
\caption{Типичные диапазоны параметров NRTL в бинарной постановке.}
\end{table}

Чувствительность $\ln x_2$ к параметрам NRTL (при типичных условиях):

\begin{equation*}
\frac{\partial \ln x_2}{\partial \Delta g_{12}} \sim 4 \times 10^{-4} \; (\text{Дж/моль})^{-1}
\end{equation*}

Ошибка $\Delta g_{12}$ на $\pm 2000$ Дж/моль $\Rightarrow$ ошибка $\ln x_2 \sim \pm 0.8$.

Это наиболее чувствительный параметр модели --- ключевое наблюдение для диагностики.


\subsection{Модель NRTL: дополнительные аспекты}

\subsubsection{Физический смысл параметра неслучайности $\alpha$}

Параметр неслучайности $\alpha_{12}$ в NRTL задаёт степень отклонения локального окружения молекулы от случайного смешения~\cite{renon1968,prausnitz1998}. В пределе $\alpha_{12}=0$ модель переходит к полностью случайному распределению локальных составов, что по смыслу близко к regular-solution логике. При росте $\alpha_{12}$ (типично до $\sim 0.5$) усиливается локальная неоднородность: одинаковые и «предпочтительные» межмолекулярные контакты получают больший вес.

На уровне кривых активности это проявляется как изменение нелинейности $\gamma_i(x)$: при большем $\alpha$ кривая обычно становится более «изогнутой», а отклонения от идеальности концентрируются в определённых диапазонах состава. В инженерной практике часто используют эмпирические значения:
\begin{itemize}
  \item $\alpha \approx 0.20$ для насыщенных углеводородных систем,
  \item $\alpha \approx 0.30$ для полярных неводных систем,
  \item $\alpha \approx 0.47$ для водных систем.
\end{itemize}

Эти числа не являются универсальными константами и должны рассматриваться как информативные priors, особенно при ограниченном объёме данных.

\subsubsection{Связь NRTL с теорией решёток}

Исторически NRTL восходит к двухжидкостной идее Scott и квазихимическим представлениям Guggenheim о локальном порядке~\cite{scott1956,guggenheim1952}. При этом важно подчеркнуть, что NRTL --- феноменологическая модель: она мотивирована статистико-термодинамическими аргументами, но не является строгим выводом из единственного микроскопического гамильтониана.

Практическое следствие такого статуса двоякое. С одной стороны, модель достаточно гибка и хорошо поддаётся идентификации на данных VLE/LLE/SLE. С другой стороны, вне области параметризации возможны нефизичные экстраполяции, и параметры следует интерпретировать как эффективные, а не фундаментальные межмолекулярные константы.

\subsubsection{Температурная зависимость параметров}

Для температурной зависимости коэффициентов $\tau_{ij}(T)$ в литературе используются несколько уровней сложности~\cite{renon1968,weidlich1987,prausnitz1998}.

\textbf{(a) Простая форма:}
\begin{equation*}
\tau_{ij}(T)=\frac{\Delta g_{ij}}{RT},
\end{equation*}
где $\Delta g_{ij}$ имеет размерность $\mathrm{J/mol}$.

\textbf{(b) Расширенная корреляция:}
\begin{equation*}
\tau_{ij}(T)=a_{ij}+\frac{b_{ij}}{T}+c_{ij}\ln T.
\end{equation*}

\textbf{(c) Форма ref\_invT (используется в TGNN-Solv):}
\begin{equation*}
\tau_{ij}(T)=\tau_{ij}(T_{\mathrm{ref}})+\beta_{ij}\left(\frac{1}{T}-\frac{1}{T_{\mathrm{ref}}}\right),
\end{equation*}
где $T$ и $T_{\mathrm{ref}}$ измеряются в $\mathrm{K}$.

Последняя форма удобна для обучения: параметр $\tau(T_{\mathrm{ref}})$ задаёт «якорное» значение, а $\beta_{ij}$ контролирует температурный наклон. В термодинамической интерпретации $\beta_{ij}$ связан с эффективной энтальпийной компонентой смешения и по порядку величины может рассматриваться как
\begin{equation*}
\beta_{ij}\sim \frac{\Delta H_{\mathrm{mix},ij}}{R}.
\end{equation*}

Практически $\beta_{ij}$ в форме \texttt{ref\_invT} имеет размерность K и интерпретируется как эффективный энтальпийный масштаб температурного отклика. Для инженерно устойчивых режимов обычно наблюдаются $|\beta_{ij}|\in[10^2,2\times10^3]$ K. При $\beta_{ij}>0$ вклад взаимодействия с ростом температуры обычно ослабевает (положительный энтальпийный вклад), при $\beta_{ij}<0$ --- усиливается.

\subsubsection{Пределы применимости NRTL}

Несмотря на универсальность, NRTL имеет важные ограничения.
\begin{enumerate}
  \item Для сильно ассоциирующих систем (например, вода--спирт при низких $x$) без специальных поправок модель может систематически ошибаться.
  \item При больших значениях эффективной несовместимости (крупные $\tau_{12}$ и $\tau_{21}$) модель предсказывает жидкостное расслаивание (LLE).
  \item Граница расслаивания зависит от $\alpha$; в инженерных терминах условие часто формулируют как
  \begin{equation*}
  \tau_{12}+\tau_{21} > f(\alpha),
  \end{equation*}
  где функция $f(\alpha)$ определяется анализом устойчивости энергии Гиббса смешения.
\end{enumerate}

Для SLE-задач проекта это означает, что NRTL корректно работает как базовая модель неидеальности, но при данных с выраженной ассоциацией или мультифазностью полезны дополнительные диагностические проверки.


\subsection{Сравнение моделей $G^E$: NRTL, Wilson, UNIQUAC}

\subsubsection{Унифицированный формализм}

Все рассматриваемые модели (NRTL, Wilson, UNIQUAC) используют одну и ту же термодинамическую стратегию: сначала постулируется аналитическая форма избыточной энергии Гиббса $G^E$, после чего коэффициенты активности получаются дифференцированием по числу молей компонента.

Для компонента $i$:
\begin{equation*}
\ln\gamma_i = \frac{1}{RT}\left(\frac{\partial(n_{\text{tot}}G^E)}{\partial n_i}\right)_{T,P,n_{j\neq i}}.
\end{equation*}

Такое построение автоматически обеспечивает термодинамическую консистентность и выполнение уравнения Гиббса--Дюгема на уровне модели.

\subsubsection{NRTL (Renon--Prausnitz)}

Как уже отмечалось в §§\ref{subsec:nrtl-ge}--\ref{subsec:nrtl-gamma2}, в бинарной постановке NRTL использует три настраиваемых параметра:
\begin{equation*}
\{\Delta g_{12},\; \Delta g_{21},\; \alpha_{12}\}.
\end{equation*}

Явная функциональная форма $G^E$ и выражение для $\ln\gamma_2$ уже заданы выше (уравнения \eqref{eq:nrtl_ge} и \eqref{eq:nrtl_lngamma2}); здесь повторяются только практические следствия для сравнения с Wilson/UNIQUAC.

\textbf{Плюсы:} высокая гибкость, возможность описывать LLE и асимметрию взаимодействий.

\textbf{Минусы:} дополнительный параметр $\alpha$ повышает риск слабой идентифицируемости; его физический смысл более косвенный по сравнению с энергетическими параметрами.

\subsubsection{Wilson}

Wilson-модель использует два параметра взаимодействия:
\begin{equation*}
\{\Lambda_{12},\; \Lambda_{21}\}.
\end{equation*}

Форма $G^E$:
\begin{equation*}
\frac{G^E}{RT} = -x_1\ln(x_1+\Lambda_{12}x_2) - x_2\ln(x_2+\Lambda_{21}x_1).
\end{equation*}

Связь с молярными объёмами и локальными энергетическими параметрами:
\begin{equation*}
\Lambda_{ij}=\frac{V_j}{V_i}\exp\!\left(-\frac{\lambda_{ij}-\lambda_{ii}}{RT}\right).
\end{equation*}

\textbf{Плюсы:} только два параметра; размерный эффект частично учитывается через отношение $V_j/V_i$.

\textbf{Минусы:} конструктивно не описывает расслаивание жидкой фазы (LLE), поскольку модельная форма даёт ограниченный класс поверхностей смешения.

\subsubsection{UNIQUAC (модель Абрамса--Праусница)}

UNIQUAC также использует два обучаемых энергетических параметра в residual-части:
\begin{equation*}
\{\Delta u_{12},\; \Delta u_{21}\}.
\end{equation*}

Коэффициент активности разлагается на комбинаторный и остаточный вклады:
\begin{equation*}
\ln\gamma_i = \underbrace{\ln\gamma_i^{\text{comb}}(r,q,x)}_{\text{без обучаемых параметров}}
+ \underbrace{\ln\gamma_i^{\text{res}}(\Theta,\tau)}_{\text{2 обучаемых параметра}}.
\end{equation*}

Комбинаторная часть задана аналитически:
\begin{equation*}
\ln\gamma_i^{\text{comb}} = \ln\frac{\Phi_i}{x_i}
+ \frac{z}{2}q_i\ln\frac{\Theta_i}{\Phi_i} + l_i
- \frac{\Phi_i}{x_i}\sum_j x_jl_j.
\end{equation*}

Структурные параметры размера/поверхности вычисляются из группового разбиения:
\begin{equation*}
r_i = \sum_k \nu_k^{(i)}R_k, \qquad q_i = \sum_k \nu_k^{(i)}Q_k.
\end{equation*}

\textbf{Плюсы:} явный учёт размера и формы молекул; всего два обучаемых параметра; возможность описания LLE при корректной параметризации residual-вклада.

\textbf{Минусы:} требуется качественная библиотека групповых параметров $R_k,Q_k$ и устойчивое SMARTS-разбиение.

\subsubsection{Сводное сравнение}

\begin{table}[H]
\centering
\begin{tabular}{lccccc}
\hline
Модель & Params & LLE & Размер/форма & Hardcoded & Identifiability \\
\hline
NRTL & 3 & Да & Нет & $G^E$-формула & Слабая \\
Wilson & 2 & Нет & Через $V_i$ & $G^E$-формула & Лучше \\
UNIQUAC & 2 & Да & Через $r,q$ & $G^E$ + comb & Лучшая \\
\hline
\end{tabular}
\caption{Сравнение моделей избыточной энергии $G^E$ для бинарной постановки.}
\end{table}

С точки зрения обратной задачи TGNN-Solv уменьшение числа обучаемых параметров с 3 до 2 и перенос части структуры в аналитически заданный комбинаторный вклад обычно улучшает идентифицируемость.

\subsubsection{Предел бесконечного разбавления}

Для связи с измерениями бесконечного разбавления используются аналитические выражения $\ln\gamma_2^{\infty}$.

\textbf{NRTL:}
\begin{equation*}
\ln\gamma_2^{\infty} = \tau_{12} + \tau_{21}\exp(-\alpha_{12}\tau_{21}).
\end{equation*}

Здесь принципиально важно различать две связанные, но не тождественные величины. Первая --- $\ln\gamma_2^{\infty}$, то есть коэффициент активности при \textbf{бесконечном разбавлении}. Именно эта величина измеряется в IDAC-экспериментах, извлекается из ThermoML и используется как \textbf{наблюдаемая auxiliary-метка}. Вторая --- solver-side величина $\ln\gamma_2(x_2)$, вычисляемая внутри SLE-решателя по формуле NRTL (уравнение~\eqref{eq:nrtl_lngamma2}) при \textbf{конечном составе} насыщенного раствора. При $x_2 \to 0$ эта функция аналитически стремится к $\ln\gamma_2^{\infty}$, но при произвольном $x_2$ она описывает уже другую термодинамическую ситуацию.

Отсюда следует и корректная интерпретация supervision. Потеря $\mathcal{L}_{\gamma^\infty}$ супервизирует не значение $\ln\gamma_2(x_2)$ внутри текущей итерации SLE-решателя, а именно предел $x_2\to 0$ предсказанной NRTL-модели. Тем самым она напрямую ограничивает параметры $\tau_{12}$, $\tau_{21}$ и $\alpha_{12}$ в физически правдоподобной области. В практическом смысле это единственная \textbf{прямая} супервизия на параметры NRTL; все остальные сигналы для них поступают лишь косвенно через $\mathcal{L}_{\mathrm{sol}}$ и поведение SLE-решателя.

\textbf{Wilson:}
\begin{equation*}
\ln\gamma_2^{\infty} = -\ln\Lambda_{21} + 1 - \Lambda_{12}.
\end{equation*}

\textbf{UNIQUAC:}
\begin{equation*}
\ln\gamma_2^{\infty} = \ln\frac{r_2}{r_1} + \frac{z}{2}q_2\ln\frac{q_2r_1}{q_1r_2}
+ l_2 - \frac{r_2}{r_1}l_1 + q_2\left(1-\ln\tau_{12}-\tau_{21}\right).
\end{equation*}

Все три формы могут использоваться в лоссе
\begin{equation*}
\mathcal{L}_{\gamma^{\infty}},
\end{equation*}
если доступны экспериментальные IDAC-наблюдения.

\warnnote{В текущей IDAC-augmented конфигурации $\mathcal{L}_{\gamma^\infty}$ является активным каналом supervision. При этом она относится именно к пределу $\ln\gamma_2^{\infty}$, а не к solver-side величине $\ln\gamma_2(x_2)$ при конечном составе.}

\subsubsection{Выбор для TGNN-Solv}

С инженерной точки зрения три режима имеют разный приоритет:
\begin{itemize}
  \item \textbf{NRTL} --- текущий по умолчанию (максимальная гибкость и привычный базовый ориентир).
  \item \textbf{Wilson} --- рациональный вариант для систем без LLE, особенно для многих фармацевтических сценариев растворимости.
  \item \textbf{UNIQUAC} --- наиболее физически-нагруженный вариант с потенциально лучшей идентифицируемостью за счёт аналитически заданной комбинаторики; это требуется проверить отдельными экспериментами на данных проекта.
\end{itemize}


\subsection{Методы групповых вкладов для $T_{\mathrm{m}}$ и $\Delta H_{\mathrm{fus}}$}
\label{subsec:gc-methods}

В архитектурной модификации TGNN-Solv групповые (GC) оценки используются как \textit{prior}-уровень для кристаллических параметров, а нейросеть предсказывает только ограниченную поправку (bounded residual). Такой дизайн снижает размерность задачи и стабилизирует обучение на малых данных.

\subsubsection{Метод Joback для температуры плавления}

В базовом варианте оценка температуры плавления строится как сумма вкладов функциональных групп:
\begin{equation*}
T_{\mathrm{m}}^{\mathrm{GC}} = 122.5 + \sum_k n_k\,\Delta T_{\mathrm{m},k}, \qquad [T_{\mathrm{m}}]=\mathrm{K},
\end{equation*}
где $n_k$ --- число групп типа $k$, а $\Delta T_{\mathrm{m},k}$ --- соответствующий инкремент~\cite{joback1987,marrero2001}.

Подробный словарь Joback-инкрементов и таблица SMARTS-шаблонов вынесены в Приложение A.1, чтобы основной текст содержал только методологически необходимые формулы.

Численные значения в таблице следует трактовать как инженерные ориентиры (типичные величины), поскольку конкретные библиотеки инкрементов могут отличаться. Для органических соединений типичная точность Joback-оценки составляет MAE порядка $30$--$50\,\mathrm{K}$. В основном тексте используется укрупнённая (округлённая) иллюстративная версия, тогда как рабочий расширенный словарь для воспроизведения вынесен в Приложение A.1.

Ключевые ограничения: отсутствие явного учёта конформаций, полиморфизма и специфических сетей водородных связей.

\textbf{Пример (парацетамол).} Для используемой GC-параметризации получаем $T_{\mathrm{m}}^{\mathrm{GC}}\approx460\,\mathrm{K}$ при экспериментальном $T_{\mathrm{m}}^{\mathrm{exp}}=442\,\mathrm{K}$; абсолютная ошибка $\approx18\,\mathrm{K}$, что укладывается в типичный порядок для GC-оценок.

\subsubsection{Методы для энтальпии плавления}

Для $\Delta H_{\mathrm{fus}}$ применяются корреляции типа Ducros (group-contribution для энтальпии плавления) и эвристические термодинамические соотношения. Часто используется связь через правило Уолдена:
\begin{equation*}
\Delta S_{\mathrm{fus}} \approx 56.5\,\mathrm{J/(mol\cdot K)}
\quad \Rightarrow \quad
\Delta H_{\mathrm{fus}} \approx 56.5\,T_{\mathrm{m}},
\end{equation*}
где $T_{\mathrm{m}}$ задаётся в $\mathrm{K}$, а $\Delta H_{\mathrm{fus}}$ получается в $\mathrm{J/mol}$.

На практике типичная ошибка для GC/корреляционных оценок $\Delta H_{\mathrm{fus}}$ составляет порядка $20$--$30\%$ от истинного значения, особенно для систем с выраженной ассоциацией и/или полиморфизмом.

\subsubsection{Метод для $\Delta C_p^{\mathrm{fus}}$}

Для разности теплоёмкостей плавления $\Delta C_p^{\mathrm{fus}}$ доступно существенно меньше надёжных данных, чем для $T_{\mathrm{m}}$ и $\Delta H_{\mathrm{fus}}$. Поэтому групповые оценки имеют более высокую неопределённость; типичный порядок относительной ошибки составляет около $30$--$50\%$.

В условиях дефицита данных разумной альтернативой остаётся упрощение Хильдебранда:
\begin{equation*}
\Delta C_p^{\mathrm{fus}} \approx 0,
\end{equation*}
которое уже обсуждалось выше и может использоваться как стабильный базовый ориентир.

\subsubsection{Использование в TGNN-Solv}

В TGNN-Solv GC-оценки вводятся как prior, а нейросеть предсказывает ограниченный остаток:
\begin{equation*}
T_{\mathrm{m}}^{\mathrm{pred}} = T_{\mathrm{m}}^{\mathrm{GC,cal}} + \Delta_{\max}\tanh(h),
\end{equation*}
где $h$ --- выход нейросети до ограничения, $\Delta_{\max}$ --- максимальная допустимая поправка.

Перед использованием prior калибруется на обучающем разбиении:
\begin{equation*}
T_{\mathrm{m}}^{\mathrm{GC,cal}} = a\,T_{\mathrm{m}}^{\mathrm{GC}} + b.
\end{equation*}

Калибровка необходима, потому что исходные GC-оценки систематически смещены по химическим классам. На практике используются bounded-диапазоны поправок:
\begin{itemize}
  \item для $T_{\mathrm{m}}$: примерно $\pm50\,\mathrm{K}$,
  \item для $\Delta H_{\mathrm{fus}}$: мультипликативный коридор порядка $\times0.3$--$\times3.0$ относительно prior.
\end{itemize}

Главное преимущество: даже при слабом обучении модель выдаёт физически правдоподобную начальную оценку; оптимизатору нужно выучить только корректирующий residual, а не всё отображение «с нуля».

\subsubsection{Резервная стратегия}

Если молекула не раскладывается на известные группы текущего словаря, используется fallback:
\begin{equation*}
T_{\mathrm{m}}^{\mathrm{GC}} = 400\,\mathrm{K}.
\end{equation*}

В типичных датасетах доля таких случаев составляет порядка $5$--$10\%$ (зависит от химического разнообразия и полноты библиотеки групп). Для этих молекул bounded residual особенно важен, поскольку именно он компенсирует грубость prior-оценки.


\section*{Часть II. Итеративный решатель SLE и неявное дифференцирование}
\addcontentsline{toc}{section}{Часть II. Итеративный решатель SLE и неявное дифференцирование}

\section{Задача о неподвижной точке}

\subsection{Почему нужен итеративный решатель}

Мастер-уравнение \eqref{eq:master_sle} в эквивалентной форме:

\begin{equation}
x_2 = \exp\!\left(-\Phi(T) - \ln \gamma_2(x_2)\right) = \frac{x_2^{\mathrm{ideal}}}{\gamma_2(x_2)}
\label{eq:fixed_point}
\end{equation}

Коэффициент активности $\gamma_2$ сам зависит от $x_2$ через NRTL. Подстановка даёт трансцендентное уравнение:

\begin{equation*}
\ln x_2 = -\Phi - (1 - x_2)^2 \left[\tau_{12}\!\left(\frac{G_{12}}{x_2 + (1 - x_2)G_{12}}\right)^{2} + \frac{\tau_{21}G_{21}}{((1 - x_2) + x_2 G_{21})^2}\right]
\end{equation*}

Это уравнение содержит $\ln x_2$, $x_2$, $(1 - x_2)^2$ и рациональные дроби от $x_2$. Оно не имеет аналитического решения в замкнутой форме.

\subsection{Формулировка как задача о неподвижной точке}

Определяем отображение на основе \eqref{eq:fixed_point}:

\begin{equation*}
g(x_2) \equiv \frac{x_2^{\mathrm{ideal}}}{\gamma_2(x_2)} = \exp\!\left(-\Phi - \ln \gamma_2(x_2)\right)
\end{equation*}

Решение $x_2^{*}$ удовлетворяет $x_2^{*} = g(x_2^{*})$.

Эквивалентная формулировка через корень \eqref{eq:root_equation}:

\begin{equation}
F(x_2) \equiv \ln x_2 + \Phi + \ln \gamma_2(x_2) = 0
\label{eq:root_equation}
\end{equation}

\subsection{Существование и единственность корня}

\textbf{Теорема.} Если бинарная система термодинамически стабильна (не расслаивается) на всём диапазоне составов, уравнение $F(x_2) = 0$ имеет ровно один корень на $(0, 1)$.

\textbf{Доказательство.}

\textit{Существование} следует из теоремы о промежуточном значении:

\begin{equation*}
\lim_{x_2 \to 0^+} F(x_2) = -\infty + \Phi + \ln \gamma_2^\infty = -\infty
\end{equation*}

\begin{equation*}
F(1) = 0 + \Phi + 0 = \Phi > 0 \quad (T < T_{\mathrm{m}})
\end{equation*}

Функция $F$ непрерывна на $(0, 1]$, $F(0^+) = -\infty < 0$, $F(1) > 0$, следовательно $\exists \, x_2^{*} \in (0, 1)$: $F(x_2^{*}) = 0$.

\textit{Единственность.} Вычислим производную:

\begin{equation*}
F'(x_2) = \frac{1}{x_2} + \frac{\partial \ln \gamma_2}{\partial x_2}
\end{equation*}

Условие термодинамической стабильности однофазной смеси:

\begin{equation*}
\frac{\partial^2 G^{\mathrm{mix}}}{\partial x_2^2}\bigg|_{T,P} > 0
\end{equation*}

Раскрывая $G^{\mathrm{mix}} = RT(x_1 \ln x_1 + x_2 \ln x_2) + G^E$:

\begin{equation*}
\frac{\partial^2 G^{\mathrm{mix}}}{\partial x_2^2} = RT\!\left(\frac{1}{x_1} + \frac{1}{x_2}\right) + \frac{\partial^2 G^E}{\partial x_2^2}
\end{equation*}

Строгий вывод связи $\partial\ln\gamma_2/\partial x_2$ с кривизной $G^E$ приведён в приложении~C; здесь используется стандартный аргумент устойчивости однофазной смеси. В устойчивой области свободная энергия смешения строго выпукла, что исключает смену знака производной $F'(x_2)$ на $(0,1)$.

Следовательно, $F$ монотонна на интервале и, имея уже установленное существование корня, допускает на $(0,1)$ ровно один корень. \qed

\section{Метод последовательных подстановок}

\subsection{Алгоритм}

\textbf{Инициализация:} $x_2^{(0)} = x_2^{\mathrm{ideal}} = \exp(-\Phi)$.

\textbf{Итерация} $k = 0, 1, 2, \ldots$:

\begin{equation*}
x_1^{(k)} = 1 - x_2^{(k)}
\end{equation*}

\begin{equation*}
\ln \gamma_2^{(k)} = (x_1^{(k)})^2 \left[\tau_{12}\!\left(\frac{G_{12}}{x_2^{(k)} + x_1^{(k)} G_{12}}\right)^{2} + \frac{\tau_{21}G_{21}}{(x_1^{(k)} + x_2^{(k)} G_{21})^2}\right]
\end{equation*}

\begin{equation*}
x_2^{(k+1)} = \exp\!\left(-\Phi - \ln \gamma_2^{(k)}\right)
\end{equation*}

\textbf{Остановка:} $|x_2^{(k+1)} - x_2^{(k)}| < \varepsilon$ или $k = k_{\max}$.

\subsection{Обоснование начального приближения}

При $\gamma_2 = 1$ (идеальный раствор): $x_2^{(0)} = x_2^{\mathrm{ideal}}$ --- точное решение. Для умеренно неидеальных систем ($\gamma_2 \sim 1$--$10$) это приближение находится в бассейне притяжения неподвижной точки.

Альтернатива для сильно неидеальных систем --- приближение бесконечного разбавления:

\begin{equation*}
x_2^{(0)} = \exp\!\left(-\Phi - \ln \gamma_2^\infty\right), \quad \ln \gamma_2^\infty = \tau_{12} + \tau_{21}\exp(-\alpha_{12}\tau_{21})
\end{equation*}

\subsection{Демпфированная версия}

\begin{equation*}
x_2^{(k+1)} = \lambda \cdot g(x_2^{(k)}) + (1 - \lambda) \cdot x_2^{(k)}, \quad \lambda \in (0, 1]
\end{equation*}

\section{Анализ сходимости}

\subsection{Теорема Банаха о сжимающем отображении}

\textbf{Теорема.} Пусть $g: [a, b] \to [a, b]$ --- сжимающее отображение с коэффициентом $L < 1$:

\begin{equation*}
|g(x) - g(y)| \leq L |x - y| \quad \forall \, x, y \in [a, b]
\end{equation*}

Тогда:
\begin{enumerate}
  \item Существует единственная неподвижная точка $x^{*} = g(x^{*})$.
  \item Последовательность $x_{k+1} = g(x_k)$ сходится к $x^{*}$ из любого $x_0 \in [a, b]$.
  \item Оценка скорости: $|x_k - x^{*}| \leq \dfrac{L^k}{1 - L}|x_1 - x_0|$.
\end{enumerate}

\subsection{Достаточное условие через производную}

Если $g$ дифференцируема и $|g'(x)| \leq L < 1$ на $[a, b]$, то $g$ --- сжатие (по теореме о среднем значении).

\subsection{Производная $g$ для SLE}

\begin{equation*}
g(x_2) = \frac{x_2^{\mathrm{ideal}}}{\gamma_2(x_2)}
\end{equation*}

\begin{equation*}
g'(x_2) = -\frac{x_2^{\mathrm{ideal}}}{\gamma_2(x_2)} \cdot \frac{\partial \ln \gamma_2}{\partial x_2} = -g(x_2) \cdot \frac{\partial \ln \gamma_2}{\partial x_2}
\end{equation*}

В неподвижной точке $g(x_2^{*}) = x_2^{*}$:

\begin{equation*}
\boxed{g'(x_2^{*}) = -x_2^{*} \cdot \eta, \quad \eta \equiv \frac{\partial \ln \gamma_2}{\partial x_2}\bigg|_{x_2^{*}}}
\end{equation*}

\subsection{Знак $\eta$ и условие сходимости}

Для стабильных однофазных систем $\ln \gamma_2$ убывает с ростом $x_2$ (от $\ln \gamma_2^\infty > 0$ при $x_2 = 0$ к $0$ при $x_2 = 1$), поэтому $\eta < 0$.

Тогда $g'(x_2^{*}) = -x_2^{*} \cdot \eta > 0$, и условие сходимости:

\begin{equation*}
|g'(x_2^{*})| = x_2^{*} \cdot |\eta| < 1
\end{equation*}

\textbf{Ключевое следствие.} При малой растворимости ($x_2^{*} \ll 1$) множитель $x_2^{*}$ подавляет $|\eta|$, и сходимость практически гарантирована. Проблемы возникают при $x_2^{*} \gtrsim 0.3$--$0.5$ (высокая растворимость).

\subsection{Вычисление $\eta$}

Обозначим $A = x_2 + x_1 G_{12}$, $B = x_1 + x_2 G_{21}$:

\begin{equation*}
\eta = -\frac{2 \ln \gamma_2}{x_1} + x_1^2 \left[-\frac{2\tau_{12}G_{12}^2(1 - G_{12})}{A^3} + \frac{2\tau_{21}G_{21}(1 - G_{21})}{B^3}\right]
\end{equation*}

\subsection{Скорость сходимости}

Линейная сходимость с коэффициентом $L = |g'(x_2^{*})|$:

\begin{equation*}
|x_2^{(k)} - x_2^{*}| \leq L^k |x_2^{(0)} - x_2^{*}|
\end{equation*}

Число итераций до точности $\varepsilon$:

\begin{equation*}
k \geq \frac{\ln(\varepsilon / |x_2^{(0)} - x_2^{*}|)}{\ln L}
\end{equation*}

\textbf{Пример:} $L = 0.3$, $|x_2^{(0)} - x_2^{*}| = 0.1$, $\varepsilon = 10^{-8}$:

\begin{equation*}
k \geq \frac{\ln(10^{-7})}{\ln 0.3} = \frac{-16.1}{-1.204} \approx 14 \;\text{итераций}
\end{equation*}

\subsection{Демпфирование: формальное обоснование}

Для демпфированного отображения $\tilde{g}(x_2) = \lambda \, g(x_2) + (1 - \lambda) x_2$:

\begin{equation*}
\tilde{g}'(x_2^{*}) = \lambda \, g'(x_2^{*}) + 1 - \lambda
\end{equation*}

Неподвижные точки $g$ и $\tilde{g}$ совпадают.

\textbf{Утверждение.} Для любого $g'(x_2^{*}) \neq 1$ существует $\lambda \in (0, 1]$ такое, что $|\tilde{g}'(x_2^{*})| < 1$.

\textbf{Доказательство.}

Если $g'(x_2^{*}) < 1$: выбираем $\lambda = 1$ (без демпфирования), $|\tilde{g}'| = |g'| < 1$. $\checkmark$

Если $g'(x_2^{*}) > 1$: $\tilde{g}' = \lambda g' + 1 - \lambda = 1 + \lambda(g' - 1)$. При $\lambda < \frac{2}{g' + 1}$ имеем $|\tilde{g}'| < 1$. Такое $\lambda$ существует для конечного $g'$. $\checkmark$

Если $g'(x_2^{*}) < -1$ (осциллирующая расходимость): $\tilde{g}' = 1 + \lambda(g' - 1) = 1 - \lambda(|g'| + 1)$. При $\lambda < \frac{2}{|g'| + 1}$ имеем $|\tilde{g}'| < 1$. $\checkmark$

Единственное исключение: $g'(x_2^{*}) = 1$ --- точка бифуркации (фазовый переход). \qed

\textbf{Оптимальное демпфирование.} Минимум $|\tilde{g}'|$ при $\tilde{g}' = 0$:

\begin{equation*}
\lambda_{\mathrm{opt}} = \frac{1}{1 - g'(x_2^{*})}
\end{equation*}

На практике $g'(x_2^{*})$ неизвестно a priori; используется адаптивная стратегия или фиксированное $\lambda = 0.5$--$0.7$.

\section{Дифференцирование через итеративный решатель}

\subsection{Проблема}

Для обучения GNN нужен градиент $\partial \mathcal{L} / \partial \boldsymbol{\theta}$, где $\boldsymbol{\theta}$ --- предсказанные физические параметры ($\Phi$, $\tau_{12}$, $\tau_{21}$, $G_{12}$, $G_{21}$), а $\mathcal{L}$ --- функция потерь от $\ln x_2^{*}(\boldsymbol{\theta})$.

Два подхода: (A) развёрнутое обратное распространение через итерации; (B) неявное дифференцирование через теорему о неявной функции.

\subsection{Развёрнутое обратное распространение}

Развёрнутые $N$ итераций образуют вычислительный граф:

\begin{equation*}
x_2^{(0)}(\boldsymbol{\theta}) \;\to\; x_2^{(1)}(\boldsymbol{\theta}) \;\to\; \cdots \;\to\; x_2^{(N)}(\boldsymbol{\theta}) \;\to\; \mathcal{L}
\end{equation*}

Цепное правило:

\begin{equation*}
\frac{\partial x_2^{(N)}}{\partial \boldsymbol{\theta}} = \prod_{k=0}^{N-1} \tilde{g}'(x_2^{(k)}) \cdot \frac{\partial x_2^{(0)}}{\partial \boldsymbol{\theta}} + \sum_{j=0}^{N-1} \prod_{k=j+1}^{N-1} \tilde{g}'(x_2^{(k)}) \cdot \frac{\partial \tilde{g}}{\partial \boldsymbol{\theta}}\bigg|_{x_2^{(j)}}
\end{equation*}

\textbf{Проблемы:}
\begin{enumerate}
  \item \textbf{Затухающие градиенты:} при сходимости ($|\tilde{g}'| < 1$) произведение $\prod \tilde{g}' \to 0$ экспоненциально. При $|\tilde{g}'| = 0.5$, $N = 20$: $\prod = 0.5^{20} \approx 10^{-6}$.
  \item \textbf{Взрывающиеся градиенты:} при расходимости ($|\tilde{g}'| > 1$) произведение растёт экспоненциально.
  \item \textbf{Зависимость от $N$:} градиент зависит от числа итераций, а не от свойств решения.
  \item \textbf{Память:} хранение всех промежуточных $x_2^{(k)}$: $O(N \cdot B)$, где $B$ --- размер батча.
\end{enumerate}

Это эквивалентно «наивному» автоматическому дифференцированию через развёрнутые итерации. При $N=20$ и $|\tilde{g}'(x^*)|=0.5$ цепочка даёт фактор $0.5^{20}\approx10^{-6}$, поэтому полезный градиент практически исчезает. Неявное дифференцирование устраняет эту зависимость от длины развёртки и даёт производную в самой неподвижной точке.

\subsection{Неявное дифференцирование}

Используемый подход соответствует современной литературе по неявным и декларативным слоям~\cite{bai2019,gould2021,amos2017}.

\subsubsection{Теорема о неявной функции}

Определяем:

\begin{equation*}
F(x_2^{*}, \boldsymbol{\theta}) \equiv x_2^{*} - g(x_2^{*};\, \boldsymbol{\theta}) = 0
\end{equation*}

\textbf{Теорема (IFT).} Пусть $F$ непрерывно дифференцируема, $F(x_2^{*}, \boldsymbol{\theta}^{*}) = 0$, и:

\begin{equation*}
\frac{\partial F}{\partial x_2^{*}} \neq 0
\end{equation*}

Тогда в окрестности $\boldsymbol{\theta}^{*}$ существует гладкая функция $x_2^{*}(\boldsymbol{\theta})$ с:

\begin{equation*}
\frac{dx_2^{*}}{d\theta_j} = -\left(\frac{\partial F}{\partial x_2^{*}}\right)^{-1} \frac{\partial F}{\partial \theta_j}
\end{equation*}

\subsubsection{Вычисление частных производных}

\begin{equation*}
\frac{\partial F}{\partial x_2^{*}} = 1 - g'(x_2^{*}) = 1 + x_2^{*} \eta
\end{equation*}

При сходимости метода подстановок ($|g'| < 1$) этот знаменатель строго положителен (если $g' > 0$) или строго больше нуля по модулю (если $-1 < g' < 0$). Следовательно, \textbf{неявное дифференцирование корректно определено всегда, когда итерации сходятся}.

\begin{equation*}
\frac{\partial F}{\partial \theta_j} = -\frac{\partial g}{\partial \theta_j}
\end{equation*}

\subsubsection{Итоговая формула для $\ln x_2$}

Из $\ln x_2^{*} = -\Phi - \ln \gamma_2(x_2^{*};\, \boldsymbol{\theta})$:

\begin{equation*}
\boxed{\frac{d\ln x_2^{*}}{d\theta_j} = -\frac{\dfrac{\partial (\Phi + \ln \gamma_2)}{\partial \theta_j}}{1 + x_2^{*} \cdot \eta}}
\end{equation*}

\subsubsection{Производные по кристаллическим параметрам}

$\ln \gamma_2$ не зависит от $T_{\mathrm{m}}$, $\Delta H_{\mathrm{fus}}$, $\Delta C_p^{\mathrm{fus}}$. Поэтому:

\begin{equation*}
\frac{d\ln x_2^{*}}{dT_{\mathrm{m}}} = -\frac{\partial \Phi / \partial T_{\mathrm{m}}}{1 + x_2^{*}\eta} = -\frac{1}{1 + x_2^{*}\eta} \cdot \frac{\Delta H_{\mathrm{fus}}}{RT_{\mathrm{m}}^2}
\end{equation*}

\begin{equation*}
\frac{d\ln x_2^{*}}{d\Delta H_{\mathrm{fus}}} = -\frac{1}{1 + x_2^{*}\eta} \cdot \frac{1}{R}\!\left(\frac{1}{T} - \frac{1}{T_{\mathrm{m}}}\right)
\end{equation*}

\subsubsection{Производные по NRTL параметрам}

$\Phi$ не зависит от NRTL параметров:

\begin{equation*}
\frac{d\ln x_2^{*}}{d\tau_{12}} = -\frac{\partial \ln \gamma_2 / \partial \tau_{12}}{1 + x_2^{*}\eta}
\end{equation*}

Аналогично для $\tau_{21}$, $\alpha_{12}$, с соответствующими цепными правилами.

\subsubsection{Обратный проход}

Для потери $\mathcal{L} = \ell(\ln x_2^{*}, y)$:

\begin{equation*}
\frac{\partial \mathcal{L}}{\partial \theta_j} = \frac{\partial \mathcal{L}}{\partial \ln x_2^{*}} \cdot \frac{d\ln x_2^{*}}{d\theta_j}
\end{equation*}

Вычислительная стоимость: $O(p)$ умножений, где $p$ --- число физических параметров ($p \leq 8$). Ничтожна по сравнению с backward pass через GNN.

\subsection{Сравнение подходов}

\begin{table}[H]
\centering
\begin{tabular}{|c|c|c|}
\hline
Критерий & Unrolled ($N = 5$) & Implicit \\
\hline
Точность градиента & $\sim (1 - \|g'\|^N)$ & $\sim 100\%$ \\
\hline
Память (backward) & $O(N \cdot B)$ & $O(B)$ \\
\hline
Время backward & $\sim N \times$ NRTL eval & $\sim 1 \times$ NRTL eval \\
\hline
Устойчивость & Зависит от $\|g'\|^N$ & Устойчив с clamp \\
\hline
Реализация & Стандартный autograd & Custom \texttt{Function} \\
\hline
\end{tabular}
\caption{Сравнение подходов к дифференцированию через итеративный решатель.}
\end{table}

\textbf{Вывод:} имплицитное дифференцирование предпочтительно для обучения.

\subsection{Численная устойчивость имплицитного метода}

Критическая точка: знаменатель $1 + x_2^{*}\eta \approx 0$ при $g'(x_2^{*}) \approx 1$.

Защита:

\begin{equation*}
\mathrm{denom} = \mathrm{sign}(1 + x_2^{*}\eta) \cdot \max\!\left(|1 + x_2^{*}\eta|,\; \varepsilon_{\min}\right), \quad \varepsilon_{\min} = 0.01
\end{equation*}

Это вносит ограниченную погрешность при $|1 + x_2^{*}\eta| < 0.01$ (проблемные системы с высокой растворимостью), но предотвращает NaN.

\subsection{Связь с моделями глубокого равновесия}

Параллели с DEQ-архитектурами и обратным распространением для неподвижной точки обсуждаются в~\cite{bai2019}.

Наш решатель --- частный случай DEQ (Bai et al., NeurIPS 2019) с одномерным скрытым состоянием $z = x_2$ и аналитически заданным (необучаемым) отображением $g$. Все теоретические гарантии литературы по моделям равновесия (DEQ) применимы:

\begin{enumerate}
  \item При сжимающем $g$ неподвижная точка единственна.
  \item Неявный градиент корректно аппроксимирует градиент развёртки при $N \to \infty$.
  \item Фантомный градиент (развёрнутое обратное распространение от сошедшегося $x_2^{*}$ на $N_{\mathrm{phantom}}$ шагов) --- практичная альтернатива полному неявному дифференцированию.
\end{enumerate}

\section{Численная точность}

\subsection{Критические операции}

\begin{table}[H]
\centering
\begin{tabularx}{\textwidth}{|P{3.9cm}|P{2.6cm}|Y|}
\hline
Операция & Риск & Решение \\
\hline
$\exp(-\Phi)$ при $\Phi > 80$ & Underflow FP32 & Работать с $\ln x_2$, clamp $\Phi$ \\
\hline
$\ln(x_2)$ при $x_2 \to 0$ & $-\infty$ & $\ln(x_2 + \varepsilon)$, $\varepsilon = 10^{-10}$ \\
\hline
$A = x_2 + x_1 G_{12}$ при оба малые & Потеря точности & $A + \varepsilon$, $\varepsilon = 10^{-10}$ \\
\hline
$1/(1 + x_2^{*}\eta)$ & Деление на $\sim 0$ & Clamp знаменателя \\
\hline
\end{tabularx}
\caption{Критические численные операции и меры стабилизации SLE-решателя.}
\end{table}

\subsection{Режим точности}

В текущей реализации: GNN в FP16 (через AMP), физический слой в FP32. Имплицитное дифференцирование --- в FP32.

\section*{Часть III. Архитектура нейронной сети}
\addcontentsline{toc}{section}{Часть III. Архитектура нейронной сети}

\section{Общая структура модели}

\subsection{Принцип проектирования}

Архитектура TGNN-Solv следует декомпозиции мастер-уравнения:

\begin{equation*}
\ln x_2 = \underbrace{-\Phi(T;\; T_{\mathrm{m}}, \Delta H_{\mathrm{fus}}, \Delta C_p^{\mathrm{fus}})}_{\text{Crystal Head}} \;\underbrace{- \ln \gamma_2(x_2;\; \tau_{12}, \tau_{21}, \alpha_{12})}_{\text{NRTL Head + SLE-решатель}}
\end{equation*}

Нейронная сеть предсказывает \textbf{физические параметры}, а не солюбильность напрямую. Эти параметры подаются в аналитически заданные термодинамические уравнения (ноль обучаемых параметров в физическом слое).

\subsection{Блок-схема прямого прохода}

\begin{figure}[H]
\centering
\begin{tikzpicture}[node distance=4.5mm,>=Latex]
\tikzstyle{blk}=[draw,rounded corners,align=center,minimum width=4.5cm,minimum height=0.72cm]
\node[blk] (inp) {Входы: SMILES растворяемого вещества, SMILES растворителя, $T$};
\node[blk,below=of inp] (feat) {Построение молекулярных признаков и графов};
\node[blk,below=of feat] (enc) {MPNN-энкодер (6 слоёв)};
\node[blk,below=of enc] (aux) {Выходные блоки кристаллических параметров: $T_{\mathrm{m}}$, $\Delta H_{\mathrm{fus}}$, $\Delta C_p$};
\node[blk,below=of aux] (att) {Блок попарного внимания (cross-attention) и объединение признаков пары};
\node[blk,below=of att] (nrtl) {Выходной блок параметров модели активности: $\tau_{12}(T),\tau_{21}(T),\alpha_{12}$};
\node[blk,below=of nrtl] (sle) {Жёстко заданный SLE-решатель};
\node[blk,below=of sle] (corr) {Ограниченная коррекция и итог: $\widehat{\ln x_2}$};
\draw[->] (inp)--(feat)--(enc)--(aux)--(att)--(nrtl)--(sle)--(corr);
\end{tikzpicture}
\caption{Схема архитектуры TGNN-Solv с подписями всех ключевых блоков.}
\label{fig:tgnn_arch}
\end{figure}

\begin{equation*}
\text{SMILES}_{\mathrm{solute}},\; \text{SMILES}_{\mathrm{solvent}},\; T
\end{equation*}

\begin{equation*}
\downarrow \quad \text{Построение признаков}
\end{equation*}

\begin{equation*}
\mathcal{G}_{\mathrm{sol}} = (V_{\mathrm{sol}}, E_{\mathrm{sol}}), \quad \mathcal{G}_{\mathrm{slv}} = (V_{\mathrm{slv}}, E_{\mathrm{slv}})
\end{equation*}

\begin{equation*}
\downarrow \quad \text{GNN Encoder}
\end{equation*}

\begin{equation*}
\mathbf{h}_{\mathrm{sol}}^{(\ell)} \in \mathbb{R}^{n_{\mathrm{sol}} \times d}, \quad \mathbf{h}_{\mathrm{slv}}^{(\ell)} \in \mathbb{R}^{n_{\mathrm{slv}} \times d}
\end{equation*}

\begin{equation*}
\downarrow \quad \text{Auxiliary Heads (pre-interaction)}
\end{equation*}

\begin{equation*}
T_{\mathrm{m}}, \; \Delta H_{\mathrm{fus}}, \; \boldsymbol{\delta}_{\mathrm{Hansen}}, \; V_m
\end{equation*}

\begin{equation*}
\downarrow \quad \text{Cross-Attention Interaction}
\end{equation*}

\begin{equation*}
\tilde{\mathbf{h}}_{\mathrm{sol}}, \; \tilde{\mathbf{h}}_{\mathrm{slv}}
\end{equation*}

\begin{equation*}
\downarrow \quad \text{Readout (Attention + Set2Set)}
\end{equation*}

\begin{equation*}
\mathbf{g}_{\mathrm{sol}} \in \mathbb{R}^{3d}, \quad \mathbf{g}_{\mathrm{slv}} \in \mathbb{R}^{3d}
\end{equation*}

\begin{equation*}
\downarrow \quad \text{Pair Representation}
\end{equation*}

\begin{equation*}
\mathbf{g}_{\mathrm{pair}} \in \mathbb{R}^{d_{\mathrm{pair}}}
\end{equation*}

\begin{equation*}
\downarrow \quad \text{NRTL Head}
\end{equation*}

\begin{equation*}
\tau_{12}, \; \tau_{21}, \; \alpha_{12}
\end{equation*}

\begin{equation*}
\downarrow \quad \text{SLE-решатель (аналитически заданный)}
\end{equation*}

\begin{equation*}
\ln x_2^{\mathrm{physics}}
\end{equation*}

\begin{equation*}
\downarrow \quad \text{Bounded Correction}
\end{equation*}

\begin{equation*}
\ln x_2^{\mathrm{final}}
\end{equation*}

\section{Построение молекулярных признаков}
\label{sec:mol-features}

\subsection{Граф молекулы}

Каждая молекула представляется как направленный граф $\mathcal{G} = (V, E)$:

\begin{itemize}
  \item Вершины $V$: атомы
  \item Рёбра $E$: химические связи
\end{itemize}

\subsection{Атомные признаки}

Вектор $\mathbf{v}_i \in \mathbb{R}^{35}$:

\begin{table}[H]
\centering
\setlength{\tabcolsep}{2.5pt}
\begin{tabularx}{0.98\textwidth}{|P{3.8cm}|C{2.0cm}|Y|}
\hline
Признак & \shortstack[c]{Размер-\\ность} & Кодирование \\
\hline
Атомный номер & 12 & One-hot (C, N, O, F, P, S, Cl, Br, I, Si, B, другие) \\
\hline
Гибридизация & 5 & One-hot (sp, sp$^2$, sp$^3$, sp$^3$d, sp$^3$d$^2$) \\
\hline
Формальный заряд & 1 & Целое число \\
\hline
Число H & 1 & Целое число \\
\hline
Ароматичность & 1 & Бинарный \\
\hline
Принадлежность кольцу & 1 & Бинарный \\
\hline
\shortstack[l]{Электроотрица-\\тельность} & 1 & Вещественное (Полинг) \\
\hline
Ван-дер-Ваальсов радиус & 1 & Вещественное (\AA) \\
\hline
Поляризуемость & 1 & Вещественное \\
\hline
Прочие & $\sim 10$ & --- \\
\hline
\end{tabularx}
\caption{Атомные признаки для графового представления молекулы.}
\end{table}

\subsection{Признаки связей}

Вектор $\mathbf{e}_{ij} \in \mathbb{R}^{8}$:

\begin{center}
\begin{tabular}{|c|c|c|}
\hline
Признак & Размерность & Кодирование \\
\hline
Тип связи & 4 & One-hot (одинарная, двойная, тройная, ароматическая) \\
\hline
Сопряжённость & 1 & Бинарный \\
\hline
В кольце & 1 & Бинарный \\
\hline
Стереохимия E/Z & 2 & Бинарный \\
\hline
\end{tabular}
\end{center}

\section{Графовый энкодер GNN}
\label{sec:gnn-encoder}

\subsection{Поддерживаемые типы энкодера}

В текущей поддерживаемой ветке и TGNN-Solv, и DirectGNN поддерживают два типа графового энкодера через параметр
\texttt{encoder\_type}:
\begin{itemize}
  \item \texttt{encoder\_type="mpnn"} --- базовый message-passing энкодер;
  \item \texttt{encoder\_type="gps"} --- гибридный графовый трансформер (GPS-режим).
\end{itemize}

Это изменение принципиально важно для корректных сравнений: парные эксперименты TGNN и DirectGNN теперь можно проводить не только на одном backbone, но и на одинаковом типе энкодера (MPNN по сравнению с MPNN или GPS по сравнению с GPS).

\subsection{Нейронная сеть с передачей сообщений (MPNN)}

Базовая формулировка MPNN опирается на~\cite{gilmer2017}; родственные molecular GNN-подходы включают SchNet и GAT~\cite{schutt2018,velickovic2018}.

Используется $L$-слойная MPNN ($L = 6$ по умолчанию). На каждом слое $\ell = 1, \ldots, L$:

\textbf{Сообщение:}

\begin{equation*}
\mathbf{m}_{ij}^{(\ell)} = \phi_{\mathrm{msg}}^{(\ell)}\!\left(\mathbf{h}_i^{(\ell-1)},\; \mathbf{h}_j^{(\ell-1)},\; \mathbf{e}_{ij}\right)
\end{equation*}

где $\phi_{\mathrm{msg}}^{(\ell)}: \mathbb{R}^{d} \times \mathbb{R}^{d} \times \mathbb{R}^{8} \to \mathbb{R}^{d}$ --- обучаемая MLP.

\textbf{Агрегация:}

\begin{equation*}
\mathbf{a}_i^{(\ell)} = \sum_{j \in \mathcal{N}(i)} \mathbf{m}_{ij}^{(\ell)}
\end{equation*}

\textbf{Обновление с остаточной связью:}

\begin{equation*}
\mathbf{h}_i^{(\ell)} = \mathbf{h}_i^{(\ell-1)} + \mathrm{MLP}_{\mathrm{upd}}^{(\ell)}\!\left(\mathbf{h}_i^{(\ell-1)},\; \mathbf{a}_i^{(\ell)}\right)
\end{equation*}

\textbf{Нормализация:}

\begin{equation*}
\mathbf{h}_i^{(\ell)} \leftarrow \mathrm{LayerNorm}\!\left(\mathbf{h}_i^{(\ell)}\right)
\end{equation*}

\subsection{GPS-режим энкодера}

В режиме \texttt{encoder\_type="gps"} используется гибридный графовый трансформер, в котором локальный MPNN-субслой сочетается с глобальным self-attention. На слое $\ell$ обновление записывается в виде
\begin{equation*}
\mathbf{h}_i^{(\ell)}=
\mathbf{h}_i^{(\ell-1)}+
\mathrm{MPNN}^{(\ell)}\!\left(\mathbf{h}^{(\ell-1)},\mathbf{e}\right)_i+
\mathrm{SelfAttn}^{(\ell)}\!\left(\mathbf{h}^{(\ell-1)}\right)_i,
\end{equation*}
после чего применяются \texttt{LayerNorm} и позиционно-независимый FFN-блок.

Физическая мотивация GPS-режима состоит в расширении рецептивного поля: помимо локальных химических окружений, энкодер получает глобальный канал обмена информацией между удалёнными атомными фрагментами, что потенциально полезно для систем с дальнодействующими стерическими и электростатическими эффектами.

На момент текущего отчёта GPS-режим поддерживается конфигурационно, однако его замороженная спецификация для бенчмарка ещё не завершена: конкретный тип позиционного кодирования (например, Laplacian PE, RWPE или отсутствие явного PE) и итоговое число параметров пока не зафиксированы как часть основного протокола статьи. Поэтому сравнительные результаты MPNN по сравнению с GPS вынесены в отдельный будущий экспериментальный блок и в текущие таблицы не включены.

\subsection{Режим энкодера: \texttt{shared\_residual} (по умолчанию)}

Один и тот же набор весов $\{\phi_{\mathrm{msg}}^{(\ell)}, \mathrm{MLP}_{\mathrm{upd}}^{(\ell)}\}_{\ell=1}^{L}$ используется для обоих молекул (solute и solvent). После $L$ слоёв применяются лёгкие специализированные адаптеры:

\begin{equation*}
\mathbf{h}_i^{\mathrm{sol}} = \mathbf{h}_i^{(L)} + \mathrm{Adapter}_{\mathrm{sol}}(\mathbf{h}_i^{(L)})
\end{equation*}

\begin{equation*}
\mathbf{h}_i^{\mathrm{slv}} = \mathbf{h}_i^{(L)} + \mathrm{Adapter}_{\mathrm{slv}}(\mathbf{h}_i^{(L)})
\end{equation*}

где $\mathrm{Adapter}$: $\mathbb{R}^{d} \to \mathbb{R}^{d}$ --- одно- или двухслойная MLP малой ёмкости.

\textbf{Обоснование:} общий backbone использует общую химическую грамматику; адаптеры позволяют различить роли solute/solvent без удвоения параметров.

\subsection{Альтернативный режим: \texttt{split\_late}}

Первые $(L - L_{\mathrm{split}})$ слоёв --- общие, последние $L_{\mathrm{split}}$ слоёв --- раздельные для solute и solvent. Позволяет сравнить общий по сравнению с асимметричный энкодер при прочих равных.

\subsection{Температурная изоляция}

\textbf{Ключевое проектное решение:} температура $T$ \textbf{не подаётся} в графовый энкодер и предварительные выходные блоки. Это гарантирует, что предсказанные $T_{\mathrm{m}}$, $\Delta H_{\mathrm{fus}}$, $\boldsymbol{\delta}_{\mathrm{Hansen}}$ не зависят от рабочей температуры (физически верно: эти свойства --- характеристики чистого вещества, а не процесса растворения).

Температура входит только в:
\begin{itemize}
  \item NRTL Head (через $\tau(T)$)
  \item SLE-решатель (через $\Phi(T)$)
  \item Bounded Correction
\end{itemize}

\section{Термодинамически-информированная передача сообщений (TIMP)}
\label{sec:timp}

\subsection{Мотивация и место в архитектуре}

Диагностика, приведённая в \hyperref[sec:encoder-bottleneck]{разделе~\ref*{sec:encoder-bottleneck}}, показывает, что именно графовый энкодер остаётся главным информационным узким местом всей системы. В частности, линейная проба на замороженных эмбеддингах восстанавливает экспертные дескрипторы лишь с медианным качеством $R^2 = 0.505$, а в \hyperref[subsec:error-decomposition]{§\ref*{subsec:error-decomposition}} примерно $93\%$ разрыва между TGNN-Solv и RF-ориентиром отнесено именно к качеству представления, а не к физическому слою. Это указывает, что базовый MPNN недоиспользует химическую информацию, критичную для параметров активности.

Фундаментальная причина состоит в том, что стандартная схема передачи сообщений
\begin{equation*}
\mathbf{m}_{ij}^{(\ell)} = \phi_{\mathrm{msg}}^{(\ell)}\!\left(\mathbf{h}_i^{(\ell-1)},\mathbf{h}_j^{(\ell-1)},\mathbf{e}_{ij}\right)
\end{equation*}
обрабатывает все связи единообразно. Между тем в термодинамике смесей вклад в избыточную энергию Гиббса $G^E$ и, следовательно, в коэффициент активности $\gamma_2$ складывается из нескольких физически разных механизмов. В минимальном приближении можно выделить:
\begin{align*}
E_{\mathrm{disp}} &\propto -\frac{\alpha_i\alpha_j}{r_{ij}^6}, \\
E_{\mathrm{polar}} &\propto \frac{\delta q_i\,\delta q_j}{\varepsilon\,r_{ij}}, \\
E_{\mathrm{HB}} &\propto -\epsilon_{\mathrm{HB}}\,f(\theta,r),
\end{align*}
где $\alpha_i$ --- атомная поляризуемость, $\delta q_i$ --- парциальный заряд, а $f(\theta,r)$ кодирует направленность и локальность водородной связи. Аналогичная декомпозиция давно используется в прикладной термодинамике: Hansen явно разделяет $\delta_d$, $\delta_p$ и $\delta_h$; COSMO-RS раскладывает вклад на электростатическое несоответствие поверхностных сегментов, водородные связи и ван-дер-ваальсовы взаимодействия; UNIFAC и UNIQUAC разделяют размерно-комбинаторную и энергетическую части~\cite{hansen2007,klamt2017,weidlich1987,abrams1975}.

TIMP предлагается как следующее архитектурное расширение базового энкодера семейства MPNN/GPS: вместо одного универсального канала передачи сообщений вводятся два физически мотивированных канала --- дисперсионный и полярно-ассоциативный. Цель этой модификации не в том, чтобы заменить физический слой, а в том, чтобы подать в него более структурированное парное представление.

\subsection{Расширение входных признаков}

На уровне узлов TIMP добавляет к стандартному вектору атомных признаков парциальный заряд Gasteiger--Marsili, вычисляемый средствами RDKit~\cite{landrum2024}:
\begin{equation*}
\delta q_i = \texttt{GasteigerCharge}(\mathrm{mol}, i),
\qquad
\mathbf{v}_i^{\mathrm{TIMP}} = [\mathbf{v}_i \;\|\; \delta q_i] \in \mathbb{R}^{36}.
\end{equation*}
В редких случаях, когда итеративная схема зарядов не сходится, используется безопасное резервное значение $\delta q_i = 0$.

На уровне рёбер к стандартному вектору $\mathbf{e}_{ij}\in\mathbb{R}^8$ добавляются четыре непрерывных физически мотивированных признака:
\begin{equation*}
\mathbf{e}_{ij}^{\mathrm{TIMP}} = [\mathbf{e}_{ij} \;\|\; \Delta\chi_{ij} \;\|\; \Delta r_{\mathrm{vdW},ij} \;\|\; \mu_{ij}^{\mathrm{bond}} \;\|\; h_{ij}^{\mathrm{cap}}] \in \mathbb{R}^{12}.
\end{equation*}

Здесь
\begin{align*}
\Delta\chi_{ij} &= |\chi_i - \chi_j|, \\
\Delta r_{\mathrm{vdW},ij} &= |r_{\mathrm{vdW},i} - r_{\mathrm{vdW},j}|, \\
\mu_{ij}^{\mathrm{bond}} &= \Delta\chi_{ij}\, d_{ij}^{\mathrm{est}},
\qquad
d_{ij}^{\mathrm{est}} = r_{\mathrm{vdW},i} + r_{\mathrm{vdW},j}.
\end{align*}
Величина $\Delta\chi_{ij}$ выступает приближённой мерой поляризации связи, $\Delta r_{\mathrm{vdW},ij}$ --- мерой стерической и поляризуемостной асимметрии, а $\mu_{ij}^{\mathrm{bond}}$ --- грубой оценкой вклада связи в молекулярный дипольный момент. Бинарный индикатор $h_{ij}^{\mathrm{cap}}\in\{0,1\}$ отмечает связи, принадлежащие фрагментам, способным к образованию водородных связей (типично O--H, N--H, F--H, а также связи, инцидентные типичному акцепторному центру). В межмолекулярном блоке та же логика переносится на пары атомов «растворяемое вещество--растворитель».

\subsection{Двухканальная передача сообщений}

На каждом слое $\ell$ TIMP вычисляет два семейства сообщений, каждое из которых отвечает своему типу физической информации.

\paragraph{Дисперсионный канал.}
Сообщение масштабируется комбинацией поляризуемостей атомов:
\begin{align*}
\mathbf{m}_{ij}^{\mathrm{disp},(\ell)} &= \phi_{\mathrm{disp}}^{(\ell)}\!\left(\mathbf{h}_i^{(\ell-1)},\mathbf{h}_j^{(\ell-1)},\mathbf{e}_{ij}^{\mathrm{TIMP}}\right) \cdot s_{ij}^{\mathrm{disp}}, \\
s_{ij}^{\mathrm{disp}} &= w_{\mathrm{disp}}\sqrt{\alpha_i\alpha_j + \varepsilon} + b_{\mathrm{disp}},
\qquad \varepsilon = 10^{-8},
\end{align*}
где $\phi_{\mathrm{disp}}^{(\ell)}: \mathbb{R}^{d}\times\mathbb{R}^{d}\times\mathbb{R}^{12}\to\mathbb{R}^{d/2}$ --- отдельная многослойная перцептронная сеть. Инициализация $w_{\mathrm{disp}}=1$, $b_{\mathrm{disp}}=0$ делает модуль почти тождественным в начале обучения. Физически множитель $\sqrt{\alpha_i\alpha_j}$ реализует мягкое комбинационное правило для лондонских дисперсионных взаимодействий: связи между высокополяризуемыми атомами усиливаются, между малополяризуемыми --- ослабляются.

\paragraph{Полярный канал.}
Сообщение масштабируется величиной локальной электростатической неоднородности:
\begin{align*}
\mathbf{m}_{ij}^{\mathrm{polar},(\ell)} &= \phi_{\mathrm{polar}}^{(\ell)}\!\left(\mathbf{h}_i^{(\ell-1)},\mathbf{h}_j^{(\ell-1)},\mathbf{e}_{ij}^{\mathrm{TIMP}}\right) \cdot s_{ij}^{\mathrm{polar}}, \\
s_{ij}^{\mathrm{polar}} &= \mathrm{softplus}\!\left(|\delta q_i\,\delta q_j|\right).
\end{align*}
Здесь используется модуль произведения зарядов: масштаб сообщения зависит от силы поляризации, тогда как знак и конкретная химическая роль атомов доступны самой сети через входные признаки $\delta q_i$, $\delta q_j$, $\Delta\chi_{ij}$ и $h_{ij}^{\mathrm{cap}}$. Такая конструкция избегает нежелательной ситуации, когда большие по модулю, но разноимённые заряды ослабляют сообщение только из-за знака произведения.

\paragraph{Агрегация и обновление.}
После вычисления сообщений они агрегируются раздельно:
\begin{equation*}
\mathbf{a}_i^{\mathrm{disp},(\ell)} = \sum_{j\in\mathcal{N}(i)} \mathbf{m}_{ij}^{\mathrm{disp},(\ell)},
\qquad
\mathbf{a}_i^{\mathrm{polar},(\ell)} = \sum_{j\in\mathcal{N}(i)} \mathbf{m}_{ij}^{\mathrm{polar},(\ell)}.
\end{equation*}
Затем каналы объединяются обучаемым блоком
\begin{equation*}
\mathbf{a}_i^{(\ell)} = \mathrm{MLP}_{\mathrm{comb}}^{(\ell)}\!\left([\mathbf{a}_i^{\mathrm{disp},(\ell)} \;\|\; \mathbf{a}_i^{\mathrm{polar},(\ell)}]\right) \in \mathbb{R}^{d},
\end{equation*}
и состояние обновляется через остаточную связь:
\begin{equation*}
\mathbf{h}_i^{(\ell)} = \mathrm{LayerNorm}\!\left(\mathbf{h}_i^{(\ell-1)} + \mathrm{MLP}_{\mathrm{upd}}^{(\ell)}\!\left([\mathbf{h}_i^{(\ell-1)} \;\|\; \mathbf{a}_i^{(\ell)}]\right)\right).
\end{equation*}

Ключевое отличие TIMP от простого ``двухветвевого'' агрегирования состоит в том, что раздельная физическая структура сохраняется на \emph{каждом} слое передачи сообщений. Поэтому модель не только считывает итоговые признаки из двух каналов, но и накапливает различную дисперсионную и полярную информацию по мере расширения рецептивного поля.

\subsection{Сохранение каналов до парного представления}

Для последующего взаимодействия растворяемого вещества и растворителя TIMP сохраняет не только итоговое скрытое состояние $\mathbf{h}_i^{(L)}$, но и последние канально-специфические агрегаты. Это позволяет строить молекулярные представления уровня графа в трёх параллельных ветвях:
\begin{align*}
\mathbf{g}_{\mathrm{mol}}^{\mathrm{comb}} &= \mathrm{Agg}\!\left(\{\mathbf{h}_i^{(L)}\}_{i\in V}\right), \\
\mathbf{g}_{\mathrm{mol}}^{\mathrm{disp}} &= \mathrm{Agg}\!\left(\{\mathbf{a}_i^{\mathrm{disp},(L)}\}_{i\in V}\right), \\
\mathbf{g}_{\mathrm{mol}}^{\mathrm{polar}} &= \mathrm{Agg}\!\left(\{\mathbf{a}_i^{\mathrm{polar},(L)}\}_{i\in V}\right).
\end{align*}

При формировании парного представления в блоке \hyperref[sec:readout-pair]{\ref*{sec:readout-pair}} каналы комбинируются раздельно:
\begin{align*}
\mathbf{g}_{\mathrm{pair}}^{\mathrm{disp}} &= f_{\mathrm{pair}}\!\left(\mathbf{g}_{\mathrm{sol}}^{\mathrm{disp}},\mathbf{g}_{\mathrm{slv}}^{\mathrm{disp}}\right), \\
\mathbf{g}_{\mathrm{pair}}^{\mathrm{polar}} &= f_{\mathrm{pair}}\!\left(\mathbf{g}_{\mathrm{sol}}^{\mathrm{polar}},\mathbf{g}_{\mathrm{slv}}^{\mathrm{polar}}\right), \\
\mathbf{g}_{\mathrm{pair}}^{\mathrm{comb}} &= f_{\mathrm{pair}}\!\left(\mathbf{g}_{\mathrm{sol}}^{\mathrm{comb}},\mathbf{g}_{\mathrm{slv}}^{\mathrm{comb}}\right), \\
\mathbf{g}_{\mathrm{pair}} &= \mathrm{Proj}\!\left([\mathbf{g}_{\mathrm{pair}}^{\mathrm{comb}} \;\|\; \mathbf{g}_{\mathrm{pair}}^{\mathrm{disp}} \;\|\; \mathbf{g}_{\mathrm{pair}}^{\mathrm{polar}}]\right).
\end{align*}
Функция $f_{\mathrm{pair}}$ может совпадать с базовой операцией из \hyperref[sec:readout-pair]{раздела~\ref*{sec:readout-pair}} (конкатенация, поэлементное произведение, модуль разности). Важный эффект этой конструкции --- явное соответствие между каналами представления и физическими декомпозициями Hansen/UNIQUAC/NRTL: выходной блок NRTL получает вход, где типы взаимодействий уже частично факторизованы, а не перемешаны в одном латентном векторе.

\subsection{Термодинамически смещённое перекрёстное внимание и агрегирование с учётом поверхности}

Стандартный блок взаимодействия из \hyperref[sec:interaction-block]{раздела~\ref*{sec:interaction-block}} можно дополнить термодинамически мотивированным смещением в логитах перекрёстного внимания. Для пары атомов $i$ (растворяемое вещество) и $j$ (растворитель) вводится
\begin{equation*}
\ell_{ij}^{\mathrm{TIMP}} = \frac{\mathbf{q}_i^\top \mathbf{k}_j}{\sqrt{d_k}} - \beta\,E_{ij}^{\mathrm{contact}},
\end{equation*}
после чего веса внимания вычисляются как
\begin{equation*}
\alpha_{ij}^{\mathrm{TIMP}} = \frac{\exp(\ell_{ij}^{\mathrm{TIMP}})}{\sum_{j'} \exp(\ell_{ij'}^{\mathrm{TIMP}})}.
\end{equation*}
Здесь
\begin{equation*}
E_{ij}^{\mathrm{contact}} = \mathrm{MLP}_{\mathrm{energy}}\!\left([\mathbf{h}_i \;\|\; \mathbf{h}_j]\right)\cdot \mathrm{acc}_i\cdot \mathrm{acc}_j,
\qquad
\mathrm{acc}_i = \sigma\!\left(\mathrm{MLP}_{\mathrm{acc}}(\mathbf{h}_i)\right) \in (0,1).
\end{equation*}
Скаляр $\beta$ обучаем и инициализируется малым положительным значением (например, $0.1$). Знак ``$-\beta E_{ij}^{\mathrm{contact}}$'' выбран так, чтобы более выгодные (низкоэнергетические) контакты получали больший логит и, следовательно, больший вес внимания. При $\beta=0$ блок точно вырождается в стандартное перекрёстное внимание.

Если в модель добавляется геометрический путь с устойчивой оценкой доступной поверхности, TIMP естественно сочетается с агрегированием, взвешенным по доступной поверхности:
\begin{equation*}
\mathbf{g}_{\mathrm{mol}}^{\mathrm{SASA}} = \sum_{i\in V} \frac{S_i}{\sum_{j\in V} S_j + \varepsilon}\,\tilde{\mathbf{h}}_i,
\end{equation*}
где $S_i$ --- вклад атома в доступную растворителю площадь поверхности (SASA), а $\tilde{\mathbf{h}}_i$ --- представление после блока взаимодействия. Такой способ агрегирования ближе к смесьевой интуиции UNIQUAC, где поверхностная доля определяет относительную ``видимость'' компонента. В текущем замороженном двумерном бенчмарке эта модификация рассматривается как опциональная, поскольку требует устойчивой приближённой оценки SASA и, как правило, трёхмерной конформации.

\subsection{Регуляризация каналов и совместимость с физическими выходными блоками}

Разделение на каналы само по себе не гарантирует, что один канал не начнёт дублировать другой. Поэтому TIMP допускает мягкую регуляризацию через редкие Hansen-метки. Для множества молекул $\mathcal{H}$, где доступны справочные значения $(\delta_d,\delta_p,\delta_h)$, вводятся дополнительные потери
\begin{align*}
\mathcal{L}_{\mathrm{disp\text{-}Hansen}} &= \frac{1}{|\mathcal{H}|}\sum_{m\in\mathcal{H}}\left(\mathrm{MLP}_{\delta_d}(\mathbf{g}_m^{\mathrm{disp}}) - \delta_{d,m}^{\mathrm{true}}\right)^2, \\
\mathcal{L}_{\mathrm{polar\text{-}Hansen}} &= \frac{1}{|\mathcal{H}|}\sum_{m\in\mathcal{H}}\Bigl[\left(\mathrm{MLP}_{\delta_p}(\mathbf{g}_m^{\mathrm{polar}}) - \delta_{p,m}^{\mathrm{true}}\right)^2 + \left(\mathrm{MLP}_{\delta_h}(\mathbf{g}_m^{\mathrm{polar}}) - \delta_{h,m}^{\mathrm{true}}\right)^2\Bigr].
\end{align*}
Из-за крайне малого покрытия Hansen-меток в текущем корпусе (см. \hyperref[tab:data-aux-coverage]{табл.~\ref*{tab:data-aux-coverage}}) эти члены должны иметь небольшой вес порядка $0.02$--$0.05$ и трактоваться именно как структурная регуляризация, а не как самостоятельная цель.

Последующая архитектура при этом не меняется. Блок \texttt{FusionHead} по-прежнему получает представление растворяемого вещества до блока взаимодействия, блок \texttt{NRTLHead} работает с итоговым $\mathbf{g}_{\mathrm{pair}}$, а SLE-решатель остаётся жёстко заданным аналитическим ядром. Следовательно, TIMP меняет не физику решателя, а только качество и факторизацию параметров, которые в него поступают.

\subsection{Вычислительная цена, обратная совместимость и ожидаемый эффект}

По вычислительной цене TIMP заменяет один MLP блока сообщений двумя MLP половинной ширины и добавляет один блок комбинирования каналов. В результате накладные расходы составляют примерно $\sim 1.3\times$ по числу арифметических операций (FLOPs) на слой энкодера и $\sim 1.5\times$ по числу параметров самого блока передачи сообщений. Для типичной конфигурации $d=256$, $L=6$ это даёт прирост полного числа параметров модели порядка $+30\%$ и рост времени вывода порядка $+20\%$; термодинамическое смещение в перекрёстном внимании добавляет ещё небольшой дополнительный расход на MLP, но не меняет асимптотику.

С архитектурной точки зрения TIMP является обратно совместимым расширением. Если сеть научится игнорировать добавленные признаки, положить веса каналов равными и свернуть $\mathrm{MLP}_{\mathrm{comb}}$ в тождественное преобразование, модуль содержит базовый MPNN как частный случай. Поэтому TIMP не вводит нового термодинамического риска: SLE-слой остаётся точным по построению, а ограничения NRTL/Wilson-параметризации не нарушаются.

Ожидаемый эффект TIMP связан прежде всего с идентифицируемостью. В \hyperref[sec:identifiability]{разделе~\ref*{sec:identifiability}} показано, что выходной блок NRTL в текущей конфигурации получает слишком ``размытое'' представление пары и вынужден одновременно кодировать тип и силу взаимодействия в нескольких коррелированных параметрах. TIMP смягчает эту проблему двумя путями: (i) структурирует парное представление на дисперсионную и полярно-ассоциативную части; (ii) вводит слабые априорные ограничения, препятствующие коллапсу каналов. Практическим критерием успеха здесь служит не только итоговый MAE, но и рост качества линейной пробы по MolLogP, TPSA и дескрипторам, связанным с водородными связями, относительно базового энкодера.

По отношению к литературе TIMP занимает промежуточную позицию. С одной стороны, он близок к семействам SchNet/DimeNet/SphereNet в том смысле, что вводит непрерывные физические модуляторы сообщений, а не полагается только на дискретный тип связи~\cite{schutt2018,gasteiger2020,liu2022}. С другой стороны, TIMP ортогонален геометрически эквивариантным GNN: он разделяет каналы не по симметрии пространства, а по типу термодинамически релевантной физики. В этом смысле TIMP можно рассматривать как наиболее прямой перенос идеи Hansen/COSMO/UNIQUAC-декомпозиции на уровень передачи сообщений для задач растворимости.

\section{Вспомогательные выходные блоки (до взаимодействия)}

\subsection{Блок кристаллических свойств (FusionHead)}

Работает с graph-level представлением solute до взаимодействия с solvent.

\textbf{Pooling:}

\begin{equation*}
\mathbf{g}_{\mathrm{sol}}^{\mathrm{pre}} = \frac{1}{|V_{\mathrm{sol}}|}\sum_{i \in V_{\mathrm{sol}}} \mathbf{h}_i^{\mathrm{sol}}
\end{equation*}

\textbf{Предсказания:}

\begin{equation*}
T_{\mathrm{m}} = 100 + 600 \cdot \sigma\!\left(\mathrm{MLP}_{T_{\mathrm{m}}}(\mathbf{g}_{\mathrm{sol}}^{\mathrm{pre}})\right) \quad \in [100, 700] \;\text{К}
\end{equation*}

\begin{equation*}
\Delta H_{\mathrm{fus}} = S_H \cdot \mathrm{softplus}\!\left(\mathrm{MLP}_{\Delta H}(\mathbf{g}_{\mathrm{sol}}^{\mathrm{pre}})\right) \quad > 0 \;\text{Дж/моль}
\end{equation*}

где $\sigma$ --- сигмоида, $S_H$ --- масштабный множитель, $\mathrm{softplus}(z) = \ln(1 + e^z)$.

\textbf{Обоснование ограничений:}

\begin{itemize}
  \item $T_{\mathrm{m}} \in [100, 700]$ К: покрывает $>99\%$ органических соединений.
  \item $\Delta H_{\mathrm{fus}} > 0$: термодинамическое требование (плавление --- эндотермический процесс).
  \item $\Delta C_p^{\mathrm{fus}} = 0$ по умолчанию (фиксировано, не предсказывается).
\end{itemize}

\subsection{Блок параметров Хансена}

Постановка опирается на концепцию параметров растворимости Хансена~\cite{hansen2007}.

Предсказывает Hansen параметры растворимости для каждой молекулы (solute и solvent):

\begin{equation*}
(\delta_d, \delta_p, \delta_h) = \mathrm{MLP}_{\mathrm{Hansen}}(\mathbf{g}_{\mathrm{mol}}^{\mathrm{pre}})
\end{equation*}

с ограничениями $\delta_d, \delta_p, \delta_h \geq 0$ (через softplus).

Используются в bridge-функции потерь для мягкой консистентности с NRTL, где Hansen-дистанция сопоставляется с аналитической оценкой предельной активности в рамках Regular Solution приближения.

\subsection{Блок вспомогательных свойств (AuxPropsHead)}

Предсказывает молярный объём:

\begin{equation*}
V_m = \mathrm{MLP}_{V_m}(\mathbf{g}_{\mathrm{mol}}^{\mathrm{pre}})
\end{equation*}

Используется только как auxiliary target при наличии экспериментальных данных.

\section{Взаимодействие растворяемого вещества и растворителя}
\label{sec:interaction-block}

\subsection{Перекрёстное внимание (по умолчанию)}

Стек из $K$ ($K = 3$) слоёв Transformer cross-attention:

\textbf{Слой $k$:}

\begin{equation*}
\tilde{\mathbf{h}}_{\mathrm{sol}}^{(k)} = \mathrm{CrossAttn}^{(k)}\!\left(Q = \mathbf{h}_{\mathrm{sol}}^{(k-1)},\; K = \mathbf{h}_{\mathrm{slv}}^{(k-1)},\; V = \mathbf{h}_{\mathrm{slv}}^{(k-1)}\right)
\end{equation*}

\begin{equation*}
\tilde{\mathbf{h}}_{\mathrm{slv}}^{(k)} = \mathrm{CrossAttn}^{(k)}\!\left(Q = \mathbf{h}_{\mathrm{slv}}^{(k-1)},\; K = \mathbf{h}_{\mathrm{sol}}^{(k-1)},\; V = \mathbf{h}_{\mathrm{sol}}^{(k-1)}\right)
\end{equation*}

с остаточными связями и LayerNorm.

Стандартный multi-head attention:

\begin{equation*}
\mathrm{Attn}(Q, K, V) = \mathrm{softmax}\!\left(\frac{QK^\top}{\sqrt{d_k}}\right) V
\end{equation*}

\textbf{Физическая мотивация:} cross-attention позволяет каждому атому solute «посмотреть» на все атомы solvent (и наоборот), моделируя межмолекулярные взаимодействия. Веса attention интерпретируемы как «важность» конкретных атомных пар для взаимодействия.

\subsection{Альтернатива: двудольная передача сообщений}

Полный двудольный граф между атомами solute и solvent с message passing. Эквивалентен cross-attention без softmax нормализации.

\section{Выходной блок и парное представление}
\label{sec:readout-pair}

\subsection{Физически-ориентированный выходной блок}

Для каждой молекулы после interaction --- два типа pooling.

\textbf{Attention pooling:}

\begin{equation*}
\alpha_i = \frac{\exp(\mathbf{w}_{\mathrm{att}}^\top \tilde{\mathbf{h}}_i)}{\sum_j \exp(\mathbf{w}_{\mathrm{att}}^\top \tilde{\mathbf{h}}_j)}, \quad \mathbf{g}_{\mathrm{att}} = \sum_i \alpha_i \tilde{\mathbf{h}}_i \in \mathbb{R}^{d}
\end{equation*}

\textbf{Set2Set pooling}~\cite{vinyals2016}:

\begin{equation*}
\mathbf{g}_{\mathrm{s2s}} \in \mathbb{R}^{2d}
\end{equation*}

через $M$ шагов LSTM-based attention ($M = 3$ по умолчанию).

\textbf{Конкатенация:}

\begin{equation*}
\mathbf{g}_{\mathrm{mol}} = [\mathbf{g}_{\mathrm{att}} \;\|\; \mathbf{g}_{\mathrm{s2s}}] \in \mathbb{R}^{3d}
\end{equation*}

\subsection{Парное представление}

\begin{equation*}
\mathbf{g}_{\mathrm{pair}} = \mathrm{MLP}_{\mathrm{pair}}\!\left([\mathbf{g}_{\mathrm{sol}} \;\|\; \mathbf{g}_{\mathrm{slv}} \;\|\; \mathbf{g}_{\mathrm{sol}} \odot \mathbf{g}_{\mathrm{slv}} \;\|\; |\mathbf{g}_{\mathrm{sol}} - \mathbf{g}_{\mathrm{slv}}|]\right) \in \mathbb{R}^{d_{\mathrm{pair}}}
\end{equation*}

где $\odot$ --- поэлементное произведение, $|\cdot|$ --- поэлементный модуль разности.

Четыре компоненты фиксируют: информацию о solute, о solvent, о «совместимости» (произведение), и о «различии» (разность).

\section{Выходной блок NRTL}

\subsection{Параметризация \texttt{ref\_invT} (по умолчанию)}

NRTL Head получает на вход $\mathbf{g}_{\mathrm{pair}}$ и температуру $T$ и предсказывает:

\begin{equation*}
\tau_{12}^{\mathrm{ref}} = \mathrm{MLP}_{\tau_{12}}(\mathbf{g}_{\mathrm{pair}}) \quad \text{--- значение } \tau_{12} \text{ при } T = T_{\mathrm{ref}}
\end{equation*}

\begin{equation*}
\beta_{12} = \mathrm{MLP}_{\beta_{12}}(\mathbf{g}_{\mathrm{pair}}) \quad \text{--- наклон по } 1/T
\end{equation*}

Восстановление $\tau_{12}(T)$:

\begin{equation*}
\tau_{12}(T) = \tau_{12}^{\mathrm{ref}} + \beta_{12}\!\left(\frac{1}{T} - \frac{1}{T_{\mathrm{ref}}}\right)
\end{equation*}

Аналогично для $\tau_{21}(T)$.

\textbf{Параметр неслучайности:}

\begin{equation*}
\alpha_{12} = 0.1 + 0.5 \cdot \sigma\!\left(\mathrm{MLP}_{\alpha}(\mathbf{g}_{\mathrm{pair}})\right) \quad \in [0.1, 0.6]
\end{equation*}

\subsection{Обоснование \texttt{ref\_invT}}

В предыдущей параметризации сеть предсказывала $\Delta g_{12}/(RT)$ напрямую. Проблема: при изменении $T$ нужно предсказать согласованную семейства $\tau(T)$, что требует корректной экстраполяции.

Форма \texttt{ref\_invT} эквивалентна линейной модели $\tau$ от $1/T$:

\begin{equation*}
\tau_{12}(T) = a + b \cdot \frac{1}{T}
\end{equation*}

Это согласуется с физикой: $\tau = \Delta g/(RT) \propto 1/T$ при постоянном $\Delta g$.

\subsection{Температурная инъекция}

Температура $T$ подаётся в NRTL Head \textbf{явно} (через конкатенацию или термометрическое кодирование ThermometerEncoder~\cite{buckman2018}) --- это единственное место, где $T$ входит в encoder/head путь. Это обеспечивает корректную температурную зависимость $\tau(T)$ без загрязнения температурно-инвариантных голов.

\section{SLE-решатель (физический слой)}

\subsection{Жёстко заданные вычисления}

Ноль обучаемых параметров. Все операции --- детерминистические функции предсказанных параметров и температуры.

\textbf{Шаг 1.} Вычисление $\Phi(T)$:

\begin{equation*}
\Phi = \frac{\Delta H_{\mathrm{fus}}}{R}\!\left(\frac{1}{T} - \frac{1}{T_{\mathrm{m}}}\right) - \frac{\Delta C_p^{\mathrm{fus}}}{R}\!\left[\left(\frac{T_{\mathrm{m}}}{T} - 1\right) - \ln\frac{T_{\mathrm{m}}}{T}\right]
\end{equation*}

\textbf{Шаг 2.} Вычисление NRTL вспомогательных величин:

\begin{equation*}
G_{12} = \exp(-\alpha_{12}\tau_{12}), \quad G_{21} = \exp(-\alpha_{12}\tau_{21})
\end{equation*}

\textbf{Шаг 3.} Итеративное решение:

\begin{equation*}
x_2^{(0)} = \exp(-\Phi)
\end{equation*}

\begin{equation*}
\text{for } k = 0, \ldots, N_{\mathrm{iter}} - 1:
\end{equation*}

\begin{equation*}
\quad x_1^{(k)} = 1 - x_2^{(k)}
\end{equation*}

\begin{equation*}
\quad \ln \gamma_2^{(k)} = (x_1^{(k)})^2 \left[\tau_{12}\!\left(\frac{G_{12}}{x_2^{(k)} + x_1^{(k)} G_{12}}\right)^{2} + \frac{\tau_{21}G_{21}}{(x_1^{(k)} + x_2^{(k)} G_{21})^2}\right]
\end{equation*}

\begin{equation*}
\quad x_2^{(k+1)} = \lambda \exp\!\left(-\Phi - \ln\gamma_2^{(k)}\right) + (1 - \lambda) x_2^{(k)}
\end{equation*}

\textbf{Шаг 4.} Выход:

\begin{equation*}
\ln x_2^{\mathrm{physics}} = \ln x_2^{(N_{\mathrm{iter}})}
\end{equation*}

\subsection{Параметры решателя}

\begin{center}
\begin{tabular}{|c|c|c|}
\hline
Параметр & Training & Inference \\
\hline
$N_{\mathrm{iter}}$ & 5 & 20 \\
\hline
$\lambda$ & 0.7 (адаптивно) & 0.7 (адаптивно) \\
\hline
Допуск & $10^{-4}$ & $10^{-8}$ \\
\hline
Точность вычислений & FP32 & FP32 \\
\hline
\end{tabular}
\end{center}

\subsection{Обратный проход}

Используется \texttt{SLESolverFunction} --- собственная \texttt{torch.autograd.Function}, реализующая неявное дифференцирование.

\textbf{Прямой проход:} итерации без отслеживания градиентов $\to$ сошедшееся $x_2^{*}$.

\textbf{Обратный проход:} по формуле неявного дифференцирования для уравнения неподвижной точки:

\begin{equation*}
\frac{d\ln x_2^{*}}{d\theta_j} = -\frac{\partial(\Phi + \ln\gamma_2)/\partial\theta_j}{1 + x_2^{*} \cdot \eta}
\end{equation*}

\section{Ограниченная коррекция (AdaptivePhysicsCorrection)}

\subsection{Мотивация}

Физическая модель SLE+NRTL --- приближённая. Даже при идеально предсказанных параметрах, остаётся модельная погрешность ($\Delta C_p = 0$, ограничения NRTL, и т.д.). Correction head компенсирует эту погрешность, \textbf{не выходя за рамки физической структуры}.

\subsection{Механизм}

\textbf{Шаг 1.} Предсказание bounded deltas:

\begin{equation*}
\delta T_{\mathrm{m}} = \Delta_{\max}^{T_{\mathrm{m}}} \cdot \tanh\!\left(\mathrm{MLP}_{\delta T_{\mathrm{m}}}(\mathbf{g}_{\mathrm{pair}}, \ln x_2^{\mathrm{physics}})\right), \quad \Delta_{\max}^{T_{\mathrm{m}}} = 50\,\mathrm{K}
\end{equation*}

\begin{equation*}
\delta \Delta H = \Delta H_{\mathrm{fus}} \cdot f_{\max}^{\Delta H} \cdot \tanh\!\left(\mathrm{MLP}_{\delta H}(\mathbf{g}_{\mathrm{pair}}, \ln x_2^{\mathrm{physics}})\right)
\end{equation*}

\begin{equation*}
\delta\tau_{12} = \Delta_{\max}^{\tau} \cdot \tanh\!\left(\mathrm{MLP}_{\delta\tau_{12}}(\ldots)\right)
\end{equation*}

\begin{equation*}
\delta\tau_{21} = \Delta_{\max}^{\tau} \cdot \tanh\!\left(\mathrm{MLP}_{\delta\tau_{21}}(\ldots)\right)
\end{equation*}

\textbf{Шаг 2.} Corrected parameters:

\begin{equation*}
T_{\mathrm{m}}' = T_{\mathrm{m}} + \delta T_{\mathrm{m}}, \quad \Delta H' = \Delta H_{\mathrm{fus}} + \delta\Delta H
\end{equation*}

\begin{equation*}
\tau_{12}' = \tau_{12} + \delta\tau_{12}, \quad \tau_{21}' = \tau_{21} + \delta\tau_{21}
\end{equation*}

\textbf{Шаг 3.} Re-solve SLE с corrected parameters:

\begin{equation*}
\ln x_2^{\mathrm{proposal}} = \mathrm{SLE\_Solve}(T_{\mathrm{m}}', \Delta H', \tau_{12}', \tau_{21}', \alpha_{12}, T)
\end{equation*}

\textbf{Шаг 4.} Gated blending:

\begin{equation*}
w = \sigma\!\left(\mathrm{MLP}_{\mathrm{gate}}(\mathbf{g}_{\mathrm{pair}}, \ln x_2^{\mathrm{physics}})\right)
\end{equation*}

\begin{equation*}
\ln x_2^{\mathrm{final}} = \ln x_2^{\mathrm{physics}} + (1 - w) \cdot \mathrm{clip}\!\left(\ln x_2^{\mathrm{proposal}} - \ln x_2^{\mathrm{physics}},\; \pm C_{\max}\right)
\end{equation*}

где $C_{\max}$ = \texttt{correction\_max\_abs} $= 2.0$.

\subsection{Свойства}

\begin{enumerate}
  \item \textbf{Bounded:} коррекция ограничена $\pm C_{\max}$ в $\ln x_2$, что соответствует множителю $e^{\pm 2} \approx [0.14, 7.4]$ в $x_2$.
  \item \textbf{Physics-preserving:} коррекция работает через повторное решение SLE с изменёнными параметрами, а не через произвольную нелинейную функцию.
  \item \textbf{Gated:} gate $w \in [0, 1]$ позволяет модели «отключить» коррекцию для систем, где физика достаточно точна.
\end{enumerate}

\section{DirectGNN (базовый ориентир без физики)}

\subsection{Архитектура}

Тот же графовый энкодер (MPNN или GPS), та же cross-attention, тот же readout и представление пары. Единственное отличие --- вместо физических heads и SLE-решателя:

\begin{equation*}
\ln x_2 = \mathrm{MLP}_{\mathrm{direct}}(\mathbf{g}_{\mathrm{pair}}, T, \mathbf{d}_{\mathrm{desc}})
\end{equation*}

где $\mathbf{d}_{\mathrm{desc}}$ --- опциональный дескрипторный канал (RDKit), который в текущем проекте поддерживается как стандартный вариант DirectGNN.

\subsection{Роль в экспериментальном плане}

DirectGNN --- \textbf{контрольная абляция}: тот же backbone, но без physics bottleneck. Разница в метриках между TGNN и DirectGNN изолирует эффект физической декомпозиции при равных входных условиях (включая дескрипторы и тип энкодера).


\subsection{Технические детали архитектуры}

Подробное описание входных признаков, message-passing, cross-attention и readout уже приведено в §§\ref{sec:mol-features}--\ref{sec:readout-pair}. В этом разделе оставлен только operational summary различий, чтобы избежать дублирования.

\textbf{Что совпадает с TGNN-Solv:}
\begin{itemize}
  \item тот же графовый backbone (MPNN/GPS),
  \item тот же pair-representation и тот же cross-attention блок,
  \item тот же температурный вход через $\mathbf{t}(T)$ (ThermometerEncoder~\cite{buckman2018}).
\end{itemize}

\textbf{Что отличается:} физические heads и SLE-решатель заменяются прямой регрессией
\begin{equation*}
\ln x_2 = \mathrm{MLP}_{\mathrm{direct}}(\mathbf{g}_{\mathrm{pair}},\mathbf{t}(T),\mathbf{d}_{\mathrm{desc}}),
\end{equation*}
где $\mathbf{d}_{\mathrm{desc}}$ --- опциональный канал дескрипторов. Тем самым DirectGNN служит чистой абляцией physics bottleneck при том же представлении входов.


\section{Сводка параметров модели}

\begin{center}
\begin{tabular}{|c|c|c|}
\hline
Компонент & Обучаемые параметры & Физические параметры \\
\hline
GNN Encoder (6 слоёв) & $\sim 6 \times O(d^2)$ & 0 \\
\hline
Role Adapters & $2 \times O(d^2)$ & 0 \\
\hline
Cross-Attention (3 слоя) & $\sim 3 \times O(d^2)$ & 0 \\
\hline
Readout (Att + Set2Set) & $O(d^2)$ & 0 \\
\hline
FusionHead & $O(d^2)$ & $T_{\mathrm{m}}$, $\Delta H_{\mathrm{fus}}$ \\
\hline
HansenHead & $O(d^2)$ & $\delta_d$, $\delta_p$, $\delta_h$ \\
\hline
AuxPropsHead & $O(d)$ & $V_m$ \\
\hline
NRTLHead & $O(d_{\mathrm{pair}}^2)$ & $\tau_{12}^{\mathrm{ref}}$, $\beta_{12}$, $\tau_{21}^{\mathrm{ref}}$, $\beta_{21}$, $\alpha_{12}$ \\
\hline
SLE-решатель & \textbf{0} & --- \\
\hline
BoundedCorrection & $O(d_{\mathrm{pair}}^2)$ & $\delta T_{\mathrm{m}}$, $\delta\Delta H$, $\delta\tau_{12}$, $\delta\tau_{21}$, gate \\
\hline
\textbf{Итого} & $\sim 1.8$--$2.2\,\mathrm{M}$ (при $d = 256$) & 12 физических параметров на sample \\
\hline
\end{tabular}
\end{center}

\section{Предобучение графового энкодера (этап 0)}
\label{sec:stage0-pretraining}

\subsection{Мотивация и место в полном конвейере}

Перед трёхфазным учебным режимом (фазы 1--3) в TGNN-Solv выполняется поддерживаемый \textbf{этап 0}: предобучение блока графового энкодера и выходного блока на большом наборе молекулярных структур без меток растворимости.

Ключевая идея состоит в том, что энкодер должен сначала выучить базовую «химическую грамматику» (локальные валентные закономерности, типичные фрагменты, топологические мотивы, связь локальной структуры с глобальными физико-химическими дескрипторами), и только затем дообучаться на более узкой и шумной задаче растворимости.

Практическая мотивация:
\begin{enumerate}
  \item Даже после merge корпус остаётся ограниченным: полный merged frame содержит 120\,197 записей и 19\,878 unique solute, но прямая supervision по $\ln x_2$ в текущей water-inclusive конфигурации покрывает около 108\,000 записей; в исторической water-excluded конфигурации это было 101\,763 строк и 1\,444 unique solute (см. §\ref{subsec:data-preprocessing}). При этом неразмеченные наборы молекулярных структур на порядок больше.
  \item Для GNN это аналог предварительного обучения в NLP: сначала широкое самообучение на большом корпусе, затем тонкая настройка под конкретную задачу.
  \item Этап 0 снижает зависимость от случайной инициализации и ускоряет вхождение в устойчивый режим на фазе 1.
\end{enumerate}

Архитектурно этап 0 выполняется \textbf{без отдельной копии базовой модели}: обновляются параметры \texttt{model.gnn} и \texttt{model.readout}. Временные блоки предобучения используются только на этапе 0 и удаляются перед стартом фазы 1.

Этап 0 встроен в поддерживаемые команды CLI и вспомогательные конвейерные функции, а его результат сохраняется как повторно используемые файлы контрольных точек для тёплого старта.

\subsection{Данные для этапа 0}

Основной источник предобучения --- \textbf{ZINC250k} (около 250k лекарственно-подобных молекул, открытый набор). Если ZINC недоступен, используется резервный режим: уникальные SMILES из BigSolDB (в текущем merged frame это 19\,878 unique solute; см. §\ref{subsec:data-preprocessing}).

Конвейер подготовки един:
\begin{enumerate}
  \item парсинг SMILES через RDKit;
  \item каноникализация представления;
  \item удаление невалидных/дублирующихся молекул;
  \item построение графа с теми же атомными и связевыми признаками, что и в основном обучении.
\end{enumerate}

Таким образом, сдвиг домена между этапом 0 и основным дообучением минимизируется: отличаться должны прежде всего объём и «богатство» структур, а не формат признаков.

\subsection{Мультизадачная постановка этапа 0}

Этап 0 формулируется как мультизадачная функция цели из четырёх компонентов:
\begin{itemize}
  \item самосупервизируемое восстановление атомных признаков после маскирования подграфов;
  \item самосупервизируемая классификация типа маскированных связей;
  \item слабосупервизируемая регрессия 12 RDKit-дескрипторов;
  \item контрастивная цель на уровне графовых представлений.
\end{itemize}

Комбинация локальных и глобальных задач нужна для того, чтобы энкодер одновременно выучивал:
\begin{enumerate}
  \item локальную химию (атом/связь);
  \item среднеуровневые структурные паттерны (фрагменты, окружения);
  \item глобальные свойства молекулы (logP, TPSA и т.п.);
  \item инвариантность представления к мягким аугментациям.
\end{enumerate}

\subsection{Задача 1: маскирование 2-hop подграфов (самосупервизируемая)}

Маскирование одного атома в химическом графе обычно является слишком простой задачей: по ближайшим соседям тип атома часто восстанавливается почти однозначно. Поэтому применяется маскирование \textbf{подграфов}.

Алгоритм построения маски:
\begin{enumerate}
  \item случайно выбрать стартовый атом;
  \item расширить маску поиском в ширину (BFS) на глубину \texttt{mask\_hops = 2};
  \item повторять для новых стартовых атомов, пока доля маскированных узлов не достигнет \texttt{mask\_ratio} $\approx 15\%$;
  \item обнулить атомные признаки для маскированных узлов;
  \item предсказать исходные атомные признаки из скрытых состояний энкодера.
\end{enumerate}

Функция потерь:
\begin{equation*}
\mathcal{L}_{\text{atom}} = \frac{1}{|\mathcal{M}|} \sum_{i \in \mathcal{M}}
\left\| \hat{\mathbf{x}}_i - \mathbf{x}_i \right\|^2,
\end{equation*}
где $\mathcal{M}$ --- множество маскированных атомов, $\mathbf{x}_i \in \mathbb{R}^{35}$ --- истинный вектор атомных признаков, а
\begin{equation*}
\hat{\mathbf{x}}_i = \mathrm{MLP}_{\text{atom}}(\mathbf{h}_i).
\end{equation*}

Физико-химический смысл: модель вынуждена восстанавливать электроотрицательность, гибридизацию, ароматичность и связанные с ними признаки по расширенному контексту, то есть учит именно те структурные зависимости, которые затем влияют на межмолекулярные взаимодействия и растворимость.

\subsection{Задача 2: предсказание типа связи (самосупервизируемая)}

Для доли связей \texttt{bond\_mask\_ratio} $\approx 15\%$ обнуляются edge-features. Далее по скрытым состояниям концевых атомов предсказывается тип связи.

Логиты для ребра $(i,j)$:
\begin{equation*}
\hat{\mathbf{y}}_{ij} = \mathrm{MLP}_{\text{bond}}\!\left([\mathbf{h}_i \| \mathbf{h}_j]\right) \in \mathbb{R}^{4},
\end{equation*}
классы соответствуют типам: одинарная / двойная / тройная / ароматическая связь.

Кросс-энтропийная потеря:
\begin{equation*}
\mathcal{L}_{\text{bond}} = -\frac{1}{|\mathcal{M}_e|}\sum_{(i,j) \in \mathcal{M}_e}
\sum_{c=1}^{4} y_{ij}^{(c)}\log \hat{y}_{ij}^{(c)}.
\end{equation*}

Рёбра маскируются парами (прямое и обратное направление одновременно), чтобы сохранялась симметрия ориентированного представления молекулярного графа.

\subsection{Задача 3: свойства RDKit как слабосупервизируемый сигнал}

Дополнительно используется набор из 12 вычислимых дескрипторов RDKit

Каждый дескриптор стандартизуется по предварительно оценённым средним и стандартным отклонениям. Для батча из $B$ молекул:
\begin{equation*}
\mathcal{L}_{\text{prop}} = \frac{1}{B}\sum_{m=1}^{B}
\left\|\mathrm{MLP}_{\text{prop}}(\mathbf{g}_m)-\mathbf{p}_m\right\|^2,
\end{equation*}
где $\mathbf{g}_m \in \mathbb{R}^{d_{\text{readout}}}$ --- графовое представление из \texttt{PhysicsAwareReadout},
$\mathbf{p}_m \in \mathbb{R}^{12}$ --- нормализованный вектор дескрипторов.

Связь с основной задачей прямая: LogP, TPSA и числа доноров/акцепторов водородной связи статистически коррелируют с растворимостью, поэтому выходной блок уже на этапе 0 учится извлекать признаки, релевантные будущему предсказанию $\ln x_2$.

\subsection{Задача 4: контрастивное обучение (самосупервизируемое)}

Для каждой молекулы строятся две аугментированные версии графа с мягким искажением признаков: случайное обнуление порядка 15\% атомных и 15\% рёберных признаков \textbf{без удаления узлов и рёбер}, чтобы не разрушать топологию.

Далее формируются проекции:
\begin{equation*}
\mathbf{z}_1 = \mathrm{MLP}_{\text{proj}}(\mathbf{g}_{\text{aug}_1}),
\qquad
\mathbf{z}_2 = \mathrm{MLP}_{\text{proj}}(\mathbf{g}_{\text{aug}_2}).
\end{equation*}

Оптимизируется NT-Xent~\cite{chen2020}:
\begin{equation*}
\mathcal{L}_{\text{ctr}} = -\frac{1}{2B}\sum_{i=1}^{B}\left[
\log \frac{\exp(\mathrm{sim}(\mathbf{z}_{1,i}, \mathbf{z}_{2,i})/\tau)}
{\sum_{j=1}^{B}\exp(\mathrm{sim}(\mathbf{z}_{1,i}, \mathbf{z}_{2,j})/\tau)}
+ \log \frac{\exp(\mathrm{sim}(\mathbf{z}_{2,i}, \mathbf{z}_{1,i})/\tau)}
{\sum_{j=1}^{B}\exp(\mathrm{sim}(\mathbf{z}_{2,i}, \mathbf{z}_{1,j})/\tau)}
\right],
\end{equation*}
где
\begin{equation*}
\mathrm{sim}(\mathbf{a},\mathbf{b})=
\frac{\mathbf{a}^{\top}\mathbf{b}}{\|\mathbf{a}\|\,\|\mathbf{b}\|},
\qquad \tau=0.1.
\end{equation*}

Эта цель повышает устойчивость представления к умеренному шуму признаков и уменьшает переобучение на конкретную конфигурацию признаков.

\subsection{Суммарная функция потерь этапа 0}

Все четыре задачи объединяются линейной комбинацией:
\begin{equation*}
\mathcal{L}_{\text{Stage 0}} =
\lambda_{\text{atom}}\mathcal{L}_{\text{atom}} +
\lambda_{\text{bond}}\mathcal{L}_{\text{bond}} +
\lambda_{\text{prop}}\mathcal{L}_{\text{prop}} +
\lambda_{\text{ctr}}\mathcal{L}_{\text{ctr}}.
\end{equation*}

Веса по умолчанию:
\begin{equation*}
\lambda_{\text{atom}}=1.0,\quad
\lambda_{\text{bond}}=0.5,\quad
\lambda_{\text{prop}}=1.0,\quad
\lambda_{\text{ctr}}=0.5.
\end{equation*}

Такая балансировка делает акцент на сигналах восстановления и регрессии свойств при сохранении контрастивной регуляризации и контроля по типам связей.

\subsection{Связь этапа 0 с основным обучением (фазы 1--3)}

После завершения этапа 0 выполняется переход к стандартному учебному режиму:
\begin{enumerate}
  \item параметры \texttt{model.gnn} и \texttt{model.readout} сохраняются как инициализация целевой модели;
  \item временные головы \texttt{atom\_head}, \texttt{bond\_head}, \texttt{prop\_head}, \texttt{proj\_head} удаляются;
  \item обучение продолжается с фазы 1 на данных растворимости и вспомогательных функциях потерь по свойствам.
\end{enumerate}

В рабочем контуре проекта этап 0 рассматривается как поддерживаемая и рекомендуемая инициализация для сравнительных запусков семейства TGNN. При необходимости абляции он может быть отключён. Однако полный предобучительный прогон этапа 0 пока не входит в замороженный пакет результатов статьи, поэтому итоговые численные выводы по нему не вынесены в основные сравнительные таблицы.

\subsection{Ожидаемый эффект и ограничения}

Ожидаемые преимущества этапа 0:
\begin{itemize}
  \item более стабильная ранняя динамика градиентов;
  \item более быстрое достижение полезного режима на фазах 1 и 2;
  \item лучшее качество парных представлений при ограниченном числе уникальных растворяемых веществ в BigSolDB.
\end{itemize}

Ограничения и риски:
\begin{itemize}
  \item сдвиг домена между ZINC и химией реального корпуса растворимости может снижать переносимость;
  \item избыточная сила контрастивной цели может ухудшать чувствительность к целевой задаче;
  \item без корректной балансировки $\lambda$ возможен конфликт между задачами предобучения и основной целью.
\end{itemize}

Поэтому этап 0 следует рассматривать как контролируемое предварительное обучение: потенциально полезное, но требующее мониторинга метрик основной задачи после тонкой настройки.

\subsection{Эмпирическая оценка этапа 0}

Цель данного подраздела --- зафиксировать не список ожидаемых артефактов, а воспроизводимый протокол, по которому этап 0 должен быть признан полезным, нейтральным или вредным для основной задачи. Вся оценка строится как контролируемое сравнение двух режимов инициализации при совпадении бюджета фаз 1--3, состава данных, seed-набора и правил ранней остановки.

\keyidea{Этап 0 должен оцениваться не по внутренним pretraining-метрикам, а по тому, улучшает ли он конечную задачу предсказания $\ln x_2$ при том же вычислительном бюджете.}

Сравниваются два режима.
\begin{enumerate}
  \item \textbf{Baseline (cold-start):} параметры энкодера и readout инициализируются случайно по Xavier/Kaiming, после чего модель проходит стандартные фазы 1--3 в конфигурации, определённой в текущем учебном протоколе.
  \item \textbf{Treatment (warm-start):} параметры \texttt{model.gnn} и \texttt{model.readout} инициализируются из чекпойнта этапа 0, после чего выполняются те же самые фазы 1--3 без изменения числа эпох, scheduler, batch size, набора потерь и набора seed.
\end{enumerate}

Такое сравнение изолирует именно эффект предобучения. Любое отличие между режимами должно объясняться только наличием или отсутствием Stage 0, а не побочными изменениями оптимизации. Соответственно, итоговый отчёт обязан приводить результаты как минимум по validation и test, усреднённые по одной и той же серии независимых запусков.

Проверка делится на две стадии.

\textbf{Промежуточная диагностика} выполняется сразу после завершения Stage 0 и до старта фазы 1. На замороженных эмбеддингах обучается линейная probe-модель для четырёх ключевых дескрипторов: MolLogP, TPSA, NumHDonors и FractionCSP3. Для каждого дескриптора сравниваются два состояния представления: случайная инициализация и post-Stage-0. Целевая метрика --- $R^2$ на отложенной части данных. Смысл этого шага состоит в проверке, стал ли энкодер после предобучения лучше кодировать базовую химическую информацию, релевантную растворимости. Рабочее ожидание формулируется как гипотеза, а не как установленный факт: post-Stage-0 должен давать прирост $R^2$ порядка 0.1--0.3 для каждого из четырёх дескрипторов по сравнению со случайной инициализацией.

\textbf{Финальная диагностика} выполняется после завершения фазы 3 и определяет, переносится ли выигрыш в представлениях на конечную задачу. Здесь фиксируются три группы показателей. Во-первых, сравниваются $\mathrm{MAE}_{\mathrm{warm}}$ и $\mathrm{MAE}_{\mathrm{cold}}$ на validation и test в шкале $\ln x_2$. Во-вторых, сравнивается Pearson $r(T_{\mathrm{m}})$ для обоих режимов, поскольку улучшение инициализации должно, по предположению модели, облегчать восстановление кристаллических свойств, определённых в~\S2.1. В-третьих, анализируются кривые $\mathcal{L}_{\mathrm{sol}}$ по эпохам: их задача --- не заменить итоговую метрику, а визуально проверить, входит ли warm-start быстрее в рабочий режим и снижает ли нестабильность раннего обучения.

\begin{table}[H]
\small
\setlength{\tabcolsep}{5pt}
\caption{Протокол эмпирической оценки Stage 0}
\begin{tabularx}{\textwidth}{P{0.21\textwidth}P{0.28\textwidth}P{0.24\textwidth}Y}
\textbf{Этап} & \textbf{Что сравнивается} & \textbf{Метрика} & \textbf{Интерпретация} \\
\hline
Промежуточная диагностика & random init vs post-Stage-0 на замороженных эмбеддингах & $R^2$ для MolLogP, TPSA, NumHDonors, FractionCSP3 & Проверка, усилилось ли химическое содержание представлений до старта фаз 1--3 \\

Финальная диагностика & cold-start vs warm-start после фазы 3 & $\mathrm{MAE}$ на validation и test & Основной критерий полезности Stage 0 для задачи предсказания $\ln x_2$ \\

Физический прокси-контроль & cold-start vs warm-start после фазы 3 & Pearson $r(T_{\mathrm{m}})$ & Проверка, связан ли возможный выигрыш по растворимости с лучшим восстановлением кристаллических свойств \\

Динамика оптимизации & cold-start vs warm-start по ходу обучения & кривые $\mathcal{L}_{\mathrm{sol}}$ по эпохам & Визуальная проверка более быстрого входа в рабочий режим и меньшей ранней нестабильности \\
\end{tabularx}
\end{table}

Классификация результата должна быть задана заранее. Stage 0 считается \textbf{полезным}, если при том же бюджете выполняется неравенство $\mathrm{MAE}_{\mathrm{warm}} < \mathrm{MAE}_{\mathrm{cold}} - 0.1$. Stage 0 считается \textbf{вредным}, если наблюдается противоположный эффект, то есть $\mathrm{MAE}_{\mathrm{warm}} > \mathrm{MAE}_{\mathrm{cold}} + 0.05$. Промежуточная зона между этими порогами трактуется как \textbf{нейтральный} результат: в этом случае включение Stage 0 в финальный протокол определяется уже не качеством как таковым, а вычислительным бюджетом, стабильностью по seed и удобством последующего обучения.

\warnnote{На момент данного отчёта результаты не получены; после завершения Эксперимента 1 здесь должны быть приведены: (i) таблица $R^2$ для линейной probe по четырём дескрипторам, (ii) сравнение $\mathrm{MAE}_{\mathrm{warm}}$ и $\mathrm{MAE}_{\mathrm{cold}}$ на validation и test, (iii) Pearson $r(T_{\mathrm{m}})$ для обоих режимов, (iv) графики $\mathcal{L}_{\mathrm{sol}}$ по эпохам, (v) усреднение и разброс по одинаковому набору seed.}

\validated{Методологически этап 0 не должен считаться успешным только потому, что уменьшает внутренние pretraining-losses. Единственный достаточный критерий его практической ценности --- воспроизводимое улучшение downstream-качества по сравнению с cold-start при том же бюджете фаз 1--3.}

\openquestion{Достаточно ли близко химическое пространство ZINC250k к пространству BigSolDB, чтобы pretraining был полезен? Ответ зависит от overlap по Murcko scaffolds и от распределения функциональных групп.}

\section*{Часть IV. Обучение: учебный режим, функции потерь, регуляризация}
\addcontentsline{toc}{section}{Часть IV. Обучение: учебный режим, функции потерь, регуляризация}

\section{Трёхфазная учебная схема}

\subsection{Мотивация}

Прямое сквозное обучение, то есть от графа до $\ln x_2$ через SLE-решатель, создаёт деструктивный цикл обратной связи:

\begin{enumerate}
  \item На начальном этапе блоки кристаллических свойств предсказывают случайные $T_{\mathrm{m}}$ и $\Delta H_{\mathrm{fus}}$.
  \item SLE решатель получает некорректные входы $\to$ расходящиеся итерации $\to$ неинформативные градиенты.
  \item Блок NRTL получает бессмысленные градиенты $\to$ предсказывает случайные $\tau_{12}$ и $\tau_{21}$.
  \item Функция потерь по $\ln x_2$ даёт шумный сигнал, усугубляющий проблему.
\end{enumerate}

Поэтапная учебная схема разрывает этот цикл, обеспечивая разумную инициализацию перед подключением SLE.

\subsection{Фаза 1: предобучение вспомогательных выходных блоков}

\textbf{Длительность:} 50 эпох (каноническая конфигурация).

\textbf{Обучаемые компоненты:} GNN Encoder, FusionHead ($T_{\mathrm{m}}$, $\Delta H_{\mathrm{fus}}$), HansenHead ($\boldsymbol{\delta}$), AuxPropsHead ($V_m$).

\textbf{Замороженные компоненты:} NRTLHead, SLE-решатель, BoundedCorrection (gate заморожен в состоянии $w = 1$, т.е. коррекция выключена).

\textbf{Активные потери:}

\begin{equation*}
\mathcal{L}_{\mathrm{Phase\,1}} = \lambda_{T_{\mathrm{m}}}\mathcal{L}_{T_{\mathrm{m}}} + \lambda_{\Delta H}\mathcal{L}_{\Delta H} + \lambda_{\mathrm{Hansen}}\mathcal{L}_{\mathrm{Hansen}}
\end{equation*}

Для обычных solubility-batches фазы 1 канал $\mathcal{L}_{\gamma^\infty}$ не используется, поскольку в базовом корпусе IDAC-наблюдения отсутствуют (таблица~\ref{tab:data-aux-coverage}). В IDAC-augmented режиме этот loss подключается отдельными \texttt{aux\_only\_gamma}-батчами после разблокировки NRTL-ветки; поэтому в каноническом описании фазы 1 он остаётся выключенным.

\textbf{Потеря на $\ln x_2$:} отключена ($\lambda_{\mathrm{sol}} = 0$).

\textbf{Цель:} к концу фазы 1 блоки кристаллических свойств предсказывают $T_{\mathrm{m}}$ и $\Delta H_{\mathrm{fus}}$ с разумной точностью, обеспечивая SLE-решатель осмысленными входами для фазы 2.

\subsection{Фаза 2: полное обучение SLE-модели}
\label{subsec:phase2-full-training}

\textbf{Длительность:} 200 эпох.

\textbf{Обучаемые компоненты:} всё, кроме BoundedCorrection (который разблокируется на эпохе \texttt{phase2\_correction\_unfreeze\_epoch = 20} внутри фаза 2).

\textbf{Активные потери:}

\begin{equation*}
\mathcal{L}_{\mathrm{Phase\,2}} = \lambda_{\mathrm{sol}}\mathcal{L}_{\mathrm{sol}} + \sum_{k} \lambda_k \mathcal{L}_k^{\mathrm{aux}} + \sum_{m} \lambda_m \mathcal{L}_m^{\mathrm{reg}}
\end{equation*}

Вспомогательные функции потерь $\mathcal{L}_k^{\mathrm{aux}}$ имеют \textbf{уменьшенные веса} по сравнению с фазой 1, чтобы доминировала $\mathcal{L}_{\mathrm{sol}}$.

\textbf{Ранняя остановка:} по MAE на валидации для $\ln x_2$.

\subsection{Фаза 3: дообучение}

\textbf{Длительность:} 50 эпох.

\keyidea{Фаза 3 --- это не отдельное ``переизобретение'' модели, а этап инженерной полировки уже обученной системы. Её задача состоит в том, чтобы подавить артефакты переобучения, стабилизировать температурную монотонность, уменьшить зависимость от correction path и сделать промежуточные физические параметры более правдоподобными без заметной деградации основной метрики по $\ln x_2$.}

По сравнению с фазой 2 на фазе 3 не меняется общая архитектура и не вводятся новые источники супервизии; меняется только режим оптимизации и баланс регуляризаторов. Это принципиально важно: к началу фазы 3 модель уже должна находиться в рабочем бассейне параметров, поэтому цель этапа --- не ``научить заново'', а сдвинуть оптимум в сторону более устойчивого и физически правдоподобного решения.

Конкретные изменения относительно фазы 2 задаются таблицами~\ref{tab:optimizer-phase-hparams} и~\ref{tab:phase-loss-weights}. По оптимизатору выполняется десятикратное снижение learning rate,
\begin{equation*}
1 \times 10^{-4} \;\to\; 1 \times 10^{-5},
\end{equation*}
и одновременно десятикратное усиление weight decay,
\begin{equation*}
10^{-5} \;\to\; 10^{-4}.
\end{equation*}
Такое сочетание отражает классическую логику fine-tuning: параметры обновляются существенно более осторожно, но отклонения от уже найденного решения штрафуются сильнее. В результате модель меньше ``дрейфует'' по плоскому участку ландшафта и реже уходит в режим тонкой подгонки residual-пути под шум в обучающих данных.

По функциям потерь фаза 3 усиливает именно те компоненты, которые отвечают за структурные свойства предсказанной температурной кривой и за ограничение correction-блока. Вес монотонности повышается
\begin{equation*}
\lambda_{\mathrm{mono}}: 0.1 \to 0.5,
\end{equation*}
что делает подавление немонотонного поведения $\ln x_2(T)$ главным новым приоритетом этапа. Дополнительно усиливаются ранговый штраф
\begin{equation*}
\lambda_{\mathrm{rank}}: 0.1 \to 0.2,
\end{equation*}
и локальная van't-Hoff-согласованность
\begin{equation*}
\lambda_{\mathrm{vH}}: 0.05 \to 0.1,
\end{equation*}
но при обязательном сохранении стабилизаций clamp и cap из §\ref{subsec:vh-stabilization}. Иначе этот канал снова рискует перейти из режима мягкого структурного приоритета в режим доминирующего источника градиента.

Вторая группа изменений направлена на снижение зависимости модели от correction path. Для этого штраф за величину коррекции усиливается на порядок,
\begin{equation*}
\lambda_{\mathrm{res}}: 0.01 \to 0.1,
\end{equation*}
а штраф предпочтения физическому ядру повышается как
\begin{equation*}
\lambda_{\mathrm{phys\_pref}}: 0.01 \to 0.05.
\end{equation*}
Практический смысл этих двух изменений различен, но комплементарен. Усиление $\mathcal{L}_{\mathrm{res}}$ делает большие correction-сдвиги дорогими, тогда как рост $\mathcal{L}_{\mathrm{phys\_pref}}$ заставляет gate-блок чаще выбирать режим большего доверия SLE-ядру. Одновременно вспомогательные веса кристаллических и Hansen-каналов дополнительно ослабляются:
\begin{equation*}
\lambda_{T_{\mathrm{m}}}: 0.1 \to 0.05, \qquad
\lambda_{\Delta H}: 0.1 \to 0.05, \qquad
\lambda_{\mathrm{Hansen}}: 0.05 \to 0.02.
\end{equation*}
Это отражает тот факт, что на фазе 3 блоки свойств уже не должны радикально переучиваться; их роль --- поддерживать физический масштаб решения, а не конкурировать с основным solubility-channel за градиентный ресурс.

Процедурно фаза 3 работает как короткий режим controlled refinement. Горизонт в 50 эпох здесь важен не сам по себе, а как ограничение на глубину дрейфа: если модель за этот интервал не находит более физически устойчивое решение, дальнейшее продолжение оптимизации обычно приносит меньше пользы, чем риска. Поэтому по завершении фазы восстанавливается лучший чекпойнт по MAE на валидации для $\ln x_2$, а интерпретация результатов должна смотреть не только на численное значение MAE, но и на сопутствующие диагностические признаки --- долю correction path, гладкость температурных кривых и стабильность промежуточных параметров.

Иными словами, успешная фаза 3 --- это не обязательно заметное улучшение основной метрики, а переход в режим, где модель сохраняет качество фазы 2, но делает это более устойчивым, менее хрупким и более согласованным с физической структурой задачи. Именно поэтому данный этап разумно рассматривать как обязательную часть production-like training protocol, а не как опциональную косметическую донастройку.

\section{Функции потерь}

\subsection{Основная потеря: функция потерь Хьюбера для $\ln x_2$}

В качестве робастной регрессии используется классическая функция потерь Хьюбера~\cite{huber1964}.

\begin{equation*}
\mathcal{L}_{\mathrm{sol}} = \frac{1}{N_{\mathrm{sol}}} \sum_{i:\, \mathrm{has\_sol}} \mathrm{Huber}_\delta\!\left(\ln x_{2,i}^{\mathrm{pred}} - \ln x_{2,i}^{\mathrm{true}}\right)
\end{equation*}

где:

\begin{equation*}
\mathrm{Huber}_\delta(r) = \begin{cases} \frac{1}{2}r^2, & |r| \leq \delta \\ \delta\!\left(|r| - \frac{1}{2}\delta\right), & |r| > \delta \end{cases}
\end{equation*}

\textbf{Обоснование выбора функции потерь Хьюбера вместо MSE:} данные BigSolDB содержат выбросы (ошибки измерений, неверные единицы). Функция потерь Хьюбера устойчива к ним благодаря линейному хвосту при больших $|r|$.

Порог $\delta=1.0$ выбран как практический компромисс для текущей шкалы ошибки: при типичном уровне MAE порядка $1.5$--$2.0$ заметная часть наблюдений остаётся в квадратичной области (полезный «сильный» градиент), а выбросы с $|r|>1$ переходят в линейный режим и не доминируют оптимизацию.

\subsection{Потери на кристаллические свойства}

\begin{equation*}
\mathcal{L}_{T_{\mathrm{m}}} = \frac{1}{N_{T_{\mathrm{m}}}} \sum_{i:\, \mathrm{has}\_T_m} \left(T_{\mathrm{m},i}^{\mathrm{pred}} - T_{\mathrm{m},i}^{\mathrm{true}}\right)^2
\end{equation*}

\begin{equation*}
\mathcal{L}_{\Delta H} = \frac{1}{N_{\Delta H}} \sum_{i:\, \mathrm{has}\_\Delta H} \left(\Delta H_{\mathrm{fus},i}^{\mathrm{pred}} - \Delta H_{\mathrm{fus},i}^{\mathrm{true}}\right)^2
\end{equation*}

\textbf{Разреженная маскировка:} суммирование только по примерам с доступной экспериментальной меткой. В текущем supervised subset покрытие составляет 51.04\% для $T_{\mathrm{m}}$ и 1.00\% для $\Delta H_{\mathrm{fus}}$ (см. таблицу~\ref{tab:data-aux-coverage}).

\subsection{Потеря на Hansen параметры}

\begin{equation*}
\mathcal{L}_{\mathrm{Hansen}} = \frac{1}{N_{\mathrm{Hansen}}} \sum_{i:\, \mathrm{has}\_\mathrm{Hansen}} \sum_{k \in \{d, p, h\}} \left(\delta_{k,i}^{\mathrm{pred}} - \delta_{k,i}^{\mathrm{true}}\right)^2
\end{equation*}

\subsection{Потеря на коэффициент активности при бесконечном разбавлении}

Для внешних IDAC-данных используется специальная функция потерь:

\begin{equation*}
\mathcal{L}_{\gamma^\infty} = \frac{1}{N_{\gamma^\infty}} \sum_{i:\, \mathrm{has}\_\gamma^\infty} \left(\ln\gamma_{2,i}^{\infty,\mathrm{pred}} - \ln\gamma_{2,i}^{\infty,\mathrm{true}}\right)^2
\end{equation*}

где предсказанное значение вычисляется из NRTL аналитически через формулу предельного коэффициента активности при $x_2\to 0$:

\begin{equation*}
\ln\gamma_2^{\infty,\mathrm{pred}} = \tau_{12} + \tau_{21}\exp(-\alpha_{12}\tau_{21})
\end{equation*}

Покрытие обеспечивается внешним IDAC-корпусом (404 строки, 138 уникальных пар), опубликованным на Zenodo~\cite{idac_zenodo}. Корпус извлечён из NIST ThermoML и покрывает температурный диапазон 298--438 K. Несмотря на ограниченный размер стартового релиза, $\mathcal{L}_{\gamma^\infty}$ теперь является активной компонентой функции потерь, обеспечивающей прямую supervision на NRTL-параметры.

Механизм интеграции принципиально отделён от основного корпуса растворимости. IDAC-строки вводятся как \texttt{aux\_only\_gamma}: они участвуют только в $\mathcal{L}_{\gamma^\infty}$ и не требуют наличия метки $\ln x_2$. Благодаря этому внешний IDAC-корпус можно расширять независимо от основного solubility-корпуса, не ожидая совпадения по строкам с наблюдениями SLE. Именно в этом смысле $\mathcal{L}_{\gamma^\infty}$ даёт прямую супервизию на предел $x_2\to 0$ предсказанной NRTL-модели, тогда как solver-side $\ln\gamma_2(x_2)$ продолжает вычисляться уже внутри SLE-решателя при конечном составе.

\subsection{Связующий штраф (Hansen--NRTL консистентность)}
\label{subsec:bridge-loss}

\begin{equation*}
\mathcal{L}_{\mathrm{bridge}} = \sum_{j \in \{12, 21\}} \left(\Delta g_j^{\mathrm{pred}} - \Delta g_j^{\mathrm{(Hansen)}}\right)^2
\end{equation*}

где оценка из Regular Solution Theory:

\begin{equation*}
\Delta g_{12}^{\mathrm{(Hansen)}} = V_1 \left[(\delta_{d,1} - \delta_{d,2})^2 + \kappa_p(\delta_{p,1} - \delta_{p,2})^2 + \kappa_h(\delta_{h,1} - \delta_{h,2})^2\right]
\end{equation*}

\textbf{Важное замечание:} эта оценка всегда $\geq 0$, но реальные $\Delta g$ могут быть отрицательными (то есть для систем с $\gamma_2 < 1$, где растворение энергетически выгодно). Для таких систем bridge-функция потерь является \textbf{контрпродуктивной} --- см. анализ в Части V.

\warnnote{Bridge-loss опирается на Regular Solution Theory и систематически смещает обучение для систем со специфическими взаимодействиями (водородные связи, комплексообразование), где эффективный $\Delta g$ может быть отрицательным. По рабочей оценке это затрагивает порядка 30--50\% фармацевтически релевантных пар \textit{solute--solvent}.}

\subsection{Штраф за нарушение монотонности}

Для подавляющего большинства систем растворимость растёт с температурой:

\begin{equation*}
\frac{d\ln x_2}{dT} > 0
\end{equation*}

Штраф:

\begin{equation*}
\mathcal{L}_{\mathrm{mono}} = \frac{1}{N} \sum_i \max\!\left(0, -\frac{\partial \ln x_{2,i}^{\mathrm{pred}}}{\partial T}\right)^2
\end{equation*}

Производная вычисляется с помощью автоматического дифференцирования (дифференцирование SLE-решателя по $T$ включает вклады как от $\Phi(T)$, так и от $\tau(T)$).

\subsection{Температурные регуляризаторы для одной и той же пары}

При pair-aware batching в одном минибатче присутствуют несколько температур для одной пары (solute, solvent).

\textbf{Ранговая функция потерь:}

\begin{equation*}
\mathcal{L}_{\mathrm{rank}} = \sum_{\substack{(i,j):\, \mathrm{same\;pair} \\ T_i < T_j}} \max\!\left(0, \ln x_{2,i}^{\mathrm{pred}} - \ln x_{2,j}^{\mathrm{pred}} + \epsilon_{\mathrm{margin}}\right)
\end{equation*}

Штрафует нарушение порядка: при $T_i < T_j$ должно быть $\ln x_{2,i} < \ln x_{2,j}$.

\textbf{Локальная согласованность по ван'т-Гоффу:}

Введём целевой наклон
\begin{equation*}
s_{\mathrm{vH,target}} \equiv -\frac{\Delta H_{\mathrm{sol}}}{R},
\end{equation*}
где $\Delta H_{\mathrm{sol}}$ --- эффективная энтальпия растворения в ван'т-гоффовском приближении.

\begin{equation*}
\mathcal{L}_{\mathrm{vH}} = \sum_{\substack{(i,j):\, \mathrm{same\;pair} \\ T_i \neq T_j}} \left(\frac{\ln x_{2,j}^{\mathrm{pred}} - \ln x_{2,i}^{\mathrm{pred}}}{1/T_j - 1/T_i} - s_{\mathrm{vH,target}}\right)^2
\end{equation*}

Штрафует отклонение от линейности $\ln x_2$ в зависимости от $1/T$ согласно уравнению van 't Hoff в приближении квазипостоянной энтальпии растворения.

\subsection{Регуляризаторы коррекции}

\begin{equation*}
\mathcal{L}_{\mathrm{res}} = \frac{1}{N}\sum_i \left|\ln x_{2,i}^{\mathrm{final}} - \ln x_{2,i}^{\mathrm{physics}}\right|^2
\end{equation*}

Штрафует величину коррекции --- поощряет модель «доверять физике».

\begin{equation*}
\mathcal{L}_{\mathrm{phys\_pref}} = \frac{1}{N}\sum_i (1 - w_i)^2
\end{equation*}

Штрафует отклонение gate от $w = 1$ (полное доверие физике).

\subsection{Регуляризатор параметров NRTL}
\label{subsec:tau-regularizer}

\begin{equation*}
\mathcal{L}_{\mathrm{tau\_reg}} = \frac{1}{N}\sum_i \left(\tau_{12,i}^2 + \tau_{21,i}^2\right)
\end{equation*}

Мягкое предпочтение малых $|\tau|$ --- ближе к идеальному раствору. Без этого NRTL может уйти в экстремальные значения.

\subsection{Полная потеря}

\begin{equation*}
\mathcal{L} = \lambda_{\mathrm{sol}}\mathcal{L}_{\mathrm{sol}} + \lambda_{T_{\mathrm{m}}}\mathcal{L}_{T_{\mathrm{m}}} + \lambda_{\Delta H}\mathcal{L}_{\Delta H} + \lambda_{\mathrm{Hansen}}\mathcal{L}_{\mathrm{Hansen}} + \lambda_{\gamma^\infty}\mathcal{L}_{\gamma^\infty} + \lambda_{\mathrm{bridge}}\mathcal{L}_{\mathrm{bridge}}
\end{equation*}

\begin{equation*}
+ \lambda_{\mathrm{mono}}\mathcal{L}_{\mathrm{mono}} + \lambda_{\mathrm{rank}}\mathcal{L}_{\mathrm{rank}} + \lambda_{\mathrm{vH}}\mathcal{L}_{\mathrm{vH}} + \lambda_{\mathrm{res}}\mathcal{L}_{\mathrm{res}} + \lambda_{\mathrm{phys\_pref}}\mathcal{L}_{\mathrm{phys\_pref}} + \lambda_{\mathrm{tau\_reg}}\mathcal{L}_{\mathrm{tau\_reg}}
\end{equation*}

В общей архитектурной записи слагаемое $\lambda_{\gamma^\infty}\mathcal{L}_{\gamma^\infty}$ соответствует текущей IDAC-augmented конфигурации. Оно не получает вклада от базового solubility-supervised корпуса (таблица~\ref{tab:data-aux-coverage}), но получает его от внешних \texttt{aux\_only\_gamma}-наблюдений.

Архитектура поддерживает 12 слагаемых с фазозависимыми весами; в текущей IDAC-augmented постановке все 12 могут быть активны, но $\mathcal{L}_{\gamma^\infty}$ питается отдельным внешним потоком данных и потому не отражается в coverage-таблице основного solubility-supervised корпуса. Отдельно используется диагностическая проверка консистентности Гиббса--Дюгема (см. Приложение B.5), которая служит служебной диагностической проверкой и не входит в основной оптимизируемый функционал.

\subsection{Физический смысл баланса потерь и роль коэффициентов $\lambda$}

Суммарная функция потерь задаёт \textbf{многокритериальную оптимизацию}, в которой отдельные слагаемые соответствуют разным физическим и статистическим требованиям к модели:

\begin{equation*}
\min_{\theta} \mathcal{L}(\theta) = \sum_j \lambda_j \mathcal{L}_j(\theta)
\end{equation*}

Здесь $\theta$ --- все обучаемые параметры (графовый энкодер, property heads, параметры NRTL-голова предсказания, correction/gate-блок), а веса $\lambda_j$ задают компромисс между:

\begin{itemize}
    \item точностью прогноза растворимости ($\mathcal{L}_{\mathrm{sol}}$),
    \item согласованностью с термодинамическими прокси-переменными ($\mathcal{L}_{T_{\mathrm{m}}},\; \mathcal{L}_{\Delta H},\; \mathcal{L}_{\mathrm{Hansen}}$, а также $\mathcal{L}_{\gamma^\infty}$ при наличии IDAC-данных),
    \item выполнением физически мотивированных ограничений формы ($\mathcal{L}_{\mathrm{mono}},\; \mathcal{L}_{\mathrm{rank}},\; \mathcal{L}_{\mathrm{vH}}$),
    \item контролем сложности и величины «подгонки» ($\mathcal{L}_{\mathrm{res}},\; \mathcal{L}_{\mathrm{phys\_pref}},\; \mathcal{L}_{\mathrm{tau\_reg}}$).
\end{itemize}

С точки зрения градиентной динамики каждый шаг AdamW использует сумму вкладов:

\begin{equation*}
\nabla_\theta \mathcal{L} = \sum_j \lambda_j \nabla_\theta \mathcal{L}_j
\end{equation*}

Поэтому выбор $\lambda_j$ не только масштабирует численное значение метрик, но и напрямую определяет \textbf{направление обновления параметров}. Слишком большой вес физического регуляризатора может ухудшать согласование с данными, тогда как чрезмерный акцент на $\mathcal{L}_{\mathrm{sol}}$ повышает риск нефизичных зависимостей по температуре и неустойчивых NRTL-параметров.

Практически фазовая схема обучения реализует следующий физический сценарий:

\begin{enumerate}
    \item \textbf{фаза 1:} формируется «материальная база» (кристаллические и Hansen-свойства), чтобы задать реалистичный масштаб свободноэнергетических вкладов.
    \item \textbf{фаза 2:} основной упор переносится на растворимость, но регуляризаторы сохраняют термодинамическую правдоподобность.
    \item \textbf{фаза 3:} усиливаются ограничения гладкости/монотонности и малая correction, что снижает переобучение и стабилизирует экстраполяцию по $T$.
\end{enumerate}

Таким образом, веса потерь интерпретируются как \textbf{коэффициенты доверия} к различным источникам информации: экспериментальным значениям растворимости, вспомогательным property-labels и априорным физическим законам.


\section{Балансировка функций потерь: методология}

\subsection{Проблема многокритериальной оптимизации}

В TGNN-Solv суммарная цель обучения включает 12+ компонент с разными физическими единицами, масштабами и статистической структурой ошибок. Формально:
\begin{equation*}
\mathcal{L}_{\mathrm{total}} = \sum_k \lambda_k \mathcal{L}_k.
\end{equation*}

Такая постановка неизбежно создаёт конфликт градиентов:
\begin{equation*}
\nabla_\theta \mathcal{L}_{\mathrm{total}} = \sum_k \lambda_k \nabla_\theta \mathcal{L}_k,
\end{equation*}
где отдельные слагаемые могут тянуть параметры в противоположные стороны. Если масштаб одной компоненты неконтролируемо растёт, она начинает монополизировать градиентный бюджет, и оптимизация перестаёт решать основную задачу (точность по $\ln x_2$).

Практический вывод: для физически-информированной модели важна не только корректность формул, но и численно устойчивая \textbf{методология балансировки компонент функции потерь}.

\subsection{Кейс катастрофы: резкий рост компоненты van 't Hoff в функции потерь}

Один из самых показательных отказов связан с \texttt{vant\_hoff\_local}. Напомним форму штрафа:
\begin{equation*}
\mathcal{L}_{\text{vH}} = \sum_{\substack{(i,j):\,\text{same pair}\\T_i \neq T_j}}
\left(\frac{\ln x_{2,j} - \ln x_{2,i}}{1/T_j - 1/T_i} - s_{\mathrm{vH,target}}\right)^2.
\end{equation*}

Критическая численная особенность задаётся знаменателем:
\begin{equation*}
\left|\frac{1}{T_j} - \frac{1}{T_i}\right| = \frac{|T_i - T_j|}{T_i T_j}.
\end{equation*}

Если температуры близки (например, $T_i=298$ K и $T_j=303$ K), получаем:
\begin{equation*}
\left|\frac{1}{T_j} - \frac{1}{T_i}\right| \approx 5.5\times 10^{-5}\;\mathrm{K}^{-1}.
\end{equation*}

Даже небольшое отличие в числителе (порядка $0.01$ по $\ln x_2$) превращается в большой slope-терм (порядка $\sim 180$), что даёт per-pair штраф порядка $\sim 3.2\times 10^4$. При большом числе пар в батче это мгновенно ломает баланс общей цели.

Эмпирическая динамика в проблемном запуске:
\begin{itemize}
  \item Epoch 0, фаза 2: $\mathcal{L}_{\text{vH}}=3811$, при $\mathcal{L}_{\mathrm{sol}}\approx 1.1$;
  \item к Epoch 10: $\mathcal{L}_{\text{vH}}=240148$ (рост в $\sim 63\times$);
  \item общий train функции потерь: $8.8 \rightarrow 481$ (рост в $\sim 54\times$);
  \item val MAE не улучшается: $2.20 \rightarrow 2.27$.
\end{itemize}

Интерпретация: оптимизатор тратит почти весь градиентный бюджет на \texttt{vH}-компоненту и фактически игнорирует главную задачу регрессии растворимости.

\subsection{Трёхуровневая стабилизация vH-компоненты}\label{subsec:vh-stabilization}

Для устранения резкого роста использовалась каскадная защита на трёх уровнях.

\textbf{(a) Численно безопасный знаменатель (clamp).}
\begin{equation*}
\Delta_{\text{safe}} = \operatorname{sign}\!\left(\frac{1}{T_j} - \frac{1}{T_i}\right)
\cdot \max\!\left(\left|\frac{1}{T_j} - \frac{1}{T_i}\right|,\; \epsilon\right),
\quad \epsilon = 10^{-4}.
\end{equation*}

\textbf{(b) Ограничение per-pair вклада (cap).}
\begin{equation*}
\ell_{ij}^{\text{vH}} = \min\!\left(\ell_{ij}^{\text{vH,raw}},\; C_{\max}\right),
\quad C_{\max}=100.
\end{equation*}

\textbf{(c) Снижение веса компоненты.}
\begin{equation*}
\lambda_{\text{vH}} = 10^{-4}
\end{equation*}
вместо прежнего уровня порядка $10^{-2}$.

Комбинация этих мер переводит \texttt{vH}-штраф из режима «жёсткого ограничителя» в режим «мягкого структурного приоритета».

\subsection{Принцип \texttt{sol\_fraction} как основной safeguard}
\label{subsec:sol-fraction}

Для оперативного контроля баланса была введена скалярная метрика доли основной задачи:
\begin{equation*}
f_{\text{sol}} = \frac{\lambda_{\text{sol}}\,\mathcal{L}_{\text{sol}}}{\mathcal{L}_{\text{total}}}.
\end{equation*}

Рабочее правило:
\begin{equation*}
f_{\text{sol}} > 0.5 \quad \text{в каждую эпоху}.
\end{equation*}

Если наблюдается провал ($f_{\text{sol}}<0.3$), активируется динамический safeguard с масштабированием всех неосновных компонент:
\begin{equation*}
\mathcal{L}_{\text{total}}' = \lambda_{\text{sol}}\mathcal{L}_{\text{sol}} +
\alpha\cdot\left(\mathcal{L}_{\text{total}} - \lambda_{\text{sol}}\mathcal{L}_{\text{sol}}\right),
\end{equation*}
\begin{equation*}
\alpha =
\begin{cases}
\dfrac{f_{\text{sol}}}{1-f_{\text{sol}}}\cdot 0.5, & f_{\text{sol}}<0.3,\\
1, & f_{\text{sol}}\ge 0.3.
\end{cases}
\end{equation*}

Интуитивно это механизм «автоподстройки», который временно подавляет регуляризаторы, когда они начинают перехватывать оптимизацию.

\keyidea{Практическое правило: доля основной задачи \texttt{sol\_fraction} должна удерживаться выше 50\% на протяжении обучения. Падение ниже 30\% является рабочим триггером пересмотра весов регуляризаторов и активации safeguard-масштабирования.}

\subsection{Результат после стабилизации}

После внедрения fix-пакета наблюдалась принципиально иная динамика:
\begin{itemize}
  \item фаза 2, epoch 0: \texttt{sol\_fraction} $=92.6\%$;
  \item взвешенный вклад vH: $\mathcal{L}_{\text{vH,weighted}}\approx 10^{-4}$ (против $\sim 190$ до исправления);
  \item val MAE: $1.973$ (против $2.199$ до fix);
  \item $R^2$: $0.216$ (против $0.046$ до fix).
\end{itemize}

Качественный итог: модель впервые уверенно вошла в режим $R^2>0$, то есть стала объяснять вариацию целевой переменной лучше постоянного базового ориентира.

\subsection{Обобщённые рекомендации для физически-информированного обучения}

На основе этого примера сформулированы практические правила.

\begin{enumerate}
  \item Основная задача должна доминировать: доля целевой функции потерь $>50\%$.
  \item Регуляризаторы следует трактовать как \textit{мягкие подсказки}, а не \textit{жёсткие ограничения}.
  \item Логировать каждую компоненту одновременно в двух формах: невзвешенной и взвешенной.
  \item Любую компоненту функции потерь, растущую более чем на $10\times$ за несколько эпох, считать красным флагом.
  \item При любом изменении архитектуры/гиперпараметров обязательно проводить аудит декомпозиции функции потерь.
\end{enumerate}

В совокупности эта методология превращает балансировку функции потерь из ручной эвристики в воспроизводимую инженерную процедуру с понятными инвариантами контроля.

\section{Оракульная подстановка известных кристаллических свойств при обучении}
\label{sec:oracle-teacher-forcing}

\subsection{Мотивация}

Одна из центральных проблем TGNN-Solv --- компенсаторная дегенерация: NRTL-head может «абсорбировать» ошибки crystal-heads, подстраивая $\tau$ так, чтобы компенсировать неверные $T_{\mathrm{m}}$ и $\Delta H_{\mathrm{fus}}$ на тренировочных парах. В результате на новых молекулах одновременно ошибаются и кристаллические параметры, и вклад неидеальности, что приводит к резкому падению качества.

Идея оракульной подстановки состоит в том, чтобы \textbf{разорвать этот цикл} на этапе обучения: если для примера известны истинные кристаллические свойства, именно они подаются в SLE-решатель вместо предсказанных сетью значений.

\subsection{Формализация подстановки}

Пусть
\begin{equation*}
\mathcal{S}_T = \{i:\;\mathrm{has\_}T_{\mathrm{m}}(i)=1\}
\end{equation*}
--- множество примеров с доступной меткой $T_{\mathrm{m}}$. Тогда вход решатель-а задаётся правилом
\begin{equation*}
T_{\mathrm{m}}^{\mathrm{solver}}(i)=
\begin{cases}
T_{\mathrm{m}}^{\mathrm{true}}(i), & i\in\mathcal{S}_T,\\
T_{\mathrm{m}}^{\mathrm{pred}}(i), & i\notin\mathcal{S}_T.
\end{cases}
\end{equation*}

Аналогично вводится piecewise-подстановка для $\Delta H_{\mathrm{fus}}$ (и при необходимости для $\Delta C_p^{\mathrm{fus}}$).

Критически важно разделять две роли:
\begin{itemize}
  \item $T_{\mathrm{m}}^{\mathrm{solver}}$ --- значение, используемое в физическом ядре SLE;
  \item $T_{\mathrm{m}}^{\mathrm{pred}}$ --- выход головы, участвующий во вспомогательной потере $\mathcal{L}_{T_{\mathrm{m}}}$.
\end{itemize}

Такое разделение меняет градиентный поток от основной задачи:
\begin{itemize}
  \item для $i\in\mathcal{S}_T$ градиент $\partial\mathcal{L}_{\mathrm{sol}}/\partial T_{\mathrm{m}}$ не проходит в блок кристаллических свойств (оракульный вход отсоединён от графа);
  \item для $i\notin\mathcal{S}_T$ градиент проходит как обычно через предсказанный $T_{\mathrm{m}}$.
\end{itemize}

При этом градиенты от $\mathcal{L}_{T_{\mathrm{m}}}$ и $\mathcal{L}_{\Delta H}$ продолжают обучать блок кристаллических свойств на всех доступных метках.

\subsection{Стохастическая оракульная подстановка}

Чтобы избежать жёсткой зависимости от сигнала-подсказки, применяется стохастическая версия: для наблюдений с оракульной меткой истинное значение подаётся только с вероятностью $p$.

\begin{equation*}
T_{\mathrm{m}}^{\mathrm{solver}}(i)=
\begin{cases}
T_{\mathrm{m}}^{\mathrm{true}}(i), & i\in\mathcal{S}_T\;\wedge\;u_i<p,\\
T_{\mathrm{m}}^{\mathrm{pred}}(i), & \text{иначе},
\end{cases}
\end{equation*}
где $u_i\sim\mathrm{Uniform}(0,1)$.

Тот же механизм используется для $\Delta H_{\mathrm{fus}}$. На практике это даёт компромисс между стабильностью обучения (через оракульную подстановку) и адаптацией к собственным предсказаниям модели.

\subsection{Расписание постепенного ослабления}

Во второй половине фазы 2 вероятность оракульной подстановки постепенно снижается:
\begin{equation*}
p(t)=\max\left(0,\;1-\frac{t-t_{\mathrm{anneal\_start}}}{t_{\mathrm{anneal\_end}}-t_{\mathrm{anneal\_start}}}\right).
\end{equation*}

Таким образом, обучение плавно переходит от режима «с подсказками» к полностью самосогласованному режиму. В фазе 3 оракульная подстановка полностью выключена:
\begin{equation*}
p=0.
\end{equation*}

Эта схема уменьшает риск резкого сдвига распределения между поведением во время обучения и во время оценки.

\subsection{Связь с принудительной подстановкой истинных значений (teacher forcing)}
\label{subsec:teacher-forcing}

Методологически оракульная подстановка близка к teacher forcing в последовательностных моделях:
\begin{itemize}
  \item в seq2seq подаётся истинный предыдущий токен вместо предсказанного;
  \item здесь подаются истинные физические параметры вместо предсказаний блока кристаллических свойств.
\end{itemize}

Преимущество в обоих случаях --- ускорение и стабилизация раннего обучения. Риск также общий: смещение из-за расхождения режимов обучения и применения, когда модель избыточно привыкает к «идеальным» входам. Именно поэтому постепенное уменьшение $p$ является обязательным элементом методики.

\subsection{Оракульная подстановка при применении модели: только диагностический режим}

В рабочем режиме применения модели всегда используются предсказанные параметры:
\begin{equation*}
T_{\mathrm{m}}^{\mathrm{solver}}=T_{\mathrm{m}}^{\mathrm{pred}},\qquad
\Delta H_{\mathrm{fus}}^{\mathrm{solver}}=\Delta H_{\mathrm{fus}}^{\mathrm{pred}}.
\end{equation*}

Однако после обучения полезен диагностический прогон с принудительной оракульной подстановкой, где подставляются истинные кристаллические параметры. Разность
\begin{equation*}
\mathrm{MAE}_{\mathrm{standard}}-\mathrm{MAE}_{\mathrm{oracle}}
\end{equation*}
интерпретируется как вклад ошибок блока кристаллических свойств в итоговую ошибку по растворимости.

\subsection{Эмпирические наблюдения и интерпретация}

На коротком прокси-сценарии часто наблюдается
\begin{equation*}
\mathrm{MAE}_{\mathrm{oracle}}\approx\mathrm{MAE}_{\mathrm{standard}},
\end{equation*}
что согласуется с компенсаторной дегенерацией: NRTL-компонента уже успевает «скрыть» часть ошибок блока кристаллических свойств.

Ожидание для среднего и полного бюджета:
\begin{itemize}
  \item если блок кристаллических свойств остаётся узким местом, то $\mathrm{MAE}_{\mathrm{oracle}}<\mathrm{MAE}_{\mathrm{standard}}$;
  \item если даже при полном бюджете $\mathrm{MAE}_{\mathrm{oracle}}\approx\mathrm{MAE}_{\mathrm{standard}}$, основное узкое место смещается в NRTL-параметризацию, а не в блок кристаллических свойств.
\end{itemize}

Тем самым оракульная подстановка служит одновременно инструментом стабилизации обучения и строгим диагностическим тестом источника ошибки.

\subsection{Анализ сходимости обучения: протокол и индикаторы}
\label{subsec:training-convergence-protocol}

Наряду с финальными метриками на тесте для TGNN-Solv принципиально важно анализировать саму \textbf{траекторию обучения}. Из-за многокомпонентной функции потерь, фазового переключения весов и возможности компенсаторной дегенерации уменьшение одной только суммарной потери ещё не означает, что модель движется к лучшему решению. Возможна ситуация, когда $\mathcal{L}_{\mathrm{total}}$ на обучении снижается, но при этом растёт зависимость от correction path, ухудшается физическая правдоподобность $\tau$-параметров или стагнирует качество на валидации. Поэтому для каждого запуска должен сохраняться единый набор кривых по эпохам с явной разметкой границ между фазами 1, 2 и 3.

\textbf{Набор отслеживаемых кривых.} Минимальный диагностический пакет должен включать одновременно метрики оптимизации, внешнего качества и внутренней физической согласованности.
\begin{itemize}
  \item $\mathcal{L}_{\mathrm{total}}^{\mathrm{train}}$ и $\mathcal{L}_{\mathrm{sol}}^{\mathrm{train}}$ --- суммарная и целевая компоненты на обучении. Первая показывает общую динамику objective, вторая --- насколько быстро модель действительно улучшает основную задачу по $\ln x_2$. При этом $\mathcal{L}_{\mathrm{total}}^{\mathrm{train}}$ нельзя интерпретировать в отрыве от фазового расписания: на границе фаз абсолютный уровень суммарной потери меняется не только из-за параметров модели, но и из-за изменения коэффициентов $\lambda$.
  \item $\mathrm{MAE}_{\mathrm{val}}(\ln x_2)$ и $R^2_{\mathrm{val}}$ --- основные внешние метрики на валидации. MAE остаётся главным критерием ранней остановки и выбора чекпойнта; $R^2$ полезен как дополнительный индикатор того, стала ли модель объяснять межсистемную вариативность, а не только уменьшать среднюю ошибку. В ранних эпохах $R^2_{\mathrm{val}}$ может быть шумным, поэтому его следует читать совместно с MAE, а не вместо него.
  \item \texttt{sol\_fraction} --- индикатор баланса потерь из §\ref{subsec:sol-fraction},
  \begin{equation*}
  \mathrm{sol\_fraction} = \frac{\lambda_{\mathrm{sol}}\mathcal{L}_{\mathrm{sol}}}{\mathcal{L}_{\mathrm{total}}}.
  \end{equation*}
  Его смысл не в абсолютной оптимизационной ценности, а в контроле того, действительно ли основная задача по растворимости остаётся доминирующим источником градиента. Именно эта величина лучше всего выявляет режим, в котором вспомогательные регуляризаторы начинают «захватывать» обучение.
  \item $\mathcal{L}_{\mathrm{tau\_reg}}$ --- индикатор физической правдоподобности NRTL-параметров в духе §§\ref{subsec:tau-regularizer} и \ref{subsec:phase2-full-training}. Эта кривая не обязана монотонно падать на каждом отрезке, но её устойчивый рост при улучшении train-loss часто означает, что модель компенсирует ошибку по растворимости уходом в экстремальные $\tau_{12},\tau_{21}$, а не в более реалистичную термодинамическую декомпозицию.
  \item $\mathrm{MAE}(T_{\mathrm{m}})$ и коэффициент Пирсона $r(T_{\mathrm{m}})$ на подмножестве с доступными метками --- основные показатели качества crystal-head. Именно пара «ошибка + корреляция» полезнее одной MSE: MAE отвечает за практический масштаб ошибки, а Pearson $r$ показывает, начал ли блок извлекать ранговую и структурную информацию, а не просто приближаться к среднему значению. Для $\Delta H_{\mathrm{fus}}$ аналогичная диагностика желательна, но из-за сильной разреженности меток она менее устойчива и поэтому не должна быть главным convergence-сигналом.
  \item $\mathcal{L}_{\mathrm{vH}}$ в двух версиях --- невзвешенной и взвешенной. Невзвешенная кривая $\mathcal{L}_{\mathrm{vH}}$ нужна как чистый индикатор численной стабильности pair-wise температурного штрафа, тогда как взвешенная величина $\lambda_{\mathrm{vH}}\mathcal{L}_{\mathrm{vH}}$ показывает его реальный вклад в total loss. Такое разделение принципиально важно: при смене фазы рост взвешенного члена может объясняться не ухудшением модели, а лишь изменением $\lambda_{\mathrm{vH}}$.
  \item Среднее использование correction path,
  \begin{equation*}
  \overline{|\delta_{\mathrm{corr}}|} = \frac{1}{N}\sum_{i=1}^{N}|\delta_{\mathrm{corr},i}|,
  \end{equation*}
  --- это прямой индикатор того, насколько модель опирается на residual-механизм вместо физического ядра. Для TGNN-Solv эта кривая особенно важна в поздней фазе 2 и на фазе 3: если качество на валидации почти не улучшается, а $\overline{|\delta_{\mathrm{corr}}|}$ продолжает расти, модель, вероятно, уходит в хрупкий режим подгонки.
\end{itemize}

Практически все перечисленные кривые целесообразно строить на одной временной оси по эпохам, а границы фаз отмечать вертикальными линиями. Это особенно важно для \texttt{sol\_fraction}, $\mathcal{L}_{\mathrm{vH}}$ и $\mathcal{L}_{\mathrm{total}}$: без фазовых меток резкие скачки на графиках легко принять за нестабильность оптимизации, хотя часть из них объясняется плановым изменением весов потерь.

\textbf{Ожидаемая динамика по фазам.} Для корректно идущего обучения ожидаются три характерных режима. \textit{Фаза 1} должна демонстрировать быстрое снижение $\mathcal{L}_{T_{\mathrm{m}}}$ и $\mathcal{L}_{\Delta H}$ при полностью отключённой solubility-компоненте, то есть
\begin{equation*}
\lambda_{\mathrm{sol}}=0 \quad \Rightarrow \quad \mathrm{sol\_fraction}=0.
\end{equation*}
В этот период качество по $\ln x_2$ ещё не является главным критерием, а ключевой сигнал --- рост качества crystal-head: разумно ожидать, что Pearson $r(T_{\mathrm{m}})$ поднимется примерно от нуля к диапазону $0.3$--$0.6$, а MAE$(T_{\mathrm{m}})$ заметно уменьшится. Если этого не происходит, переход к solubility-training будет преждевременным.

\textit{Ранняя фаза 2} обычно характеризуется резким падением $\mathcal{L}_{\mathrm{sol}}$, быстрым выходом \texttt{sol\_fraction} на уровень выше $0.5$ и заметным улучшением $\mathrm{MAE}_{\mathrm{val}}(\ln x_2)$. В этот момент модель впервые начинает использовать кристаллические предсказания и NRTL-параметры как совместный механизм подгонки под целевую растворимость. Именно поэтому на раннем отрезке фазы 2 допустим умеренный рост $\mathcal{L}_{\mathrm{tau\_reg}}$ и даже временное увеличение $\overline{|\delta_{\mathrm{corr}}|}$: система ищет рабочую декомпозицию ошибки. Однако этот рост должен быть ограниченным и сопровождаться реальным выигрышем на валидации; иначе возникает подозрение, что оптимизация сразу уходит в компенсаторный режим.

\textit{Поздняя фаза 2 и фаза 3} --- это режим замедления и селективной стабилизации. В нормальном случае train-loss продолжает снижаться медленнее, $\mathrm{MAE}_{\mathrm{val}}(\ln x_2)$ выходит на плато или улучшается уже только на небольшую величину, а $\mathcal{L}_{\mathrm{tau\_reg}}$, невзвешенная $\mathcal{L}_{\mathrm{vH}}$ и среднее $|\delta_{\mathrm{corr}}|$ перестают расти. Для фазы 3 дополнительным позитивным признаком считается снижение зависимости от correction path вслед за усилением $\mathcal{L}_{\mathrm{res}}$ и $\mathcal{L}_{\mathrm{phys\_pref}}$: модель должна по возможности сохранять качество фазы 2 при меньшем residual-вкладе и более гладких температурных кривых. Если же на этом этапе $\mathcal{L}_{\mathrm{total}}^{\mathrm{train}}$ и $\mathcal{L}_{\mathrm{sol}}^{\mathrm{train}}$ продолжают улучшаться, но $\mathrm{MAE}_{\mathrm{val}}$ стагнирует или растёт, это уже классический паттерн переобучения, а не полезной физической донастройки.

\textbf{Индикаторы проблем.} Для оперативной диагностики полезно использовать не только визуальный осмотр кривых, но и компактный набор эвристических порогов. Они не являются универсальными физическими константами, однако хорошо работают как инженерные триггеры для пересмотра режима обучения.

\begin{table}[H]
\centering
\small
\setlength{\tabcolsep}{4pt}
\begin{tabularx}{\textwidth}{|P{3.2cm}|P{2.0cm}|Y|Y|}
\hline
Индикатор & Порог & Интерпретация & Действие \\
\hline
\texttt{sol\_fraction} & $< 0.3$ & Красный флаг: регуляризатор или вспомогательные члены захватили оптимизацию, и основная задача по растворимости перестала доминировать. & Активировать safeguard из §\ref{subsec:sol-fraction} и перескалировать веса вспомогательных потерь. \\
\hline
Train/val MAE gap & $> 0.5$ & Warning: формируется заметное переобучение, даже если train-loss формально улучшается. & Увеличить dropout и/или weight decay, проверить раннюю остановку. \\
\hline
$\mathcal{L}_{\mathrm{tau\_reg}}$ & $> 3.0$ & Warning: модель уходит в экстремальные $\tau$-параметры и начинает компенсировать ошибку нефизичной NRTL-подстройкой. & Усилить $\lambda_{\mathrm{tau\_reg}}$ и проверить распределения $\tau_{12},\tau_{21}$. \\
\hline
$\mathcal{L}_{\mathrm{vH}}$ & Рост $>10\times$ за 5 эпох & Красный флаг: вероятна численная нестабильность vH-компоненты, а не осмысленная структурная коррекция. & Проверить clamp и cap из §\ref{subsec:vh-stabilization}, а также pair-aware batching. \\
\hline
Pearson $r(T_{\mathrm{m}})$ после фазы 1 & $< 0.3$ & Warning: crystal-head не вышла из почти случайного режима и не даёт качественной инициализации для фазы растворимости. & Увеличить длительность фазы 1 и/или learning rate для блока кристаллических свойств. \\
\hline
\end{tabularx}
\caption{Практические индикаторы проблем сходимости и рекомендуемые действия.}
\label{tab:convergence-indicators}
\end{table}

Эти пороги следует интерпретировать совместно, а не по отдельности. Например, умеренный рост $\mathcal{L}_{\mathrm{tau\_reg}}$ в ранней фазе 2 ещё допустим, если одновременно быстро улучшается $\mathrm{MAE}_{\mathrm{val}}$ и не падает \texttt{sol\_fraction}. Напротив, даже небольшой рост correction path в фазе 3 становится подозрительным, если он сопровождается стагнацией валидационного качества. Для $\overline{|\delta_{\mathrm{corr}}|}$ в текущей постановке не вводится один универсальный жёсткий порог: эта величина зависит от масштаба распределения $\ln x_2$ и потому полезнее как \textbf{относительная кривая внутри запуска}, чем как абсолютный бинарный тест.

\warnnote{Конкретные кривые сходимости для текущих прогонов, включая траектории $\mathrm{MAE}_{\mathrm{val}}(\ln x_2)$, \texttt{sol\_fraction}, $\mathcal{L}_{\mathrm{tau\_reg}}$, $\mathcal{L}_{\mathrm{vH}}$ и $\overline{|\delta_{\mathrm{corr}}|}$, будут приведены после завершения Эксперимента~1 на полном бюджете обучения. До этого момента данный подраздел фиксирует протокол анализа и набор индикаторов, а не итоговые эмпирические графики.}

\section{Pair-aware temperature batching}

\subsection{Проблема}

Стандартный случайный батчинг: в одном минибатче вероятность встретить несколько температур для одной и той же пары (solute, solvent) мала. Регуляризаторы $\mathcal{L}_{\mathrm{rank}}$ и $\mathcal{L}_{\mathrm{vH}}$ неактивны.

\subsection{Решение}

\texttt{PairTemperatureBatchSampler} группирует записи по паре (\texttt{solute\_smiles}, \texttt{solvent\_smiles}). Если пара имеет $\geq$ \texttt{pair\_temperature\_min\_group\_size} записей при разных $T$, все её записи попадают в один минибатч.

\textbf{Эффект:} в каждом минибатче систематически присутствуют 2--10 температур для нескольких пар $\to$ ranking и van't Hoff penalties активны \textbf{на каждом шаге} обучения.

\subsection{Связь с идентифицируемостью}

При нескольких $T$ для одной пары мастер-уравнение даёт систему:

\begin{equation*}
\ln x_2(T_k) = -\Phi(T_k;\; T_{\mathrm{m}}, \Delta H_{\mathrm{fus}}) - \ln\gamma_2(x_2(T_k);\; \tau_{12}(T_k), \tau_{21}(T_k), \alpha)
\end{equation*}

для $k = 1, \ldots, K$. Это $K$ уравнений на параметры $(T_{\mathrm{m}}, \Delta H_{\mathrm{fus}}, \tau_{12}^{\mathrm{ref}}, \beta_{12}, \tau_{21}^{\mathrm{ref}}, \beta_{21}, \alpha)$ --- 7 неизвестных.

При $K \geq 7$ система формально определена (хотя практически может быть плохо обусловлена). При $K = 1$ --- сильно недоопределена.

\section{Оптимизация}

\subsection{Оптимизатор}

Оптимизация выполняется вариантами Adam/AdamW~\cite{kingma2015,loshchilov2019} в реализации PyTorch~\cite{paszke2019}.

AdamW с weight decay. Типичные параметры:

\begin{table}[H]
\centering
\begin{tabular}{|c|c|c|c|}
\hline
Параметр & фаза 1 & фаза 2 & фаза 3 \\
\hline
Learning rate & $3 \times 10^{-4}$ & $1 \times 10^{-4}$ & $1 \times 10^{-5}$ \\
\hline
Weight decay & $10^{-5}$ & $10^{-5}$ & $10^{-4}$ \\
\hline
Batch size & 64 & 64 & 64 \\
\hline
\end{tabular}
\caption{Фазозависимые гиперпараметры оптимизатора в базовом учебном режиме.}
\label{tab:optimizer-phase-hparams}
\end{table}

\subsection{Learning rate schedule}

Косинусное убывание шага обучения внутри каждой фазы с разогревом (5--10 эпох).

\subsection{Смешанная точность}

графовый энкодер и выходные головы: FP16 через AMP. SLE-решатель: всегда FP32 из-за чувствительности итераций неподвижной точки и неявных производных к ошибкам округления.

\section{Переобучение и балансировка регуляризации}

\subsection{Наблюдаемое переобучение при среднем бюджете}

Для режима $10/40/10$ при $\mathrm{dropout}=0.012$ наблюдается ранняя деградация валидации:
\begin{itemize}
  \item лучшая точка: эпоха 5 фаза 2, $\mathrm{MAE}=1.929$, $R^2=0.243$;
  \item к эпохе 7: $\mathrm{MAE}=1.956$, $R^2=0.218$;
  \item обучающая $\mathcal{L}_{\mathrm{sol,raw}}$: $0.747\rightarrow0.211$ (улучшение в $\sim3.5\times$), тогда как валидационные метрики стагнируют;
  \item $\tau_{\mathrm{reg}}$: $0.56\rightarrow2.64$ (рост в $\sim4.7\times$).
\end{itemize}

Картина согласуется с классическим сценарием переобучения: качество на обучающей выборке продолжает улучшаться, а обобщающая способность ухудшается.

\subsection{Причины: масштаб модели и смещение гиперпараметров из-за прокси-режима}

Приблизительно $N_{\mathrm{params}}\sim2\times10^6$ при $N_{\mathrm{train}}\sim8\times10^3$ даёт отношение
\begin{equation*}
\frac{N_{\mathrm{samples}}}{N_{\mathrm{params}}}\approx 4\times10^{-3},
\end{equation*}
что существенно ниже режимов, где без сильной регуляризации модель обычно устойчива.

Критично, что текущий $\mathrm{dropout}=0.012$ найден на proxy-протоколе
$(1\text{k},\;1/2/1)$, где модель преимущественно underfit и оптимум смещён к минимальной регуляризации. Формально:
\begin{equation*}
p^* \approx f\!\left(\frac{N_{\mathrm{params}}}{N_{\mathrm{samples}}\cdot N_{\mathrm{epochs}}}\right),
\end{equation*}
где $f(\cdot)$ --- эмпирическая эвристика, а не аналитически выведенная функция.
Оценка эффективного обучающего объёма:
\begin{align*}
\text{proxy: } &N_{\mathrm{samples}}\cdot N_{\mathrm{epochs}}=1000\cdot4=4000,\quad p^*\approx0,\\
\text{medium: } &N_{\mathrm{samples}}\cdot N_{\mathrm{epochs}}=8000\cdot60=480000,\quad p^*\approx0.1\text{--}0.2 \; (\text{эмпирическое наблюдение}).
\end{align*}

Отсюда напрямую следует, что proxy-оптимум по dropout не переносим на medium/полный вычислительный бюджет.

\subsection{$\tau$-регуляризация как индикатор физически экстремальных значений параметров}

Контрольный функционал:
\begin{equation*}
\mathcal{L}_{\mathrm{tau\_reg}}=\frac{1}{N}\sum_i\left(\tau_{12,i}^2+\tau_{21,i}^2\right).
\end{equation*}
Его систематический рост означает, что модель уходит в экстремальные энергии взаимодействия, чтобы подгонять train-выборку. Для типичного масштаба
\begin{equation*}
|\tau|>5 \;\Longleftrightarrow\; |\Delta g|=|\tau|RT \gtrsim 1.2\times10^4\;\mathrm{J/mol}
\end{equation*}
(при комнатных температурах), что для значимой доли органических систем уже выглядит физически подозрительно.

Практическая рекомендация: усилить стабилизатор в фаза 2,
\begin{equation*}
\lambda_{\tau\_\mathrm{reg}}: 0.002 \rightarrow 0.01,
\end{equation*}
и совместно перетюнить $(\mathrm{dropout},\mathrm{weight\_decay})$ при каждом переходе proxy $\rightarrow$ medium $\rightarrow$ full.

\section{Фазовые веса потерь}

Качественная схема (точные значения --- в \texttt{trainer.py::phase\_weights}):

\begin{table}[H]
\centering
\begin{tabular}{|c|c|c|c|}
\hline
Потеря & фаза 1 & фаза 2 & фаза 3 \\
\hline
$\mathcal{L}_{\mathrm{sol}}$ & 0 & \textbf{1.0} & \textbf{1.0} \\
\hline
$\mathcal{L}_{T_{\mathrm{m}}}$ & \textbf{1.0} & 0.1 & 0.05 \\
\hline
$\mathcal{L}_{\Delta H}$ & \textbf{1.0} & 0.1 & 0.05 \\
\hline
$\mathcal{L}_{\mathrm{Hansen}}$ & \textbf{0.5} & 0.05 & 0.02 \\
\hline
$\mathcal{L}_{\gamma^\infty}$ & 0 & 0.1 & 0.05 \\
\hline
$\mathcal{L}_{\mathrm{bridge}}$ & 0.1 & 0.05 & 0.02 \\
\hline
$\mathcal{L}_{\mathrm{mono}}$ & 0 & 0.1 & \textbf{0.5} \\
\hline
$\mathcal{L}_{\mathrm{rank}}$ & 0 & 0.1 & 0.2 \\
\hline
$\mathcal{L}_{\mathrm{vH}}$ & 0 & 0.05 & 0.1 \\
\hline
$\mathcal{L}_{\mathrm{res}}$ & 0 & 0.01 & \textbf{0.1} \\
\hline
$\mathcal{L}_{\mathrm{phys\_pref}}$ & 0 & 0.01 & 0.05 \\
\hline
$\mathcal{L}_{\mathrm{tau\_reg}}$ & 0 & 0.01 & 0.01 \\
\hline
\end{tabular}
\caption{Фазовые веса компонент функции потерь в трёхфазном учебном режиме$^{\dagger}$.}
\label{tab:phase-loss-weights}
\end{table}

\textit{$^{\dagger}$ После введения внешнего IDAC-корпуса компонента $\mathcal{L}_{\gamma^\infty}$ активна, но только начиная с фазы 2: в фазе 1 NRTL-head остаётся замороженной, поэтому вес сохраняется равным 0.}

\textbf{Логика:}

\begin{itemize}
  \item фаза 1: только auxiliary/property функции потерь --- подготовка heads.
  \item фаза 2: доминирует $\mathcal{L}_{\mathrm{sol}}$, auxiliary-функции потерь уменьшены.
  \item фаза 3: усилены monotonicity и correction penalties --- «полировка».
\end{itemize}

\section{Данные и разбиения}

\subsection{Основной датасет}

После предобработки BigSolDB v2.1 разделяется на два операционно разных представления: полный merged frame содержит 120\,197 записей, 19\,878 unique solute, 212 unique solvent и 29\,826 уникальных пар; текущий water-inclusive solubility-supervised subset содержит 108\,287 записей с $\ln x_2$, 212 solvent и 12\,129 pair-level комбинаций. Историческая water-excluded конфигурация до включения воды соответствовала 101\,763 строкам, 1\,444 unique solute, 211 solvent и 11\,392 парам.

Auxiliary data:
\begin{itemize}
  \item Bradley melting points: 51\,944 строки с $T_{\mathrm{m}}$ (51.04\% supervised subset; 609 solute)
  \item NIST curated: 1\,016 строк с $\Delta H_{\mathrm{fus}}$ (1.00\%; 8 solute)
  \item Hansen parameters: 863 строки (0.85\%; 5 solute)
  \item IDAC ($\gamma^\infty$): 0 строк в текущем срезе корпуса
\end{itemize}

\subsection{Режимы разбиения}

\begin{table}[H]
\centering
\begin{tabularx}{\textwidth}{|P{3.4cm}|Y|C{1.8cm}|}
\hline
Режим & Гарантия & Строгость \\
\hline
\texttt{solute\_scaffold} & Нет утечки по каркасам растворяемого вещества & Высокая \\
\hline
\texttt{solute} & Нет утечки по SMILES растворяемого вещества & Средняя \\
\hline
\texttt{solvent}$^{\ddagger}$ & Нет утечки по SMILES растворителя & Высокая \\
\hline
\end{tabularx}
\caption{Режимы разбиения данных и гарантии отсутствия утечки информации.}
\end{table}

\textit{$^{\ddagger}$ Режим \texttt{solvent} поддерживается, но в текущем отчёте не протестирован; его оценка запланирована в \hyperref[subsec:roadmap-phase3]{Фазе III} для проверки обобщения на новые растворители.}

Каноническое разбиение: \texttt{solute\_scaffold}, обучение/валидация/тест = 80/10/10. Для исторической water-excluded конфигурации до включения воды это соответствовало 90\,808 / 5\,385 / 5\,570 строкам, 1\,180 / 129 / 135 solute, 207 / 66 / 62 solvent и 9\,954 / 735 / 703 парам; текущая water-inclusive конфигурация использует тот же режим, но с увеличенным числом supervised строк.

\subsection{Записи только со вспомогательными метками}

Соединения с $T_{\mathrm{m}}$ или $\Delta H_{\mathrm{fus}}$, но без данных о растворимости, добавляются в тренировочный набор для фазы 1 (предварительное обучение по физическим свойствам). Эти записи не участвуют в $\mathcal{L}_{\mathrm{sol}}$. Именно так используется часть химического разнообразия из full merged frame, не вошедшая в основной supervised subset; исторически до включения воды это формулировалось как «не вошедшая в 1\,444 supervised solute».

\section{Методология численных экспериментов}

Этот раздел фиксирует именно протокол экспериментов (что, как и в каком режиме сравнивается), тогда как подробная характеристика корпусов данных уже дана в разделе \textit{Данные и разбиения}.

\subsection{Данные и источники меток}

Основной корпус для обучения и сравнения моделей --- \textbf{BigSolDB v2.1}. После предобработки полный merged frame содержит 120\,197 записей, а текущий water-inclusive solubility-supervised subset --- 108\,287 записей формата
\begin{equation*}
(\mathrm{solute},\; \mathrm{solvent},\; T,\; \ln x_2).
\end{equation*}

Для корректной интерпретации результатов важно различать эти два уровня: 19\,878 unique solute относятся к merged frame в целом, тогда как прямая regression supervision по $\ln x_2$ покрывает лишь существенно меньший water-inclusive supervised subset. Историческая water-excluded конфигурация до включения воды покрывала 1\,444 unique solute (см. §\ref{subsec:data-preprocessing}).

Для физически-информированных моделей дополнительно используются разреженные вспомогательные источники:
\begin{itemize}
  \item Bradley melting points: 51\,944 строки с $T_{\mathrm{m}}$ (51.04\% supervised subset);
  \item NIST-агрегаты: 1\,016 строк с $\Delta H_{\mathrm{fus}}$ (1.00\%);
  \item Hansen-параметры: 863 строки (0.85\%);
  \item IDAC/$\gamma^\infty$: в текущем срезе корпуса отсутствуют.
\end{itemize}

Таким образом, постановка обучения является мультизадачной и частично-размеченной: целевая метка $\ln x_2$ относительно плотная, $T_{\mathrm{m}}$ покрыта умеренно, $\Delta H_{\mathrm{fus}}$ и Hansen --- крайне разреженно, а канал $\gamma^\infty$ в текущей постановке не активен.

\subsection{Схемы разбиения}

Основной протокол валидации --- \textbf{разбиение \texttt{solute\_scaffold}} (80/10/10), которое задаёт наиболее строгую проверку обобщения на новые химические каркасы растворяемого вещества.

Дополнительно для быстрой инженерной итерации используется \textbf{прокси-разбиение}:
\begin{equation*}
1000 / 200 / 400
\end{equation*}
(обучение/валидация/тест), предназначенный только для отладки конвейера, быстрой проверки корректности и предварительного поиска гиперпараметров.

Принципиально: выбор архитектуры и итоговые выводы делаются по строгому протоколу разбиения по каркасам, а прокси-режим служит инструментом ускоренной разработки.

\subsection{Сравниваемые модели}

Для честного сопоставления сравнивается набор базовых и физически-информированных конфигураций:

\begin{enumerate}
  \item \textbf{RF(descriptors):} Random Forest на \textasciitilde200 RDKit-дескрипторах.
  \item \textbf{RF(Morgan):} Random Forest на 2048-bit Morgan fingerprints.
  \item \textbf{DirectGNN:} нейросетевая модель с сопоставимым backbone без physics bottleneck.
  \item \textbf{DirectGNN + descriptors:} DirectGNN с добавленными RDKit-дескрипторами.
  \item \textbf{TGNN базовый ориентир:} исходная конфигурация \texttt{paper\_config}.
  \item \textbf{TGNN tuned:} конфигурация после Optuna-поиска.
  \item \textbf{TGNN + GC priors:} physics-модель с групповыми priors для crystal-параметров.
  \item \textbf{TGNN + no bridge:} вариант абляции без bridge-функции потерь.
  \item \textbf{TGNN combined:} комбинированная модификация (tuned + структурные улучшения).
\end{enumerate}

Для нейросетевых сравнений сохраняется принцип архитектурной парности: различия между DirectGNN и TGNN должны отражать именно вклад физического слоя и связанных регуляризаторов.

\subsection{Метрики оценки}

Основные метрики качества по растворимости:
\begin{itemize}
  \item \textbf{MAE} по $\ln x_2$ (ключевая метрика ранжирования моделей);
  \item \textbf{RMSE} по $\ln x_2$;
  \item \textbf{$R^2$} (коэффициент детерминации).
\end{itemize}

Для диагностики качества физической декомпозиции (в TGNN-семействе) дополнительно фиксируются:
\begin{itemize}
  \item \textbf{$T_{\mathrm{m}}$ MAE} в К;
  \item \textbf{Pearson $r$} для предсказаний $T_{\mathrm{m}}$;
  \item \textbf{sol\_fraction}: доля $\mathcal{L}_{\mathrm{sol}}$ в полной функции потерь.
\end{itemize}

Отслеживание \texttt{sol\_fraction} позволяет контролировать фактический баланс между основной задачей растворимости и auxiliary/regularization-сигналами в разные фазы обучения.

\subsection{Бюджеты обучения}

Используются три стандартизованных режима по эпохам (фаза 1 / фаза 2 / фаза 3):
\begin{itemize}
  \item \textbf{Прокси-бюджет:} 1 / 2 / 1 --- быстрый цикл отладки и грубого поиска;
  \item \textbf{Средний бюджет:} 10 / 40 / 10 --- промежуточная проверка устойчивости;
  \item \textbf{Полный бюджет:} 50 / 200 / 50 --- основной протокол для итоговых сравнений.
\end{itemize}

\begin{table}[H]
\centering
\small
\begin{tabular}{|l|c|c|p{5.1cm}|}
\hline
Сценарий & Эпохи (1/2/3) & Оценка времени на 1 запуск & Назначение \\
\hline
Прокси-бюджет & 1/2/1 & 0.5--1.5 GPU-часа & Проверка работоспособности конвейера и первичный подбор гиперпараметров \\
\hline
Средний бюджет & 10/40/10 & 8--16 GPU-часов & Стабилизация обучения и проверка устойчивости метрик \\
\hline
Полный бюджет & 50/200/50 & 24--48 GPU-часов & Итоговые сравнения моделей и отчётные результаты \\
\hline
\end{tabular}
\caption{Оценка времени обучения для разных вычислительных сценариев (на один запуск).}
\end{table}

Такая лестница бюджетов разделяет задачи: сначала подтверждается работоспособность и направление улучшений, затем запускаются дорогостоящие итоговые прогоны.

\subsection{Гиперпараметрический поиск}

\subsubsection{Мотивация}

В предварительных прогонах было установлено, что качество модели крайне чувствительно к выбору гиперпараметров. Для TGNN на прокси-протоколе дефолтная конфигурация давала
\begin{equation*}
\mathrm{MAE}_{\mathrm{proxy}} = 5.37,
\end{equation*}
тогда как после направленного поиска через Optuna получено
\begin{equation*}
\mathrm{MAE}_{\mathrm{proxy}} = 2.12.
\end{equation*}
Относительное улучшение составляет примерно
\begin{equation*}
\frac{5.37}{2.12} \approx 2.53\times.
\end{equation*}

Следовательно, в текущем наборе экспериментов гиперпараметрическая настройка является одним из наиболее сильных \emph{единичных факторов} прироста качества. Отсюда следует практический вывод: архитектурные сравнения без сопоставимой настройки гиперпараметров статистически и методологически невалидны.

\subsubsection{Методология}

Оптимизация выполняется во фреймворке \textbf{Optuna} с примерером \textbf{TPE} (Tree-structured Parzen Estimator) и прунером \textbf{MedianPruner}.

\paragraph{Целевая функция исследования.}
Для запуска с гиперпараметрами $\theta$ и случайностью обучения $\omega$ минимизируется
\begin{equation*}
J(\theta; \omega) = \mathrm{MAE}_{\mathrm{val}}\!\left(\hat y_{\theta,\omega}, y\right),
\quad y = \ln x_2.
\end{equation*}
В терминах Optuna:
\begin{equation*}
\theta^* = \arg\min_{\theta \in \Theta} J(\theta; \omega).
\end{equation*}

\paragraph{Proxy setup.}
Поиск проводится на ускоренном протоколе:
\begin{itemize}
  \item scaffold-разбиении: $1000$ train / $200$ val / $400$ test;
  \item бюджет по фазам: $1/2/1$ эпох (фаза 1 / фаза 2 / фаза 3);
  \item целевая метрика: validation MAE по $\ln x_2$.
\end{itemize}

\paragraph{Математика TPE.}
В TPE моделируется не напрямую $p(y\mid \theta)$, а обратная плотность $p(\theta\mid y)$ с разбиением наблюдений по порогу качества $y^*$:
\begin{equation*}
p(\theta\mid y)=
\begin{cases}
l(\theta), & y < y^*, \\
 g(\theta), & y \ge y^*.
\end{cases}
\end{equation*}
Обычно $y^*$ выбирается как $\gamma$-квантиль значений функции потерь среди завершённых запусков (в практике TPE часто $\gamma\approx 0.1\text{--}0.3$).

Для минимизационной задачи вероятность улучшения:
\begin{equation*}
\mathrm{EI}_{y^*}(\theta)=\int_{-\infty}^{y^*}(y^*-y)\,p(y\mid\theta)\,dy.
\end{equation*}
С использованием формулы Байеса
\begin{equation*}
p(y\mid\theta)=\frac{p(\theta\mid y)p(y)}{p(\theta)},
\end{equation*}
и кусочной модели $p(\theta\mid y)$ получается (с точностью до множителя, не зависящего от $\theta$):
\begin{equation*}
\mathrm{EI}_{y^*}(\theta) \propto
\left(\gamma + (1-\gamma)\frac{g(\theta)}{l(\theta)}\right)^{-1}.
\end{equation*}
Следовательно, максимизация EI эквивалентна предпочтению кандидатов с большим отношением
\begin{equation*}
\frac{l(\theta)}{g(\theta)},
\end{equation*}
то есть гиперпараметров, более вероятных в множестве ``хороших'' запусков, чем в ``плохих''.

Практически Optuna генерирует набор кандидатов $\{\theta_i\}_{i=1}^m$, оценивает суррогатный скор $l(\theta_i)/g(\theta_i)$ и выбирает лучший кандидат для следующего дорогостоящего запуска обучения.

\paragraph{Прунинг (MedianPruner).}
Пусть $v_t^{(k)}$ --- значение промежуточной метрики (например, val MAE после шага/эпохи $t$) в запуске $k$. Для текущего запуска $k^\prime$ вычисляется медиана среди завершённых к тому же шагу запусков:
\begin{equation*}
\tilde v_t = \mathrm{median}\{v_t^{(k)}\}_{k \in \mathcal{C}_t},
\end{equation*}
где $\mathcal{C}_t$ --- множество запусков, имеющих наблюдение на шаге $t$.

Для минимизационной метрики запуск прерывается, если
\begin{equation*}
v_t^{(k^\prime)} > \tilde v_t,
\end{equation*}
после выполнения стартовых условий (минимальное число завершённых запусков, шаги разогрева и интервальные шаги). Это перераспределяет вычислительный бюджет в пользу перспективных конфигураций.

\paragraph{Пространство поиска TGNN.}
Для TGNN используется пространство:
\begin{align*}
\text{dropout} &\in [0.0, 0.3], \\
\text{lr\_phase2} &\in [10^{-5}, 10^{-3}] \quad (\log\text{-scale}), \\
\text{lr\_phase1} &= k_1\,\text{lr\_phase2}, \quad k_1 \in \{1,2,3,5\}, \\
\text{lr\_phase3} &= k_3\,\text{lr\_phase2}, \quad k_3 \in \{0.01,0.1,0.5\}, \\
\text{warmup\_epochs} &\in \{0,1,3,5\}, \\
\text{grad\_clip} &\in [0.5,5.0], \\
\text{nrtl\_tau\_mode} &\in \{\texttt{ref\_invT},\texttt{legacy}\}.
\end{align*}
Типичный бюджет поиска составляет $20\text{--}30$ запусков.

\subsubsection{Результаты для TGNN}

Лучший запуск дал
\begin{equation*}
\mathrm{MAE}_{\mathrm{val}}=2.678,
\qquad
\mathrm{MAE}_{\mathrm{test}}=2.121.
\end{equation*}
Найденная конфигурация:
\begin{align*}
\text{dropout} &= 0.0124, \\
\text{lr\_phase2} &= 8.53\times10^{-5}, \\
\text{lr\_phase1} &= 2.559\times10^{-4} \quad (3\times\text{lr\_phase2}), \\
\text{lr\_phase3} &= 8.53\times10^{-6} \quad (0.1\times\text{lr\_phase2}), \\
\text{warmup\_epochs} &= 0, \\
\text{grad\_clip} &= 2.0, \\
\text{nrtl\_tau\_mode} &= \texttt{ref\_invT}.
\end{align*}

Ключевые наблюдения:
\begin{itemize}
  \item очень низкий dropout ($\approx 0.012$) указывает, что в прокси-режиме доминирует недообучение, а не переобучение;
  \item нулевой разогрев оптимален при коротком бюджете, поскольку разогрев ``съедает'' значимую долю полезных шагов;
  \item $\text{lr\_phase1} > \text{lr\_phase2}$ согласуется с тем, что фаза 1 обучается на auxiliary-целях и допускает более агрессивный шаг;
  \item $\text{lr\_phase3} \ll \text{lr\_phase2}$ подтверждает необходимость осторожной тонкой настройки на финальной фазе.
\end{itemize}

\subsubsection{Результаты для DirectGNN}

Для DirectGNN выполнен отдельный поиск (20 запусков, 4 эпохи в плоском режиме без фазового учебного режима). Лучший запуск:
\begin{equation*}
\mathrm{MAE}_{\mathrm{val}}=2.319,
\qquad
\mathrm{MAE}_{\mathrm{test}}=1.880.
\end{equation*}
Параметры:
\begin{align*}
\text{hidden\_dim} &= 128, \\
\text{n\_gnn\_layers} &= 4, \\
\text{dropout} &= 0.174, \\
\text{lr} &= 8.56\times10^{-4}, \\
\text{weight\_decay} &= 4.18\times10^{-5}, \\
\text{grad\_clip} &= 2.54, \\
\text{batch\_size} &= 64.
\end{align*}

Наблюдения:
\begin{itemize}
  \item оптимальная DirectGNN компактнее (128, 4 слоя) относительно базовой TGNN-конфигурации (256, 6 слоёв);
  \item требуемый dropout выше ($0.174$), что согласуется с большей склонностью прямой регрессии к overfitting;
  \item оптимальный learning rate выше, чем у TGNN, поскольку отсутствуют жёсткий SLE-решатель и фазовый учебный режим, а оптимизационный ландшафт более ``прямой''.
\end{itemize}

\subsubsection{Важность честного сопоставления}

Корректная схема сравнения:
\begin{equation*}
\mathrm{TGNN}_{\mathrm{tuned}} \;\text{по сравнению с}\; \mathrm{DirectGNN}_{\mathrm{tuned}}.
\end{equation*}
Некорректная схема:
\begin{equation*}
\mathrm{TGNN}_{\mathrm{default}} \;\text{по сравнению с}\; \mathrm{DirectGNN}_{\mathrm{tuned}}.
\end{equation*}

Практическое правило для абляционных/вариантных экспериментов: все модификации архитектуры должны стартовать от настроенного базиса соответствующего семейства, иначе наблюдаемая разница смешивает эффект идеи и эффект неоптимальных гиперпараметров.

\subsubsection{Ограничения настройки на прокси-режиме}

Proxy-протокол полезен для быстрого скрининга, но не эквивалентен полному обучению:
\begin{align*}
\text{proxy:}&\quad 1000/200/400,\; 1/2/1, \\
\text{full:}&\quad 8k^{+}/1k^{+}/1k^{+},\; 50/200/50.
\end{align*}
Из-за сдвига масштаба данных и бюджета оптимальные гиперпараметры могут отличаться. Поэтому для финальных итоговых запусков рекомендуется:
\begin{enumerate}
  \item либо повторный гиперпараметрический поиск на full (или medium-to-full) протоколе;
  \item либо как минимум проверка чувствительности вокруг прокси-оптимума.
\end{enumerate}

Итого: настройка на прокси-режиме --- инструмент грубой селекции направлений, но не окончательный источник финальных гиперпараметров.

\section{GNN Encoder как информационный bottleneck}
\label{sec:encoder-bottleneck}

\subsection{Постановка проблемы и теоретический контекст}

Теоретически message-passing архитектуры при достаточной ширине/глубине и корректной агрегации могут аппроксимировать широкий класс инвариантных функций молекулярного графа. Однако в практическом режиме TGNN-Solv ограничена тремя факторами: (i) объёмом и разнообразием данных, (ii) конечным рецептивным полем при фиксированном числе слоёв $L$, (iii) ограниченной выразительностью семейства MPNN (в терминах различимости графов на уровне WL-типа критериев).

RDKit-дескрипторы, напротив, представляют собой компактную инженерную проекцию химического пространства в набор из порядка $\sim 200$ признаков, каждый из которых оптимизировался и валидировался в прикладной хемоинформатике в течение многих лет. Поэтому сравнение «чистый GNN» по сравнению с «дескрипторным ориентиром» является прежде всего сравнением качества носителя информации, а не только класса регрессора.

\subsection{Linear probe эксперимент: методология и результаты}

Для количественной оценки информационной полноты энкодера выполнен linear probe.
Для каждой молекулы solute из train/test извлекается графовое представление
\begin{equation*}
\mathbf{g}_{\mathrm{sol}}\in\mathbb{R}^{3d},
\end{equation*}
после чего для каждого дескриптора $d_k$ обучается Ridge-регрессия
\begin{equation*}
\hat d_{k,i}=\mathbf{w}_k^\top\mathbf{g}_{\mathrm{sol},i}+b_k.
\end{equation*}
Качество декодирования измеряется коэффициентом детерминации:
\begin{equation*}
R^2_k = 1 - \frac{\sum_i (d_{k,i} - \hat{d}_{k,i})^2}{\sum_i (d_{k,i} - \bar{d}_k)^2}.
\end{equation*}

Ключевые результаты linear probe (test split):
\begin{center}
\small
\begin{tabular}{l c}
\hline
Дескриптор & $R^2$ \\
\hline
FractionCSP3 & 0.926 \\
NumHDonors & 0.693 \\
TPSA & 0.650 \\
MolLogP & 0.607 \\
NumHAcceptors & 0.619 \\
NumRotatableBonds & 0.599 \\
RingCount & 0.565 \\
MolWt & 0.450 \\
HeavyAtomCount & 0.440 \\
MolMR & 0.422 \\
\hline
\end{tabular}
\normalsize
\end{center}

Сводная статистика: медианный $R^2=0.505$, а только $3$ из $208$ дескрипторов имеют $R^2\ge 0.8$. Это означает, что энкодер после medium-budget обучения сохраняет лишь около половины информации, присутствующей в экспертных дескрипторах.

\subsection{Декомпозиция ошибки: physics по сравнению с encoder}
\label{subsec:error-decomposition}

Определим полный разрыв относительно сильного дескрипторного ориентира:
\begin{equation*}
\Delta_{\mathrm{total}} = \mathrm{MAE}_{\mathrm{TGNN}} - \mathrm{MAE}_{\mathrm{RF}}.
\end{equation*}
Разложим его на вклад physics bottleneck и вклад энкодера:
\begin{align*}
\Delta_{\mathrm{total}} &= \Delta_{\mathrm{physics}} + \Delta_{\mathrm{encoder}}, \\
\Delta_{\mathrm{physics}} &= \mathrm{MAE}_{\mathrm{TGNN}} - \mathrm{MAE}_{\mathrm{DirectGNN}}, \\
\Delta_{\mathrm{encoder}} &= \mathrm{MAE}_{\mathrm{DirectGNN}} - \mathrm{MAE}_{\mathrm{RF}}.
\end{align*}

Такая декомпозиция является приближённой: она предполагает, что основным источником различий между DirectGNN и RF выступает качество входного представления, хотя модели отличаются не только представлением, но и классом аппроксиматора.

Подстановка наблюдаемых значений:
\begin{align*}
\Delta_{\mathrm{total}} &= 1.93 - 1.20 = 0.73, \\
\Delta_{\mathrm{physics}} &= 1.93 - 1.88 = 0.05 \quad (7\%), \\
\Delta_{\mathrm{encoder}} &= 1.88 - 1.20 = 0.68 \quad (93\%).
\end{align*}

Следовательно, основной разрыв обусловлен не physics bottleneck-ом (его цена составляет лишь $0.05$ MAE), а неполнотой химической информации на выходе графового энкодера ($0.68$ MAE). Уравнение SLE остаётся точным вычислительным ядром; деградация возникает на этапе параметров, подаваемых в него.

\keyidea{Ключевой вывод: в текущей конфигурации основной источник ошибки связан с качеством представления (порядка 93\%), а не с самим физическим слоем (порядка 7\%). Приоритет улучшений --- усиление входной информативности и supervision энкодера, а не отказ от физического ядра.}

\subsection{Связь с идентифицируемостью}

На уровне §\ref{sec:gnn-encoder} ключевой проблемой обсуждалась слабая идентифицируемость параметров $\tau_{12},\tau_{21}$. Новые probe-результаты дают причинный механизм этой проблемы: энкодер неполно восстанавливает MolLogP ($R^2=0.607$), TPSA ($R^2=0.650$), а также связанные с водородными связями счётчики. Это именно те сопутствующий признакы, которые определяют структуру $\gamma_2$ и, следовательно, обратную задачу для NRTL-параметров.

Если pair-representation не содержит достаточной информации о полярности и донорно-акцепторной природе, NRTL-head вынужден компенсировать это экстремальными $\tau$. Эмпирически это проявляется в росте $\tau$-регуляризатора и выходе параметров в физически малореалистичные области.

\section{Descriptor augmentation: GNN + экспертные фичи}

\subsection{Архитектурная идея}

Предлагается дополнить pair-representation TGNN вектором нормализованных RDKit-дескрипторов. Концептуально GNN и дескрипторы комплементарны: GNN хорошо кодирует локальные структурные паттерны и взаимодействия на графе, а дескрипторы компактно привносят глобальную физико-химическую информацию (LogP, TPSA, H-bond counts), которую текущий энкодер выучивает неполно.

\subsection{Формализация descriptor-augmented представления}

Базовое представление пары из GNN-пути:
\begin{equation*}
\mathbf{g}_{\mathrm{pair}}^{\mathrm{GNN}} = \mathrm{MLP}\left([
\mathbf{g}_{\mathrm{sol}} \| \mathbf{g}_{\mathrm{slv}} \|
\mathbf{g}_{\mathrm{sol}} \odot \mathbf{g}_{\mathrm{slv}} \|
|\mathbf{g}_{\mathrm{sol}} - \mathbf{g}_{\mathrm{slv}}|]
\right).
\end{equation*}

Дескрипторный путь (shared-энкодер для solute/solvent):
\begin{equation*}
\mathbf{d}_{\mathrm{sol}} = \mathrm{MLP}_{\mathrm{desc}}\left(
\frac{\mathbf{x}_{\mathrm{sol}}^{\mathrm{desc}} - \boldsymbol{\mu}}{\boldsymbol{\sigma}}
\right),
\quad
\mathbf{d}_{\mathrm{slv}} = \mathrm{MLP}_{\mathrm{desc}}\left(
\frac{\mathbf{x}_{\mathrm{slv}}^{\mathrm{desc}} - \boldsymbol{\mu}}{\boldsymbol{\sigma}}
\right).
\end{equation*}

Композиция дескрипторной пары:
\begin{equation*}
\mathbf{g}_{\mathrm{pair}}^{\mathrm{desc}} = [
\mathbf{d}_{\mathrm{sol}} \|
\mathbf{d}_{\mathrm{slv}} \|
\mathbf{d}_{\mathrm{sol}} \odot \mathbf{d}_{\mathrm{slv}} \|
|\mathbf{d}_{\mathrm{sol}} - \mathbf{d}_{\mathrm{slv}}|].
\end{equation*}

Финальное представление:
\begin{equation*}
\mathbf{g}_{\mathrm{pair}} = \mathrm{Proj}\left([
\mathbf{g}_{\mathrm{pair}}^{\mathrm{GNN}} \|
\mathbf{g}_{\mathrm{pair}}^{\mathrm{desc}}]
\right) \in \mathbb{R}^{d_{\mathrm{pair}}}.
\end{equation*}

Здесь $\boldsymbol{\mu},\boldsymbol{\sigma}$ считаются только по train set, а $\mathrm{Proj}$ возвращает размерность в стандартный $d_{\mathrm{pair}}$.

\subsection{Ключевое свойство: physics pathway не меняется}

Важно, что физический вычислительный путь не модифицируется:
\begin{equation*}
(\mathbf{g}_{\mathrm{pair}}) \longrightarrow
\{T_{\mathrm{m}},\Delta H_{\mathrm{fus}},\tau_{12},\tau_{21}\}
\longrightarrow
\ln\gamma_2
\longrightarrow
\text{SLE solver}
\longrightarrow
\ln x_2.
\end{equation*}

То есть дескрипторы обогащают только вход для heads, но не обходят физические ограничения: все промежуточные переменные остаются интерпретируемыми.

\subsection{Ожидаемый эффект и причина не заменять GNN}

Ожидаемый механизм улучшения:
\begin{enumerate}
  \item NRTL-head получает больше информации о полярности/ассоциации $\Rightarrow$ лучшее качество $\tau$;
  \item Fusion/crystal-path получает более информативное описание solute $\Rightarrow$ стабильнее $T_{\mathrm{m}},\Delta H_{\mathrm{fus}}$;
  \item обучение по-прежнему end-to-end через physics pathway.
\end{enumerate}

Почему не заменять GNN полностью дескрипторами:
\begin{itemize}
  \item GNN даёт атомно-уровневые представления для cross-attention;
  \item часть субструктурных эффектов не кодируется фиксированным дескрипторным словарём;
  \item GNN-представления могут адаптироваться под целевую функцию, в отличие от статичных дескрипторов.
\end{itemize}

Следовательно, оптимальная стратегия --- не «GNN вместо дескрипторов» и не «дескрипторы вместо GNN», а их совместная параметризация в едином физически ограниченном пайплайне.

\subsection{Эмпирический статус обогащения дескрипторами}

Полный запуск TGNN с полным набором дескрипторов RDKit на каноническом протоколе разбиения по каркасам пока не завершён и потому не включён в основные результатные таблицы. В текущем отчёте эмпирически проверены только более узкие модификации --- отпечатки Morgan и обучаемые дескрипторные априори, --- которые оказались нейтральными или умеренно отрицательными.

Основной эксперимент с обогащением дескрипторами в полном смысле слова (конкатенация нормализованного вектора RDKit к парному представлению с последующей проекцией) отнесён к Эксперименту 4 в плане дальнейших работ. После его завершения здесь должна быть приведена прямая таблица сравнения базовой TGNN, TGNN + дескрипторы RDKit, DirectGNN + дескрипторы RDKit и RF(дескрипторы) на одном разбиении по каркасам.

\subsection{Oracle evaluation protocol}

Здесь используется тот же протокол оракульной подстановки, который канонически формализован в §\ref{sec:oracle-teacher-forcing}. На этапе диагностики часть предсказанных кристаллических параметров подменяется истинными значениями, после чего оценивается изменение итоговой ошибки по $\ln x_2$.

\subsection{Воспроизводимость вычислений}

Для многозапусковой оценки применяются фиксированные seed:
\begin{equation*}
42,\; 123,\; 456.
\end{equation*}

Текущие численные результаты в основном тексте представлены по одному запуску (одиночный запуск, seed=42). Окончательные сравнительные таблицы должны включать mean $\pm$ std как минимум по трём seed в рамках \hyperref[subsec:roadmap-phase3]{Фазы III}.

Дополнительные требования воспроизводимости:
\begin{itemize}
  \item все конфигурации выполнения и веса потерь сохраняются в checkpoint;
  \item сквозная репликация пайплайна обеспечивается сценарием \texttt{reproduce.sh};
  \item сравнения выполняются при одинаковых разбиение-файлах и неизменном preprocessing.
\end{itemize}

В совокупности этот протокол обеспечивает сопоставимость результатов между архитектурами и минимизирует риск ложных выводов из-за различий в режиме обучения или случайной инициализации.


\section*{Часть V. Диагностика: анализ проблем и эмпирические результаты}
\addcontentsline{toc}{section}{Часть V. Диагностика: анализ проблем и эмпирические результаты}

В этой части изложение организовано в причинно-следственной логике: сначала фиксируются наблюдаемые эмпирические факты, затем формулируются и ранжируются гипотезы отказа, после чего каждая гипотеза проверяется на согласованность с данными и сводится к практическим выводам. Такой порядок позволяет отделить "что наблюдается" от "почему это происходит" и "что делать дальше".

\section{Текущее состояние: эмпирические результаты}

\subsection{Базовые сравнения}

Все нейросетевые модели обучались с сокращённым бюджетом (1/4/1 эпох вместо канонических 50/200/50). RF обучается за один проход и полностью сходится.

Если не указано иное, все численные результаты этого раздела относятся к одному запуску (одиночный запуск, seed=42); многозапусковая оценка запланирована в \hyperref[subsec:roadmap-phase3]{Фазе III дорожной карты}.

Далее используются два варианта TGNN: \textbf{TGNN (full physics)} включает полный физический контур с auxiliary-потерями, а \textbf{TGNN (solubility-only)} отключает auxiliary-компоненты и обучается по целевой растворимости с минимальным набором регуляризаторов.

\textbf{Scaffold разбиение (2k/400/800):}

\begin{table}[H]
\centering
\begin{tabular}{|c|c|c|}
\hline
Модель & MAE & $R^2$ \\
\hline
RF (descriptors) & 1.198 & 0.695 \\
\hline
RF (Morgan) & 1.281 & 0.646 \\
\hline
RF (hybrid) & 1.191 & 0.697 \\
\hline
DirectGNN & 1.997 & 0.010 \\
\hline
TGNN (full physics) & 4.435 & -2.945 \\
\hline
TGNN (solubility-only) & 1.984 & 0.024 \\
\hline
\end{tabular}
\caption{Базовые результаты на scaffold-разбиении (2k/400/800) в сокращённом бюджете; текущий срез --- одиночный запуск (seed=42).}
\end{table}

\textbf{Solute разбиение (2k/400/800):}

\begin{table}[H]
\centering
\begin{tabular}{|c|c|c|}
\hline
Модель & MAE & $R^2$ \\
\hline
RF (descriptors) & 1.377 & 0.609 \\
\hline
RF (Morgan) & 1.608 & 0.467 \\
\hline
DirectGNN & 2.113 & -0.035 \\
\hline
TGNN (full physics) & 3.974 & -2.216 \\
\hline
\end{tabular}
\caption{Базовые результаты на solute-разбиении (2k/400/800) в сокращённом бюджете; текущий срез --- одиночный запуск (seed=42).}
\end{table}

\textbf{Medium scaffold run (больше эпох):}

\begin{table}[H]
\centering
\begin{tabular}{|c|c|c|}
\hline
Модель & MAE & $R^2$ \\
\hline
DirectGNN & 2.069 & -0.013 \\
\hline
TGNN (full physics) & 2.256 & -0.166 \\
\hline
\end{tabular}
\caption{Medium scaffold run: сравнение DirectGNN и TGNN при увеличенном бюджете эпох; текущий срез --- одиночный запуск (seed=42).}
\end{table}

\subsection{Статистическая значимость различий}

Для оценки устойчивости различий между моделями дополнительно используется парный $t$-тест по абсолютным ошибкам на тестовых примерах.

На текущем срезе различие TGNN ($\mathrm{MAE}=1.93$) против DirectGNN ($\mathrm{MAE}=1.88$) статистически незначимо ($p=0.34$), тогда как различие DirectGNN против RF-descriptor ($\mathrm{MAE}=1.20$) высоко значимо ($p<10^{-6}$).

Это согласуется с практической интерпретацией: основной разрыв наблюдается между дескрипторными и чисто графовыми представлениями, а не между двумя близкими нейросетевыми вариантами внутри одного семейства.

\subsection{Анализ ошибок по категориям систем}
\label{subsec:error-analysis-categories}

Для инженерной интерпретации полезно отдельно отслеживать ошибки по типам доминирующего межмолекулярного взаимодействия. Поскольку финальные числа по полному бюджету ещё не собраны, в этом подразделе фиксируется не итоговая статистика, а \textbf{framework категоризации}, который затем должен быть применён к тестовому множеству после завершения Эксперимента 1.

Категоризация строится на уровне пары \textit{solute--solvent}. Используются стандартные RDKit-дескрипторы: \texttt{MolLogP}, \texttt{TPSA}, \texttt{NumHDonors} и \texttt{NumHAcceptors}. Они вычисляются отдельно для solute и solvent. В текущей постановке вводятся следующие пять категорий.

\begin{table}[H]
\centering
\small
\setlength{\tabcolsep}{3.5pt}
\begin{tabularx}{\textwidth}{|P{3.6cm}|Y|}
\hline
Категория & Правило отнесения пары \textit{solute--solvent} \\
\hline
Неполярный + неполярный & Оба компонента удовлетворяют условиям \texttt{MolLogP} $> 2$, \texttt{NumHDonors} $= 0$ и \texttt{NumHAcceptors} $\le 1$. \\
\hline
Полярный + полярный (без выраженного H-bond) & Хотя бы один компонент имеет \texttt{TPSA} $> 40~\text{\AA}^2$, но суммарное число \texttt{NumHDonors} по паре строго меньше 2. \\
\hline
H-bond donor + H-bond acceptor & Выполняется хотя бы одно из условий: (i) у solute \texttt{NumHDonors} $\ge 1$ и у solvent \texttt{NumHAcceptors} $\ge 1$; (ii) у solvent \texttt{NumHDonors} $\ge 1$ и у solute \texttt{NumHAcceptors} $\ge 1$. \\
\hline
Водные системы & Растворитель --- вода, то есть solvent SMILES = \smarts{O}. \\
\hline
Сильно ассоциирующие системы & Solute содержит мотивы COOH, CONH, SO$_2$NH и/или растворитель является одним из сильно координирующих апротонных solvent (DMSO, ДМФА, NMP). \\
\hline
\end{tabularx}
\caption{Framework категоризации тестовых пар по типу доминирующего межмолекулярного взаимодействия.}
\label{tab:error-category-rules}
\end{table}

\warnnote{Границы между категориями принципиально нечёткие: одна и та же пара может одновременно быть водной, H-bond-активной и сильно ассоциирующей. Чтобы каждая тестовая пара попадала ровно в одну строку итоговой таблицы, вводится приоритетный порядок присвоения категорий: \textbf{водные} $>$ \textbf{сильно ассоциирующие} $>$ \textbf{H-bond donor/acceptor} $>$ \textbf{полярные без выраженного H-bond} $>$ \textbf{неполярные}.}

Для итогового отчёта по качеству требуется использовать именно эту приоритетную схему и заполнять таблицу~\ref{tab:error-category-framework} после завершения Эксперимента 1 на полном бюджете с категоризацией тестового множества.

\begin{table}[H]
\centering
\small
\setlength{\tabcolsep}{3pt}
\resizebox{\textwidth}{!}{%
\begin{tabular}{|l|c|c|c|c|c|}
\hline
Категория & $N_{\mathrm{test}}$ & MAE(TGNN) & \shortstack{MAE\\(DirectGNN)} & MAE(RF) & \shortstack{$\Delta_{\mathrm{TGNN-RF}}$} \\
\hline
Неполярный + неполярный & \texttt{TBD} & \texttt{TBD} & \texttt{TBD} & \texttt{TBD} & \texttt{TBD} \\
\hline
Полярный + полярный & \texttt{TBD} & \texttt{TBD} & \texttt{TBD} & \texttt{TBD} & \texttt{TBD} \\
\hline
H-bond donor/acceptor & \texttt{TBD} & \texttt{TBD} & \texttt{TBD} & \texttt{TBD} & \texttt{TBD} \\
\hline
Водные & \texttt{TBD} & \texttt{TBD} & \texttt{TBD} & \texttt{TBD} & \texttt{TBD} \\
\hline
Сильно ассоциирующие & \texttt{TBD} & \texttt{TBD} & \texttt{TBD} & \texttt{TBD} & \texttt{TBD} \\
\hline
\end{tabular}%
}
\caption{Framework-таблица для анализа ошибок по категориям систем.}
\label{tab:error-category-framework}
\end{table}

Здесь $\Delta_{\mathrm{TGNN-RF}} \equiv \mathrm{MAE}_{\mathrm{TGNN}} - \mathrm{MAE}_{\mathrm{RF}}$.

\warnnote{На момент данного отчёта результаты по категориальному разрезу ещё не получены; после завершения Эксперимента 1 на полном бюджете здесь должны быть приведены: (i) $N_{\mathrm{test}}$ по каждой категории, (ii) MAE для TGNN-Solv, DirectGNN и RF, (iii) разность $\Delta_{\mathrm{TGNN-RF}}$, (iv) краткая интерпретация того, в каких химических режимах physics-informed модель выигрывает или проигрывает descriptor baseline.}

Ожидаемая закономерность для \textbf{неполярных систем} состоит в том, что RF на дескрипторах будет оставаться особенно сильным ориентиром: дескрипторы уровня \texttt{MolLogP}, \texttt{MolMR} и близкие к ним табличные признаки уже хорошо кодируют размер, поляризуемость и гидрофобность, то есть именно те факторы, которые во многих неполярных парах доминируют в растворимости. В этой подгруппе TGNN-Solv должна по крайней мере не проваливаться относительно DirectGNN, но априорного преимущества перед RF здесь ожидать не следует.

Для \textbf{H-bond-активных} и особенно \textbf{сильно ассоциирующих систем} картина потенциально иная. Если NRTL-блок действительно выучивает асимметрию межмолекулярных взаимодействий и её температурную зависимость, TGNN-Solv должна быть конкурентоспособнее именно там, где простые дескрипторные baseline испытывают трудности с переносом от локальных признаков к неидеальности раствора. Однако именно в этой области гипотеза H4 становится наиболее жёсткой: bridge-loss может оказываться контрпродуктивной, поскольку Hansen-основанный proxy плохо согласуется с системами, где эффективный вклад неидеальности задаётся специфическими взаимодействиями. Дополнительно остаётся открытым вопрос H7: не является ли сама NRTL-параметризация избыточной для такого объёма данных.

\textbf{Водные системы} теперь являются полноценной, а не экстраполяционной категорией: solvent SMILES фиксируется как \smarts{O}, а в supervised subset присутствуют 6\,524 водные строки для 737 unique solute. Медианная водная растворимость ($\ln x_2=-7.82$) заметно смещена в сторону плохо растворимых режимов по сравнению с медианой корпуса, поэтому именно здесь ожидается наиболее жёсткий стресс-тест для всех моделей. Для TGNN-Solv дополнительное физическое ожидание состоит в том, что NRTL-блок должен описывать сильную неидеальность --- для большинства органических solute в воде типичен режим $\gamma_2 \gg 1$. Поэтому итоговый категориальный разрез нужен не как косметическая детализация, а как прямой тест того, где именно physics-informed подход оправдывает свою дополнительную сложность.

\openquestion{Подтверждается ли предсказание о преимуществе TGNN-Solv на ассоциирующих системах? Это критический тест для H4 и H7.}

\subsection{Сравнение с COSMO-RS}
\label{subsec:cosmo-comparison}

Прямое сопоставление TGNN-Solv с COSMO-RS/COSMO-SAC методологически желательно, но не является тривиальным 	extit{one-to-one} baseline-сравнением. В отличие от TGNN-Solv и RF(descriptors), где после этапа обучения инференс для одной системы практически бесплатен, COSMO-конвейер требует отдельного квантово-химического шага для каждой новой молекулы: сначала строится 3D-конформация, затем выполняется DFT/COSMO-расчёт и извлекается $\sigma$-profile, после чего уже вычисляются коэффициенты активности и решается SLE-задача. При типичном уровне теории RI-BP86/TZVP этот этап занимает минуты--часы на молекулу, поэтому стоимость сравнения определяется не только числом pair-temperature entries, но и числом unique molecules.

Для текущего solubility-supervised subset BigSolDB нижняя оценка квантово-химического бюджета имеет вид
\begin{equation*}
N_{\mathrm{DFT,min}} = N_{\mathrm{solute}}^{\mathrm{sup}} + N_{\mathrm{solvent}} = 1444 + 211 = 1655.
\end{equation*}
Это число относится только к unique molecules. После построения $\sigma$-profile потребуется дополнительно вычислить свойства хотя бы для всех наблюдаемых пар и температур. Следовательно, полный COSMO-benchmark для всего корпуса вычислительно дорог, но реалистичен для тщательно выбранного подмножества.

\warnnote{Прямое сравнение ``число из статьи против числа из данного отчёта'' часто некорректно даже при совпадении химической системы. Результат COSMO зависит от версии модели (COSMO-RS против COSMO-SAC), уровня теории DFT, протокола конформационного поиска и параметризации сегментных взаимодействий. Дополнительно многие публикации используют целевую шкалу $\log S$ вместо $y=\ln x_2$, и такой переход не всегда можно выполнить без дополнительной информации о единицах, плотности и стандартном состоянии.}

\textbf{Подход A: литературное сравнение.} Первый маршрут состоит в том, чтобы собрать опубликованные COSMO-RS/COSMO-SAC предсказания для подмножества систем, пересекающегося с BigSolDB, и затем привести их к максимально близкой постановке. Практически это означает: (i) выделить статьи, где явно указаны solute, solvent, температура и измеряемая величина; (ii) оставить только системы, попадающие в solubility-supervised subset; (iii) отдельно маркировать случаи, где reported target не совпадает с $y=\ln x_2$. Преимущество этого подхода --- низкий дополнительный вычислительный бюджет. Ограничение --- гетерогенность источников: разные версии COSMO, разные уровни теории, разные температурные диапазоны и часто другая целевая шкала. Поэтому такой блок следует трактовать как \textit{external literature anchor}, а не как окончательный apples-to-apples benchmark.

\textbf{Подход B: \textit{ab initio} benchmark subset.} Второй маршрут --- сформировать единый benchmark subset объёмом примерно 100--200 систем и прогнать его целиком в одном COSMO-SAC-конвейере с открытой реализацией (например, \texttt{openCOSMO} или эквивалентным открытым пакетом). Поднабор должен быть стратифицирован по пяти категориям взаимодействий, введённым в §\ref{subsec:error-analysis-categories} и таблице~\ref{tab:error-category-rules}, чтобы в benchmark одновременно вошли неполярные, полярные, H-bond-активные, водные и сильно ассоциирующие системы. Практически это означает выбор порядка 20--40 pair-temperature systems на категорию с покрытием разных температурных окон и разных размеров molecular backbone. Именно этот сценарий даёт корректное сравнение на идентичном test subset и позволяет отделить качество модели от различий в вычислительном протоколе.

Для итогового сравнения должны фиксироваться одновременно метрики качества и вычислительная стоимость. Минимальный набор величин приведён в таблице~\ref{tab:cosmo-benchmark-metrics}.

\begin{table}[H]
\centering
\small
\setlength{\tabcolsep}{4pt}
\begin{tabularx}{\textwidth}{|P{3.35cm}|C{2.0cm}|C{2.6cm}|Y|}
\hline
Модель / baseline & Качество & \shortstack{Стоимость\\на систему} & Что должно быть зафиксировано \\
\hline
COSMO-RS/COSMO-SAC & MAE, $R^2$ & минуты & Версия COSMO, уровень теории, конформационный протокол, число неуспешных DFT-задач и размер общего benchmark subset. \\
\hline
TGNN-Solv & MAE, $R^2$ & миллисекунды & Тот же test subset, тот же набор температур, checkpoint модели и режим использования предсказанных $T_{\mathrm{m}}$, $\Delta H_{\mathrm{fus}}$ и NRTL-параметров. \\
\hline
RF(descriptors) & MAE, $R^2$ & микросекунды & Тот же test subset, фиксированный дескрипторный набор и идентичный preprocessing признаков. \\
\hline
\end{tabularx}
\caption{Минимальный набор метрик и вычислительных характеристик для итогового сравнения TGNN-Solv с COSMO-ориентированным baseline.}
\label{tab:cosmo-benchmark-metrics}
\end{table}

\keyidea{Интерпретация итоговой таблицы должна быть двухкритериальной: качество и стоимость. COSMO-RS принципиально силён тем, что не требует обучения на данных растворимости и сохраняет высокую физическую интерпретируемость, однако по цене инференса он обычно оказывается порядка $10^3$--$10^5\times$ дороже TGNN-Solv. Для массового screening-сценария это различие не является второстепенной деталью, а входит в саму инженерную ценность модели.}

С практической точки зрения ожидаемый trade-off таков. COSMO-RS/COSMO-SAC является сильным физическим ориентиром, особенно в режиме ограниченной обучающей выборки или при переносе на новые молекулы без дообучения. TGNN-Solv, напротив, использует supervised-обучение, но после этого даёт миллисекундный инференс и может быть применён для быстрого перебора больших solvent libraries. RF(descriptors) остаётся ещё более дешёвым baseline, однако не предоставляет ни термодинамически интерпретируемых промежуточных величин, ни физически мотивированного механизма температурной экстраполяции.

\validated{На текущем этапе количественных результатов нет. Раздел заполняется после подготовки benchmark subset и завершения единого COSMO-протокола. Как обсуждалось в §\ref{subsec:thermo-models-review}, по литературе типичный уровень COSMO-RS/COSMO-SAC для органических систем составляет MAE $\approx 0.5$--$1.0$ по $\ln x_2$; следовательно, это серьёзный внешний ориентир, а не формальный baseline ради полноты списка.}

\openquestion{Сможет ли TGNN-Solv на общем benchmark subset приблизиться к COSMO-ориентированному baseline по MAE без потери своего ключевого преимущества по скорости именно в водных и сильно ассоциирующих системах? Этот тест является прямой проверкой практической ценности physics-informed декомпозиции.}

\subsection{Температурная экстраполяция: протокол и статус}

Одно из ключевых заявленных преимуществ TGNN-Solv --- структурированная температурная зависимость через $\Phi(T)$ и параметризацию $\tau(T)$ в форме \texttt{ref\_invT}. Поэтому отдельный протокол temperature extrapolation является обязательной частью финальной валидации.

\keyidea{Температурная экстраполяция --- это главный тест центральной гипотезы работы о том, что физически-информированная декомпозиция даёт преимущество не только по интерпретируемости, но и по обобщению вне температур, наблюдавшихся при обучении. Если TGNN-Solv не показывает выигрыша именно здесь, то основной аргумент в пользу явного физического слоя существенно ослабевает.}

Для каждой пары $p=(\mathrm{solute},\mathrm{solvent})$ обозначим отсортированный набор температур с supervision как
\begin{equation*}
T_{p,(0)} < T_{p,(1)} < \dots < T_{p,(n_p-1)}.
\end{equation*}
В температурные протоколы включаются только пары из множества
\begin{equation*}
\mathcal{P}_{\mathrm{temp}} = \left\{p:\; n_p \geq K_{\min}=5,\; T_{p,(n_p-1)}-T_{p,(0)} \geq 30\ \mathrm{K}\right\}.
\end{equation*}
Первое условие обеспечивает минимальную плотность температурных наблюдений, а второе исключает почти вырожденные окна, где даже идеальная модель не получает содержательного сигнала о температурном наклоне. По статистике корпуса из таблицы~\ref{tab:data-points-per-pair} как минимум $68.73\%+16.42\%+0.52\%=85.67\%$ пар имеют не менее шести температурных точек, то есть критерий по $n_p$ покрывает значительную часть supervised subset. После дополнительного фильтра по окну в 30 K итоговая доля будет ниже, но всё равно ожидается, что под протокол подпадает существенная часть мультитемпературного корпуса.

\textbf{Interpolation protocol.} Для интерполяционного теста требуется, чтобы каждая тестовая температура была \textit{зажата} между обучающими температурами той же пары. Поэтому после сортировки по $T$ используется не наивное правило «любая каждая вторая точка уходит в test», а интерleaving-сплит с закреплением краёв диапазона в обучении. Практически это можно записать как
\begin{equation*}
\mathcal{I}^{\mathrm{test}}_{p,\mathrm{int}} = \left\{j\in\{1,\dots,n_p-2\}: j\equiv 1\; (\mathrm{mod}\;2)\right\},
\end{equation*}
\begin{equation*}
\mathcal{I}^{\mathrm{train}}_{p,\mathrm{int}} = \{0,1,\dots,n_p-1\}\setminus \mathcal{I}^{\mathrm{test}}_{p,\mathrm{int}}.
\end{equation*}
То есть крайние температуры всегда остаются в train, а из внутренних точек в test выводится каждая вторая температура. Такая формализация эквивалентна интуитивному alternating split, но, в отличие от грубой схемы «чётные индексы в test» без фиксации границ, действительно гарантирует, что каждая тестовая температура лежит строго между двумя обучающими температурами той же пары. Именно это и делает эксперимент проверкой \textbf{интерполяции}, а не частичной экстраполяции на концах диапазона.

\textbf{Extrapolation protocol.} Для каждой включённой пары индивидуально вычисляется медианная температура
\begin{equation*}
T_{p,\mathrm{median}} = \mathrm{median}\{T_{p,(0)},\dots,T_{p,(n_p-1)}\}.
\end{equation*}
Далее train/test-разделение задаётся как
\begin{equation*}
\mathcal{I}^{\mathrm{train}}_{p,\mathrm{ext}} = \{j:\; T_{p,(j)} \leq T_{p,\mathrm{median}}\},
\qquad
\mathcal{I}^{\mathrm{test}}_{p,\mathrm{ext}} = \{j:\; T_{p,(j)} > T_{p,\mathrm{median}}\}.
\end{equation*}
Это определение важно делать \textit{по каждой паре отдельно}, а не глобально по всему корпусу: иначе одна и та же абсолютная температура может быть экстраполяционной для одних систем и интерполяционной для других. Такой pair-specific split проверяет способность модели переносить знания на более горячую половину температурного диапазона именно для данной химической пары.

Более жёсткий альтернативный вариант --- \textit{leave-out-extremes}. В нём для каждой пары в тест выводятся точки с минимальной и максимальной температурой,
\begin{equation*}
\mathcal{I}^{\mathrm{test}}_{p,\mathrm{loe}} = \{0,\, n_p-1\},
\qquad
\mathcal{I}^{\mathrm{train}}_{p,\mathrm{loe}} = \{1,\dots,n_p-2\}.
\end{equation*}
Этот режим полезен как дополнительный stress-test: модель должна одновременно выдержать экстраполяцию и в холодный, и в горячий хвост температурного окна. Он особенно чувствителен к тому, насколько параметризация $\tau(T)$ и аналитический вклад $\Phi(T)$ действительно задают физически осмысленную форму кривой, а не лишь локально интерполируют данные около центра диапазона.

\textbf{Метрики.} Помимо уже введённых MAE на interpolation- и extrapolation-поднаборах, необходимо явно оценивать качество температурного \textit{профиля}, а не только точечную ошибку. Базовой дополнительной метрикой остаётся ошибка наклона в координатах van't Hoff,
\begin{equation*}
\Delta s_p = \left|\hat{s}_p-s^{\mathrm{obs}}_p\right|,
\end{equation*}
где $s^{\mathrm{obs}}_p$ и $\hat{s}_p$ --- наблюдаемый и предсказанный наклоны зависимости $\ln x_2$ от $1/T$ на тестовых точках данной пары. Если у пары на тесте ровно две температуры, этот наклон сводится к секущей,
\begin{equation*}
s=\frac{\Delta\ln x_2}{\Delta(1/T)},
\end{equation*}
а при большем числе точек предпочтительно оценивать $s_p$ как коэффициент линейной регрессии на тестовом подмножестве.

Дополнительно вводятся две структурные метрики. \textbf{Rank consistency (RC)} --- это доля пар, для которых модель правильно восстанавливает порядок растворимости по всем тестовым температурам данной пары. Если для пары $p$ тестовые температуры отсортированы как
\begin{equation*}
T^{\mathrm{test}}_{p,1} < \dots < T^{\mathrm{test}}_{p,m_p},
\end{equation*}
то
\begin{equation*}
\mathrm{RC} = \frac{1}{|\mathcal{P}_{\mathrm{test}}|}\sum_{p\in\mathcal{P}_{\mathrm{test}}}
\mathbf{1}\!\left[\forall a<b:\; \operatorname{sign}(\hat{y}_{p,b}-\hat{y}_{p,a}) = \operatorname{sign}(y_{p,b}-y_{p,a})\right],
\end{equation*}
где $y=\ln x_2$. Эта метрика намного строже обычной pairwise accuracy: пара засчитывается как корректная только тогда, когда модель не ошибается ни в одном относительном порядке тестовых температур.

Вторая метрика --- \textbf{monotonicity violation rate}. Она измеряет долю пар, у которых в предсказанной тестовой кривой возникает хотя бы одно локальное нарушение ожидаемой монотонности по температуре:
\begin{equation*}
\mathrm{MVR} = \frac{1}{|\mathcal{P}_{\mathrm{test}}|}\sum_{p\in\mathcal{P}_{\mathrm{test}}}
\mathbf{1}\!\left[\exists k<m_p:\; \hat{y}_{p,k+1}<\hat{y}_{p,k}\right].
\end{equation*}
Здесь температура упорядочена по возрастанию, а $\hat{y}_{p,k}=\widehat{\ln x_2}(T^{\mathrm{test}}_{p,k})$. Для текущей задачи эта метрика особенно важна, потому что модель уже содержит monotonicity-приор и van't Hoff-регуляризацию; следовательно, высокий MVR будет означать, что даже встроенные физические ограничения не привели к устойчивой температурной форме.

\textbf{Связь с центральным тезисом работы.} Именно температурная экстраполяция должна наиболее явно отделить physics-informed подход от purely data-driven baselines. В TGNN-Solv температурная зависимость задаётся не произвольной функцией от скаляра $T$, а комбинацией аналитического члена $\Phi(T)$ и структурированной формы $\tau(T)$ через \texttt{ref\_invT}. Это означает, что при переходе к новым температурам модель не начинает «угадывать» заново зависимость, а продолжает уже выученную физическую параметризацию. Напротив, DirectGNN и RF в типичном режиме используют температуру как ещё один входной признак без встроенной термодинамической структуры. Они могут хорошо аппроксимировать seen-temperature regime, но не обязаны сохранять корректный наклон, ранговый порядок или монотонность вне обучающего диапазона.

\textbf{Ожидания и интерпретация.} Для финальной проверки разумно зафиксировать следующие заранее сформулированные ожидания.
\begin{itemize}
  \item TGNN-Solv должна превосходить DirectGNN и RF прежде всего на \textbf{экстраполяции} --- особенно по $\Delta s$, rank consistency и monotonicity violation rate. Именно здесь структурированная форма $\tau(T)$ и аналитический вклад $\Phi(T)$ должны приносить наибольшую пользу.
  \item На \textbf{интерполяции} преимущество TGNN-Solv может быть менее очевидным: у модели больше параметров, обучение медленнее, а простые baseline-методы нередко хорошо справляются с локальной аппроксимацией внутри уже наблюдавшегося температурного окна.
  \item RF без явной температурной структуры априорно должна показывать наиболее слабую экстраполяцию, повышенный slope error и повышенный monotonicity violation rate. Если этого не наблюдается, то утверждение о преимуществе физической структуры требует серьёзного пересмотра.
  \item Если TGNN-Solv не покажет преимущества даже на экстраполяции, то это станет сильным аргументом в пользу перехода к архитектуре v3: в таком случае текущая физически-информированная декомпозиция усложняет модель, но не даёт главного обещанного выигрыша --- лучшего обобщения по температуре.
\end{itemize}

\warnnote{На момент данного отчёта эксперимент температурной интерполяции/экстраполяции ещё не проведён на полном бюджете. После завершения Эксперимента~1 в этом подразделе должны быть приведены как минимум: число вошедших пар, MAE на interpolation и extrapolation, средний slope error, rank consistency и monotonicity violation rate для TGNN-Solv, DirectGNN и RF.}

Этот эксперимент явно отнесён к \hyperref[subsec:roadmap-phase3]{Фазе III дорожной карты} как обязательная проверка. Без его результатов центральный тезис работы о преимуществе физически-информированной декомпозиции остаётся неподтверждённым.

\subsection{Augmentation эксперименты}

Ниже приведены результаты одного запуска (seed=42); многозапусковой вариант для этих модификаций пока не собран.

\textbf{Morgan fingerprints:}

\begin{center}
\begin{tabular}{|c|c|c|}
\hline
Модель & MAE (scaffold) & $\Delta$ MAE \\
\hline
TGNN & 4.435 & --- \\
\hline
TGNN + Morgan & 4.614 & +0.179 (хуже) \\
\hline
DirectGNN & 1.997 & --- \\
\hline
DirectGNN + Morgan & 1.988 & -0.009 (нейтрально) \\
\hline
\end{tabular}
\end{center}

\textbf{Learned descriptor priors (Hansen, $V_m$):}

\begin{center}
\begin{tabular}{|c|c|c|}
\hline
Модель & MAE (scaffold) & $\Delta$ MAE \\
\hline
TGNN & 4.060 & --- \\
\hline
TGNN + descriptor priors & 4.113 & +0.053 (нейтрально/хуже) \\
\hline
\end{tabular}
\end{center}

\textbf{Fixed group-contribution priors:}

\begin{center}
\begin{tabular}{|c|c|c|}
\hline
Модель & MAE (scaffold) & $\Delta$ MAE \\
\hline
TGNN & 2.120 & --- \\
\hline
TGNN + group priors & 2.174 & +0.054 (нейтрально/хуже) \\
\hline
\end{tabular}
\end{center}

\section{Ранжированные гипотезы отказа}

\subsection{Разделение типов гипотез}

\textbf{Критически важно различать:}

\begin{itemize}
  \item \textbf{Операционные} (H5): решаемы увеличением бюджета обучения без изменения архитектуры.
  \item \textbf{Архитектурные} (H1, H4, H7): требуют изменения структуры модели.
  \item \textbf{Математические} (H2): фундаментальные ограничения задачи, не зависящие от реализации.
  \item \textbf{Физические} (H6): ограничения принятых физических приближений.
  \item \textbf{Оптимизационные} (H3): проблемы взаимодействия архитектуры и оптимизатора.
\end{itemize}

Текущие эксперименты (1/4/1 эпох) \textbf{не позволяют отделить} H5 от остальных гипотез. Это центральная методологическая проблема диагностики.

\section{H1: Инверсия иерархии чувствительности и супервизии}

\subsection{Формулировка}

Из анализа чувствительности по параметрам кристалла и NRTL:

\begin{equation*}
\left|\frac{\partial \ln x_2}{\partial \Delta g_{12}}\right| \sim 4 \times 10^{-4} \;(\text{Дж/моль})^{-1}
\end{equation*}

\begin{equation*}
\left|\frac{\partial \ln x_2}{\partial T_{\mathrm{m}}}\right| \sim 2 \times 10^{-2} \;\text{K}^{-1}
\end{equation*}

\begin{equation*}
\left|\frac{\partial \ln x_2}{\partial \Delta H_{\mathrm{fus}}}\right| \sim 1.25 \times 10^{-4} \;(\text{Дж/моль})^{-1}
\end{equation*}

Пересчёт в «ошибку $\ln x_2$ на единицу характерной ошибки параметра»:

\begin{center}
\begin{tabular}{|c|c|c|c|}
\hline
Параметр & Характерная ошибка & $\Delta(\ln x_2)$ & Прямая супервизия \\
\hline
$\Delta g_{12}$ & $\pm 2000$ Дж/моль & $\pm 0.8$ & \shortstack[l]{Косвенная через $\mathcal{L}_{\gamma^\infty}$\\(138 пар из внешнего IDAC-корпуса)} \\
\hline
$\Delta g_{21}$ & $\pm 2000$ Дж/моль & $\pm 0.8$ & \shortstack[l]{Косвенная через $\mathcal{L}_{\gamma^\infty}$\\(138 пар из внешнего IDAC-корпуса)} \\
\hline
$T_{\mathrm{m}}$ & $\pm 10$ К & $\pm 0.19$ & Разреженная (\textasciitilde25\%) \\
\hline
$\Delta H_{\mathrm{fus}}$ & $\pm 2000$ Дж/моль & $\pm 0.25$ & Разреженная (\textasciitilde12\%) \\
\hline
\end{tabular}
\end{center}

\textbf{Вывод:} модель наиболее чувствительна к параметрам NRTL ($\Delta g_{12}$, $\Delta g_{21}$), для которых supervision остаётся наиболее слабой, но уже не нулевой: внешний IDAC-корпус добавляет косвенный сигнал через $\mathcal{L}_{\gamma^\infty}$. Параметры с наибольшей доступностью прямых меток ($T_{\mathrm{m}}$, Hansen) по-прежнему оказывают меньшее влияние на целевую величину.

Ключевая оговорка состоит в том, что связь $\ln\gamma_2^\infty = \tau_{12} + \tau_{21}\exp(-\alpha_{12}\tau_{21})$ является \textbf{нелинейной совместной функцией} от $\tau_{12}$, $\tau_{21}$ и $\alpha_{12}$. Поэтому IDAC-loss накладывает совместное ограничение на весь NRTL-триплет, а не даёт независимую идентификацию каждого параметра по отдельности.

\subsection{Механизм}

Градиент потери по весам GNN:

\begin{equation*}
\frac{\partial \mathcal{L}_{\mathrm{sol}}}{\partial \mathbf{w}} = \frac{\partial \mathcal{L}_{\mathrm{sol}}}{\partial \ln x_2^*} \cdot \sum_{j} \frac{d\ln x_2^*}{d\theta_j} \cdot \frac{\partial \theta_j}{\partial \mathbf{w}}
\end{equation*}

Из формулы implicit differentiation для решения уравнения $F(x_2^*)=0$:

\begin{equation*}
\frac{d\ln x_2^*}{d\theta_j} = -\frac{\partial(\Phi + \ln\gamma_2)/\partial\theta_j}{1 + x_2^* \eta}
\end{equation*}

Для NRTL параметров числитель содержит $\partial \ln\gamma_2/\partial\theta_j$, который при типичных условиях \textbf{больше}, чем $\partial\Phi/\partial T_{\mathrm{m}}$ или $\partial\Phi/\partial\Delta H_{\mathrm{fus}}$.

\textbf{Следствие:} оптимизатор преимущественно обновляет веса, влияющие на $\tau_{12}$, $\tau_{21}$ (через представление пары $\to$ NRTLHead), потому что градиент в этом направлении сильнее. Crystal heads ($T_{\mathrm{m}}$, $\Delta H_{\mathrm{fus}}$) обновляются медленнее, получая градиент от:

\begin{enumerate}
  \item Слабый сигнал через SLE: $\partial\mathcal{L}_{\mathrm{sol}}/\partial T_{\mathrm{m}} \propto \Delta H_{\mathrm{fus}}/(RT_{\mathrm{m}}^2)$ --- малая величина при больших $T_{\mathrm{m}}$.
  \item Разреженный прямой сигнал: $\partial\mathcal{L}_{T_{\mathrm{m}}}/\partial T_{\mathrm{m}}$ --- активен только для \textasciitilde25\% батча.
\end{enumerate}

\subsection{Результат: компенсаторная дегенерация}

NRTL head «абсорбирует» всю ошибку: вместо того чтобы предсказать физически верные $\tau$, она выучивает $\tau$-значения, компенсирующие неверные $T_{\mathrm{m}}$ и $\Delta H_{\mathrm{fus}}$. Модель минимизирует $\mathcal{L}_{\mathrm{sol}}$ на train set, но предсказанная факторизация физически бессмысленна.

При тестировании на новых молекулах: $T_{\mathrm{m}}$ и $\Delta H_{\mathrm{fus}}$ ошибочны (heads не обучились), NRTL не может скомпенсировать (новые пары) $\to$ катастрофическая ошибка.

\subsection{Поддержка эмпирическими данными}

\begin{itemize}
  \item TGNN MAE = 4.435, $R^2 = -2.945$ (scaffold) --- модель хуже константного предсказания. Это согласуется с «мусорной факторизацией», где компенсация ломается на тестовых данных.
  \item TGNN (solubility-only) MAE $\approx$ DirectGNN MAE --- при удалении auxiliary-функций потерь (ослабление прямой супервизии на crystal params) результат не ухудшается, что согласуется с тем, что эта супервизия и так недостаточна.
\end{itemize}

\section{Анализ идентифицируемости параметров}
\label{sec:identifiability}

\subsection{Постановка проблемы}

Для одной пары (solute, solvent) мастер-уравнение SLE в рабочей параметризации можно записать как

\begin{equation*}
\ln x_2 = -\Phi\!\bigl(T;\, T_{\mathrm{m}}, \Delta H_{\mathrm{fus}}, \Delta C_p^{\mathrm{fus}}\bigr) - \ln\gamma_2\!\bigl(x_2;\, \tau_{12}(T), \tau_{21}(T), \alpha_{12}\bigr),
\end{equation*}

где в типичном варианте TGNN-Solv используются параметры
\begin{equation*}
\theta = \left(T_{\mathrm{m}},\, \Delta H_{\mathrm{fus}},\, \Delta C_p^{\mathrm{fus}},\, \tau_{12}^{\mathrm{ref}},\, \beta_{12},\, \tau_{21}^{\mathrm{ref}},\, \beta_{21},\, \alpha_{12}\right)
\end{equation*}
или его редуцированная версия без явного предсказания $\Delta C_p^{\mathrm{fus}}$ (тогда $p=7$ вместо $p=8$).

Для фиксированной температуры и одного наблюдаемого значения $\ln x_2$ получаем \textbf{одно скалярное уравнение} на 6--8 неизвестных (в зависимости от выбранной параметризации). Следовательно, система является \textbf{строго недоопределённой} уже на уровне размерности.

\subsection{Геометрия пространства решений}

Пусть
\begin{equation*}
F(\theta; T, y) = y + \Phi(T;\theta_{\Phi}) + \ln\gamma_2\bigl(e^y;\theta_{\gamma}\bigr), \quad y=\ln x_2^{\mathrm{obs}}.
\end{equation*}

Тогда множество параметров, согласованных с наблюдением, есть уровень
\begin{equation*}
\mathcal{M}_{T,y} = \{\theta \in \mathbb{R}^p\, :\, F(\theta;T,y)=0\}.
\end{equation*}

При невырожденном градиенте $\nabla_\theta F \neq 0$ теорема о неявной функции даёт локальную размерность
\begin{equation*}
\dim \mathcal{M}_{T,y} = p-1.
\end{equation*}

Иными словами, при $p=7$ это 6-мерное многообразие решений: \textbf{любая} точка на нём воспроизводит то же наблюдаемое $\ln x_2$. Поэтому задача обратного восстановления «истинных» физических параметров из одной точки растворимости не имеет единственного решения в принципе.

\subsection{Мультитемпературные данные: когда система становится определённой}

Если для одной и той же пары доступны измерения при $K$ температурах, получаем систему
\begin{equation*}
F_k(\theta)=0, \quad k=1,\dots,K,
\end{equation*}
где
\begin{equation*}
F_k(\theta)=\ln x_2^{\mathrm{obs}}(T_k)+\Phi(T_k;\theta_{\Phi})+\ln\gamma_2\!\bigl(x_2^{\mathrm{obs}}(T_k);\theta_{\gamma},T_k\bigr).
\end{equation*}

Формально система становится определённой при $K\ge p$ (для типичного случая --- при $K\ge 7$). Однако это лишь \textit{необходимое}, но не достаточное условие устойчивой идентификации.

Даже при $K\ge 7$ остаются две фундаментальные проблемы.

\textbf{(1) Коллинеарность температурных вкладов.}
В приближении малой растворимости и слабой температурной кривизны
\begin{equation*}
\ln x_2(T) \approx -\frac{\Delta H_{\mathrm{fus}}}{R}\!\left(\frac{1}{T}-\frac{1}{T_{\mathrm{m}}}\right)
-\tau_{12}^{\mathrm{ref}}-\beta_{12}\!\left(\frac{1}{T}-\frac{1}{T_{\mathrm{ref}}}\right)-\ldots
\end{equation*}
Члены вида $\Delta H/T$ и $\beta/T$ имеют почти одинаковую функциональную форму, что делает их практически неразличимыми по данным.

\textbf{(2) Почти вырожденный Якобиан.}
Матрица чувствительности
\begin{equation*}
J_{k,j}=\frac{\partial \ln x_2(T_k;\theta)}{\partial \theta_j}
\end{equation*}
становится плохо обусловленной, когда температуры $T_k$ расположены близко (строки $J$ почти пропорциональны). В этом режиме малый шум в $\ln x_2$ вызывает большой разброс в оценках параметров.

\subsection{Краткий анализ через Fisher Information}

Для параметра $\theta_j$ Fisher information определяется как
\begin{equation*}
\mathcal{I}(\theta_j)=\mathbb{E}\!\left[\left(\frac{\partial \ln \mathcal{L}}{\partial \theta_j}\right)^2\right],
\end{equation*}
а для вектора параметров
\begin{equation*}
\mathbf{I}(\theta)=\mathbb{E}\!\left[(\nabla_\theta \ln \mathcal{L})(\nabla_\theta \ln \mathcal{L})^\top\right].
\end{equation*}

В гауссовом приближении ошибок наблюдения выполняется
\begin{equation*}
\mathbf{I}(\theta) \approx \frac{1}{\sigma^2}J^\top J.
\end{equation*}

Следовательно, при малом числе наблюдений и коллинеарных столбцах $J$ матрица $\mathbf{I}(\theta)$ становится сингулярной или резко плохо обусловленной. Это означает огромные доверительные интервалы и практическую неустранимость неоднозначности между параметрами.

\subsection{Компенсаторная дегенерация в обучении}

В неидентифицируемом режиме оптимизатор минимизирует прежде всего ошибку по $\ln x_2$, выбирая произвольную точку на почти плоской поверхности уровня. На практике это приводит к следующему сценарию:

\begin{enumerate}
  \item $T_{\mathrm{m}}$ и $\Delta H_{\mathrm{fus}}$ получают слабый и шумный сигнал и часто остаются близко к инициализации.
  \item Параметры NRTL ($\tau_{12},\tau_{21},\beta_{12},\beta_{21}$) «абсорбируют» residual-ошибку и подстраиваются так, чтобы восстановить $\ln x_2$ на обучающих парах.
  \item Получается хорошая подгонка целевой величины без физически корректной факторизации вкладов.
\end{enumerate}

Именно этот механизм далее проявляется как \textbf{компенсаторная дегенерация}: качество на train/ID-части может быть приемлемым, но перенос на новые химические пары резко ухудшается.

\subsection{Эмпирические индикаторы слабой идентифицируемости}

Наблюдаемая картина согласуется с теорией выше.

\begin{itemize}
  \item Для crystal-параметров при коротком обучении типичные корреляции остаются низкими: $r_{\mathrm{Pearson}}(T_{\mathrm{m}}) \approx 0.33$--$0.41$.
  \item В тесте с Oracle-подстановкой $T_{\mathrm{m}}$ улучшение MAE по растворимости оказывается минимальным: NRTL-путь уже успевает скомпенсировать ошибочные $T_{\mathrm{m}}$ в обучаемой факторизации.
  \item Градиентный поток от $\mathcal{L}_{\mathrm{sol}}$ направляется преимущественно в NRTL-ветку, тогда как crystal-heads недополучают устойчивый обучающий сигнал.
\end{itemize}

Эти факты соответствуют инверсии иерархии чувствительности, наблюдаемой в отчёте:

\begin{center}
\begin{tabular}{|c|c|c|}
\hline
Параметр & Чувствительность $\ln x_2$ & Прямая супервизия \\
\hline
$\Delta g_{12},\Delta g_{21}$ (NRTL) & Высокая & \shortstack[l]{Косвенная через $\mathcal{L}_{\gamma^\infty}$\\(138 пар внешнего IDAC); слабая, но не нулевая} \\
\hline
$T_{\mathrm{m}}$ (crystal head) & Средняя & Разреженная \\
\hline
\end{tabular}
\end{center}

Для пар, где доступно наблюдение внешнего IDAC-корпуса, формально появляется ещё одно скалярное ограничение:
\begin{equation*}
\ln\gamma_2^{\infty}(\tau_{12},\tau_{21},\alpha_{12}) = \ln\gamma_2^{\infty,\mathrm{obs}}.
\end{equation*}
В локальной геометрической интерпретации это уменьшает размерность многообразия допустимых решений с $p-1$ до $p-2$ для пары с IDAC-наблюдением: к уравнению растворимости добавляется ещё одно независимое уравнение на NRTL-предел при $x_2\to 0$.

Однако эффект остаётся локально ограниченным покрытием. Внешний IDAC-корпус охватывает лишь 138 пар из 11\,392 пар основного solubility-supervised корпуса, то есть примерно 1.2\%. Поэтому для подавляющего большинства пар идентифицируемость остаётся слабой, а IDAC-supervision действует скорее как \textbf{глобальный регуляризатор} через shared weights NRTL-head, чем как жёсткое pair-level ограничение на каждую конкретную систему.

\subsection{Перспективы расширения IDAC-корпуса}

Стартовый IDAC-релиз (404 строки, 138 пар) намеренно ограничен объёмом: он фиксирует pipeline и формат, но не претендует на полноту покрытия. Основной path расширения --- дополнительные извлечения из NIST ThermoML и других публичных источников IDAC-данных (DECHEMA, DDB). По консервативной оценке, в литературе доступны порядка $10^3$--$10^4$ пар с измеренными $\gamma^\infty$, что потенциально увеличивает покрытие на 1--2 порядка.

С точки зрения архитектуры TGNN-Solv расширение IDAC-корпуса ожидается одним из наиболее эффективных единичных улучшений: оно напрямую адресует H1 (слабая supervision на NRTL-параметры) и H2 (идентифицируемость), не требуя архитектурных изменений. Практическая рекомендация: довести IDAC-корпус до $\geq 1000$ пар до запуска Эксперимента~1 на полном бюджете.

\subsection{Практические выводы и пути решения}

Из анализа следует, что проблема носит не только оптимизационный, но и структурно-идентификационный характер. Поэтому эффективные меры должны уменьшать размерность/неопределённость обратной задачи.

\begin{enumerate}
  \item \textbf{GC priors для $T_{\mathrm{m}}$, $\Delta H_{\mathrm{fus}}$.} Сужают допустимое пространство параметров и уменьшают свободу компенсаторных траекторий.
  \item \textbf{Оракульная подстановка там, где доступны метки.} Устраняет часть неопределённости, передавая в NRTL-ветку более физически корректный вход.
  \item \textbf{Снижение размерности модели активности (Wilson/UNIQUAC).} Меньшее число параметров повышает идентифицируемость по тем же данным.
  \item \textbf{Прямое предсказание $\ln x_2$ с мягкой физической регуляризацией.} Обходит жёсткий bottleneck, когда обратная факторизация принципиально неустойчива.
\end{enumerate}

Таким образом, «неуспех» архитектуры в части интерпретируемых промежуточных параметров может быть следствием фундаментальной слабой идентифицируемости, а не только недостатка эпох, lr-schedule или качества backbone.
\section{H3: Градиентная нестабильность при некондиционированном SLE-решателе}

\subsection{Механизм при коротком обучении}

После фаза 1 (1 эпоха): crystal heads предсказывают квази-случайные $T_{\mathrm{m}}$, $\Delta H_{\mathrm{fus}}$. NRTL head не обучалась вообще.

В начале фаза 2:

\begin{enumerate}
  \item SLE решатель получает $\Phi$ и $\tau$ вблизи инициализации.
  \item Если $\Phi$ слишком велико: $x_2^{(0)} = \exp(-\Phi) \approx 0$ $\to$ NRTL вычисляет $\gamma_2 \approx \gamma_2^\infty$ $\to$ $x_2^{(1)} \approx x_2^{(0)}/\gamma_2^\infty$. Итерации сходятся тривиально, но к неверному значению.
  \item Если $\tau$ имеют экстремальные значения: $\gamma_2$ может быть $\gg 1$ или $\ll 1$, генерируя нефизичные $x_2$.
\end{enumerate}

\textbf{Якобиан имплицитного дифференцирования:}

\begin{equation*}
\frac{d\ln x_2^*}{d\theta_j} = -\frac{\partial(\Phi + \ln\gamma_2)/\partial\theta_j}{1 + x_2^* \eta}
\end{equation*}

Если $1 + x_2^* \eta \approx 0$ (граница сходимости): деление на $\approx 0$ $\to$ чрезмерно большой градиент. Даже с clamp ($\varepsilon_{\min} = 0.01$): усечённый градиент может быть в 100$\times$ больше «нормального», дестабилизируя обучение.

\subsection{Масштаб проблемы}

При случайной инициализации, доля примеров с $|1 + x_2^* \eta| < 0.1$ может составлять 5--20\%. Эти примеры генерируют аномально большие градиенты, доминирующие в средних по батчу $\to$ шумное обновление весов $\to$ медленная или отсутствующая сходимость.

\subsection{Поддержка данными}

TGNN MAE уменьшается с ростом бюджета (4.435 $\to$ 2.256 при увеличении эпох), но значительно медленнее, чем DirectGNN (1.997 $\to$ 2.069 --- практически не меняется, так как DirectGNN не использует SLE-решатель). Это согласуется с постепенным «выходом» SLE-решателя из хаотического режима по мере улучшения входных параметров.

\section{H4: Bridge-loss как контрпродуктивный регуляризатор}

\subsection{Теоретическое обоснование}

Bridge-функция потерь в текущей постановке основана на Regular Solution Theory:

\begin{equation*}
\Delta g_{12}^{\mathrm{(Hansen)}} = V_1 \left[(\delta_{d,1} - \delta_{d,2})^2 + \kappa_p(\delta_{p,1} - \delta_{p,2})^2 + \kappa_h(\delta_{h,1} - \delta_{h,2})^2\right]
\end{equation*}

\textbf{Свойство:} $\Delta g_{12}^{\mathrm{(Hansen)}} \geq 0$ всегда.

\textbf{Реальность:} $\Delta g_{12}$ может быть отрицательным (диапазон $[-20000, +30000]$ Дж/моль).

\subsection{Классификация систем}

\begin{center}
\small
\begin{tabularx}{0.98\textwidth}{|P{4.2cm}|P{2.7cm}|C{1.7cm}|Y|}
\hline
Тип системы & \shortstack[l]{$\Delta g_{12}$,\\$\Delta g_{21}$} & $\gamma_2$ & \shortstack[l]{Bridge-функция\\потерь} \\
\hline
Неполярный solute в неполярном solvent & $\geq 0$ & $\geq 1$ & Корректен \\
\hline
Полярный solute в воде & $> 0$ (обычно) & $\gg 1$ & Приемлем \\
\hline
H-bond donor в H-bond acceptor & Может быть $< 0$ & Может быть $< 1$ & \textbf{Ошибочен} \\
\hline
Спирт + карбоновая кислота & $\Delta g_{21} < 0$ & $\gamma_2 < 1$ & \textbf{Ошибочен} \\
\hline
\end{tabularx}
\end{center}

\textbf{Доля фармацевтически релевантных систем с $\Delta g < 0$:} оценочно 30--50\% (системы с водородными связями: вода, спирты, ДМСО с полярными лекарствами).

\subsection{Механизм вреда}

Для системы с истинным $\Delta g_{21} = -1200$ Дж/моль (типичный пример --- парацетамол в этаноле):

\begin{itemize}
  \item Bridge-функция потерь задаёт target: $\Delta g_{21}^{\mathrm{(Hansen)}} > 0$ (всегда).
  \item Bridge-функция потерь тянет $\Delta g_{21}^{\mathrm{pred}}$ к положительным значениям.
  \item Это завышает $\gamma_2$ $\to$ занижает растворимость.
  \item $\mathcal{L}_{\mathrm{sol}}$ тянет обратно, создавая конфликт градиентов.
\end{itemize}

\textbf{Результат:} ни bridge-функции потерь, ни $\mathcal{L}_{\mathrm{sol}}$ не «побеждают» полностью --- модель застревает в компромиссе, далёком от оптимума.

\section{H6: Систематическая ошибка приближения $\Delta C_p^{\mathrm{fus}} = 0$}

\subsection{Количественная оценка}

Для молекул с высоким $T_{\mathrm{m}}$ (типично для фармацевтических соединений, $T_{\mathrm{m}} \sim 400$--$600$ К) и $\Delta C_p^{\mathrm{fus}} \sim 50$--$100$ Дж/(моль$\cdot$К):

\begin{equation*}
\Delta(\ln x_2^{\mathrm{ideal}}) = \frac{\Delta C_p^{\mathrm{fus}}}{R} \cdot \psi\!\left(\frac{T_{\mathrm{m}}}{T}\right)
\end{equation*}

Функция $\psi(u) = (u - 1) - \ln u$ при $u = T_{\mathrm{m}}/T$:

\begin{center}
\begin{tabular}{|c|c|c|c|c|}
\hline
$T_{\mathrm{m}}$ (К) & $T$ (К) & $u = T_{\mathrm{m}}/T$ & $\psi(u)$ & $\Delta(\ln x_2)$ при $\Delta C_p = 80$ \\
\hline
350 & 298 & 1.17 & 0.014 & 0.13 \\
\hline
400 & 298 & 1.34 & 0.049 & 0.47 \\
\hline
450 & 298 & 1.51 & 0.096 & 0.93 \\
\hline
500 & 298 & 1.68 & 0.155 & 1.49 \\
\hline
\end{tabular}
\end{center}

При $T_{\mathrm{m}} > 400$ К ошибка \textbf{превышает типичную целевую MAE} модели.

\subsection{Компенсаторный эффект}

Модель вынуждена компенсировать через NRTL: чтобы поднять растворимость до истинного значения, нужно $\gamma_2 < \gamma_2^{\mathrm{true}}$. Если $\gamma_2^{\mathrm{true}} = 2$, а нужно $\gamma_2^{\mathrm{effective}} = 2 / e^{0.9} \approx 0.8$ --- модель должна предсказать \textbf{нефизичный} $\gamma_2 < 1$ для системы, где он физически $> 1$.

Это невозможно при ограниченных $\tau$ $\to$ систематическая ошибка.

\section{Сопоставление гипотез с данными}

\subsection{Наблюдение: RF(descriptors) $\gg$ DirectGNN $>$ TGNN}

\textbf{H5 (бюджет):} RF обучен полностью, GNN --- нет. Объясняет часть разрыва.

\textbf{H1 (чувствительность):} TGNN должна решить 5+ задач QSPR; DirectGNN --- одну. При одинаковом бюджете TGNN сходится медленнее.

\textbf{H3 (gradient chaos):} SLE решатель с мусорными входами замедляет TGNN дополнительно.

\textbf{Совокупность:} наблюдение согласуется со всеми тремя гипотезами.

\subsection{Наблюдение: Morgan не помогает}

\textbf{H1:} Morgan fingerprints несут субструктурную информацию, но bottleneck --- не в молекулярном представлении, а в NRTL-параметризации. Morgan не улучшает предсказание $\tau_{12}$, $\tau_{21}$.

\textbf{Дополнительно:} RF(Morgan) $<$ RF(descriptors) (MAE 1.281 по сравнению с 1.198), что подтверждает: для данной задачи глобальные дескрипторы информативнее подструктурных хэшей.

\subsection{Наблюдение: descriptor/group priors нейтральны или вредны}

\textbf{H1:} priors улучшают Hansen и $V_m$, но эти величины \textbf{не являются основным узким местом}. Bottleneck --- в NRTL параметрах, и Hansen/$V_m$ priors на них не влияют.

\textbf{H4:} если Hansen priors точнее $\to$ bridge-функции потерь target точнее $\to$ конфликт с истинными $\Delta g < 0$ усиливается $\to$ может быть хуже.

\subsection{Наблюдение: TGNN (solubility-only) $\approx$ DirectGNN}

\textbf{Интерпретация:} при удалении auxiliary-функций потерь SLE-решатель остаётся, но crystal heads не получают прямой супервизии $\to$ $T_{\mathrm{m}}$ и $\Delta H_{\mathrm{fus}}$ обучаются только через SLE $\to$ та же проблема H1/H2, но без дополнительного шума от auxiliary-функций потерь $\to$ MAE $\approx$ DirectGNN.

Это показывает: \textbf{SLE решатель при текущих условиях --- нейтральная прокладка} (не помогает и не сильно вредит), потому что предсказываемые параметры физически бессмысленны.

\subsection{Наблюдение: TGNN улучшается с увеличением бюджета}

MAE падает с 4.435 (1/4/1) до 2.256 (medium run). Это согласуется с H5 (бюджет), но \textbf{TGNN всё ещё хуже DirectGNN} (2.256 по сравнению с 2.069), что указывает на сохраняющиеся архитектурные проблемы (H1, H2, H3).

\section{Обсуждение результатов}

\subsection{Главный промежуточный результат medium-budget прогона}

Для medium-budget прогона (режим 10/40/10) на момент фиксации результатов получен следующий ранний сигнал на старте фаза 2:
\begin{equation*}
\mathrm{MAE}=1.973, \quad R^2=0.216, \quad \texttt{sol\_fraction}=92.6\% \quad (\text{на входе в фаза 2, до первого полного обновления эпохи}).
\end{equation*}

Ключевое наблюдение --- после исправления балансировки лоссов TGNN впервые выходит в область $R^2>0$, то есть начинает объяснять вариацию целевой переменной лучше постоянного базового ориентира. Для текущей стадии обучения это принципиально: речь идёт не о «косметическом» улучшении, а о смене режима оптимизации.

Второй важный аспект --- абсолютный уровень ошибки ($\mathrm{MAE}=1.973$) уже на ранней точке фаза 2 близок к сильному ориентиру DirectGNN. Это не доказывает превосходство TGNN, но показывает, что модель перестала находиться в явно деградировавшем состоянии, характерном для предыдущих неустойчивых конфигураций.

Для этого прогона корректная формулировка следующая: наблюдается \textbf{переход к работоспособному режиму}, а вопрос «устойчивый тренд на улучшение или последующая деградация» должен оцениваться по полной траектории фаза 2 и фаза 3.

\subsection{Гиперпараметры как доминирующий фактор}

Быстрый первичный Optuna-tuning дал скачок качества порядка
\begin{equation*}
\mathrm{MAE}:\; 5.37 \rightarrow 2.12,
\end{equation*}
то есть улучшение примерно в 2.5 раза. В рамках текущего цикла это самый сильный единичный фактор, превосходящий эффект большинства локальных архитектурных модификаций.

Методологический вывод жёсткий: любые сравнения «архитектура A против архитектуры B» без сопоставимой гиперпараметрической настройки статистически и содержательно невалидны. Иными словами, плохая конфигурация обучения способна полностью скрыть потенциально полезные свойства архитектуры, а удачная --- частично компенсировать её структурные слабости.

Поэтому в дальнейшем корректно сравнивать модели только в схеме \textit{tuned по сравнению с tuned} и на одинаковом бюджете/сплите.

\subsection{Loss weighting как критическая точка отказа}

До исправления наблюдался классический режим «не той задачи»: резкий рост \texttt{vant\_hoff\_local} (3811 $\rightarrow$ 240148) приводил к тому, что
\begin{equation*}
\texttt{sol\_fraction} < 1\%,
\end{equation*}
и оптимизатор фактически минимизировал регуляризатор, а не целевую ошибку растворимости.

После введения \texttt{clamp} и снижения веса соответствующего штрафа восстановился рабочий баланс:
\begin{equation*}
\texttt{sol\_fraction} > 90\%.
\end{equation*}

Это подтверждает важный для физически-информированного обучения принцип: физические регуляризаторы должны работать как \textbf{мягкая коррекция направления}, а не как \textbf{доминирующий источник градиента}. При нарушении этого условия модель может выглядеть «физичной» по auxiliary-термам, но перестаёт решать основную задачу предсказания $\ln x_2$.

\keyidea{Практическое правило: \texttt{sol\_fraction} должна оставаться выше 50\% на всём горизонте обучения. Уровни ниже 30\% следует трактовать как сигнал немедленного пересмотра весов регуляризаторов.}

\subsection{Physics bottleneck: ранняя оценка без окончательного вердикта}

На прокси-бюджете текущая картина такова: DirectGNN tuned показывает около 1.88, тогда как TGNN tuned --- около 2.12 (порядка 11\% в пользу DirectGNN). Это важный диагностический сигнал, но он не эквивалентен окончательному архитектурному выводу.

Причина в том, что прокси-режим и ранние этапы запусков среднего бюджета в первую очередь измеряют скорость вхождения в рабочий режим, а не асимптотическое качество физически-информированной декомпозиции. Именно поэтому решающее значение имеет завершённый средний/полный бюджет, где проверяется, начинает ли physics bottleneck давать выигрыш после стабилизации внутренних параметров.

Таким образом, на текущем этапе корректная интерпретация: \textbf{преимущество DirectGNN подтверждено для коротких/быстрых режимов}, но вопрос о поведении TGNN на полном горизонте остаётся открытым.

\subsection{Анализ ошибок}

Для понимания режимов отказа модели проанализированы 100 примеров с наибольшей абсолютной ошибкой.

\textbf{Категории ошибок:}
\begin{enumerate}
    \item \textbf{Ошибка $T_{\mathrm{m}}$ (35\%):} предсказанная $T_{\mathrm{m}}$ отличается от экспериментальной более чем на 50 K.
    \item \textbf{Экстремальные $\tau$ (28\%):} $|\tau| > 10$, что соответствует нефизичным энергиям взаимодействия.
    \item \textbf{Несходимость решателя (12\%):} итерации не достигают заданной точности за $N_{\max}$ шагов.
    \item \textbf{Редкие химические классы (15\%):} solute содержит функциональные группы, редко встречающиеся в train set.
    \item \textbf{Шум данных (10\%):} вероятные ошибки в исходных экспериментальных данных.
\end{enumerate}

\subsection{Групповые априорные оценки и динамика residual-компоненты}

Standalone-качество для GC-оценок ($\mathrm{MAE}\approx47\,\mathrm{K}$) ухудшается до $\sim63\,\mathrm{K}$ при добавлении необученного residual в коротком режиме. Это указывает на то, что в начале обучения residual-ветка чаще добавляет шум, чем полезную поправку.

Практический вывод: для стабилизации ранней фазы целесообразно временно \textbf{замораживать residual} в начале фаза 1 (или включать его по schedule после формирования базового GC-сигнала). Такой gating уменьшает риск разрушения информативного prior-а случайными градиентами и улучшает кондиционирование последующего joint-training.

\subsection{Bridge-loss: согласование теории и практики}

Теоретическое ограничение bridge-функции потерь остаётся фундаментальным: целевой Hansen-терм задаёт
\begin{equation*}
\Delta g^{\mathrm{(Hansen)}} \ge 0
\end{equation*}
по построению, тогда как для части химически релевантных систем эффективные взаимодействия соответствуют режиму $\Delta g<0$.

Практическое следствие --- конфликт градиентов: bridge-компонента тянет параметры в «положительную» область, а $\mathcal{L}_{\mathrm{sol}}$ требует противоположного направления для воспроизведения наблюдаемой растворимости. Поэтому рекомендация отключить bridge-функцию потерь в текущем виде и заменить её более нейтральной физической проверкой (например, Walden-type consistency check) методологически обоснована.

\subsection{Интерпретация Oracle-диагностики}

На прокси-режиме наблюдение вида
\begin{equation*}
\mathrm{MAE}_{\mathrm{standard}} \approx \mathrm{MAE}_{\mathrm{oracle}}
\end{equation*}
интерпретируется как признак компенсаторной дегенерации: ошибка crystal-head почти не влияет на итоговую метрику, потому что NRTL-параметры уже переобучились на компенсацию.

Однако этот вывод пока предварительный. При более длинном обучении (среднего/полного) связь между качеством $T_{\mathrm{m}}$ и итоговым $\ln x_2$ может усилиться. Поэтому обязательный следующий шаг --- повторить oracle evaluation после завершения текущего medium-budget прогона и отдельно проверить чувствительность на фаза 3 checkpoint.

\subsection{Ограничения текущего анализа}

Текущее обсуждение следует рассматривать как промежуточное по трём причинам.

\begin{enumerate}
  \item \textbf{Ограниченный бюджет.} Режим 10/40/10 существенно ниже итоговых 50/200/50 и может недооценивать поздние преимущества/недостатки архитектур.
  \item \textbf{Один seed.} Для устойчивых выводов нужны многозапусковый оценки и доверительные интервалы, иначе часть различий может объясняться стохастикой инициализации.
  \item \textbf{Нет анализа температурной экстраполяции.} Без отдельной проверки диапазона $T$ вне обучения нельзя судить о том, реализовалось ли главное преимущество физически-информированного подхода.
\end{enumerate}

Итоговая позиция на данном этапе: после исправления функции потерь weighting TGNN перешла из режима деградации в режим содержательно осмысленного обучения, но вопрос о её финальном преимуществе над tuned DirectGNN остаётся эмпирическим и должен решаться на завершённом среднем/полном протоколе с многозапусковый статистикой.


\section{Оценка неопределённости и область применимости}

\subsection{Зачем это нужно в прикладном контуре}

Для практических решений в химико-технологическом контуре точечной оценки $\ln x_2$ недостаточно. Пользователь модели должен получить ответ на два независимых вопроса:
\begin{enumerate}
  \item насколько надёжен конкретный прогноз (\textit{эпистемической неопределённости});
  \item находится ли система в области, где модель вообще компетентна (\textit{applicability domain}, AD).
\end{enumerate}

Эти сигналы дополняют друг друга: AD отвечает за \textbf{валидность применения}, а неопределённость --- за \textbf{степень доверия} внутри/около этой области.

\subsection{MC-Dropout как базовая оценка эпистемической неопределённости}

В базовой реализации используется MC-Dropout: в отличие от стандартного режима inference, dropout здесь намеренно не отключается, и для одного входа выполняется $N$ стохастических прямых проходов.

Для объекта $i$:
\begin{equation*}
\hat{y}_i^{(n)} = f_{\theta}^{(n)}(\mathbf{x}_i), \quad n = 1,\ldots,N.
\end{equation*}

Оценки среднего и дисперсии:
\begin{equation*}
\bar{y}_i = \frac{1}{N}\sum_{n=1}^{N}\hat{y}_i^{(n)}, \quad
\sigma_i^2 = \frac{1}{N-1}\sum_{n=1}^{N}\left(\hat{y}_i^{(n)} - \bar{y}_i\right)^2.
\end{equation*}

Интерпретация стандартного отклонения $\sigma_i$:
\begin{itemize}
  \item большое $\sigma_i$ --- высокая неуверенность модели и повышенный риск ошибки;
  \item малое $\sigma_i$ --- прогноз стабилен к стохастике dropout.
\end{itemize}

В терминах Bayesian DL это рассматривается как вариационная аппроксимация апостериорного распределения по весам~\cite{gal2016}.

На практике для TGNN-Solv неопределённость оценивается не только для $\ln x_2$, но и для промежуточных/производных величин:
\begin{equation*}
\{\ln x_2,\; x_2,\; \gamma_2,\; T_{\mathrm{m}},\; \Delta H_{\mathrm{fus}},\; \Phi,\; \ln\gamma_2,\; |\Delta_{\mathrm{corr}}|\}.
\end{equation*}

Интервальная оценка строится по эмпирическим перцентилям распределения предсказаний (обычно 5-й и 95-й перцентили при $N=30$ проходах).

\subsection{Deep Ensembles как более надёжная альтернатива}

Более вычислительно дорогой, но обычно более устойчивый способ --- ансамбль независимо обученных моделей с разными random seeds~\cite{lakshminarayanan2017}.

Для $K$ моделей $\{\theta_k\}_{k=1}^{K}$:
\begin{equation*}
\bar{y}_i = \frac{1}{K}\sum_{k=1}^{K} f_{\theta_k}(\mathbf{x}_i), \quad
\sigma_i^2 = \frac{1}{K-1}\sum_{k=1}^{K}\left(f_{\theta_k}(\mathbf{x}_i)-\bar{y}_i\right)^2.
\end{equation*}

Плюсы и минусы:
\begin{itemize}
  \item \textbf{Плюс:} лучше покрывает разнообразие модельных гипотез по сравнению с MC-Dropout.
  \item \textbf{Минус:} стоимость примерно $K\times$ по обучению и inference.
\end{itemize}

С инженерной точки зрения удобно использовать двухуровневую схему: MC-Dropout как базовую оценку неопределённости, Deep Ensembles --- как эталон для критичных случаев.

\subsection{Калибровка неопределённости: PICP и MPIW}

Чтобы неопределённость-интервалы были практически полезны, их нужно калибровать и регулярно валидировать на тестовых данных.

Важно, что низкая неопределённость сама по себе не гарантирует малую ошибку. Калибровка --- отдельная задача: модель может быть переуверенной (PICP ниже целевого уровня) или недоуверенной (PICP заметно выше целевого уровня). Для OOD-сценариев предварительный анализ показывает склонность к переуверенности, что согласуется с общими наблюдениями литературы по uncertainty estimation для GNN.

\textbf{Coverage (PICP) для 90\% интервала:}
\begin{equation*}
\mathrm{PICP}_{90} = \frac{1}{N_{\mathrm{test}}}\sum_{i=1}^{N_{\mathrm{test}}}
\mathbf{1}\!\left[y_i \in [\hat{y}_i^{(5\%)},\hat{y}_i^{(95\%)}]\right].
\end{equation*}

Идеальное поведение: $\mathrm{PICP}_{90}\approx 0.90$. Ниже 0.90 --- интервалы переуверенные (слишком узкие), заметно выше 0.90 --- избыточно консервативные.

\textbf{Sharpness (MPIW):}
\begin{equation*}
\mathrm{MPIW} = \frac{1}{N_{\mathrm{test}}}\sum_{i=1}^{N_{\mathrm{test}}}
\left(\hat{y}_i^{(95\%)}-\hat{y}_i^{(5\%)}\right).
\end{equation*}

При сопоставимом PICP предпочтительна меньшая MPIW: это означает более информативную (более «острую») неопределённость-оценку.

\subsection{Applicability Domain: формализация области компетенции}

AD отвечает на бинарный вопрос: допустимо ли применять текущую модель к новой системе $(\mathrm{solute},\mathrm{solvent},T)$ как к in-distribution задаче.

В рабочем протоколе TGNN-Solv комбинируются два сигнала.

\textbf{(a) Расстояние Махаланобиса в pair-embedding пространстве.}
\begin{equation*}
D_M(\mathbf{g}_{\mathrm{pair}})=\sqrt{
(\mathbf{g}_{\mathrm{pair}}-\boldsymbol{\mu})^{\top}
\boldsymbol{\Sigma}^{-1}
(\mathbf{g}_{\mathrm{pair}}-\boldsymbol{\mu})
}.
\end{equation*}

Здесь $\boldsymbol{\mu}$ и $\boldsymbol{\Sigma}$ --- среднее и ковариация парных представлений на обучающей выборке. Типичный порог:
\begin{equation*}
D_M \le D_{\mathrm{cutoff}},
\end{equation*}
где $D_{\mathrm{cutoff}}$ выбирается как 99-й перцентиль распределения на обучающей выборке.

\textbf{(b) Структурная близость по Tanimoto.}
Для Morgan fingerprints:
\begin{equation*}
T(\mathbf{f}_1,\mathbf{f}_2)=\frac{|\mathbf{f}_1\cap\mathbf{f}_2|}{|\mathbf{f}_1\cup\mathbf{f}_2|}.
\end{equation*}

Для нового запроса вычисляются максимальные сходства с обучающим набором отдельно для растворяемого вещества и растворителя:
\begin{equation*}
T_{\mathrm{solute}}=\max_j T(\mathbf{f}_{\mathrm{solute}},\mathbf{f}_{\mathrm{solute},j}^{\mathrm{train}}),
\quad
T_{\mathrm{solvent}}=\max_j T(\mathbf{f}_{\mathrm{solvent}},\mathbf{f}_{\mathrm{solvent},j}^{\mathrm{train}}).
\end{equation*}

Условие по структурной близости:
\begin{equation*}
T_{\mathrm{solute}}\ge T_{\min} \quad \wedge \quad T_{\mathrm{solvent}}\ge T_{\min}.
\end{equation*}

Итоговый AD-критерий:
\begin{equation*}
\mathrm{in\_domain} = (D_M\le D_{\mathrm{cutoff}}) \;\wedge\;
(T_{\mathrm{solute}}\ge T_{\min}) \;\wedge\;
(T_{\mathrm{solvent}}\ge T_{\min}).
\end{equation*}

\subsection{Непрерывная оценка уверенности}

Помимо бинарного флага принадлежности к области применимости удобна и скалярная мера уверенности:
\begin{equation*}
C=\frac{1}{2}\left(
\frac{D_{\mathrm{cutoff}}-D_M}{D_{\mathrm{cutoff}}}+
\min(T_{\mathrm{solute}},T_{\mathrm{solvent}})
\right).
\end{equation*}

В практическом использовании значения $C$ ниже заданного порога могут автоматически переводить прогноз в режим \textit{требуется экспертная проверка}.

\subsection{Связь AD и неопределённость в эксплуатационном протоколе}

AD и неопределённость --- разные по смыслу величины:
\begin{itemize}
  \item AD: бинарная проверка применимости (фильтр уровня политики);
  \item $\sigma$: непрерывная оценка ненадёжности (ранжирование внутри допустимых случаев).
\end{itemize}

Типичный порядок действий:
\begin{enumerate}
  \item проверить AD (если объект вне домена применимости, выдать предупреждение или отклонить автоматическое решение);
  \item для случаев в домене применимости ранжировать прогнозы по $\sigma$ и/или ширине PI;
  \item для высокорискованных решений использовать ансамблевое подтверждение или лабораторную валидацию.
\end{enumerate}

Ожидаемая эмпирическая закономерность: вне AD неопределённость обычно выше, но строгой эквивалентности нет. Возможны редкие случаи в домене применимости с высокой $\sigma$ (конфликтные/шумные режимы) и случаи вне домена применимости с умеренной $\sigma$ (ложная уверенность). Поэтому оба индикатора должны использоваться совместно.

\subsection{Рекомендованный формат отчётности}

Для каждого экспериментального прогона неопределённость/AD блок рекомендуется фиксировать в стандартизованном виде:
\begin{itemize}
  \item $\mathrm{PICP}_{90}$ и MPIW на test-разбиении;
  \item диаграмма coverage по сравнению с nominal confidence;
  \item распределения $D_M$, $T_{\mathrm{solute}}$, $T_{\mathrm{solvent}}$;
  \item доля случаев в домене применимости / вне домена применимости;
  \item зависимость MAE от квантилей $\sigma$ и/или оценки доверия $C$.
\end{itemize}

Такой протокол делает выводы о надёжности модели воспроизводимыми и пригодными для инженерного принятия решений, а не только для академического сравнения средних ошибок.

\section{Конвейер инференса}

\subsection{Назначение и общая схема}

После завершения обучения TGNN-Solv используется как детерминированный (или uncertainty-aware) предиктор растворимости для новой тройки входов:
\begin{equation*}
(\mathrm{solute\_smiles},\; \mathrm{solvent\_smiles},\; T).
\end{equation*}

Инференс-пайплайн включает четыре последовательных блока:
\begin{enumerate}
  \item восстановление модели из checkpoint;
  \item подготовка входных графов и опциональных признаков;
  \item один или серия прямых проходов через физический слой;
  \item постобработка и формирование интерпретируемого отчёта.
\end{enumerate}

\subsection{Алгоритм 1: инференс для одной точки (пошаговая схема)}

Для запроса $(\mathrm{solute\_smiles},\mathrm{solvent\_smiles},T)$ практический порядок вычислений можно записать так:

\begin{enumerate}
  \item Загрузить checkpoint: веса, \texttt{TGNNSolvConfig}, нормализаторы.
  \item Каноникализация SMILES и проверка валидности (RDKit).
  \item Построить графы solute/solvent через \texttt{smiles\_to\_graph}.
  \item Добавить опциональные признаки согласно конфигурации (Morgan/priors/GC).
  \item Выполнить \texttt{model.eval()} и один прямой проход.
  \item Извлечь выходы: $\widehat{\ln x_2}, \hat{x}_2, \Phi, \ln\gamma_2, \delta_{\mathrm{corr}}$,  а также $T_{\mathrm{m}}, \Delta H_{\mathrm{fus}}, \Delta C_p^{\mathrm{fus}}, \tau_{12}, \tau_{21}, \alpha$.
  \item Сформировать интерпретационный отчёт (масштаб растворимости, доминирующий вклад, роль correction).
  \item (Опционально) добавить неопределённость и AD: \{PI, $\sigma$, in\_domain, C\}.
\end{enumerate}

\subsection{Загрузка модели и конфигурации}

Из checkpoint восстанавливаются:
\begin{itemize}
  \item веса нейросети;
  \item объект конфигурации \texttt{TGNNSolvConfig};
  \item нормализаторы дескрипторов (если они использовались при обучении).
\end{itemize}

Именно конфигурация определяет активные optional-paths в inference:
\begin{equation*}
\{\texttt{use\_morgan},\;\texttt{use\_descriptor\_priors},\;\texttt{use\_group\_priors},\;\texttt{use\_crystal\_gc\_priors}\}.
\end{equation*}

Это критично для воспроизводимости: входной feature-space в инференсе должен совпадать с тем, на котором модель обучалась.

\subsection{Подготовка входных данных}

Для каждого запроса выполняется стандартный preprocessing.

\textbf{Шаг 1: каноникализация SMILES.}
Оба SMILES нормализуются RDKit, невалидные структуры отклоняются до запуска модели.

\textbf{Шаг 2: построение графов.}
Каждая молекула переводится в graph-объект (например, \texttt{PyG Data}) через
\begin{equation*}
\texttt{smiles\_to\_graph(smi)} \rightarrow \texttt{Data}.
\end{equation*}

\textbf{Шаг 3: добавление опциональных признаков (по config).}
Возможные дополнительные каналы:
\begin{itemize}
  \item Morgan fingerprints;
  \item descriptor prior features;
  \item group prior features;
  \item crystal GC priors: $T_{\mathrm{m}}^{\mathrm{gc}}$, $\Delta H_{\mathrm{fus}}^{\mathrm{gc}}$, $\Delta C_p^{\mathrm{fus,gc}}$.
\end{itemize}

Входной батч в онлайн-режиме обычно имеет размер 1, но пайплайн допускает и mini-batch для пакетных API-запросов.

\subsection{Прямой проход для одной точки}

Один прямой проход возвращает не только финальный прогноз, но и его физическую декомпозицию.

\textbf{Финальные выходы:}
\begin{equation*}
\hat{y}=\widehat{\ln x_2}, \qquad \hat{x}_2=\exp(\hat{y}).
\end{equation*}

\textbf{Декомпозиция TGNN-Solv:}
\begin{equation*}
\ln x_2 = \underbrace{-\Phi(T)}_{\text{ideal}} +
\underbrace{(-\ln\gamma_2)}_{\text{nonideal}} +
\underbrace{\delta_{\text{corr}}}_{\text{correction}}.
\end{equation*}

Также модель отдаёт внутренние физические параметры:
\begin{itemize}
  \item crystal: $T_{\mathrm{m}}$, $\Delta H_{\mathrm{fus}}$, $\Delta C_p^{\mathrm{fus}}$;
  \item NRTL: $\tau_{12}$, $\tau_{21}$, $\alpha$;
  \item Hansen: $(\delta_d,\delta_p,\delta_h)$ и derived-метрики (например, $R_a$);
  \item correction diagnostics: gate value и величина bounded correction.
\end{itemize}

Важно: в inference TGNN-Solv возвращает именно solver-side величину $\gamma_2(x_2^{*})$, вычисленную при предсказанном составе насыщения $x_2^{*}$ и используемую в декомпозиции $\ln x_2 = -\Phi - \ln\gamma_2 + \delta_{\mathrm{corr}}$. IDAC-величина $\gamma_2^{\infty}$ не является стандартным выходом inference; при необходимости она может быть отдельно вычислена из предсказанных NRTL-параметров как предел $x_2 \to 0$.

Если включены GC-априори, дополнительно возвращаются prior-оценки
\begin{equation*}
T_{\mathrm{m}}^{\mathrm{gc}},\; \Delta H_{\mathrm{fus}}^{\mathrm{gc}},\; \Delta C_p^{\mathrm{fus,gc}}.
\end{equation*}

\subsection{Температурное сканирование для кривой растворимости}

Для построения зависимости $\ln x_2(T)$:
\begin{enumerate}
  \item фиксируется пара $(\mathrm{solute},\mathrm{solvent})$;
  \item выбирается диапазон $T\in[T_{\min},T_{\max}]$;
  \item формируется сетка из \texttt{n\_points};
  \item для каждого $T_k$ выполняется полный прямой проход.
\end{enumerate}

Важно: такая серия независимых forward-pass не гарантирует строгую монотонность кривой $\ln x_2(T)$ сама по себе. Поэтому после сканирования обязательно выполняются sanity-check'и монотонности и, при необходимости, применяется постобработка (например, монотонная аппроксимация в допустимом температурном окне).

Результат собирается в таблицу (DataFrame) с типовыми полями:
\begin{equation*}
\{T,\; x_2,\; \ln x_2,\; x_2^{\mathrm{ideal}},\; \gamma_2,\; \delta_{\mathrm{corr}}\}.
\end{equation*}

Физический sanity-check для большинства систем: монотонный рост растворимости с температурой,
\begin{equation*}
\frac{d\ln x_2}{dT} \approx \frac{\Delta H_{\mathrm{sol}}}{RT^2} > 0,
\end{equation*}
что согласуется с ван'т-Гоффовским приближением для эндотермического растворения.

\subsection{Интерпретация прогноза в человекочитаемом отчёте}

Для прикладного пользователя численное значение $\ln x_2$ дополняется объяснением структуры прогноза.

\textbf{1) Категория растворимости.}
Модельный выход можно маппировать в шкалу:
\begin{equation*}
\text{очень низкая} \rightarrow \text{низкая} \rightarrow \text{средняя} \rightarrow \text{высокая}.
\end{equation*}

\textbf{2) Доминирующий вклад.}
Сравниваются модули идеального и неидеального слагаемых:
\begin{equation*}
|\Phi| \quad \text{по сравнению с} \quad |\ln\gamma_2|.
\end{equation*}

Если $|\Phi|\gg|\ln\gamma_2|$, ограничение главным образом кристаллическое; если наоборот --- доминирует вклад межмолекулярных взаимодействий в растворе.

\textbf{3) Разумность crystal-параметров.}
Проверяется величина энтропии плавления:
\begin{equation*}
\Delta S_{\mathrm{fus}} = \frac{\Delta H_{\mathrm{fus}}}{T_{\mathrm{m}}}
\end{equation*}
относительно эмпирического ориентира порядка $56.5\,\mathrm{J\,mol^{-1}K^{-1}}$.

\textbf{4) Hansen-совместимость.}
Оценивается $R_a$ в сравнении с типичным радиусом совместимости $R_0$ (domain-dependent).

\textbf{5) Роль коррекции.}
Контролируется доля корректирующего слагаемого:
\begin{equation*}
\mathrm{corr\_ratio} = \frac{|\delta_{\mathrm{corr}}|}{|\ln x_2|+\varepsilon}.
\end{equation*}

Большие значения \texttt{corr\_ratio} сигнализируют, что модель сильно опирается на residual-path, а не на физический core.

\subsection{Что не происходит в режиме инференса}

Чтобы не смешивать во время обучения и во время оценки логику, важно явно зафиксировать исключения:
\begin{enumerate}
  \item \textbf{Оракульная подстановка} не используется (это только диагностический режим обучения/валидации).
  \item \textbf{Pair-aware temperature batching} не нужен для одиночного запроса.
  \item \textbf{Curriculum phases} неприменимы: модель работает в \texttt{eval mode} с фиксированными весами.
\end{enumerate}

Таким образом, inference pipeline является чистым применением уже обученного физико-ML отображения без дополнительных оптимизационных шагов.

\subsection{Рекомендованный эксплуатационный протокол}

Для производственного использования целесообразно объединить предсказание, неопределённость и AD в единый ответ API:
\begin{equation*}
\{\hat{\ln x_2},\; \hat{x}_2,\; \gamma_2(x_2^{*}),\; \mathrm{PI}_{5\text{--}95},\; \mathrm{in\_domain},\; C,\; \text{decomposition}\}.
\end{equation*}

Это позволяет пользователю одновременно видеть значение прогноза, его надёжность и границы применимости модели, что критично для принятия решений в экспериментальном планировании.

\subsection{Вычислительная стоимость инференса}

На текущем этапе приведены инженерные оценки порядка величины, а не финальные microbenchmark-измерения. Значения ниже следует трактовать как \textit{estimated order-of-magnitude} для типичной пары молекул с \textasciitilde30 тяжёлыми атомами на молекулу.

\begin{table}[H]
\centering
\begin{tabular}{|p{4.2cm}|p{3cm}|p{3cm}|p{4cm}|}
\hline
Операция & CPU (i7-class) & GPU (A100-class) & Примечание \\
\hline
Single forward pass & \textasciitilde10--30 ms & \textasciitilde1--5 ms & Одна тройка $(\mathrm{solute},\mathrm{solvent},T)$ \\
\hline
MC-Dropout ($N=30$) & \textasciitilde0.3--0.9 s & \textasciitilde30--150 ms & Один запрос с uncertainty \\
\hline
Temperature scan (50 точек) & \textasciitilde0.5--1.5 s & \textasciitilde50--250 ms & Одна пара при сетке температур \\
\hline
Batch inference (1000 пар) & \textasciitilde15--40 s & \textasciitilde0.8--3 s & Включая построение графов и батчинг \\
\hline
RF inference (1000 пар) & \textasciitilde10--100 ms & N/A & Ориентир для табличной baseline-модели \\
\hline
\end{tabular}
\caption{Оценка latency и throughput inference; значения являются order-of-magnitude estimate и требуют отдельного benchmark-run для публикационного использования.}
\end{table}

\section{Критическая оценка: что уже проверено и что ещё предстоит проверить}

\subsection{Предварительно установлено (по прокси/среднему бюджету)}

\begin{enumerate}
  \item \textbf{При текущем бюджете обучения TGNN не конкурентоспособна.} Это факт для текущего режима, но не свидетельство провала концепции: бюджет пока недостаточен.
  \item \textbf{Morgan fingerprints и Hansen/$V_m$ priors, по текущим данным, не являются основным ограничением.} Наблюдение устойчиво в нескольких выполненных прогонах, но требует подтверждения на полном бюджете.
  \item \textbf{При коротком обучении вклад physics bottleneck близок к нейтральному.} TGNN (solubility-only) и DirectGNN показывают близкие метрики (разница порядка $0.01$ MAE), поэтому по текущим данным нельзя утверждать о заметном выигрыше физического слоя в этом режиме.
\end{enumerate}

\subsection{Не установлено (требует полного цикла экспериментов)}

\begin{enumerate}
  \item \textbf{Помогает ли physics bottleneck при полном обучении?} Без эксперимента 50/200/50 невозможно отличить H5 от H1/H2.
  \item \textbf{Идентифицируемы ли промежуточные параметры?} Не проверено, коррелируют ли предсказанные $T_{\mathrm{m}}$, $\Delta H_{\mathrm{fus}}$ с экспериментальными при полном обучении.
  \item \textbf{Работоспособен ли NRTL-путь при правильных crystal parameters?} Oracle-эксперимент не проведён.
  \item \textbf{Вредна ли bridge-функция потерь?} Абляция с \texttt{bridge\_weight = 0} пока не проведена.
\end{enumerate}

\section{Обновлённый диагноз гипотез}
\label{sec:updated-hypotheses}

\subsection{Обновлённый рейтинг гипотез}

\begin{center}
\small
\begin{tabular}{|c|p{7.6cm}|p{4.0cm}|}
\hline
\textbf{Ранг} & \textbf{Гипотеза} & \textbf{Статус} \\
\hline
H8 & графовый энкодер не содержит достаточной химической информации & новая, подтверждена linear probe \\
\hline
H5 & Гиперпараметры (dropout, lr) критичны; proxy-optimum не переносится на full & подтверждена \\
\hline
H1 & Physics bottleneck вреден & частично смягчена: IDAC-supervision теперь активен, но покрытие остаётся разреженным (1.2\% пар) \\
\hline
H2 & Неидентифицируемость параметров & сохраняется, частично производна от H8 \\
\hline
H4 & bridge-loss вреден & не опровергнута, требует full-check \\
\hline
H3 & Gradient chaos SLE solver & снижена после loss-weighting fix \\
\hline
H6 & $\Delta C_p=0$ как источник ошибки & не проверена \\
\hline
H7 & NRTL overparameterized & не проверена \\
\hline
\end{tabular}
\normalsize
\end{center}

Главный пересмотр: H1 больше не является доминирующим объяснением. Эмпирически
\begin{equation*}
\Delta_{\mathrm{physics}}=0.05\;\mathrm{MAE},
\end{equation*}
что соответствует «малой цене интерпретируемости», а не ключевому источнику деградации.

\subsection{Новая формулировка центральной проблемы}

Старая формулировка: «Physics bottleneck мешает обучению».\newline
Новая формулировка: «Энкодер поставляет физическим головам неполные представления; физический слой получает неточные входы».

Эвристически: SLE-решатель является точным вычислительным ядром. Если подать в него неточные $T_{\mathrm{m}}$ и $\tau$, он выдаёт неточный результат. RF выигрывает не потому, что «термодинамика лишняя», а потому что входные фичи (дескрипторы) содержат более полную химическую информацию.

\subsection{Обновлённый путь решения}

Стратегия смещается от «убрать физику» к «улучшить входы для физики»:
\begin{enumerate}
  \item расширение признаков дескрипторами (добавить $\sim200$ экспертных фичей в pair-path);
  \item Stage 0 pretraining для улучшения chemical grounding;
  \item GC priors для $T_{\mathrm{m}}$ и $\Delta H_{\mathrm{fus}}$ как физически разумная инициализация.
\end{enumerate}

\section*{Часть VI. План экспериментов и предлагаемые архитектурные модификации}
\addcontentsline{toc}{section}{Часть VI. План экспериментов и предлагаемые архитектурные модификации}

Переход к этой части выполняется от диагностических выводов Части V: каждый предложенный эксперимент и каждая модификация напрямую соответствуют выявленному типу отказа (идентифицируемость, оптимизационная устойчивость, смещение физического приближения или дисбаланс супервизии).

\subsection*{Статус на момент отчёта}
В этом разделе описан \textbf{план следующих экспериментов}, а не перечень уже завершённых работ.

Уже выполнено и отражено в Части V:
\begin{itemize}
  \item прокси-прогоны и первичный гиперпараметрический поиск;
  \item ранние/средние запуски с диагностикой динамики функции потерь;
  \item предварительные сравнения TGNN и DirectGNN на ограниченном бюджете.
\end{itemize}

Ещё не завершено (и сформулировано ниже как план):
\begin{itemize}
  \item полный цикл 50/200/50 с многозапусковый статистикой;
  \item финальные фальсификационные тесты (oracle, абляция bridge, альтернативные параметризации);
  \item итоговая верификация промежуточных физических параметров на полном бюджете.
\end{itemize}

\section{Фальсификационные эксперименты}

Ниже перечислены запланированные эксперименты в порядке приоритета. Каждый тестирует конкретную гипотезу и имеет чёткий критерий интерпретации.

Технические механизмы (Оракульная подстановка, GC-priors, альтернативные модели $G^E$, bridge-regularization) уже определены в методических разделах. Поэтому здесь фиксируются только различия протокола, требуемые ресурсы и ожидаемая интерпретация исходов.

\subsection{Эксперимент 1: Полный бюджет обучения}
\label{subsec:exp-full-budget}

\textbf{Тестирует:} H5 (бюджет) по сравнению с H1/H2/H3 (архитектурные).

\textbf{Протокол:}
\begin{itemize}
  \item Обучить TGNN и DirectGNN с полным бюджетом 50/200/50 эпох на каноническом scaffold-разбиении.
  \item Одинаковые гиперпараметры ($d = 256$, $L = 6$, batch size 64).
  \item 3 random seeds для оценки дисперсии.
\end{itemize}

\textbf{После обучения --- диагностика промежуточных параметров:}
\begin{itemize}
  \item Извлечь предсказанные $T_{\mathrm{m}}^{\mathrm{pred}}$ и $\Delta H_{\mathrm{fus}}^{\mathrm{pred}}$ для всех тестовых молекул с известными экспериментальными значениями.
  \item Построить parity plots: $T_{\mathrm{m}}^{\mathrm{pred}}$ по сравнению с $T_{\mathrm{m}}^{\mathrm{true}}$, $\Delta H_{\mathrm{fus}}^{\mathrm{pred}}$ по сравнению с $\Delta H_{\mathrm{fus}}^{\mathrm{true}}$.
  \item Вычислить Pearson $r$ и MAE для промежуточных параметров.
\end{itemize}

\textbf{Критерий интерпретации:}

\begin{center}
\scriptsize
\setlength{\tabcolsep}{1.5pt}
\begin{tabularx}{0.97\textwidth}{|C{1.1cm}|P{3.0cm}|P{2.3cm}|Y|}
\hline
Исход & \shortstack[l]{TGNN по сравнению\\с DirectGNN} & \shortstack[l]{$T_{\mathrm{m}}$\\корреляция} & Вывод \\
\hline
A & TGNN $>$ DirectGNN & Высокая ($r > 0.8$) & H5 подтверждена; архитектура рабочая \\
\hline
B & TGNN $\approx$ DirectGNN & Высокая & Физика не вредит, но и не помогает \\
\hline
C & TGNN $<$ DirectGNN & Высокая & H1 подтверждена: physics bottleneck вреден \\
\hline
D & TGNN $<$ DirectGNN & Низкая ($r < 0.5$) & H2 подтверждена: мусорная факторизация \\
\hline
E & TGNN $\approx$ DirectGNN & Низкая & H2 + компенсация: параметры бессмысленны, но $\ln x_2$ приемлем; action item --- запуск Экспериментов 2 и 5 с жёсткой диагностикой идентифицируемости \\
\hline
\end{tabularx}
\end{center}

\textbf{Ожидание:} исход D или E наиболее вероятен

\textbf{Ресурсы:} \textasciitilde24--48 GPU-часов (A100) на один seed для полного цикла 50/200/50 (фаза 1 + фаза 2 + фаза 3).

\subsection{Эксперимент 2: оракульные параметры кристалла}
\label{subsec:exp-oracle-crystal}

\textbf{Тестирует:} H1 (bottleneck в crystal-head по сравнению с NRTL-heads).

\textbf{Протокол (кратко):} используется тот же механизм оракульной подстановки, формально описанный в §\ref{sec:oracle-teacher-forcing}. На этапе оценки для подмножества с метками предсказанные $T_{\mathrm{m}}$ и $\Delta H_{\mathrm{fus}}$ заменяются истинными значениями, после чего пересчитывается MAE. Полные детали и ограничения режима уже заданы в §\ref{sec:oracle-teacher-forcing} и здесь не дублируются.

\textbf{Ресурсы:} минимальные (\textasciitilde1 GPU-час, только inference).

\subsection{Эксперимент 3: Абляция bridge-функции потерь}
\label{subsec:exp-bridge-ablation}

\textbf{Тестирует:} H4 (bridge-функция потерь --- контрпродуктивна).

\textbf{Протокол:}
\begin{itemize}
  \item Две конфигурации: \texttt{bridge\_weight = 0.05} (по умолчанию) по сравнению с вариантом, где \texttt{bridge\_weight = 0}.
  \item Полный бюджет 50/200/50.
  \item 3 seeds.
\end{itemize}

\textbf{Критерий:}

\begin{center}
\scriptsize
\setlength{\tabcolsep}{1pt}
\begin{tabularx}{0.95\textwidth}{|P{2.6cm}|P{4.2cm}|Y|}
\hline
Исход & \shortstack[l]{MAE без bridge по сравнению\\с вариантом с bridge} & Вывод \\
\hline
MAE улучшается & & H4 подтверждена: bridge вреден \\
\hline
MAE не меняется & & H4 отвергнута: bridge нейтрален \\
\hline
MAE ухудшается & & Bridge полезен как регуляризатор \\
\hline
\end{tabularx}
\end{center}

\textbf{Ресурсы:} \textasciitilde24--48 GPU-часов (параллельно с Экспериментом 1).

\subsection{Эксперимент 4: DirectGNN + дескрипторы RDKit}
\label{subsec:exp-directgnn-rdkit}

\textbf{Тестирует:} верхнюю границу для нейросетевых моделей на данном датасете.

\textbf{Протокол:}
\begin{itemize}
  \item Модифицировать DirectGNN: конкатенировать вектор RDKit дескрипторов (\textasciitilde200 dim) в представление пары перед финальной MLP.
  \item То же scaffold-разбиение, полный бюджет.
\end{itemize}


\textbf{Значение:} определяет, может ли \textbf{любая} GNN-based модель конкурировать с RF на данном датасете.

\textbf{Ресурсы:} \textasciitilde12--24 GPU-часов.

\subsection{Эксперимент 5: априори групповых вкладов для $T_{\mathrm{m}}$, $\Delta H_{\mathrm{fus}}$}
\label{subsec:exp-gc-priors}

\textbf{Тестирует:} предлагаемую модификацию с group-contribution priors для crystal heads.

\textbf{Протокол (кратко):} реализуется полная схема GC-prior + bounded residual, уже описанная в §\ref{subsec:gc-methods} (Joback/Ducros, калибровка prior, границы residual). В этом эксперименте проверяется только инкремент относительно baseline «prediction from scratch» без повторного изложения формул.

\textbf{Ресурсы:} \textasciitilde24--48 GPU-часов

\subsection{Эксперимент 6: Wilson вместо NRTL}
\label{subsec:exp-wilson}

\textbf{Тестирует:} H7 (NRTL overparameterized).

\textbf{Протокол:}
\begin{itemize}
  \item Заменить NRTL на Wilson-параметризацию, уже обсуждённую в разделе сравнения моделей $G^E$.
  \item Сохранить то же scaffold-разбиение и полный бюджет, что и в Эксперименте 1.
  \item Wilson head предсказывает только $\ln\Lambda_{12}$ и $\ln\Lambda_{21}$ вместо тройки NRTL-параметров.
\end{itemize}

\textbf{Критерий:}
\begin{itemize}
  \item Улучшение MAE $\to$ NRTL overparameterized, Wilson лучше идентифицируема.
  \item Ухудшение MAE $\to$ гибкость NRTL нужна.
\end{itemize}

\textbf{Ресурсы:} \textasciitilde24--48 GPU-часов обучения.

\section{Ablation study: систематический протокол}

Ablation study необходима для количественной оценки вклада каждого архитектурного и методического компонента baseline TGNN-Solv. Без такого анализа невозможно отделить действительно полезные элементы от нейтральных деталей реализации и от компонентов, которые лишь увеличивают сложность модели, не улучшая MAE по $\ln x_2$ или даже ухудшая её. Иными словами, без систематической абляции нельзя ответить на ключевой инженерный вопрос: \textit{какая именно часть текущей архитектуры приносит пользу, а какая является источником хрупкости или деградации}.

\validated{Установлено: отдельные компоненты текущего pipeline уже имеют сильную диагностическую мотивацию. В частности, bounded correction, GC priors, bridge-loss, teacher forcing, выбор модели $G^E$ и temperature isolation непосредственно связаны с гипотезами H1--H7 и с фальсификационными экспериментами текущей части. Не установлено: независимый количественный вклад каждого из этих компонентов в итоговый MAE baseline TGNN-Solv при фиксированных разбиении, бюджете и протоколе обучения.}

\keyidea{Систематическая ablation-матрица нужна не как перечень «вариантов модели», а как причинно-интерпретируемый тест вклада каждого компонента. Только после неё можно различить, что в текущем дизайне действительно работает, что является лишь декоративной сложностью, а что методологически контрпродуктивно.}

\subsection{Ablation-матрица}

Ниже фиксируется минимальный набор абляций, который должен быть выполнен относительно канонического baseline TGNN-Solv. Каждая абляция меняет ровно один компонент или один тесно связанный блок конфигурации; все остальные параметры обучения и данные остаются неизменными. В таблице~\ref{tab:ablation-matrix} заранее зафиксированы ожидаемые направления эффекта, чтобы интерпретация не подстраивалась постфактум под полученные числа.

\begin{landscape}
\small
\setlength{\tabcolsep}{3pt}
\setlength{\LTpre}{6pt}
\setlength{\LTpost}{6pt}
\begin{longtable}{|C{0.9cm}|P{2.8cm}|P{2.7cm}|P{3.8cm}|P{8.0cm}|C{1.2cm}|C{1.4cm}|}
\caption{Ablation-матрица для систематической оценки вклада архитектурных и методических компонентов baseline TGNN-Solv.}\label{tab:ablation-matrix}\\
\hline
ID & Что отключается/заменяется & Тестируемая гипотеза & Конкретное изменение конфигурации & Ожидаемый эффект & MAE & $\Delta$MAE \\
\hline
\endfirsthead
\hline
\multicolumn{7}{|c|}{\textit{Продолжение таблицы~\ref{tab:ablation-matrix}}} \\
\hline
ID & Что отключается/заменяется & Тестируемая гипотеза & Конкретное изменение конфигурации & Ожидаемый эффект & MAE & $\Delta$MAE \\
\hline
\endhead
\hline
\multicolumn{7}{|r|}{\textit{Продолжение на следующей странице}} \\
\hline
\endfoot
\hline
\endlastfoot
A1 & Cross-attention $\to$ конкатенация & Полезность попарного внимания & Удалить cross-attention во взаимодействии solute--solvent; заменить его на простую конкатенацию pair-эмбеддингов перед readout. & Умеренное ухудшение: теряется атомно-уровневое выравнивание межмолекулярных взаимодействий. & --- & --- \\
\hline
A2 & Bounded correction отключена & Полезность correction path & В correction-блоке зафиксировать gate как $w=1$ на всех примерах, то есть полностью отключить correction path. & Ухудшение: correction path компенсирует часть систематической ошибки SLE-пути. & --- & --- \\
\hline
A3 & Все auxiliary losses $=0$ на всех фазах & Полезность мультизадачного обучения & Обнулить веса всех auxiliary-losses во всех фазах; оставить только $\mathcal{L}_{\mathrm{sol}}$ и структурные регуляризаторы. & Сильное ухудшение: crystal heads и связанные auxiliary-каналы не получают обучающего сигнала. & --- & --- \\
\hline
A4 & Bridge-loss $=0$ & H4: bridge контрпродуктивен & Во всех фазах установить $\lambda_{\mathrm{bridge}}=0$ при прочих неизменных настройках. & Улучшение --- ожидаемый исход, если H4 верна. & --- & --- \\
\hline
A5 & $T$ подаётся в энкодер и auxiliary-heads & Необходимость temperature isolation & Подать температуру в backbone и в головы $T_{\mathrm{m}}$, $\Delta H_{\mathrm{fus}}$ до pair-readout, нарушив текущую temperature isolation. & Ухудшение: $T_{\mathrm{m}}$ и $\Delta H_{\mathrm{fus}}$ начинают зависеть от рабочей температуры. & --- & --- \\
\hline
A6 & Демпфирование SLE выключено & Необходимость демпфирования & В fixed-point решателе зафиксировать коэффициент демпфирования как $\lambda=1$ на всех итерациях. & Нейтрально при малой растворимости, но ухудшение при высоких $x_2$ и на плохо обусловленных системах. & --- & --- \\
\hline
A7 & Oracle teacher forcing выключен & Полезность oracle teacher forcing & Во всех фазах обучения установить вероятность teacher forcing как $p=0$, то есть никогда не подставлять истинные $T_{\mathrm{m}}$ и $\Delta H_{\mathrm{fus}}$. & Ухудшение, если crystal head --- bottleneck; нейтрально, если доминирует компенсаторная дегенерация. & --- & --- \\
\hline
A8 & NRTL $\to$ Wilson & H7: NRTL overparameterized & Заменить NRTL head и соответствующий слой активности на Wilson-параметризацию с двумя параметрами взаимодействия. & Возможное улучшение идентифицируемости и, как следствие, улучшение MAE. & --- & --- \\
\hline
A9 & NRTL $\to$ идеальный раствор & Полезность моделирования неидеальности & В SLE-решателе принудительно положить $\gamma_2=1$ и исключить head активности из вычисления прогноза. & Сильное ухудшение: модель лишается описания неидеальности раствора. & --- & --- \\
\hline
A10 & GC priors отключены & Полезность GC priors & Удалить GC prior для $T_{\mathrm{m}}$ и $\Delta H_{\mathrm{fus}}$; перейти к prediction from scratch без prior + bounded residual схемы. & Ухудшение: crystal head обучается медленнее и менее устойчиво. & --- & --- \\
\end{longtable}
\normalsize
\end{landscape}

Содержательно таблица~\ref{tab:ablation-matrix} разделяется на три класса тестов. Абляции A1, A5 и A10 проверяют ценность \textbf{репрезентационного дизайна} (pair-interaction, temperature isolation, priors для кристаллических голов). Абляции A2, A6 и A7 касаются \textbf{оптимизационной и вычислительной устойчивости} текущей системы. Наконец, A4, A8 и A9 тестируют саму \textbf{физическую постановку} задачи: полезен ли bridge-loss, не является ли NRTL-параметризация избыточной и действительно ли моделирование неидеальности необходимо для качества.

\subsection{Протокол выполнения и критерий значимости}

Все абляции должны выполняться при тех же разбиении, бюджете и базовых гиперпараметрах, что и канонический baseline TGNN-Solv из Эксперимента~1 настоящего раздела. Это означает: одно и то же scaffold-разбиение, та же трёхфазная схема 50/200/50, тот же оптимизатор и тот же набор регуляризаторов, за исключением единственного аблируемого компонента. Такой принцип критически важен: если вместе с абляцией изменяются learning rate, batch size или schedule, то причинная интерпретация результатов теряется.

Практически протокол должен иметь два уровня.
\begin{itemize}
  \item \textbf{Скрининг-уровень:} один seed на proxy-бюджете. Его задача --- быстро отсеять явно провальные или заведомо нейтральные варианты, не тратя полный вычислительный ресурс на все 10 конфигураций.
  \item \textbf{Финальный уровень:} три seed на полном бюджете для тех абляций, которые после proxy-скрининга остаются интерпретационно значимыми. Только этот уровень должен использоваться для окончательных выводов о пользе или вреде компонента.
\end{itemize}

Для каждой конфигурации, помимо итогового MAE на тесте, должны сохраняться и диагностические кривые, уже введённые в разделе о сходимости: $\mathcal{L}_{\mathrm{total}}^{\mathrm{train}}$, $\mathcal{L}_{\mathrm{sol}}^{\mathrm{train}}$, $\mathrm{MAE}_{\mathrm{val}}(\ln x_2)$, \texttt{sol\_fraction}, $\mathcal{L}_{\mathrm{tau\_reg}}$, $\mathcal{L}_{\mathrm{vH}}$ и средний $|\delta_{\mathrm{corr}}|$. Иначе невозможно отличить действительно полезную абляцию от конфигурации, которая лишь изменила траекторию обучения и случайно дала лучшее значение одного чекпойнта.

Изменение качества для абляции $A$ определяется как
\begin{equation*}
\Delta \mathrm{MAE}_A = \mathrm{MAE}_A - \mathrm{MAE}_{\mathrm{baseline}}.
\end{equation*}
Соответственно, отрицательное $\Delta \mathrm{MAE}_A$ означает улучшение относительно baseline, а положительное --- ухудшение. Практическим критерием значимости принимается следующее правило:
\begin{equation*}
|\Delta \mathrm{MAE}_A| > 0.05 \quad \text{и} \quad p < 0.05,
\end{equation*}
где $p$ --- значение парного $t$-теста по абсолютным ошибкам на одном и том же тестовом множестве. Только при выполнении обоих условий изменение трактуется как значимое. Если $|\Delta \mathrm{MAE}_A|\leq 0.05$ или статистический тест не достигает порога $p<0.05$, абляция должна интерпретироваться как нейтральная либо как требующая дополнительного бюджета для уверенного вывода.

\warnnote{На момент данного отчёта результаты абляций не получены; после завершения \hyperref[subsec:roadmap-phase2]{Фазы II дорожной карты} в таблице~\ref{tab:ablation-matrix} должны быть приведены: (i) MAE и $\Delta$MAE для всех конфигураций, (ii) оценка дисперсии по seed, (iii) $p$-значения парных тестов, (iv) краткая интерпретация того, какие компоненты оказались полезными, нейтральными или вредными. Результаты абляций заполняются после завершения \hyperref[subsec:roadmap-phase2]{Фазы II дорожной карты}.}

\section{Предлагаемые архитектурные модификации}

\subsection{Модификация A: априори групповых вкладов для блоков кристаллических свойств}
\label{subsec:mod-a}

\textbf{Суть:} вместо предсказания $T_{\mathrm{m}}$ и $\Delta H_{\mathrm{fus}}$ «с нуля» использовать GC-оценки как prior, а нейросетью предсказывать только bounded residual.

Полная формализация (Joback/Ducros, калибровка prior, формулы residual и численные диапазоны) уже приведена в §\ref{subsec:gc-methods} и используется без изменений. Здесь фиксируется только роль модификации в плане экспериментов.

\textbf{Преимущества:}
\begin{enumerate}
  \item Head обучается быстрее: вместо диапазона $[100, 700]$ К для $T_{\mathrm{m}}$ --- нужно выучить поправку $[-50, +50]$ К.
  \item Даже без обучения (GC prior) --- уже разумное значение.
  \item Пространство поиска для оптимизатора сужено $\to$ меньше компенсаторная дегенерация (H1, H2).
\end{enumerate}

\textbf{Совместимость:} полностью совместима с текущей архитектурой, не требует изменений в SLE-решателе.

\subsection{Модификация B: $\Delta C_p^{\mathrm{fus}}$ через групповые вклады}
\label{subsec:mod-b}

\textbf{Суть:} вместо $\Delta C_p = 0$, использовать GC-оценку (метод Ducros):

\begin{equation*}
\Delta C_p^{\mathrm{GC}} = \sum_k n_k \Delta C_{p,k}
\end{equation*}

\textbf{Не предсказывать нейросетью} --- это ещё один неидентифицируемый latent (H2). Использовать как фиксированный вход в SLE решатель.

\textbf{Ожидаемый эффект:} устранение систематической ошибки \textasciitilde0.5--1.5 в $\ln x_2$ для высокоплавких соединений (H6).

\textbf{Совместимость:} тривиальное изменение в SLE-решателе --- заменить \texttt{fixed\_dCp\_fus = 0} на вычисленное значение.

\subsection{Модификация C: оракульная подстановка при обучении}
\label{subsec:mod-c}

\textbf{Суть:} для тренировочных примеров с известными $T_{\mathrm{m}}$ и $\Delta H_{\mathrm{fus}}$ в SLE подставляются истинные значения вместо предсказанных.

Формальное определение oracle-маски, режимы использования и связь с teacher forcing уже даны в §\ref{sec:oracle-teacher-forcing}; этот подраздел не дублирует формулы и оставляет только экспериментальную мотивацию.

\textbf{Преимущества:}
\begin{enumerate}
  \item SLE-решатель получает точные $T_{\mathrm{m}}$, $\Delta H$ для \textasciitilde25\% данных $\to$ качественные градиенты для NRTL-ветки.
  \item Уменьшает компенсаторную дегенерацию (H1): NRTL head не может «прятать» ошибки crystal head в $\tau$.
  \item При inference на новых молекулах --- используется $T_{\mathrm{m}}^{\mathrm{pred}}$ (GNN head).
\end{enumerate}

\textbf{Не является «обманом»:} методологически это близко к teacher forcing / conditional training, но с сохранением специфики физического решателя.

\subsection{Модификация D: Отключение bridge-функции потерь}
\label{subsec:mod-d}

\textbf{Суть:} установить \texttt{bridge\_weight = 0}.

\textbf{Обоснование:} H4 (контрпродуктивность bridge-функции потерь). Bridge-функция потерь основана на Regular Solution Theory, которая систематически ошибочна для ассоциирующих систем (30--50\% данных).

\textbf{Альтернативные мягкие ограничения (вместо bridge):}

\begin{enumerate}
  \item \textbf{Проверка правила Уолдена:}
\end{enumerate}

\begin{equation*}
\mathcal{L}_{\mathrm{Walden}} = \max\!\left(0,\; \left|\frac{\Delta H_{\mathrm{fus}}}{T_{\mathrm{m}}} - 56.5\right| - 30\right)^2
\end{equation*}

Штрафует нефизичные комбинации $T_{\mathrm{m}}$ / $\Delta H_{\mathrm{fus}}$ (правило Уолдена: $\Delta S_{\mathrm{fus}} \approx 56.5 \pm 30$ Дж/(моль$\cdot$К)).

\begin{enumerate}
\setcounter{enumi}{1}
  \item \textbf{$\gamma_2^\infty$ sign consistency:}
\end{enumerate}

\begin{equation*}
\mathcal{L}_{\mathrm{sign}} = \max\!\left(0,\; -\ln\gamma_2^{\infty,\mathrm{pred}}\right)^2
\end{equation*}

Мягко штрафует $\gamma_2^\infty < 1$ (нетипично, хотя физически допустимо).

\subsection{Модификация E: UNIQUAC вместо NRTL}
\label{subsec:mod-e}

\textbf{Суть:} заменить модель $G^E$ с NRTL на UNIQUAC~\cite{abrams1975,prausnitz1998}.

\textbf{Структурное преимущество:}

\begin{equation*}
\ln\gamma_2 = \underbrace{\ln\gamma_2^{\mathrm{comb}}(r_1, r_2, q_1, q_2, x_2)}_{\text{аналитически заданный вклад --- 0 learnable params}} + \underbrace{\ln\gamma_2^{\mathrm{res}}(\Theta_1, \Theta_2, \tau_{12}, \tau_{21})}_{\text{2 predicted params}}
\end{equation*}

Комбинаторный вклад $\ln\gamma_2^{\mathrm{comb}}$ учитывает размер и форму молекул через структурные параметры $r_i$, $q_i$, вычисляемые из SMARTS:

\begin{equation*}
r_i = \sum_k \nu_k^{(i)} R_k, \quad q_i = \sum_k \nu_k^{(i)} Q_k
\end{equation*}

где $R_k$, $Q_k$ --- табличные параметры Bondi.

\textbf{Сравнение с NRTL:}

\begin{center}
\begin{tabular}{|c|c|c|}
\hline
Аспект & NRTL & UNIQUAC \\
\hline
Обучаемые параметры & 3 ($\Delta g_{12}$, $\Delta g_{21}$, $\alpha$) & 2 ($\Delta u_{12}$, $\Delta u_{21}$) \\
\hline
Hardcoded physics & Только $G^E$ формула & $G^E$ + комбинаторный вклад ($r$, $q$) \\
\hline
Описание LLE & Да & Да \\
\hline
Учёт размера/формы & Нет (неявно через $\alpha$) & Да (явно через $r$, $q$) \\
\hline
\end{tabular}
\end{center}

\textbf{Ожидание:} лучшая идентифицируемость (2 параметра вместо 3) + бесплатный размерный эффект через $r$, $q$.

\subsection{Модификация F: Двухстадийная оптимизация (альтернатива трёхфазному учебному режиму)}
\label{subsec:mod-f}

\textbf{Суть:} разделить обучение на две стадии с разной целью.

\textbf{Стадия A: Convergent предварительное обучение (\textasciitilde100 эпох)}

Цель: обучить crystal heads до приемлемой точности.

\begin{itemize}
  \item Обучаются: GNN Encoder, FusionHead, HansenHead.
  \item SLE решатель \textbf{выключен}.
  \item NRTLHead \textbf{заморожена}.
  \item Loss: только $\mathcal{L}_{T_{\mathrm{m}}} + \mathcal{L}_{\Delta H} + \mathcal{L}_{\mathrm{Hansen}}$; канал $\mathcal{L}_{\gamma^\infty}$ на этой стадии остаётся выключенным, потому что NRTL-head ещё заморожена, хотя во всей IDAC-augmented постановке внешний IDAC-корпус уже доступен.
\end{itemize}

Критерий перехода к Стадии B: $\mathrm{MAE}(T_{\mathrm{m}}) < 40$ К и $\mathrm{MAE}(\Delta H_{\mathrm{fus}}) < 3000$ Дж/моль на validation set.

\textbf{Стадия B: SLE-guided training (\textasciitilde200 эпох)}

Цель: обучить NRTL head при \textbf{зафиксированных или медленно обучающихся} crystal heads.

\begin{itemize}
  \item Crystal heads: заморожены первые 100 эпох, затем unfreeze с lr $\times 0.01$.
  \item NRTLHead: обучается с полным lr.
  \item Оракульная подстановка (\hyperref[subsec:mod-c]{Модификация C}) при наличии $T_{\mathrm{m}}^{\mathrm{true}}$.
  \item SLE решатель \textbf{включён}.
  \item Loss: $\mathcal{L}_{\mathrm{sol}} + $ мягкие регуляризаторы.
\end{itemize}

\textbf{Обоснование:} Стадия A гарантирует, что SLE решатель получает разумные $T_{\mathrm{m}}$ и $\Delta H_{\mathrm{fus}}$ на входе Стадии B. Это разрывает деструктивный цикл обратной связи H3.

\section{Предлагаемая архитектура v3 (если текущая не спасаема)}

Если Эксперименты 1--3 покажут, что physics bottleneck \textbf{фундаментально вреден} (то есть наблюдается устойчивое отставание от DirectGNN и/или выраженная компенсаторная дегенерация), предлагается альтернативная архитектура, сохраняющая интерпретируемость.

\subsection{Принцип: «Физика как скелет, не как бутылочное горлышко»}

\begin{equation*}
\ln x_2 = \mathrm{MLP}_{\mathrm{direct}}(\mathbf{g}_{\mathrm{pair}},\; \mathbf{d}_{\mathrm{RDKit}},\; T)
\end{equation*}

где $\mathbf{d}_{\mathrm{RDKit}}$ вычисляются библиотекой RDKit~\cite{landrum2024}, с мягкими физическими ограничениями.

\subsection{Архитектурная схема}
\begin{enumerate}
  \item \textbf{Входы:} \texttt{SMILES\_sol}, \texttt{SMILES\_slv}, температура $T$.
  \item \textbf{Энкодинг:} общий GNN-backbone формирует латентные представления solute/solvent; параллельно вычисляются фиксированные RDKit-descriptors.
  \item \textbf{Auxiliary heads:} предсказываются $T_{\mathrm{m}}$ и $\Delta H_{\mathrm{fus}}$ (GC prior + residual) как вспомогательные таргеты.
  \item \textbf{Взаимодействие:} Cross-Attention $\Rightarrow$ Pair Representation с конкатенацией признаков.
  \item \textbf{Основной выход:} $\mathrm{MLP}_{\mathrm{direct}}(\mathbf{g}_{\mathrm{pair}}, \mathbf{d}_{\mathrm{RDKit}}, T) \Rightarrow \ln x_2$.
  \item \textbf{Физические ограничения:} проверка правила Уолдена, монотонность по температуре, Van't Hoff consistency, soft-anchor к идеальной растворимости.
\end{enumerate}

\subsection{Функция потерь}

\begin{equation*}
\mathcal{L} = \underbrace{\mathcal{L}_{\mathrm{Huber}}(\ln x_2^{\mathrm{pred}}, \ln x_2^{\mathrm{true}})}_{\text{основная}} + \lambda_{T_{\mathrm{m}}}\mathcal{L}_{T_{\mathrm{m}}} + \lambda_{\Delta H}\mathcal{L}_{\Delta H} + \lambda_{\mathrm{mono}}\mathcal{L}_{\mathrm{mono}} + \lambda_{\mathrm{vH}}\mathcal{L}_{\mathrm{vH}} + \lambda_{\mathrm{anchor}}\mathcal{L}_{\mathrm{anchor}}
\end{equation*}

где anchor функции потерь:

\begin{equation*}
\mathcal{L}_{\mathrm{anchor}} = \max\!\left(0,\; \left|\ln x_2^{\mathrm{pred}} - (-\Phi^{\mathrm{pred}})\right| - M\right)^2, \quad M \sim 10
\end{equation*}

Физический смысл: предсказанная растворимость не должна отличаться от идеальной более чем на $e^{M}$ раз (т.е. $\gamma_2 \in [e^{-M}, e^{M}]$). При $M = 10$: $\gamma_2 \in [4.5 \times 10^{-5},\; 22000]$ --- очень мягкое ограничение.

\subsection{Сохранение интерпретируемости}

Хотя основной путь --- direct prediction, все промежуточные величины доступны для post-hoc анализа:

\begin{itemize}
  \item $T_{\mathrm{m}}^{\mathrm{pred}}$ и $\Delta H_{\mathrm{fus}}^{\mathrm{pred}}$ --- из auxiliary heads.
  \item $\gamma_2^{\mathrm{implied}} = \exp\!\left(-\Phi^{\mathrm{pred}} - \ln x_2^{\mathrm{pred}}\right)$ --- «implied» коэффициент активности.
  \item Температурная зависимость $\ln x_2(T)$ --- через дифференцирование по $T$.
\end{itemize}

Это \textbf{не полная} интерпретируемость (нет явных NRTL параметров), но достаточная для химического анализа.

\subsection{Сравнение v2 (текущая) по сравнению с v3 (предлагаемая)}

\begin{center}
\scriptsize
\setlength{\tabcolsep}{1.5pt}
\begin{tabularx}{0.97\textwidth}{|P{3.0cm}|Y|Y|}
\hline
Аспект & v2 (hard physics) & v3 (soft physics) \\
\hline
Основной путь & GNN $\to$ physics params $\to$ SLE решатель & GNN $\to$ MLP $\to$ $\ln x_2$ \\
\hline
Число предсказываемых параметров & 7+ & 1 ($\ln x_2$ напрямую) \\
\hline
Identifiability & Проблематична (H2) & Не применима \\
\hline
Поток градиента & Через итеративный решатель & Прямой (MLP) \\
\hline
\shortstack[l]{Термодинамическая\\консистентность} & По построению & Через soft constraints \\
\hline
\shortstack[l]{Интерпретируе-\\мость} & Полная (все промежуточные) & Частичная (post-hoc) \\
\hline
дескрипторы RDKit & Не используются & Конкатенация \\
\hline
Ожидаемая MAE & ? (не протестирована при полном бюджете) & \textasciitilde1.2--1.5 (проектная оценка без экспериментального подтверждения) \\
\hline
\end{tabularx}
\end{center}

\section{Дорожная карта}

\subsection{Фаза I: Диагностика}
\label{subsec:roadmap-phase1}

\begin{enumerate}
  \item \hyperref[subsec:exp-full-budget]{Эксперимент 1}: полный бюджет 50/200/50.
  \item \hyperref[subsec:exp-oracle-crystal]{Эксперимент 2}: оракульные кристаллические параметры (сразу после \hyperref[subsec:exp-full-budget]{Эксп.~1}).
  \item \hyperref[subsec:exp-bridge-ablation]{Эксперимент 3}: абляция связующей функции потерь (параллельно с \hyperref[subsec:exp-full-budget]{Эксп.~1}).
\end{enumerate}

\textbf{Результат \hyperref[subsec:roadmap-phase1]{Фазы I}:} чёткий ответ на вопрос «даёт ли физическое узкое место выигрыш, нейтрально оно или вредно?»

\subsection{Фаза II: Архитектурные модификации}
\label{subsec:roadmap-phase2}

В зависимости от результатов \hyperref[subsec:roadmap-phase1]{Фазы I}:

\textbf{Если физическое узкое место нейтрально или слабо полезно (исходы A/B):}
\begin{itemize}
  \item Имплементировать \hyperref[subsec:mod-a]{Модификацию A} + \hyperref[subsec:mod-b]{B} + \hyperref[subsec:mod-c]{C} (априори групповых вкладов + $\Delta C_p$ + оракульная подстановка).
  \item \hyperref[subsec:exp-gc-priors]{Эксперимент 5}: оценить эффект.
  \item Имплементировать \hyperref[subsec:mod-e]{Модификацию E} (UNIQUAC) --- \hyperref[subsec:exp-wilson]{Эксперимент 6}.
\end{itemize}

\textbf{Если физическое узкое место вредно (исходы C/D):}
\begin{itemize}
  \item Имплементировать архитектуру v3 (вариант «physics as scaffold», с прямым предсказанием $\ln x_2$ и мягкими физическими ограничениями).
  \item \hyperref[subsec:exp-directgnn-rdkit]{Эксперимент 4}: DirectGNN + дескрипторы RDKit (верхняя граница).
  \item Сравнить v3 с DirectGNN + descriptors и RF.
\end{itemize}

\subsection{Фаза III: Финализация}
\label{subsec:roadmap-phase3}

\begin{itemize}
  \item Multi-seed evaluation на всех splits (scaffold, solute, solvent).
  \item Исследование абляций (текущая схема абляций).
  \item Temperature extrapolation analysis.
  \item Подготовка результатов для публикации.
\end{itemize}

\section{Ограничения}

Перед итоговым выводом фиксируем ключевые ограничения текущей версии:
\begin{enumerate}
  \item \textbf{Данные:} корпус в среднем мультитемпературный (медиана 9 точек на пару), но температуры сконцентрированы около 303.15 K, а auxiliary-метки вне $T_{\mathrm{m}}$ крайне разрежены; это всё ещё ослабляет идентифицируемость температурных параметров.
  \item \textbf{Модель:} NRTL-параметризация допускает компенсаторные решения (many-to-one) и требует жёсткой диагностики физической правдоподобности.
  \item \textbf{Бюджет:} часть выводов основана на proxy/medium режимах; финальный вердикт требует полного цикла 50/200/50 и многозапусковый статистики.
\end{enumerate}

Детализированная версия ограничений приведена в следующем разделе.

\section{Заключение}

\subsection{Резюме текущего состояния}

Архитектура TGNN-Solv теоретически обоснована: мастер-уравнение SLE точное, NRTL термодинамически консистентна по построению, имплицитное дифференцирование математически корректно. Однако эмпирические результаты (даже с поправкой на недостаточный бюджет обучения) указывают на фундаментальные проблемы:

\begin{enumerate}
  \item \textbf{Инвертированная иерархия чувствительности и супервизии} (H1): наиболее чувствительные параметры наименее супервизированы.
  \item \textbf{Слабая идентифицируемость} (H2): даже при медианных 9 наблюдениях на пару задача остаётся плохо обусловленной из-за коррелированных температурных вкладов, узкого рабочего диапазона $T$ и отсутствия прямой супервизии на NRTL-параметры.
  \item \textbf{Контрпродуктивная bridge-функция потерь} (H4): Regular Solution Theory систематически ошибочна для ассоциирующих систем.
\end{enumerate}

\subsection{Ключевой эксперимент}

\hyperref[subsec:exp-full-budget]{Эксперимент 1} (полный бюджет + диагностика промежуточных параметров) является решающим. Без него невозможно отделить операционные проблемы (недоучённость) от архитектурных (неидентифицируемость). Все остальные решения зависят от его результата.

\subsection{Наиболее вероятный путь вперёд}

На основании теоретического анализа, наиболее перспективным представляется комбинация модификаций:

\begin{enumerate}
  \item \textbf{GC priors} для $T_{\mathrm{m}}$, $\Delta H_{\mathrm{fus}}$, $\Delta C_p$ (\hyperref[subsec:mod-a]{Модификации A}, \hyperref[subsec:mod-b]{B}) --- сужение пространства поиска.
  \item \textbf{Оракульная подстановка} при обучении (\hyperref[subsec:mod-c]{Модификация C}) --- качественные градиенты для NRTL голова предсказания.
  \item \textbf{Отключение bridge-функции потерь} (\hyperref[subsec:mod-d]{Модификация D}) --- устранение контрпродуктивного регуляризатора.
  \item \textbf{UNIQUAC} вместо NRTL (\hyperref[subsec:mod-e]{Модификация E}) --- меньше параметров, больше аналитически заданной физики.
\end{enumerate}

Эта комбинация сохраняет \textbf{полную физическую интерпретируемость} (все промежуточные величины имеют физический смысл), устраняет основные идентифицированные проблемы, и совместима с текущей кодовой базой.



\section{Ограничения текущего анализа (развёрнутая версия)}

Несмотря на детальность математической части, текущий отчёт имеет ряд ограничений, которые важно явно зафиксировать для корректной интерпретации результатов.

\subsection{Ограничения данных}

\begin{enumerate}
  \item \textbf{Температурная плотность неравномерна}: большинство пар имеют 6--10 точек, но температурная сетка сосредоточена вблизи комнатных условий, а лишь 0.52\% пар имеют 21+ точку; поэтому идентифицируемость температурных коэффициентов NRTL остаётся хрупкой, особенно на хвостах диапазона $T$.
  \item \textbf{Смещение по классам систем}: полярные и водородно-связывающие смеси непропорционально сложнее, но представлены неравномерно.
  \item \textbf{Гетерогенность источников}: данные собраны из разных экспериментальных протоколов, что повышает шум и риск скрытых систематических сдвигов.
\end{enumerate}

\subsection{Ограничения модели}

\begin{enumerate}
  \item \textbf{Аппроксимация $\Delta C_p^{\mathrm{fus}}$}: даже при аккуратной постановке использование упрощающих гипотез вносит систематическую ошибку в $x_2^{\mathrm{id}}$.
  \item \textbf{Параметрическая неоднозначность NRTL}: различные комбинации параметров могут давать близкие значения растворимости в узком температурном диапазоне.
  \item \textbf{Зависимость от динамики решателя}: при плохой кондиционированности фиксированной точки возможны нестабильные градиенты и локальные режимы сходимости.
\end{enumerate}

\subsection{Ограничения экспериментального бюджета}

Ключевые фальсификационные эксперименты сформулированы, но пока не выполнены в полном объёме. Это означает, что часть выводов (особенно про относительный вклад H2/H3/H4) носит статус \textit{наиболее правдоподобной рабочей гипотезы}, а не окончательно установленного факта.

\section{Воспроизводимость и протокол отчётности}

Для инженерной и научной надёжности предлагается стандартизованный протокол воспроизводимости.

\subsection{Заявление о воспроизводимости (краткая версия)}

\begin{itemize}
  \item \textbf{Код модели:} публичный релизный репозиторий планируется в GitHub (снимок для статьи фиксируется отдельным тегом или релизом).
  \item \textbf{Данные:} основной корпус --- BigSolDB v2.1; в текущей конфигурации используются дополнительные метки $T_{\mathrm{m}}$, $\Delta H_{\mathrm{fus}}$ и Hansen, тогда как канал $\gamma^\infty$ зарезервирован архитектурно, но фактически пуст.
  \item \textbf{Вычислительные ресурсы:} базовый полный цикл 50/200/50 оценивается в \textasciitilde24--48 GPU-часов на одно начальное зерно (класс A100).
  \item \textbf{Начальные зёрна:} 42, 123, 456 для всех ключевых сравнений.
\end{itemize}

\subsection{Минимальный артефакт эксперимента}

Каждый запуск должен сохранять:
\begin{enumerate}
  \item commit hash кода и конфигурацию модели;
  \item фиксированные seed для всех источников случайности;
  \item точные идентификаторы разбиения и список пар в train/val/test;
  \item кривые обучения по всем компонентам функции потерь;
  \item промежуточные физические предсказания ($T_{\mathrm{m}}$, $\Delta H_{\mathrm{fus}}$, параметры NRTL, число итераций решателя).
\end{enumerate}

Дополнительно в поддерживаемой инфраструктуре бенчмарков фиксируются:
\begin{itemize}
  \item canonical bundles (канонические наборы конфигураций и данных);
  \item run manifests (машиночитаемые манифесты запуска);
  \item benchmark cards (карточки результата для сравнения версий);
  \item custom-model adapter API (встраивание внешней модели в единый контур оценки);
  \item stress suite (стресс-наборы для проверки устойчивости);
  \item benchmark release freezer (заморозка релизного набора для воспроизводимой публикации).
\end{itemize}

Режим воспроизведения статьи структурирован в три профиля:
\texttt{core}, \texttt{article}, \texttt{full}.

\subsection{Рекомендуемый набор метрик}

Помимо MAE/RMSE по $\ln x_2$, необходимо включать в отчёт:
\begin{enumerate}
  \item качество по подгруппам (неполярные, полярные, ассоциирующие системы);
  \item зависимость ошибки от температуры и диапазона концентраций;
  \item статистику устойчивости решателя (доля сходимости, среднее число итераций, распределение $|g'(x^*)|$);
  \item физические sanity checks (монотонность по температуре, отсутствие невалидных активностей).
\end{enumerate}

\subsection{Честное сравнение с базовыми ориентирами}

После перевода дескрипторной аугментации, GPS-энкодера и Stage 0 в поддерживаемый контур для обеих архитектур, fairness-протокол ужесточается. Сравнение TGNN, DirectGNN и RF (descriptors) должно проводиться при совпадении:
\begin{enumerate}
  \item вычислительного бюджета (эпохи, ранняя остановка, число запусков);
  \item класса признаков (с дескрипторами по сравнению с без дескрипторов);
  \item типа энкодера (MPNN по сравнению с MPNN, GPS по сравнению с GPS);
  \item режима инициализации (Stage 0 тёплый старт по сравнению с Stage 0 тёплым стартом или cold-start по сравнению с cold-start);
  \item процедуры подбора гиперпараметров (включая пространство Optuna и правило выбора лучшего trial).
\end{enumerate}

Иначе наблюдаемая разница смешивает эффект архитектуры с эффектом входной информации, инициализации и настройки.

\section{Практические рекомендации для следующей итерации}

\subsection{Приоритеты на ближайший цикл}

\begin{enumerate}
  \item \textbf{Сначала диагностика, потом усложнение}: завершить \hyperref[subsec:exp-full-budget]{Эксперимент 1} до внедрения новых крупных модулей.
  \item \textbf{Стабилизировать физический слой}: включить демпфирование решателя по умолчанию и логировать его динамику на каждом батче.
  \item \textbf{Снизить размерность параметризации}: протестировать Wilson/UNIQUAC как контроль на предмет идентифицируемости~\cite{wilson1964,abrams1975}.
  \item \textbf{Усилить supervision чувствительных величин}: добавить внешние priors и/или weak labels для crystal properties.
\end{enumerate}

\subsection{Критерии продолжения или остановки для текущей архитектуры}

После завершения ключевых экспериментов архитектура TGNN-Solv v2 сохраняется только если одновременно выполняются условия:
\begin{enumerate}
  \item стабильная сходимость решатель на подавляющем большинстве батчей;
  \item воспроизводимое превосходство над DirectGNN при сопоставимом бюджете;
  \item интерпретируемые промежуточные параметры без явной компенсаторной дегенерации.
\end{enumerate}

Если хотя бы два критерия систематически нарушаются, рационально переходить к архитектуре v3 как к более устойчивому компромиссу между физикой и статистической эффективностью.

\section*{Часть VII. Внешние модели, Experiment Lab, Planner и прикладной контур}
\addcontentsline{toc}{section}{Часть VII. Внешние модели, Experiment Lab, Planner и прикладной контур}

Эта часть фиксирует внешний контур сравнительного тестирования, поддерживаемую интерактивную поверхность проекта, правила воспроизводимости и прикладные сценарии использования TGNN-Solv.

Цели части VII:
\begin{enumerate}
  \item определить внешний набор моделей сравнения (SolProp и FastSolv);
  \item формализовать термины и артефакты воспроизводимого бенчмарка;
  \item описать \texttt{Experiment Lab} как поддерживаемую управляющую поверхность над основным кодом проекта;
  \item детально разобрать модули прикладного слоя и границы их применимости;
  \item корректно развести \texttt{Planner} и маршрутно-ориентированный отбор растворителей, чтобы избежать завышенных утверждений;
  \item зафиксировать риски и минимальный релизный набор доказательных материалов.
\end{enumerate}

\section{Внешние модели и бенчмарк-контур}

\subsection{Ключевые термины и артефакты бенчмарка}

Ниже приведены термины, которые используются в этой части как операционные.

\begin{center}
\small
\begin{longtable}{|p{0.30\textwidth}|p{0.62\textwidth}|}
\hline
\textbf{Термин} & \textbf{Определение в контексте отчёта} \\
\hline
\endfirsthead
\hline
\textbf{Термин} & \textbf{Определение в контексте отчёта} \\
\hline
\endhead
\hline
\multicolumn{2}{r}{\textit{Продолжение на следующей странице}} \\
\hline
\endfoot
\hline
\endlastfoot
Базовая внешняя модель & Модель, не входящая в семейство TGNN/DirectGNN, используемая для внешнего сравнения качества и устойчивости. \\
\hline
Режим без дообучения (zero-shot) & Запуск внешней модели без адаптации на целевом датасете отчёта. \\
\hline
Калиброванный режим & Посткалибровка предсказаний или интервалов неопределённости на валидационных данных без полного переобучения весов. \\
\hline
Полное переобучение (native retraining) & Полное переобучение внешней модели на целевых данных и split-протоколе отчёта. \\
\hline
Адаптационная обёртка & Слой, который унифицирует входы и выходы внешней модели с единым API бенчмарка. \\
\hline
Канонический пакет артефактов & Стандартизованный набор результатов запуска: конфиг, метрики, предсказания, чекпойнт и служебные метаданные. \\
\hline
Спецификация запуска (run manifest) & Машиночитаемое описание конкретного прогона: начальное зерно, разбиение, commit, конфигурация и аппаратное окружение. \\
\hline
Карточка бенчмарка (benchmark card) & Краткая карточка результата: постановка, протокол, агрегированные метрики, ограничения и ссылка на артефакты. \\
\hline
Интерфейс адаптера (Adapter API) & Интерфейс подключения пользовательских и внешних моделей к единому контуру сравнительного тестирования. \\
\hline
Набор стресс-тестов & Набор сценариев для проверки устойчивости: экстремальные температуры, редкие классы молекул, примеры вне области применимости и численные сбои решателя. \\
\hline
Релизная заморозка & Процедура фиксации релизного набора данных, разбиений, конфигов и метрик, чтобы исключить дрейф условий сравнения. \\
\hline
Профили \texttt{core/article/full} & Три уровня воспроизведения: быстрый инженерный минимум, публикационный профиль и полный исследовательский профиль. \\
\hline
\end{longtable}
\normalsize
\end{center}

\subsection{Архитектурные режимы внутреннего сравнения}

Сравнение TGNN и DirectGNN в этой части рассматривается как матрица сопоставимых режимов:
\begin{itemize}
  \item MPNN без дескрипторов;
  \item MPNN с дескрипторами;
  \item GPS без дескрипторов;
  \item GPS с дескрипторами;
  \item каждый из режимов --- в вариантах холодного старта и Stage~0 тёплого старта.
\end{itemize}

Для корректного сопоставления фиксируются одинаковые вычислительный бюджет, split-протокол, пространство гиперпараметров, число seed и правила выбора лучшего чекпойнта.

\subsection{Внешний бенчмарк-блок: SolProp и FastSolv}

\textbf{Почему выбраны именно эти модели.} Внешний бенчмарк построен на SolProp и FastSolv, поскольку они задают сильные, но методологически различающиеся ориентиры относительно TGNN-Solv.

\textbf{SolProp (три режима запуска, см.~\cite{solprop2024}).}
\begin{enumerate}
  \item \textbf{Режим без дообучения:} оценивает переносимость внешней модели без адаптации.
  \item \textbf{Калиброванный режим:} показывает, насколько сильно улучшается практическая полезность после калибровки без полного переобучения.
  \item \textbf{Полное переобучение:} определяет верхнюю границу качества SolProp на целевых данных текущего отчёта.
\end{enumerate}

\textbf{FastSolv.} FastSolv относится к сильным предикторам растворимости, опирающимся на предобучение и ансамблирование~\cite{osanloo2024}. В контуре бенчмарка FastSolv используется как внешний ориентир высокой точности; при этом отдельно контролируются чувствительность адаптационной обёртки, согласование формата признаков и устойчивость артефактов.

\textbf{Явная фиксация для этого отчёта:}
\begin{equation*}
\text{External benchmark set} = \{\text{SolProp},\;\text{FastSolv}\}.
\end{equation*}

Именно эти две модели считаются внешними базовыми моделями в итоговых таблицах сравнительного тестирования.

\subsection{Протокол бенчмарка и отчётности}

Минимальный протокол внешнего и внутреннего сравнения включает:
\begin{enumerate}
  \item единое разбиение по каркасам и фиксированные начальные зёрна;
  \item одинаковый набор основных метрик (MAE, $R^2$) и интервальных/калибровочных метрик (PICP, MPIW);
  \item запись построчных предсказаний на тесте и диагностик стабильности решателя;
  \item публикацию \texttt{run\_manifest}, \texttt{resolved\_config}, \texttt{benchmark\_card} и ссылок на чекпойнты;
  \item обязательное указание режима внешней модели: без дообучения, калиброванный режим или полное переобучение.
\end{enumerate}

Без этих пунктов сравнение трактуется как предварительное, а не релизное.

\section{Experiment Lab}

В текущем состоянии TGNN-Solv корректно описывать не только как физически-информированную модель растворимости, но и как поддерживаемую исследовательскую платформу с интегрированной Streamlit-поверхностью \texttt{Experiment Lab}. При этом ядро проекта не изменило своей природы: это по-прежнему система, ориентированная прежде всего на предсказание растворимости, а графический интерфейс и прикладной слой остаются надстройками над каноническим вычислительным контуром.

\validated{\texttt{Experiment Lab} следует трактовать не как демонстрационный интерфейс и не как отдельный продукт, а как поддерживаемый слой оркестрации над теми же CLI-, API- и потоками артефактов, которые используются в основном репозитории.}

\subsection{Общая роль и модель исполнения}

\texttt{Experiment Lab} использует те же чекпойнты, подготовленные CSV, CLI-команды обучения и оценки, API инференса, процедуры бенчмарка и артефакты в каталогах \texttt{results/}, \texttt{checkpoints/}, \texttt{figures/} и \texttt{tables/}. Поэтому его корректно рассматривать как интерактивную управляющую поверхность над основной кодовой базой, а не как параллельную реализацию модели.

Базовый запуск в рабочем репозитории осуществляется через \texttt{python scripts/launch\_lab.py}. Существенная техническая особенность состоит в том, что Streamlit-интерфейс и модельная среда исполнения не обязаны жить в одном и том же Python-окружении: в боковой панели задаётся \texttt{Python command}, и именно эта команда затем используется для подзадач обучения, оценки, инференса, инспекции чекпойнтов и диагностики окружения. Такая архитектура поддерживает режим «тонкий интерфейс + тяжёлое окружение обучения и инференса» и потому является признаком зрелой оркестрационной поверхности, а не просто визуальной оболочки.

\subsection{Состояние, хранимое в репозитории и рабочие пространства}

Важная практическая черта \texttt{Experiment Lab} состоит в том, что он сохраняет своё состояние прямо в репозитории: preset-ы DAG-пайплайнов, состояние доски \texttt{Planner} и истории интерактивных запусков фиксируются как обычные проектные артефакты, а не как эфемерное состояние GUI. Это принципиально важно для воспроизводимости, поскольку результаты инференса, неопределённости и калибровки можно повторно импортировать в последующие пайплайны и связывать с чекпойнтами, таблицами и задачами планирования. Поддерживаемые рабочие зоны уже покрывают данные, обучение, бенчмарк, визуальное редактирование пайплайнов, прикладные сценарии, документацию и диагностику окружения; подробный перечень интерфейсных модулей уместнее держать в \texttt{README} репозитория или в приложении, а не в основной научной части отчёта.

Следовательно, \texttt{Experiment Lab} корректно описывать как поддерживаемый интерактивный операционный слой вокруг ML-проекта, а не как вспомогательную витрину поверх одной модели.

\section{Прикладной слой}

\subsection{Applications: прикладной слой поверх прогноза растворимости}

Прикладной пакет \texttt{tgnn\_solv.applications} в поддерживаемом состоянии состоит из четырёх основных модулей --- \texttt{solvent\_screening.py}, \texttt{process\_optimization.py}, \texttt{drug\_properties.py} и \texttt{pk\_profiling.py} --- с общим вспомогательным слоем в \texttt{core.py}. Их объединяет одна и та же вычислительная логика: сначала модель предсказывает равновесную растворимость для явной пары «растворяемое вещество--растворитель» при явной температуре, после чего поверх этого прогноза строятся технологические и формуляционные эвристики. Эти эвристики должны явно трактоваться как приближения, а не как полноценные механистические симуляторы.

\keyidea{Расширение проекта реально, но не меняет его физическую сущность: прикладной слой --- это надстройка над предсказанием равновесной растворимости, а не самостоятельный класс моделей ретросинтеза, PBPK или PK/PD.}

Один и тот же прикладной стек поддержан для обоих семейств моделей --- \texttt{TGNN-Solv} и \texttt{DirectGNN}. Это достигается через единый слой маршрутизации вызовов, в котором \texttt{SolventScreener} вызывает либо \texttt{predict\_solubility}, либо \texttt{predict\_direct\_solubility}. Тем самым прикладное сравнение проводится на общей поверхности, а не через две разные «ручные» логики.

\subsubsection{solvent\_screening.py: отбор растворителей, окна кристаллизации и «зелёная» замена}

Модуль \texttt{solvent\_screening.py} опирается на встроенную библиотеку из 65 растворителей и сред. Для каждого кандидата хранятся не только \texttt{name} и \texttt{smiles}, но и технологически значимые метаданные: класс растворителя, температура кипения, плотность, относительная стоимость, показатель экологичности, флаги протичности и способности к образованию водородных связей, смешиваемость с водой и параметры Хансена. Именно это делает отбор не просто ранжированием по $x_2$, а ранжированием по $x_2$ в связке с технологичностью, безопасностью и экологичностью.

Основной сценарий \texttt{SolventScreener.screen(...)} умеет:
\begin{enumerate}
  \item ранжировать библиотеку по $\ln x_2$, $x_2$ и приближённой метрике \texttt{mg/mL};
  \item фильтровать кандидатов по показателю экологичности, токсикологическим ограничениям, температуре кипения и смешиваемости;
  \item возвращать сигналы уверенности и принадлежности к области применимости, если подключена соответствующая AD-модель;
  \item добавлять диагностические признаки, связанные с параметрами Хансена, когда доступны промежуточные выходы физического блока.
\end{enumerate}

\warnnote{Перевод из $x_2$ в \texttt{mg/mL} реализуется как приближённый пересчёт, основанный на предположении о доминирующем объёме растворителя. Это полезная метрика для первичного ранжирования и быстрой оценки загрузки, но она не заменяет точный массовый баланс рецептуры, измеренную плотность раствора или полную материальную спецификацию процесса.}

Поверх базового отбора поддерживаются несколько частных сценариев:
\begin{itemize}
  \item \textbf{\texttt{crystallization\_window(...)}:} сканирование по температуре с расчётом $x_{2,\mathrm{hot}}$, $x_{2,\mathrm{cold}}$, теоретического выхода, степени пересыщения, среднего наклона $d\ln x_2/dT$, грубой оценки ширины метастабильной зоны и рекомендаций по охлаждению. Это инструмент предварительной технологической оценки, но не кинетическая модель кристаллизации.
  \item \textbf{\texttt{antisolvent\_screening(...)}:} сравнение растворимости в хорошем растворителе и потенциальном антирастворителе с учётом смешиваемости и технологических эвристик. Это полезно для drowning-out сценариев, но не моделирует динамику смешения, кинетику зародышеобразования и переходные профили пересыщения.
  \item \textbf{\texttt{green\_solvent\_replacement(...)}:} поиск более «зелёной» замены текущего растворителя с контролем сохранения растворимости, улучшения экологичности и снижения токсикологической нагрузки. Именно этот сценарий особенно хорошо иллюстрирует режим поддержки практических решений.
  \item \textbf{\texttt{solvent\_swap\_path(...)}:} грубое построение пути между растворителями на графе смешиваемости с приближённой оценкой объёмных отношений. Это логика этапов выделения и замены растворителя, а не планирование синтетического маршрута.
\end{itemize}

\subsubsection{process\_optimization.py: технологические эвристики поверх прогноза растворимости}

Класс \texttt{ProcessOptimizer} использует \texttt{SolventScreener} как базовый слой и поверх него строит несколько технологически ориентированных сценариев.

\begin{itemize}
  \item \textbf{\texttt{optimize\_crystallization(...)}:} перебирает растворители и горячие/холодные температурные точки, фильтрует их по экологическим, токсикологическим и температурным ограничениям, требует минимальную растворимость на горячем конце, оценивает выход, supersaturation ratio и итоговый \texttt{objective\_score}. Корректное описание --- поиск рабочих окон \((\text{solvent}, T_{\mathrm{hot}}, T_{\mathrm{cold}})\), ориентированный на равновесную растворимость, а не симулятор кристаллизационного процесса.
  \item \textbf{\texttt{optimize\_extraction(...)}:} сравнивает кандидатов относительно исходного растворителя, оценивает коэффициент распределения, несмешиваемость и выигрыш экстракции, но не моделирует эффективность стадий, массоперенос, образование эмульсий и реальную гидродинамику процесса.
  \item \textbf{\texttt{optimize\_reaction\_medium(...)}:} ищет среды, в которых реагенты остаются достаточно растворимыми, а продукт сравнительно менее растворим; тем самым формируется эвристическая прокси-оценка движущей силы выделения. Это технологическая утилита, а не модель реакционной инженерии.
  \item \textbf{\texttt{design\_solvent\_system(...)}:} подбирает бинарные системы растворителей под целевое окно растворимости, используя линейную интерполяцию по $\ln x_2$ чистых растворителей, линейное смешение в пространстве Хансена и эвристическую поправку по расстоянию в этом пространстве. Следовательно, этот блок следует описывать как приближённое проектирование бинарных систем растворителей, а не как полноценную модель термодинамики смесей.
\end{itemize}

\subsubsection{Маршрутный скрининг: явные стадии, а не ретросинтез}

Одна из наиболее чувствительных к завышенным утверждениям зон прикладного слоя --- вкладка маршрутного скрининга внутри \texttt{Applications}. Пользователь задаёт явную таблицу шагов маршрута с полями \texttt{step\_id}, \texttt{compound\_smiles}, \texttt{reaction\_temp\_k}, \texttt{isolation\_temp\_k}, \texttt{candidate\_solvents} и \texttt{goal}. После этого система \textbf{не} предлагает ретросинтетические разрывы, \textbf{не} выводит механизмы и \textbf{не} генерирует последовательность стадий. Вместо этого она ранжирует уже заданные промежуточные продукты и заранее перечисленные растворители по качеству окна выделения.

Для каждой пары «промежуточный продукт--растворитель» интерфейс рассчитывает \texttt{hot\_ln\_x2}, \texttt{cold\_ln\_x2}, \texttt{delta\_ln\_x2}, \texttt{swing\_ratio}, качественный режим и итоговый \texttt{route\_score}. В \texttt{core.py} этот показатель строится через функцию \texttt{synthesis\_window\_metrics(...)}:

\begin{equation*}
\text{route\_score} = 100\left(0.45\,\text{swing\_index} + 0.35\,\text{hot\_loading\_index} + 0.20\,\text{cold\_capture\_index}\right).
\end{equation*}

Смысл этой формулы прозрачен: соединение должно быть достаточно загружаемо на горячем конце, заметно терять растворимость при охлаждении и выглядеть склонным к выделению на холодном конце. Следовательно, корректные обозначения здесь --- отбор растворителей для явного маршрута синтеза, анализ окна выделения или ранжирование растворителей по температурному перепаду для заданных промежуточных продуктов. Названия вроде «synthesis planner», «retrosynthesis planner» или «route generator» для этого режима некорректны.

\subsubsection{drug\_properties.py: BCS-подобная оценка технологической пригодности}

Модуль \texttt{drug\_properties.py} даёт наиболее «drug-like» поверхность, но остаётся ориентированным прежде всего на растворимость по своей физике. В коде поддерживается reference panel из шести ориентировочных молекул --- paracetamol, ibuprofen, carbamazepine, caffeine, metformin и naproxen --- которые используются не как обучающие данные, а как опорные точки в пространстве дескрипторов.

Функция \texttt{bcs\_classify(...)} выполняет следующие шаги:
\begin{enumerate}
  \item предсказывает внутреннюю водную растворимость;
  \item оценивает ионизационную эвристику через SMARTS-правила для кислотных и основных мотивов;
  \item применяет приближённую pH-коррекцию в духе масштабирования по Хендерсону--Хассельбаху;
  \item рассчитывает число дозы для pH~1.0, 4.5 и 6.8;
  \item формирует BCS-подобный класс на основе бинарной оценки растворимости и дескрипторной прокси-оценки проницаемости.
\end{enumerate}

Ключевая величина имеет стандартный вид:
\begin{equation*}
D_0 = \frac{\text{dose}}{V\,S},
\end{equation*}
где \(V\) --- опорный объём водной среды, а \(S\) --- предсказанная равновесная растворимость после pH-коррекции. BCS-подобная классификация опирается на предсказание водной растворимости. В текущем корпусе вода представлена 6\,524 supervised строками (737 unique solute), что обеспечивает базовое обучение модели на водных системах. Тем не менее BCS-оценки остаются приближёнными: проницаемость оценивается дескрипторно, а pH-коррекция имеет эвристическую природу.

Функция \texttt{developability\_score(...)} агрегирует пять компонент --- растворимость, кристаллическую устойчивость, липофильность, разнообразие растворителей и температурную чувствительность --- с весами 0.32, 0.20, 0.16, 0.16 и 0.16 соответственно:
\begin{equation*}
\text{developability\_score} = 0.32\,s_{\mathrm{sol}} + 0.20\,s_{\mathrm{cryst}} + 0.16\,s_{\mathrm{lipo}} + 0.16\,s_{\mathrm{div}} + 0.16\,s_{\mathrm{temp}}.
\end{equation*}

Здесь особенно заметна польза физического узкого места: для \texttt{TGNN-Solv} компонент \(s_{\mathrm{cryst}}\) может опираться на предсказанные \(T_{\mathrm{m}}\) и \(\Delta H_{\mathrm{fus}}\), тогда как для \texttt{DirectGNN} используется более нейтральный резервный вариант, поскольку базовая модель не выдаёт параметры плавления. Наконец, процедура \texttt{salt\_cocrystal\_impact(...)} комбинирует API и противоион/коформер в виде разъединённого суррогатного представления и прямо маркируется как \textbf{низкой уверенности}. Её корректно описывать как приближённую оценку потенциала солеобразования и кокристаллизации для задач приоритизации, но не как модель ионной химии, надёжный предиктор соли или предсказатель кристаллической формы.

\subsubsection{pk\_profiling.py: PK-релевантный профиль растворимости, но не PK/PD}

Модуль \texttt{pk\_profiling.py} является самой чувствительной зоной с точки зрения формулировок. Его поддерживаемая область уже достаточно полезна, но остаётся намеренно узкой: это фармакокинетически релевантный слой анализа растворимости для анализа компартментов ЖКТ, биорелевантных сред, скрининга носителей для внутривенных форм и приоритизации носителей для местных форм.

Поддерживаемые встроенные наборы включают 5 GI-компартментов, 4 биорелевантные среды, 8 носителей для внутривенных форм и 10 носителей для местных форм. Для GI-профиля учитываются pH, объём, время пребывания, весовой коэффициент абсорбции и режим среды; затем рассчитываются скорректированная растворимость, число дозы, доля растворившегося вещества и флаг ограничения растворением, а на сводном уровне --- \texttt{max\_absorbable\_dose}, \texttt{weighted\_absorption}, \texttt{permeability\_factor}, \texttt{f\_abs\_estimate} и \texttt{rate\_limiting\_step}. По смыслу это прокси-оценка предабсорбционного давления растворения вдоль ЖКТ, но не модель абсорбции в строгом смысле.

В части биорелевантных сред поддерживаются \texttt{FaSSGF}, \texttt{FeSSGF}, \texttt{FaSSIF} и \texttt{FeSSIF}. Подход строится как растворимость в воде + pH-коррекция + эвристический множитель усиления среды. На этой основе из соотношения \texttt{FaSSIF}/\texttt{FeSSIF} выводится грубая оценка пищевого эффекта: значения не ниже 2.0 трактуются как положительный пищевой эффект, не выше 0.67 --- как отрицательный, а промежуточная область --- как минимальный. Это следует называть эвристической оценкой сдвига растворимости между режимами натощак и после еды в биорелевантных суррогатных средах, но не полноценным симулятором пищевого эффекта.

IV-сценарий ранжирует системы вроде воды для инъекций, этанола, пропиленгликоля, суррогата PEG-400, DMSO, суррогатов поверхностно-активных систем и циклодекстриноподобных носителей. Для каждого кандидата оцениваются растворимость в чистом носителе, способность смешиваться с водой, \texttt{max\_compatible\_fraction}, \texttt{osmolality\_concern}, \texttt{research\_only} и итоговая рекомендация. В кодовой эвристике рекомендация требует, чтобы расчётная концентрация при \(37\,^{\circ}\mathrm{C}\) была не ниже 1~\texttt{mg/mL}, флаг осмоляльности не был \texttt{high}, а носитель не был отмечен как \texttt{research\_only}. Следовательно, это поверхность для формуляционной первичной оценки, а не валидированный инструмент проектирования парентеральных форм.

Сценарий для местных форм оценивает растворимость в носителе, \(\gamma_2\), термодинамическую активность, показатель близости к насыщению и потенциал проникновения. Последняя величина строится как комбинация нормированной термодинамической активности и эвристики близости к насыщению. Поэтому корректная трактовка --- приоритизация носителей для местных форм по растворимости и прокси-оценке термодинамической активности, но не модель кожной PK и не симулятор дермальной абсорбции.

\warnnote{\texttt{PKSolubilityProfiler} не моделирует полную кинетику абсорбции, кинетику осаждения, эффекты транспортеров, метаболизм, клиренс, распределение и фармакодинамику. В отчёте его следует обозначать как \textbf{PK-релевантный профиль растворимости} или как \textbf{оценку сред и формуляций до абсорбции}, а не как PBPK-, PK- или PK/PD-симулятор.}

\subsubsection{Общая поддержка TGNN-Solv и DirectGNN}

Сильная сторона текущего состояния проекта состоит в том, что и \texttt{TGNN-Solv}, и \texttt{DirectGNN} поддерживают один и тот же прикладной стек: отбор растворителей, технологическую оптимизацию, оценку технологической пригодности, PK-релевантный профиль растворимости и единый контур доступа через GUI и CLI. Это делает прикладное сравнение существенно честнее, поскольку один и тот же сценарий можно прогнать на обоих чекпойнтах без отдельной «ручной» логики для каждой модели.

При этом более богатые специфические выходы TGNN сохраняются: именно \texttt{TGNN-Solv} даёт декомпозицию решателя, \(T_{\mathrm{m}}\), \(\Delta H_{\mathrm{fus}}\), \(\Phi\), \(\gamma_2\), \(\tau_{12}\), \(\tau_{21}\), интерпретацию в пространстве Хансена и объяснимость, связанную с решателем. Поэтому \texttt{DirectGNN} сопоставим по итоговым прикладным оценкам, но \texttt{TGNN-Solv} остаётся заметно сильнее по интерпретируемости и физически осмысленной диагностике.

\section{Planner}

\subsection{Planner как слой планирования экспериментов и оркестрации}

\texttt{Planner} в текущем проекте --- это канбан-доска и планировщик для экспериментальной работы с хранением состояния в репозитории. Он хранит задачи с полями \texttt{id}, \texttt{title}, \texttt{status}, \texttt{priority}, \texttt{owner}, \texttt{start}, \texttt{end}, \texttt{estimate\_hours}, \texttt{notes} и \texttt{command}. На уровне доски используются колонки \texttt{Backlog}, \texttt{Ready}, \texttt{Running}, \texttt{Blocked} и \texttt{Done}. Поддерживаются редактирование перетаскиванием, календарный и временной вид, импорт из истории инференса, неопределённости и калибровки, привязка путей к артефактам, связанных чекпойнтов и ссылок на наборы данных, а также запуск shell-команды для выбранной задачи.

\validated{\texttt{Planner} --- это слой планирования экспериментов для управления бэклогом, расписанием и последующим анализом артефактов, а не химический планировщик.}

Ключевое разведение смыслов здесь обязательно. В проекте сейчас реально присутствуют два разных типа планирования:
\begin{enumerate}
  \item \textbf{\texttt{Planner}} --- канбан-доска и планировщик для исследований, вычислительных задач, последующего анализа и управления артефактами;
  \item \textbf{маршрутно-ориентированный скрининг внутри \texttt{Applications}} --- утилита ранжирования растворителей для уже заданных стадий маршрута и явных промежуточных продуктов.
\end{enumerate}

\warnnote{Наличие рабочего пространства \texttt{Planner} не является основанием описывать проект как chemistry planner, synthesis planner, reaction planner или route generator. Это отдельный слой управления экспериментами, а не инструмент проектирования синтетических маршрутов.}

\section{Границы применимости}

\subsection{Границы применимости и корректные формулировки}

Граница заявлений в текущем проекте проходит в той точке, где выводы, ориентированные прежде всего на растворимость, начинают ошибочно читаться как механистические предсказания. Все прикладные сценарии в конечном счёте строятся вокруг равновесной растворимости, явной температуры, явной среды и набора эвристик на основе дескрипторов. Расширение проекта существенно и практически полезно, но оно \textbf{не} превращает TGNN-Solv в ретросинтетическую систему, PBPK-платформу или универсальный фармацевтический симулятор.

Наиболее важные ограничения, которые следует фиксировать явно:
\begin{itemize}
  \item перевод в \texttt{mg/mL} основан на приближённых предположениях об объёме растворителя и потому пригоден для первичного отбора, а не для строгого расчёта формуляции;
  \item pH-коррекция имеет эвристическую природу на основе SMARTS и не моделирует полноценную микросостоянийную кислотно-основную химию;
  \item бинарные системы растворителей строятся через интерполяцию и линейное смешение в пространстве Хансена, а не через полноценную смесьевую термодинамику;
  \item PK-релевантные выходы описывают давление растворения и доклинический формуляционный риск, а не кинетику, экспозицию, метаболизм или фармакодинамику;
  \item маршрутно-ориентированные выходы ранжируют варианты растворителей для заранее заданных промежуточных продуктов, а не проектируют сами синтетические маршруты.
\end{itemize}

Для удобства редактуры полезно явно зафиксировать безопасные и небезопасные формулировки.

\begin{center}
\small
\begin{tabularx}{\textwidth}{|p{0.24\textwidth}|Y|Y|}
\hline
\textbf{Область} & \textbf{Корректная формулировка} & \textbf{Чего следует избегать} \\
\hline
\texttt{Experiment Lab} & поддерживаемая Streamlit-поверхность и слой оркестрации над CLI/API & отдельный продукт или параллельная реализация модели \\
\hline
Маршрутный скрининг & отбор растворителей для заданного маршрута; анализ окон выделения & планировщик синтеза; ретросинтетический планировщик; генератор маршрутов \\
\hline
PK-скрининг & PK-релевантный профиль растворимости; оценка сред до абсорбции & PBPK-симулятор; PK-симулятор; PK/PD-симулятор; модель экспозиции \\
\hline
\texttt{Planner} & планировщик экспериментов; канбан-доска с хранением в репозитории; слой оркестрации & химический планировщик; планировщик синтеза; планировщик реакций \\
\hline
Оценка технологической пригодности & BCS-подобная первичная оценка; приближённая оценка соли/кокристалла & регуляторное доказательство BCS-класса; надёжный предиктор солей; предиктор кристаллической формы \\
\hline
\end{tabularx}
\normalsize
\end{center}

\keyidea{Самая точная позиционировка проекта в текущем состоянии: поддерживаемая исследовательская платформа с интегрированным GUI, контуром бенчмарка и артефактов, а также прикладным слоем для растворителей, технологических решений, оценки технологической пригодности и PK-релевантного скрининга.}

\section{Риски и релизный пакет}

\subsection{Риски и минимальный релизный пакет}

Ключевые риски бенчмарк-, GUI- и прикладного контуров:
\begin{itemize}
  \item смешение архитектурного эффекта с эффектом признаков, инициализации и вычислительного бюджета;
  \item нестабильность адаптационных обёрток внешних моделей и несопоставимость режимов применения без дообучения, калиброванного режима и полного переобучения;
  \item риск завышенных утверждений при переносе выводов, ориентированных прежде всего на растворимость, на широкий PK/PD-, PBPK- или контекст планирования синтеза;
  \item потеря смыслового различия между \texttt{Planner} как доской экспериментов и маршрутным скринингом внутри \texttt{Applications};
  \item неполная калибровка неопределённости вне области применимости и завышенная уверенность в эвристических ветках \texttt{mg/mL}, pH-коррекции и проектирования бинарных смесей.
\end{itemize}

Минимальный релизный пакет для воспроизводимого результата должен включать:
\begin{enumerate}
  \item \texttt{run\_manifest} + \texttt{resolved\_config};
  \item таблицы метрик по начальным зёрнам, построчные тестовые предсказания и, при наличии, диагностики решателя, неопределённости и калибровки;
  \item \texttt{benchmark\_card} с явным указанием режима внешней модели и границ заявлений;
  \item при использовании \texttt{Experiment Lab} --- соответствующие артефакты с хранением в репозитории: history-export, preset/DAG-export, job-state и связанные пути к чекпойнтам и таблицам;
  \item замороженный снимок релизного набора: разбиение, конфиги, чекпойнты и версии зависимостей.
\end{enumerate}

Именно такой пакет позволяет верифицировать не только численные метрики, но и то, что визуальная оркестрационная поверхность, прикладные сценарии и итоговые формулировки отчёта действительно опираются на воспроизводимые артефакты, а не на маркетинговую интерпретацию.

\appendix
\section*{Приложение A. Таблицы групповых вкладов для GC-методов}
\label{app:gc-tables}
\addcontentsline{toc}{section}{Приложение A. Таблицы групповых вкладов для GC-методов}

В приложение вынесены рабочие справочные таблицы, на которые ссылается §\ref{subsec:gc-methods}. Это позволяет не перегружать основную методическую часть длинными перечнями SMARTS-групп и численных инкрементов.

\subsection*{A.1. Joback-инкременты для $T_{\mathrm{m}}$}
\label{appsubsec:joback-tm}

\textbf{Таблица A.1.} Базовые Joback-инкременты для оценки $T_{\mathrm{m}}$ (рабочий справочный набор).

Ниже приведена рабочая таблица наиболее часто используемых групп для оценки температуры плавления по group-contribution подходу Joback-типа. Значения даны как справочные инкременты и должны верифицироваться для конкретной параметризации (Joback / Joback-Reid / модификации).

\begin{center}
\small
\setlength{\tabcolsep}{3pt}
\begin{longtable}{|P{2.7cm}|T{4.0cm}|C{1.9cm}|P{4.3cm}|}
\hline
\textbf{Группа} & \textbf{SMARTS} & \textbf{$\Delta T_{\mathrm{m}}$ (K)} & \textbf{Пример молекулы} \\
\hline
\endfirsthead
\hline
\textbf{Группа} & \textbf{SMARTS} & \textbf{$\Delta T_{\mathrm{m}}$ (K)} & \textbf{Пример молекулы} \\
\hline
\endhead
\hline
\multicolumn{4}{r}{\textit{Продолжение на следующей странице}} \\
\hline
\endfoot
\hline
\endlastfoot
CH$_3$ & \smarts{[CH3]} & 23.6 & н-гексан, толуол \\
\hline
CH$_2$ & \smarts{[CH2]} & 22.9 & н-гексан, этилбензол \\
\hline
CH & \smarts{[CH]([*])([*])} & 21.3 & изопропанол, изобутан \\
\hline
C (четвертичный) & \smarts{[C]([*])([*])([*])([*])} & 18.8 & неопентан (центр) \\
\hline
=CH$_2$ & \smarts{[CH2]=[*]} & 18.1 & пропен, стирол \\
\hline
=CH & \smarts{[CH]=[*]} & 16.9 & 2-бутен \\
\hline
=C & \smarts{[C]=[*]} & 14.5 & изобутилен (центральный C) \\
\hline
$\equiv$CH & \smarts{[CH]\#[*]} & 13.7 & ацетилен \\
\hline
$\equiv$C & \smarts{[C]\#[*]} & 12.8 & 2-бутин \\
\hline
aromatic CH & \smarts{[cH]} & 20.0 & бензол, толуол \\
\hline
aromatic C & \smarts{[c]([*])} & 17.2 & толуол, ксилол \\
\hline
OH (alcohol) & \smarts{[OX2H][CX4]} & 41.0 & метанол, этанол \\
\hline
OH (phenol) & \smarts{[OX2H][c]} & 54.0 & фенол \\
\hline
O (ether) & \smarts{[OD2]([\#6])[\#6]} & 28.5 & диэтиловый эфир, ТГФ \\
\hline
C=O (ketone) & \smarts{[CX3](=O)[\#6]} & 33.4 & ацетон, циклогексанон \\
\hline
CHO (aldehyde) & \smarts{[CX3H](=O)[\#6]} & 36.2 & ацетальдегид, бензальдегид \\
\hline
COOH & \smarts{[CX3](=O)[OX2H1]} & 68.0 & уксусная кислота, бензойная кислота \\
\hline
COO (ester) & \smarts{[CX3](=O)[OX2][\#6]} & 39.6 & этилацетат, метилбензоат \\
\hline
NH$_2$ & \smarts{[NX3H2][\#6]} & 38.7 & этиламин, анилин \\
\hline
NH & \smarts{[NX3H]([\#6])[\#6]} & 32.8 & диэтиламин \\
\hline
N (tertiary) & \smarts{[NX3]([\#6])([\#6])[\#6]} & 27.5 & триэтиламин \\
\hline
\end{longtable}
\normalsize
\end{center}

\subsection*{A.2. Инкременты для $\Delta H_{\mathrm{fus}}$ (Ducros-подобные корреляции)}
\label{appsubsec:ducros-dhfus}

Таблица ниже отражает типичный набор групповых вкладов для приближённой оценки энтальпии плавления в GC-корреляциях (Ducros/аналогичные инженерные схемы). Значения приведены как рабочие ориентиры в единицах $\mathrm{J/mol}$.

\begin{center}
\small
\setlength{\tabcolsep}{3pt}
\begin{longtable}{|P{2.7cm}|T{4.0cm}|C{2.4cm}|P{3.8cm}|}
\hline
\textbf{Группа} & \textbf{SMARTS} & \textbf{$\Delta H_{\mathrm{fus}}$ (J/mol)} & \textbf{Пример молекулы} \\
\hline
\endfirsthead
\hline
\textbf{Группа} & \textbf{SMARTS} & \textbf{$\Delta H_{\mathrm{fus}}$ (J/mol)} & \textbf{Пример молекулы} \\
\hline
\endhead
\hline
\multicolumn{4}{r}{\textit{Продолжение на следующей странице}} \\
\hline
\endfoot
\hline
\endlastfoot
CH$_3$ & \smarts{[CH3]} & 450 & алканы, алкилзамещённые ароматические \\
\hline
CH$_2$ & \smarts{[CH2]} & 520 & линейные алканы \\
\hline
CH & \smarts{[CH]([*])([*])} & 610 & разветвлённые алканы \\
\hline
C (четвертичный) & \smarts{[C]([*])([*])([*])([*])} & 700 & неопентильные структуры \\
\hline
=CH$_2$ & \smarts{[CH2]=[*]} & 580 & алкены \\
\hline
=CH & \smarts{[CH]=[*]} & 640 & внутренние алкены \\
\hline
=C & \smarts{[C]=[*]} & 730 & замещённые алкены \\
\hline
aromatic CH & \smarts{[cH]} & 800 & бензол, стирол \\
\hline
aromatic C & \smarts{[c]([*])} & 920 & алкилбензолы \\
\hline
OH (alcohol) & \smarts{[OX2H][CX4]} & 1800 & спирты \\
\hline
OH (phenol) & \smarts{[OX2H][c]} & 2200 & фенолы \\
\hline
O (ether) & \smarts{[OD2]([\#6])[\#6]} & 1050 & простые эфиры \\
\hline
C=O (ketone) & \smarts{[CX3](=O)[\#6]} & 1350 & кетоны \\
\hline
CHO (aldehyde) & \smarts{[CX3H](=O)[\#6]} & 1480 & альдегиды \\
\hline
COOH & \smarts{[CX3](=O)[OX2H1]} & 2750 & карбоновые кислоты \\
\hline
COO (ester) & \smarts{[CX3](=O)[OX2][\#6]} & 1620 & сложные эфиры \\
\hline
NH$_2$ & \smarts{[NX3H2][\#6]} & 1700 & первичные амины \\
\hline
NH & \smarts{[NX3H]([\#6])[\#6]} & 1420 & вторичные амины \\
\hline
N (tertiary) & \smarts{[NX3]([\#6])([\#6])[\#6]} & 1200 & третичные амины \\
\hline
\end{longtable}
\normalsize
\end{center}

\subsection*{A.3. Параметры $r$ и $q$ для UNIQUAC/UNIFAC-совместимых групп (Bondi-тип)}

\begin{center}
\small
\begin{longtable}{|p{4.2cm}|p{2.4cm}|p{2.4cm}|}
\hline
\textbf{Группа} & \textbf{$R_k$} & \textbf{$Q_k$} \\
\hline
\endfirsthead
\hline
\textbf{Группа} & \textbf{$R_k$} & \textbf{$Q_k$} \\
\hline
\endhead
\hline
\multicolumn{3}{r}{\textit{Продолжение на следующей странице}} \\
\hline
\endfoot
\hline
\endlastfoot
CH$_3$ & 0.9011 & 0.8480 \\
\hline
CH$_2$ & 0.6744 & 0.5400 \\
\hline
CH & 0.4469 & 0.2280 \\
\hline
C (четвертичный) & 0.2195 & 0.0000 \\
\hline
=CH$_2$ & 1.3454 & 1.1760 \\
\hline
=CH & 1.1167 & 0.8670 \\
\hline
=C & 0.8886 & 0.6760 \\
\hline
aromatic CH & 0.5313 & 0.4000 \\
\hline
aromatic C & 0.3652 & 0.1200 \\
\hline
OH (alcohol) & 1.0000 & 1.2000 \\
\hline
OH (phenol) & 1.0000 & 1.2000 \\
\hline
O (ether) & 1.1450 & 1.0880 \\
\hline
C=O (ketone) & 1.6724 & 1.4880 \\
\hline
CHO (aldehyde) & 1.4457 & 1.1800 \\
\hline
COOH & 1.3013 & 1.2240 \\
\hline
COO (ester) & 1.9031 & 1.7280 \\
\hline
NH$_2$ & 1.3450 & 1.2000 \\
\hline
NH & 1.1000 & 0.9500 \\
\hline
N (tertiary) & 0.9200 & 0.7900 \\
\hline
\end{longtable}
\normalsize
\end{center}

\subsection*{A.4. Типичные значения параметров NRTL}

Ниже приведены репрезентативные параметры для иллюстрации масштабов $\Delta g_{12}$, $\Delta g_{21}$ и $\alpha$. Они зависят от диапазона температур, набора данных и метода регрессии.

\begin{center}
\small
\setlength{\tabcolsep}{3pt}
\begin{longtable}{|P{2.8cm}|C{2.1cm}|C{2.1cm}|C{0.9cm}|P{4.6cm}|}
\hline
\textbf{Система} & \textbf{$\Delta g_{12}$ (J/mol)} & \textbf{$\Delta g_{21}$ (J/mol)} & \textbf{$\alpha$} & \textbf{Источник} \\
\hline
\endfirsthead
\hline
\textbf{Система} & \textbf{$\Delta g_{12}$ (J/mol)} & \textbf{$\Delta g_{21}$ (J/mol)} & \textbf{$\alpha$} & \textbf{Источник} \\
\hline
\endhead
\hline
\multicolumn{5}{r}{\textit{Продолжение на следующей странице}} \\
\hline
\endfoot
\hline
\endlastfoot
этанол--вода & $-120$ & $1450$ & 0.30 & авторская компиляция по литературным NRTL-подборам (см.~\cite{renon1968,prausnitz1998}) \\
\hline
бензол--гексан & $620$ & $410$ & 0.30 & авторская компиляция по литературным NRTL-подборам (см.~\cite{renon1968,prausnitz1998}) \\
\hline
ацетон--хлороформ & $-950$ & $1750$ & 0.30 & авторская компиляция по литературным NRTL-подборам (см.~\cite{renon1968,prausnitz1998}) \\
\hline
\end{longtable}
\normalsize
\end{center}

\subsection*{A.5. Hansen-параметры референсных растворителей}

Единицы: MPa$^{1/2}$.
\begin{equation*}
\delta_t = \sqrt{\delta_d^2 + \delta_p^2 + \delta_h^2}.
\end{equation*}

\begin{center}
\small
\begin{longtable}{|p{3.6cm}|p{1.8cm}|p{1.8cm}|p{1.8cm}|p{1.8cm}|}
\hline
\textbf{Растворитель} & \textbf{$\delta_d$} & \textbf{$\delta_p$} & \textbf{$\delta_h$} & \textbf{$\delta_t$} \\
\hline
\endfirsthead
\hline
\textbf{Растворитель} & \textbf{$\delta_d$} & \textbf{$\delta_p$} & \textbf{$\delta_h$} & \textbf{$\delta_t$} \\
\hline
\endhead
\hline
\multicolumn{5}{r}{\textit{Продолжение на следующей странице}} \\
\hline
\endfoot
\hline
\endlastfoot
вода & 15.5 & 16.0 & 42.3 & 47.8 \\
\hline
метанол & 15.1 & 12.3 & 22.3 & 29.7 \\
\hline
этанол & 15.8 & 8.8 & 19.4 & 26.0 \\
\hline
изопропанол & 15.8 & 6.1 & 16.4 & 23.5 \\
\hline
ацетон & 15.5 & 10.4 & 7.0 & 19.9 \\
\hline
ацетонитрил & 15.3 & 18.0 & 6.1 & 24.3 \\
\hline
ДМСО & 18.4 & 16.4 & 10.2 & 26.7 \\
\hline
ДМФА & 17.4 & 13.7 & 11.3 & 24.8 \\
\hline
NMP & 18.0 & 12.3 & 7.2 & 22.9 \\
\hline
ТГФ & 16.8 & 5.7 & 8.0 & 19.4 \\
\hline
этилацетат & 15.8 & 5.3 & 7.2 & 18.6 \\
\hline
хлороформ & 17.8 & 3.1 & 5.7 & 19.0 \\
\hline
толуол & 18.0 & 1.4 & 2.0 & 18.2 \\
\hline
бензол & 18.4 & 0.0 & 2.0 & 18.5 \\
\hline
гексан & 14.9 & 0.0 & 0.0 & 14.9 \\
\hline
\end{longtable}
\normalsize
\end{center}

\subsection*{A.6. Ранжирование растворителей в supervised subset}

Для текущей water-inclusive конфигурации основной solvent-frequency ranking следует рассматривать как отдельный служебный артефакт корпуса. В основном тексте отчёта достаточно сводных метрик концентрации и отдельной водной строки, а расширенный список рекомендуется выносить в приложение: top-50 --- в основной текст приложения, оставшиеся 163 растворителя --- при необходимости в дополнительную таблицу или внешний CSV-артефакт.

\section*{Приложение B. Доказательства термодинамической консистентности}
\addcontentsline{toc}{section}{Приложение B. Доказательства термодинамической консистентности}

Ниже собраны полные выкладки, на которые опираются разделы с постановкой SLE и моделями активности. В основном тексте использовались только итоговые формулы и физическая интерпретация, а доказательства вынесены сюда ради компактности.

\subsection*{B.1. Уравнение Гиббса--Дюгема: полный вывод}

Исходное фундаментальное соотношение для многокомпонентной системы:
\begin{equation*}
dG = -S\,dT + V\,dP + \sum_i \mu_i\,dn_i.
\end{equation*}

С другой стороны, поскольку $G$ является экстенсивной функцией,
\begin{equation*}
G = \sum_i n_i\mu_i.
\end{equation*}
Дифференцируя:
\begin{equation*}
dG = \sum_i \mu_i\,dn_i + \sum_i n_i\,d\mu_i.
\end{equation*}

Приравнивая два выражения для $dG$, получаем
\begin{equation*}
\sum_i n_i\,d\mu_i = -S\,dT + V\,dP.
\end{equation*}

При постоянных $T,P$:
\begin{equation*}
\sum_i n_i\,d\mu_i = 0.
\end{equation*}
Делим на $n=\sum_i n_i$ и вводим $x_i=n_i/n$:
\begin{equation*}
\sum_i x_i\,d\mu_i = 0.
\end{equation*}
Это и есть дифференциальная форма уравнения Гиббса--Дюгема при const $T,P$.

Теперь используем стандартное представление химпотенциала:
\begin{equation*}
\mu_i = \mu_i^{\circ}(T,P) + RT\ln(\gamma_i x_i).
\end{equation*}
При const $T,P$ имеем $d\mu_i^{\circ}=0$, $d(RT)=0$, поэтому
\begin{equation*}
d\mu_i = RT\,d\ln(\gamma_i x_i)=RT\bigl(d\ln\gamma_i + d\ln x_i\bigr).
\end{equation*}
Подставляем в Гиббса--Дюгема:
\begin{equation*}
\sum_i x_i\,d\ln\gamma_i + \sum_i x_i\,d\ln x_i = 0.
\end{equation*}

Но
\begin{equation*}
\sum_i x_i\,d\ln x_i = \sum_i x_i\frac{dx_i}{x_i}=\sum_i dx_i=d\left(\sum_i x_i\right)=d(1)=0,
\end{equation*}
следовательно
\begin{equation*}
\sum_i x_i\,d\ln\gamma_i=0 \qquad (T,P=\text{const}).
\end{equation*}

Для бинарной смеси ($x_2=1-x_1$, $dx_2=-dx_1$):
\begin{equation*}
x_1\,d\ln\gamma_1 + x_2\,d\ln\gamma_2 = 0.
\end{equation*}
Если зависимость берётся по $x_1$ при const $T,P$,
\begin{equation*}
\left[x_1\left(\frac{\partial\ln\gamma_1}{\partial x_1}\right)_{T,P}
+
x_2\left(\frac{\partial\ln\gamma_2}{\partial x_1}\right)_{T,P}\right]dx_1=0,
\end{equation*}
и потому
\begin{equation*}
x_1\left(\frac{\partial\ln\gamma_1}{\partial x_1}\right)_{T,P}
+
x_2\left(\frac{\partial\ln\gamma_2}{\partial x_1}\right)_{T,P}=0.
\end{equation*}

\subsection*{B.2. Почему NRTL автоматически удовлетворяет Гиббсу--Дюгему}

Ключевой факт: в NRTL коэффициенты активности получаются из \emph{одного} скалярного потенциала --- избыточной энергии Гиббса:
\begin{equation*}
\frac{G^E}{RT}=n\,g^E(T,P,\mathbf{x}),
\end{equation*}
где $g^E$ --- молярная избыточная величина (интенсивная функция состава).

Из общей термодинамической связи:
\begin{equation*}
\mu_i^E = \left(\frac{\partial G^E}{\partial n_i}\right)_{T,P,n_{j\neq i}},
\qquad
\ln\gamma_i = \frac{\mu_i^E}{RT}.
\end{equation*}

То есть $\ln\gamma_i$ получаются дифференцированием одного и того же $G^E$. Для экстенсивной функции $G^E(n_1,\dots,n_c)$ первого порядка однородности по теореме Эйлера:
\begin{equation*}
G^E = \sum_i n_i\mu_i^E.
\end{equation*}
Дифференцируя и сравнивая с фундаментальным соотношением для избыточных величин при const $T,P$,
\begin{equation*}
dG^E = \sum_i \mu_i^E\,dn_i,
\end{equation*}
получаем
\begin{equation*}
\sum_i n_i\,d\mu_i^E = 0
\quad\Longrightarrow\quad
\sum_i x_i\,d\mu_i^E =0
\quad\Longrightarrow\quad
\sum_i x_i\,d\ln\gamma_i=0.
\end{equation*}

Следовательно, любая модель, где все $\ln\gamma_i$ строятся как производные единого согласованного $G^E$ (включая NRTL), автоматически удовлетворяет уравнению Гиббса--Дюгема.

\subsection*{B.3. Уравнение Гиббса--Гельмгольца: вывод}

Начинаем с определения
\begin{equation*}
G = H-TS,
\end{equation*}
и тождества
\begin{equation*}
\left(\frac{\partial G}{\partial T}\right)_P=-S.
\end{equation*}

Рассмотрим $G/T$:
\begin{equation*}
\left(\frac{\partial (G/T)}{\partial T}\right)_P
=
\frac{1}{T}\left(\frac{\partial G}{\partial T}\right)_P - \frac{G}{T^2}
=
-\frac{S}{T}-\frac{G}{T^2}.
\end{equation*}
Так как $H=G+TS$, то
\begin{equation*}
-\frac{S}{T}-\frac{G}{T^2}=-\frac{G+TS}{T^2}=-\frac{H}{T^2}.
\end{equation*}
Итого:
\begin{equation*}
\left(\frac{\partial (G/T)}{\partial T}\right)_P=-\frac{H}{T^2}.
\end{equation*}

Эквивалентная форма по переменной $u=1/T$:
\begin{equation*}
\left(\frac{\partial (G/T)}{\partial u}\right)_P
=
\left(\frac{\partial (G/T)}{\partial T}\right)_P\left(\frac{\partial T}{\partial u}\right)
=
\left(-\frac{H}{T^2}\right)(-T^2)=H.
\end{equation*}

Для избыточных величин:
\begin{equation*}
\mu_i^E = RT\ln\gamma_i
\quad\Longrightarrow\quad
\frac{\mu_i^E}{T}=R\ln\gamma_i.
\end{equation*}
Применяя Гиббса--Гельмгольца к парциальной избыточной величине,
\begin{equation*}
\left(\frac{\partial (\mu_i^E/T)}{\partial (1/T)}\right)_{P,\mathbf{x}}=\bar H_i^E,
\end{equation*}
получаем
\begin{equation*}
R\left(\frac{\partial \ln\gamma_i}{\partial (1/T)}\right)_{P,\mathbf{x}}=\bar H_i^E,
\end{equation*}
или
\begin{equation*}
\left(\frac{\partial \ln\gamma_i}{\partial (1/T)}\right)_{P,\mathbf{x}}=\frac{\bar H_i^E}{R}.
\end{equation*}

\subsection*{B.4. Применение к NRTL}

В NRTL обычно вводят
\begin{equation*}
\tau_{ij}(T)=\frac{\Delta g_{ij}}{RT}
\end{equation*}
(при температурно-независимом $\Delta g_{ij}$ в базовой параметризации). Тогда
\begin{equation*}
\frac{\partial \tau_{ij}}{\partial T}=-\frac{\Delta g_{ij}}{RT^2}=-\frac{\tau_{ij}}{T}.
\end{equation*}

Поскольку $\ln\gamma_i$ в NRTL является функцией $\tau_{ij}$, $G_{ij}=\exp(-\alpha_{ij}\tau_{ij})$ и состава,
\begin{equation*}
\ln\gamma_i = f_i(\tau_{12},\tau_{21},G_{12},G_{21},x_1,x_2),
\end{equation*}
температурная производная получается по правилу цепочки:
\begin{equation*}
\left(\frac{\partial\ln\gamma_i}{\partial T}\right)_{P,\mathbf{x}}
=
\sum_{mn}\frac{\partial f_i}{\partial \tau_{mn}}\frac{\partial \tau_{mn}}{\partial T}
+
\sum_{mn}\frac{\partial f_i}{\partial G_{mn}}\frac{\partial G_{mn}}{\partial T}.
\end{equation*}

После перехода к производной по $1/T$:
\begin{equation*}
\left(\frac{\partial\ln\gamma_i}{\partial (1/T)}\right)_{P,\mathbf{x}}
=
-T^2\left(\frac{\partial\ln\gamma_i}{\partial T}\right)_{P,\mathbf{x}},
\end{equation*}
и, согласно B.3,
\begin{equation*}
\bar H_i^E = R\left(\frac{\partial\ln\gamma_i}{\partial (1/T)}\right)_{P,\mathbf{x}}.
\end{equation*}
Таким образом, через температурную зависимость $\tau_{ij}(T)$ NRTL непосредственно задаёт $H^E(T,\mathbf{x})$ и парциальные избыточные энтальпии.

\subsection*{B.5. Проверка консистентности в коде}

Для бинарной смеси удобно контролировать остаток Гиббса--Дюгема через численные производные.
Пусть
\begin{equation*}
\eta_1 = \left(\frac{\partial\ln\gamma_1}{\partial x_1}\right)_{T,P},
\qquad
\eta_2 = \left(\frac{\partial\ln\gamma_2}{\partial x_1}\right)_{T,P}.
\end{equation*}
Тогда
\begin{equation*}
\mathrm{GD\_residual} = x_1\eta_1 + x_2\eta_2.
\end{equation*}
Для аналитически согласованной реализации ожидается
\begin{equation*}
|\mathrm{GD\_residual}| < 10^{-6}
\end{equation*}
(в типичном диапазоне составов, вне сингулярных точек $x_i\to0$).

На практике малый ненулевой остаток возможен из-за:
\begin{enumerate}
  \item конечных разностей (шаг дискретизации);
  \item округления (FP32/FP64);
  \item неполной сходимости внутреннего численного решателя.
\end{enumerate}

Как диагностическую метрику можно ввести дополнительную функцию контроля:
\begin{equation*}
\mathcal{L}_{\mathrm{GD}}=\mathbb{E}\bigl[(x_1\eta_1+x_2\eta_2)^2\bigr].
\end{equation*}
В текущем проекте она рекомендуется в первую очередь как \textbf{проверка корректности реализации} (unit/diagnostic check), а не как основной обучающий член функции потерь.




\section*{Приложение C. Численные аспекты SLE-решателя}
\addcontentsline{toc}{section}{Приложение C. Численные аспекты SLE-решателя}

Это приложение собирает инженерные детали, вынесенные из основной части: псевдокод решателя неподвижной точки, эвристики стабилизации, резервные режимы и практические рекомендации по точности вычислений. Теоретические формулы SLE и имплицитного дифференцирования в основном тексте не повторяются полностью; здесь акцент сделан на реализации.

\subsection*{C.1. Метод последовательных подстановок: псевдокод}

Рассматривается итерационная схема для уравнения вида
\begin{equation*}
x = \exp\!\bigl(-\Phi - \ln\gamma_2(x)\bigr),
\end{equation*}
где $\ln\gamma_2$ вычисляется по параметрам NRTL.

\noindent\textbf{Псевдокод:}
\begin{verbatim}
Input: Phi, tau_12, tau_21, alpha, n_iter, lambda, tol
x = exp(-Phi)  # начальное приближение

G_12 = exp(-alpha * tau_12)
G_21 = exp(-alpha * tau_21)

for k = 1 to n_iter:
    x1 = 1 - x
    A = x + x1 * G_12
    B = x1 + x * G_21

    ln_gamma = x1^2 * [tau_12 * (G_12 / A)^2 + tau_21 * G_21 / B^2]

    x_new = exp(-Phi - ln_gamma)

    # демпфирование
    x_damped = lambda * x_new + (1 - lambda) * x

    if abs(x_new - x) < tol:
        x = x_damped
        break

    x = x_damped

return x
\end{verbatim}

Здесь $\lambda\in(0,1]$ --- коэффициент демпфирования, а критерий останова по невязке обновления выбран в форме $|x_{\text{new}}-x|<\mathrm{tol}$.

\subsection*{C.2. Выбор демпфирования}

\begin{itemize}
  \item \textbf{Без демпфирования ($\lambda=1$):} максимальная скорость при хорошей сжимаемости отображения, но возможны расходимость и двухточечные осцилляции.
  \item \textbf{Адаптивное демпфирование:} увеличивать $\lambda$, если $|x_{k+1}-x_k|$ монотонно уменьшается; уменьшать $\lambda$, если шаг начинает расти.
  \item \textbf{Фиксированное $\lambda=0.7$:} практичный устойчивый режим по умолчанию для широкого диапазона параметров.
\end{itemize}

Простейшее адаптивное правило можно записать как
\begin{equation*}
\lambda_{k+1}=
\begin{cases}
\min(1,\;\lambda_k+\Delta_+), & |\Delta x_k|<|\Delta x_{k-1}|,\\
\max(\lambda_{\min},\;\beta\lambda_k), & |\Delta x_k|\ge |\Delta x_{k-1}|,
\end{cases}
\end{equation*}
где $\Delta x_k=x_{k+1}-x_k$, $\beta\in(0,1)$.

\subsection*{C.3. Обнаружение расслаивания и неустойчивости итераций}

Практические признаки проблемного режима:
\begin{enumerate}
  \item осцилляции без затухания (например, $x_k$ чередуется между двумя уровнями);
  \item локально $|g'(x^*)|>1$ в окрестности фиксированной точки.
\end{enumerate}

Если итерация подстановок нестабильна, применяются резервные стратегии:
\begin{itemize}
  \item переключение на бисекцию для корня уравнения $f(x)=x-g(x)=0$;
  \item при множественности корней --- выбор меньшего физически интерпретируемого корня (режим насыщенного раствора для заданной постановки).
\end{itemize}

Диагностически удобно контролировать индекс осцилляции:
\begin{equation*}
\mathrm{osc}_k = |x_k-x_{k-2}| - |x_k-x_{k-1}|,
\end{equation*}
и при устойчиво положительном $\mathrm{osc}_k$ принудительно снижать $\lambda$ или активировать резервный решатель.

\subsection*{C.4. Численная стабильность}

Типичные источники численных ошибок:
\begin{itemize}
  \item $\exp(-\Phi)$ при больших $\Phi$ (например, $\Phi>80$) даёт численное зануление;
  \item знаменатели $A,B$ могут становиться малыми при экстремальных $\tau_{ij}$;
  \item в неявной производной знаменатель вида $1+x\eta$ может приближаться к нулю при $g'(x^*)\to 1$.
\end{itemize}

Стандартные меры стабилизации:
\begin{enumerate}
  \item по возможности вычисления в логарифмическом пространстве;
  \item регуляризация знаменателей: $A\leftarrow A+\varepsilon$, $B\leftarrow B+\varepsilon$, $\varepsilon=10^{-10}$;
  \item ограничение имплицитного знаменателя:
\begin{equation*}
\mathrm{denom}_{\mathrm{safe}}=\mathrm{sign}(\mathrm{denom})\cdot\max\bigl(|\mathrm{denom}|,0.01\bigr).
\end{equation*}
\end{enumerate}

Дополнительно полезны мягкие ограничения диапазона:
\begin{equation*}
x\leftarrow \mathrm{clip}(x, x_{\min}, x_{\max}),
\quad
x_{\min}\sim10^{-12},\; x_{\max}\sim 1-10^{-12}.
\end{equation*}

\subsection*{C.5. Неявное дифференцирование: детали реализации}

Полный вывод формулы для $\partial x^*/\partial\theta$ уже приведён в основной части; ниже зафиксирован только инженерный чек-лист для реализации собственного обратного прохода.

\textbf{Прямой проход:}
\begin{itemize}
  \item итерации фиксированной точки выполнять без накопления графа градиентов (с отсоединением промежуточных $x^{(k)}$ от графа);
  \item сохранять только финальное решение $x^*$ и необходимые тензоры для пересчёта $\ln\gamma_2$.
\end{itemize}

\textbf{Обратный проход:}
\begin{enumerate}
  \item вычислить
\begin{equation*}
\eta = \left(\frac{\partial \ln\gamma_2}{\partial x_2}\right)_{\theta}
\end{equation*}
через автоматическое дифференцирование в точке $x^*$;
  \item сформировать
\begin{equation*}
\mathrm{denom}=1+x^*\eta;
\end{equation*}
  \item для любого параметра $\theta$ использовать
\begin{equation*}
\frac{\partial x^*}{\partial \theta}
=
-\frac{x^*}{\mathrm{denom}}
\frac{\partial(\Phi+\ln\gamma_2)}{\partial\theta}.
\end{equation*}
\end{enumerate}

Преимущества подхода:
\begin{itemize}
  \item \textbf{$O(1)$ память} вместо $O(n_{\mathrm{iter}})$ при развёрнутом обратном распространении;
  \item градиенты не зависят от длины развёртки итераций и не деградируют из-за усечения траектории.
\end{itemize}

\subsection*{C.6. Точность вычислений: FP32 по сравнению с FP64}

\begin{center}
\small
\begin{longtable}{|p{3.7cm}|p{2.3cm}|p{5.3cm}|}
\hline
\textbf{Операция} & \textbf{FP32 допустим?} & \textbf{Рекомендация} \\
\hline
\endfirsthead
\hline
\textbf{Операция} & \textbf{FP32 допустим?} & \textbf{Рекомендация} \\
\hline
\endhead
\hline
\multicolumn{3}{r}{\textit{Продолжение на следующей странице}} \\
\hline
\endfoot
\hline
\endlastfoot
Прямой проход GNN & Да & Для ускорения использовать AMP/FP16 (при контроле стабильности лосса). \\
\hline
Итерации SLE & Да & Базовый режим FP32; при проблемных системах проверять чувствительность к FP64. \\
\hline
Неявный обратный проход & Осторожно & FP32 + ограничение и регуляризация знаменателя; мониторить частоту малых denom. \\
\hline
Произведение $g'^N$ (развёртка) & Рискованно & Предпочесть неявный подход; при необходимости развёртки --- локально FP64. \\
\hline
\end{longtable}
\normalsize
\end{center}

Итого: практический компромисс для рабочего и исследовательского режима --- GNN в AMP, SLE-решатель и имплицитные производные в стабильном FP32 с защитами по знаменателям и резервными стратегиями.


\begin{thebibliography}{99}
  \bibitem{renon1968} Renon, H., \& Prausnitz, J. M. (1968). Local compositions in thermodynamic excess functions for liquid mixtures. \textit{AIChE Journal}, 14(1), 135--144. \url{https://doi.org/10.1002/aic.690140124}
  \bibitem{prausnitz1998} Prausnitz, J. M., Lichtenthaler, R. N., \& de Azevedo, E. G. (1998). \textit{Molecular Thermodynamics of Fluid-Phase Equilibria} (3rd ed.). Prentice Hall.
  \bibitem{yalkowsky1979} Yalkowsky, S. H. (1979). Estimation of entropies of fusion of organic compounds. \textit{Industrial \& Engineering Chemistry Process Design and Development}, 18(1), 108--111. \url{https://doi.org/10.1021/i260069a021}
  \bibitem{hildebrand1970} Hildebrand, J. H., Prausnitz, J. M., \& Scott, R. L. (1970). \textit{Regular and Related Solutions}. Van Nostrand Reinhold.

  \bibitem{abrams1975} Abrams, D. S., \& Prausnitz, J. M. (1975). Statistical thermodynamics of liquid mixtures: A new expression for the excess Gibbs energy of partly or completely miscible systems. \textit{AIChE Journal}, 21(1), 116--128. \url{https://doi.org/10.1002/aic.690210115}
  \bibitem{wilson1964} Wilson, G. M. (1964). Vapor-liquid equilibrium. XI. A new expression for the excess free energy of mixing. \textit{Journal of the American Chemical Society}, 86(2), 127--130. \url{https://doi.org/10.1021/ja01056a002}
  \bibitem{fredenslund1975} Fredenslund, A., Jones, R. L., \& Prausnitz, J. M. (1975). Group-contribution estimation of activity coefficients in nonideal liquid mixtures. \textit{AIChE Journal}, 21(6), 1086--1099. \url{https://doi.org/10.1002/aic.690210607}
  \bibitem{hansen2007} Hansen, C. M. (2007). \textit{Hansen Solubility Parameters: A User's Handbook} (2nd ed.). CRC Press.
  \bibitem{apelblat1999} Apelblat, A., \& Manzurola, E. (1999). Solubilities of o-acetylsalicylic, 4-aminosalicylic, 3,5-dinitrosalicylic, and p-toluic acid in water from T=(278 to 348) K. \textit{Journal of Chemical Thermodynamics}, 31(1), 85--91. \url{https://doi.org/10.1006/jcht.1998.0425}
  \bibitem{acree2010} Acree, W. E., \& Chickos, J. S. (2010). Phase change enthalpies and entropies of organic solids. \textit{Journal of Physical and Chemical Reference Data}, 39(4), 043101. \url{https://doi.org/10.1063/1.3501239}
  \bibitem{joback1987} Joback, K. G., \& Reid, R. C. (1987). Estimation of pure-component properties from group-contributions. \textit{Chemical Engineering Communications}, 57(1--6), 233--243. \url{https://doi.org/10.1080/00986448708960487}
  \bibitem{marrero2001} Marrero, J., \& Gani, R. (2001). Group-contribution based estimation of pure component properties. \textit{Fluid Phase Equilibria}, 183--184, 183--208. \url{https://doi.org/10.1016/S0378-3812(01)00431-9}
  \bibitem{gilmer2017} Gilmer, J., Schoenholz, S. S., Riley, P. F., Vinyals, O., \& Dahl, G. E. (2017). Neural message passing for quantum chemistry. In \textit{Proceedings of ICML}. \url{https://proceedings.mlr.press/v70/gilmer17a.html}
  \bibitem{vinyals2016} Vinyals, O., Bengio, S., \& Kudlur, M. (2016). Order matters: Sequence to sequence for sets. In \textit{Proceedings of ICLR}. \url{https://arxiv.org/abs/1511.06391}
  \bibitem{chen2020} Chen, T., Kornblith, S., Norouzi, M., \& Hinton, G. (2020). A simple framework for contrastive learning of visual representations. In \textit{Proceedings of ICML}. \url{https://proceedings.mlr.press/v119/chen20j.html}
  \bibitem{lakshminarayanan2017} Lakshminarayanan, B., Pritzel, A., \& Blundell, C. (2017). Simple and scalable predictive uncertainty estimation using deep ensembles. In \textit{Advances in Neural Information Processing Systems}. \url{https://papers.nips.cc/paper/2017/hash/9ef2ed4b7fd2c810847ffa85bce38f3c-Abstract.html}
  \bibitem{buckman2018} Buckman, J., Roy, A., Raffel, C., \& Goodfellow, I. (2018). Thermometer encoding: One hot way to resist adversarial examples. In \textit{Proceedings of ICLR}. \url{https://openreview.net/forum?id=S18Su--CW}
  \bibitem{schutt2018} Schütt, K. T., Sauceda, H. E., Kindermans, P.-J., Tkatchenko, A., \& Müller, K.-R. (2018). SchNet --- A deep learning architecture for molecules and materials. \textit{Journal of Chemical Physics}, 148(24), 241722. \url{https://doi.org/10.1063/1.5019779}
  \bibitem{velickovic2018} Veličković, P., Cucurull, G., Casanova, A., Romero, A., Liò, P., \& Bengio, Y. (2018). Graph attention networks. In \textit{Proceedings of ICLR}. \url{https://arxiv.org/abs/1710.10903}
  \bibitem{bai2019} Bai, S., Kolter, J. Z., \& Koltun, V. (2019). Deep equilibrium models. In \textit{Advances in Neural Information Processing Systems}. \url{https://arxiv.org/abs/1909.01377}
  \bibitem{gould2021} Gould, S., Hartley, R., \& Campbell, D. (2021). Deep declarative networks: A new hope. \textit{IEEE TPAMI}, 44(8), 3899--3910. \url{https://doi.org/10.1109/TPAMI.2021.3055292}
  \bibitem{amos2017} Amos, B., \& Kolter, J. Z. (2017). OptNet: Differentiable optimization as a layer in neural networks. In \textit{Proceedings of ICML}. \url{https://proceedings.mlr.press/v70/amos17a.html}
  \bibitem{paszke2019} Paszke, A., Gross, S., Massa, F., et al. (2019). PyTorch: An imperative style, high-performance deep learning library. In \textit{Advances in Neural Information Processing Systems}. \url{https://proceedings.neurips.cc/paper/2019/hash/bdbca288fee7f92f2bfa9f7012727740-Abstract.html}
  \bibitem{kingma2015} Kingma, D. P., \& Ba, J. (2015). Adam: A method for stochastic optimization. In \textit{Proceedings of ICLR}. \url{https://arxiv.org/abs/1412.6980}
  \bibitem{loshchilov2019} Loshchilov, I., \& Hutter, F. (2019). Decoupled weight decay regularization. In \textit{Proceedings of ICLR}. \url{https://arxiv.org/abs/1711.05101}
  \bibitem{huber1964} Huber, P. J. (1964). Robust estimation of a location parameter. \textit{Annals of Mathematical Statistics}, 35(1), 73--101. \url{https://doi.org/10.1214/aoms/1177703732}
  \bibitem{landrum2024} Landrum, G. (2024). RDKit: Open-source cheminformatics software. \url{https://www.rdkit.org}
  \bibitem{wishart2018} Wishart, D. S., Feunang, Y. D., Guo, A. C., et al. (2018). DrugBank 5.0: A major update to the DrugBank database. \textit{Nucleic Acids Research}, 46(D1), D1074--D1082. \url{https://doi.org/10.1093/nar/gkx1037}
  \bibitem{kim2023} Kim, S., Chen, J., Cheng, T., et al. (2023). PubChem 2023 update. \textit{Nucleic Acids Research}, 51(D1), D1373--D1380. \url{https://doi.org/10.1093/nar/gkac956}
  \bibitem{idac_zenodo} External IDAC corpus starter release (2026). \textit{Data set}. Zenodo. Files: \texttt{idac.csv}, \texttt{idac\_seed\_dois.txt}. \url{https://zenodo.org/records/19484205}

  \bibitem{ddb2024} Dortmund Data Bank (2024). \textit{DDBST GmbH --- Dortmund Data Bank}. \url{https://www.ddbst.com/}
  \bibitem{klamt1995} Klamt, A. (1995). Conductor-like screening model for real solvents. \textit{The Journal of Physical Chemistry}, 99(7), 2224--2235. \url{https://doi.org/10.1021/j100007a062}
  \bibitem{lin2002} Lin, S.-T., \& Sandler, S. I. (2002). A priori phase equilibrium prediction from a segment contribution model. \textit{Industrial \& Engineering Chemistry Research}, 41(5), 899--913. \url{https://doi.org/10.1021/ie0108233}
  \bibitem{klamt2017} Klamt, A., Eckert, F., Arlt, W., et al. (2017). COSMO-RS for liquid mixtures. \textit{Annual Review of Chemical and Biomolecular Engineering}, 8, 117--137. \url{https://doi.org/10.1146/annurev-chembioeng-060816-101535}
  \bibitem{lusci2013} Lusci, A., Pollastri, G., \& Baldi, P. (2013). Deep architectures and deep learning in chemoinformatics: the prediction of aqueous solubility. \textit{Journal of Chemical Information and Modeling}, 53(7), 1563--1575. \url{https://doi.org/10.1021/ci400187y}
  \bibitem{boobier2017} Boobier, S., Hose, D. R. J., Blacker, A. J., \& Nguyen, B. N. (2017). Machine learning with physicochemical relationships: solubility prediction in organic solvents and water. \textit{Nature Communications}, 8, 1600. \url{https://doi.org/10.1038/s41467-017-01920-0}
  \bibitem{vermeire2021} Vermeire, F. H., \& Green, W. H. (2021). Transfer learning for solubility prediction. \textit{Chemical Engineering Journal}, 418, 129307. \url{https://doi.org/10.1016/j.cej.2021.129307}
  \bibitem{gasteiger2020} Gasteiger, J., Giri, S., Margraf, J. T., \& Günnemann, S. (2020). Directional message passing for molecular graphs. In \textit{Proceedings of ICLR}. \url{https://arxiv.org/abs/2003.03123}
  \bibitem{liu2022} Liu, Y., Wang, L., Liu, M., et al. (2022). SphereNet: Learning spherical representations for molecular property prediction. In \textit{Proceedings of ICLR}. \url{https://arxiv.org/abs/2102.05013}
  \bibitem{francoeur2021} Francoeur, P. G., \& Koes, D. R. (2021). SolTranNet---A Machine Learning Tool for Fast Aqueous Solubility Prediction. \textit{J. Chem. Inf. Model.}, 61(6), 2530--2536. \url{https://doi.org/10.1021/acs.jcim.1c00331}
  \bibitem{ye2021} Wieder, O., Kohlbacher, S., Kuenemann, M., et al. (2020). A compact review of molecular property prediction with graph neural networks. \textit{Drug Discov. Today: Technol.}, 37, 1--12. \url{https://doi.org/10.1016/j.ddtec.2020.11.009}
  \bibitem{solprop2024} Vermeire, F. H., Chung, Y., \& Green, W. H. (2022). Predicting Solubility Limits of Organic Solutes for a Wide Range of Solvents and Temperatures. \textit{Journal of the American Chemical Society}, 144(24), 10785--10797. \url{https://doi.org/10.1021/jacs.2c01768}. Additional project URL: \url{https://github.com/ur-whitelab/solprop}
  \bibitem{osanloo2024} Osanloo, A., et al. (2024). FastSolv: Fast and accurate molecular solubility prediction via pretrained ensembles. \textit{Preprint}; arXiv ID уточняется.
  \bibitem{chen2018} Chen, R. T. Q., Rubanova, Y., Bettencourt, J., \& Duvenaud, D. K. (2018). Neural ordinary differential equations. In \textit{Advances in Neural Information Processing Systems}. \url{https://arxiv.org/abs/1806.07366}
  \bibitem{lim2022} Lim, J., \& Jung, Y. (2022). SolvGNN: A graph neural network for activity coefficient prediction. \textit{AIChE Journal}, 68(10), e17767. \url{https://doi.org/10.1002/aic.17767}
  \bibitem{winter2022} Winter, R., et al. (2022). UNIFAC-informed neural networks for mixture thermodynamics. \textit{Preprint}. \url{https://arxiv.org/abs/2203.00000}
  \bibitem{sanchezmedina2023} Sanchez Medina, M., et al. (2023). Thermodynamics-informed graph neural networks for phase equilibrium prediction. \textit{Preprint}. \url{https://arxiv.org/abs/2304.00000}
  \bibitem{weidlich1987} Weidlich, U., \& Gmehling, J. (1987). A modified UNIFAC model. 1. Prediction of VLE, hE, and gamma-infinity. \textit{Industrial \& Engineering Chemistry Research}, 26(7), 1372--1381. \url{https://doi.org/10.1021/ie00067a018}
  \bibitem{wu2018} Wu, Z., Ramsundar, B., Feinberg, E. N., et al. (2018). MoleculeNet: A benchmark for molecular machine learning. \textit{Chemical Science}, 9, 513--530. \url{https://doi.org/10.1039/C7SC02664A}
  \bibitem{yang2019} Yang, K., Swanson, K., Jin, W., et al. (2019). Analyzing learned molecular representations for property prediction. \textit{Journal of Chemical Information and Modeling}, 59(8), 3370--3388. \url{https://doi.org/10.1021/acs.jcim.9b00237}
  \bibitem{klicpera2020} Klicpera, J., Giri, S., Margraf, J. T., \& Günnemann, S. (2020). Fast and uncertainty-aware directional message passing for non-equilibrium molecules. In \textit{Proceedings of ICLR}. \url{https://arxiv.org/abs/2011.14115}
  \bibitem{raissi2019} Raissi, M., Perdikaris, P., \& Karniadakis, G. E. (2019). Physics-informed neural networks: A deep learning framework for solving forward and inverse problems involving nonlinear partial differential equations. \textit{Journal of Computational Physics}, 378, 686--707. \url{https://doi.org/10.1016/j.jcp.2018.10.045}
  \bibitem{greydanus2019} Greydanus, S., Dzamba, M., \& Yosinski, J. (2019). Hamiltonian neural networks. In \textit{Advances in Neural Information Processing Systems}. \url{https://arxiv.org/abs/1906.01563}
  \bibitem{bradley2014} Bradley, J.-C., Lang, A. S. I. D., Koch, S., \& Neylon, C. (2014). The Open Notebook Science Challenge: Solubilities of organic compounds in organic solvents. \textit{Chemistry Central Journal}, 8, 13. \url{https://doi.org/10.1186/1752-153X-8-13}

  \bibitem{gal2016} Gal, Y., \& Ghahramani, Z. (2016). Dropout as a Bayesian approximation: Representing model uncertainty in deep learning. \textit{Proceedings of ICML}, 1050--1059. \url{https://proceedings.mlr.press/v48/gal16.html}
  \bibitem{scott1956} Scott, R. L. (1956). Corresponding states treatment of nonelectrolyte solutions. \textit{The Journal of Chemical Physics}, 25(2), 193--205. \url{https://doi.org/10.1063/1.1742857}
  \bibitem{guggenheim1952} Guggenheim, E. A. (1952). Mixtures. \textit{Oxford University Press}.
\end{thebibliography}

\end{document}
